% main.tex — artigo IEEE ICIT 2027 (IEEEtran conference), Agente E, 24/10/2026.
% Dois formatos controlados por \ifwip: regular (6 pp, \wipfalse) e WiP (4 pp, \wiptrue).
% NENHUM número é digitado neste arquivo: todo número vem de numeros.tex (macros \V…),
% figuras.tex (caminhos) ou das tabelas geradas em ../tables/latex/.
% Categorias de evidência nas macros: \VEst… estimativa analítica | \VPc… medição no PC |
% \VBuild… build do firmware (idf.py size) | \VChip… medição no ESP32 (pending sem a bancada).
% Não compilado aqui (BLOQUEIOS B4); conferido por src/python/gerar_artigo_checar.py.
\documentclass[conference]{IEEEtran}
\IEEEoverridecommandlockouts
\usepackage{cite}
\usepackage{amsmath,amssymb}
\usepackage{graphicx}
\usepackage{booktabs}
\usepackage{url}

\newif\ifwip
% WiP: compile com  pdflatex -jobname=main_wip "\def\usewip{1}% main.tex — artigo IEEE ICIT 2027 (IEEEtran conference), Agente E, 24/10/2026.
% Dois formatos controlados por \ifwip: regular (6 pp, \wipfalse) e WiP (4 pp, \wiptrue).
% NENHUM número é digitado neste arquivo: todo número vem de numeros.tex (macros \V…),
% figuras.tex (caminhos) ou das tabelas geradas em ../tables/latex/.
% Categorias de evidência nas macros: \VEst… estimativa analítica | \VPc… medição no PC |
% \VBuild… build do firmware (idf.py size) | \VChip… medição no ESP32 (pending sem a bancada).
% Não compilado aqui (BLOQUEIOS B4); conferido por src/python/gerar_artigo_checar.py.
\documentclass[conference]{IEEEtran}
\IEEEoverridecommandlockouts
\usepackage{cite}
\usepackage{amsmath,amssymb}
\usepackage{graphicx}
\usepackage{booktabs}
\usepackage{url}

\newif\ifwip
% WiP: compile com  pdflatex -jobname=main_wip "\def\usewip{1}% main.tex — artigo IEEE ICIT 2027 (IEEEtran conference), Agente E, 24/10/2026.
% Dois formatos controlados por \ifwip: regular (6 pp, \wipfalse) e WiP (4 pp, \wiptrue).
% NENHUM número é digitado neste arquivo: todo número vem de numeros.tex (macros \V…),
% figuras.tex (caminhos) ou das tabelas geradas em ../tables/latex/.
% Categorias de evidência nas macros: \VEst… estimativa analítica | \VPc… medição no PC |
% \VBuild… build do firmware (idf.py size) | \VChip… medição no ESP32 (pending sem a bancada).
% Não compilado aqui (BLOQUEIOS B4); conferido por src/python/gerar_artigo_checar.py.
\documentclass[conference]{IEEEtran}
\IEEEoverridecommandlockouts
\usepackage{cite}
\usepackage{amsmath,amssymb}
\usepackage{graphicx}
\usepackage{booktabs}
\usepackage{url}

\newif\ifwip
% WiP: compile com  pdflatex -jobname=main_wip "\def\usewip{1}% main.tex — artigo IEEE ICIT 2027 (IEEEtran conference), Agente E, 24/10/2026.
% Dois formatos controlados por \ifwip: regular (6 pp, \wipfalse) e WiP (4 pp, \wiptrue).
% NENHUM número é digitado neste arquivo: todo número vem de numeros.tex (macros \V…),
% figuras.tex (caminhos) ou das tabelas geradas em ../tables/latex/.
% Categorias de evidência nas macros: \VEst… estimativa analítica | \VPc… medição no PC |
% \VBuild… build do firmware (idf.py size) | \VChip… medição no ESP32 (pending sem a bancada).
% Não compilado aqui (BLOQUEIOS B4); conferido por src/python/gerar_artigo_checar.py.
\documentclass[conference]{IEEEtran}
\IEEEoverridecommandlockouts
\usepackage{cite}
\usepackage{amsmath,amssymb}
\usepackage{graphicx}
\usepackage{booktabs}
\usepackage{url}

\newif\ifwip
% WiP: compile com  pdflatex -jobname=main_wip "\def\usewip{1}\input{main}"
\ifdefined\usewip\wiptrue\else\wipfalse\fi

\input{numeros.tex}
\input{numeros_aud.tex}
\input{figuras.tex}

\begin{document}

\title{Porting a Harmonic-Sum Rotation Tracker for Acoustic Motor-State Monitoring to a
Microcontroller: Stage-Wise Verified C Port, Measured Cost, and What Two Public Datasets
Allow Us to Conclude\ifwip\ (Work in Progress)\fi}

% TODO(autores): nomes e afiliações a preencher pelos autores.
\author{\IEEEauthorblockN{[Authors to be completed]}
\IEEEauthorblockA{[Affiliation]}}

\maketitle

\begin{abstract}
We port to portable C a pipeline that classifies the operating state of small electric motors
from sound by borrowing ideas from bass-guitar transcription: a harmonic-sum rotation tracker
decoded by a Viterbi search, \VParNFeat{} physical features, cosine-to-prototype scoring and an
online hidden Markov model (HMM). The port is checked against the Python reference stage by
stage, on \VEqCases{} windows and \VEqSeqCombos{} sequence configurations, with tolerances
frozen before the formal comparison. In double precision all \VEqStagesTotal{} stage checks
pass. In float32, the build intended for the microcontroller, all \VEqRealCases{} windows of real
audio pass every stage; on \VEqFloatBadCases{} synthetic constant-offset-plus-noise windows the
envelope-modulation peak lies at the float32 noise floor and the modulation, feature and z-score
checks fail, while the class decision is identical on \VEqFloatDecisionPass{} of \VEqCases{}
windows. On an ESP32 without PSRAM the reference memory layout (estimated at
\VEstWsFloatThreeMb{}\,MB for a \VParSegIdmt{}-s window) does not fit. A named low-memory
variant fits \VParSegBldc{}-s windows, and on the chip its outputs agree with the reference within
the frozen tolerances on all \VChipEqCases{} windows compared; measured on the chip it is
\VChipRtfBldc{} times slower than real time, dominated by a table-free harmonic sum in
software-emulated double precision, and the \VParSegIdmt{}-s window does not fit. The pipeline does not outperform simple baselines:
on the IDMT dataset every method, including a two-number level baseline, saturates because
class, motor and recording coincide; on the A2212 dataset classical features reach
\VDiagBalClsBldcFile{} balanced accuracy against \VDiagBalDecBldcFile{} for the pipeline, and
with a motor held out all methods drop. The contribution is therefore the verified embedded
implementation, its cost, an interpretable rotation estimate, and an explicit account of what
these datasets can and cannot support.
\end{abstract}

\begin{IEEEkeywords}
acoustic condition monitoring, embedded signal processing, ESP32, rotation tracking,
Viterbi decoding, reproducibility, data leakage
\end{IEEEkeywords}

% ======================================================================
\section{Introduction}

Acoustic monitoring of electric motors is attractive because a microphone is cheap,
non-intrusive and needs no electrical connection to the machine. Several recent works report
near-perfect accuracy on two public datasets of small laboratory motors, the
IDMT-ISA-ELECTRIC-ENGINE set~\cite{grollmisch2019} and the A2212 brushless set~\cite{estacio2023},
often with deep networks~\cite{melcnn2025,ertarin2025} or large hand-crafted feature banks~\cite{etnext2026,tekin2025}.
Two aspects receive little attention. First, most of these results split recordings into
segments before the train/test split, so neighbouring segments of one recording fall on both
sides. Second, few works measure the cost on a microcontroller; an exception is the edge
classifier of Strantzalis et al.~\cite{strantzalis2022}, while other TinyML studies stop at a model
size and never run on hardware~\cite{almaini2026}.

This paper starts from a pipeline that we call Decifrae (p11). It reuses components of a
music-transcription system: the motor is treated as a ``one-string instrument'' whose
fundamental (the rotation frequency) is tracked through a sum over its harmonics and decoded with
a Viterbi search that has a silence state and an octave-jump cost, as a bass transcriber would do.
The analogy only motivates the \emph{representation}; we do not claim that motor sound is a
musical note, nor that music-transcription performance transfers to machines.

The question we answer is whether this pipeline fits and runs on an ESP32-class
microcontroller with an implementation that is demonstrably the same computation as the Python
reference, and what the two public datasets allow us to conclude about it. In short, it fits and
runs on an ESP32 without PSRAM only for \VParSegBldc{}-s windows, through a named low-memory
variant, and \VChipRtfBldc{} times slower than real time on that board.
\ifwip
We contribute a C port verified stage by stage against the reference with tolerances frozen before
the formal comparison, a cost account that keeps estimate, PC measurement, firmware build and
on-chip measurement apart, the tracked rotation as an interpretable output, and an evaluation
under three protocols showing that the pipeline does not beat simple baselines and that the
datasets cannot separate operating state from motor identity.
\else
Our contributions are:
\begin{itemize}
\item a portable C port without dynamic allocation, verified against the Python reference stage
by stage with tolerances frozen before the formal comparison, in double and in float32, with
every residual discrepancy listed;
\item a cost account that keeps four categories apart: analytical estimate, measurement on a PC,
firmware build (\texttt{idf.py size}), and measurement on an ESP32 without PSRAM, including the
low-memory variant needed there;
\item an interpretable intermediate output, the tracked rotation frequency, together with a check
of how far it can be trusted on each dataset;
\item an evaluation of the pipeline and of simple baselines on the same windows under three
protocols (random windows, grouped by recording or file, and motor held out), showing that the
pipeline does not outperform the baselines and that the datasets cannot separate operating
state from motor identity.
\end{itemize}
\fi
Three distinctions are kept throughout. We classify the \emph{current operating state} of a
motor; we do not predict future failures. Both datasets contain small laboratory motors, not
industrial machines. And an evaluation in which the test motor was also seen in training says
nothing about generalisation to a new motor.

% ======================================================================
\section{Related Work}

\emph{Datasets and splits.} Grollmisch et al.~\cite{grollmisch2019} released the IDMT set with
an official split (training on clean recordings, testing on recordings with added noise) and a
frame-level network. Each class was recorded on a different motor, so class and physical unit
coincide. Later works on IDMT and A2212 report near-perfect accuracy with mel-spectrogram
CNNs~\cite{melcnn2025}, image encodings~\cite{ertarin2025}, wavelet pattern features with
SVM/kNN~\cite{etnext2026} and wavelet statistics with bagged trees~\cite{tekin2025}; their splits
operate on segments rather than on recordings. The A2212 data paper~\cite{estacio2023} reports
no classifier; it documents that the bearing fault exists on a single motor.
\ifwip\else
The classical-feature baseline used here follows the feature list described in the abstract of
Vanaraj et al.~\cite{vanaraj2026}; we only read the abstract, so it is an inspired
re-implementation, not a reproduction. Glowacz~\cite{glowacz2018} classified a commutator motor
with a nearest-mean rule on selected spectral lines, the closest precedent of our
cosine-to-prototype decision.
\fi

\ifwip
\emph{Rotation tracking without a tachometer} has been reviewed on vibration
data~\cite{peeters2019}, and Viterbi decoding over a time--frequency representation predates our
work~\cite{xi2015}; we add its verified embedded implementation on microphone audio.
\else
\emph{Rotation tracking without a tachometer.} Peeters et al.~\cite{peeters2019} review and
compare tacholess speed estimators on vibration data, including multi-order probabilistic
approaches that combine several harmonics. Viterbi decoding over a time--frequency
representation for speed estimation predates our work~\cite{xi2015}, and harmonic weighting has
been proposed to suppress weak orders~\cite{yoo2023}. Wheat et al.~\cite{wheat2025} build
features on a harmonic grid given a known operating speed. Our tracker is therefore not new as an
idea; what we add is its embedded, verified implementation on microphone audio.
\fi

\ifwip\else
\emph{Embedded and anomaly-oriented work.} Strantzalis et al.~\cite{strantzalis2022} deploy a
small CNN on a Cortex-M4 and measure latency, RAM and response time to state changes on IDMT,
but with a random split of overlapping spectrograms. Almaini et al.~\cite{almaini2026} quantise an
MFCC network trained on urban sounds without deploying it. Unsupervised anomaly detection on
machine sounds (autoencoder baselines of the DCASE challenges) addresses novelty rather than state
classification and is not compared here.
\fi

% ======================================================================
\section{Method}
\label{sec:method}

Fig.~\ref{fig:arq} summarises the pipeline. Each analysis window (\VParSegIdmt{}\,s for IDMT,
\VParSegBldc{}\,s for A2212, both at \VParSrKhz{}\,kHz) is RMS-normalised and processed
independently; only the HMM carries state across windows.

\begin{figure}[t]
\centering
\includegraphics[width=\columnwidth]{\FigArq}
\caption{System overview. Everything inside the dashed box is the C port that runs on the
microcontroller; audio currently reaches the chip over UART (processor-in-the-loop) or from its
flash. Working memory: analytical estimate of the reference layout, and the pool peak of the
low-memory variant measured on the ESP32.}
\label{fig:arq}
\end{figure}

\emph{Rotation tracker.} A short-time Fourier transform with a symmetric Hann window of
\VParNfft{} samples and hop \VParHop{} gives magnitude frames $|X_t(k)|$. For each of
$J=\VParJ{}$ candidate rotation frequencies $f_j$ on a logarithmic grid from \VParFmin{} to
\VParFmax{}\,Hz (step \VParStepCents{}~cents), the saliency
$s_t(j)=\sum_{h=1}^{K}|X_t(h f_j/\Delta f)|$ sums $K=\VParNharm{}$ harmonics with linear
interpolation between bins, and is normalised per frame to $\tilde e_t(j)\in[0,1]$. The emission
cost is $-\log(\VParWEnergy{}\,\tilde e_t(j)+\VParEmisEps{})$; a silence state receives a cost
derived from the frame energy relative to the loudest frame of the window. A Viterbi search
allows jumps of up to \VParMaxJump{} states (\VParMaxJumpCents{} cents) with a cost proportional
to the jump, octave jumps of $\pm$\VParOctStates{} states with a fixed cost, and transitions into
and out of silence. The decoded path gives $f_{rot}(t)$.

\emph{Features and decision.} From the path and the mean power spectrum we compute \VParNFeat{}
features\ifwip\else{} (Table~\ref{tab:params})\fi: the median and inter-quartile range of $f_{rot}$, the silent-frame
fraction, the saliency along the path, a harmonic-energy ratio, the relative energy of four groups
of harmonic orders, the frequency and prominence of the main envelope-modulation peak (Hilbert
envelope, Butterworth low-pass, decimation to \VParEnvRateHz{}\,Hz, Welch spectrum), the spectral
centroid and the spectral flatness. Features are z-scored with training statistics and compared
by cosine with one prototype per class; a softmax with temperature \VParTempCos{} gives class
probabilities. For streams, an HMM with a fixed change penalty and an extra ``no decision''
state is run as a causal forward filter, one window at a time.

\ifwip\else
\input{../tables/latex/parametros_features.tex}
\fi

\ifwip
\emph{C port.} The C code follows the Python line by line, compiles unchanged for the PC
and the ESP-IDF, processes the spectrogram frame by frame, includes a mixed-radix FFT for the
full-length Hilbert transform, and uses a caller-provided arena instead of dynamic allocation.
The real type is float32 on the microcontroller and double for verification. Any change of
arithmetic or memory layout is a named variant: \texttt{p11\_c\_sosd} keeps the envelope filter
state in double (more precise, not an approximation), and \texttt{\VChipVariant}, the one run on
our board without PSRAM, takes its memory from a pool of internal DRAM and IRAM blocks, drops the
harmonic table (positions recomputed in double) and computes the Hilbert FFT in four steps. An
ESP-DSP option (\texttt{p11\_c\_espdsp}) was neither compared nor measured and is left out.
\else
\emph{C port.} The C code follows the Python line by line and compiles unchanged for the PC and
for the ESP-IDF. It never stores the full spectrogram: each frame is windowed, transformed by a
real FFT, accumulated into the mean power spectrum and reduced to its saliency row. The Viterbi
search needs the whole window (the silence cost uses the maximum frame energy of the window), so
the normalised saliency and the back-pointers (short integers) are kept; the Hilbert envelope needs an FFT
of the full window length, which is not a power of two, so the port includes a mixed-radix FFT.
All buffers live in a caller-provided arena whose size is computed by the same layout function that
assigns the pointers. The real type is selected at compile time (float32 for the microcontroller,
double for verification). Index decisions that change integers (grid, medians, rounding of
harmonic positions) are computed in double in both builds, reproducing the reference's
round-half-to-even and first-index tie-breaking. Any change of arithmetic is a \emph{named
variant}, compared under the same tolerances. \texttt{p11\_c\_sosd} keeps the envelope filter
state in double, as SciPy does (more precise, not an approximation). \texttt{\VChipVariant}, run on
our board without PSRAM, calls the same core functions but carves its memory phase by phase from
a pool of internal DRAM and IRAM blocks (IRAM, word access only, holds only float vectors such as
the FFT twiddles), drops the harmonic-interpolation table (bin positions recomputed in double for
every frame) and computes the Hilbert FFT in four steps, the same DFT in another order. An ESP-DSP
option (\texttt{p11\_c\_espdsp}) was neither compared nor measured and is left out.
\fi

% ======================================================================
\section{Experimental Setup}
\label{sec:setup}

\emph{Data and protocols.} We use the \texttt{train/} part of IDMT (the official test set was not
opened) and all of A2212\ifwip\else{} (Table~\ref{tab:protocols})\fi. All methods are evaluated on exactly the
same windows and folds. On IDMT, the random-window protocol leaks by construction, and the
``blocks'' protocol holds out contiguous blocks of each recording; since there is one recording
per class, it still cannot separate class from motor. On A2212 we compare random windows, a
grouped split by file, and leave-one-motor-out; because the bearing fault exists only on motor
M4, we also report leave-one-motor-out without the M4 fold.

\ifwip\else\input{../tables/latex/bases_protocolos.tex}\fi

\emph{Baselines.} A level baseline uses only the RMS and peak of the window with logistic
regression, once on RMS-normalised audio (only the crest factor remains) and once on raw audio. A
classical baseline (p04) uses \VParNClassicFeat{} frame statistics (zero-crossing rate, RMS,
spectral centroid and MFCC means and deviations) with a soft-voting ensemble, inspired
by~\cite{vanaraj2026}.

\emph{Metrics.} We report balanced accuracy, macro-F1, macro one-vs-rest ROC-AUC from continuous
scores, and per-class recall. Confidence intervals (\VDataCiLevel{}\%) come from a bootstrap that
resamples groups (recording blocks or files), stratified by class, with \VDataNBoot{} resamples;
windows of one group are never treated as independent. For the motor-out protocols the
independent unit is the motor, of which there are only \VDataBldcMotors{}, so no honest interval
exists at that level; the file-level interval shown is optimistic.

\ifwip
\emph{Equivalence protocol.} Golden vectors of every stage were produced by the
instrumented reference on \VEqCases{} windows (\VEqGoldClean{} clean, \VEqGoldNoisy{} noisy,
\VEqGoldEdge{} edge cases) and \VEqSeqCombos{} HMM sequence configurations. Per-stage tolerances,
derived from rounding analysis and an independent emulation, were frozen (\texttt{\VEqTolCommit{}})
before the formal comparison, although the C port had already been debugged against the same
golden vectors by its developer; discrete outputs must agree exactly unless a tie is demonstrated.
\else
\emph{Equivalence protocol.} The instrumented Python reference produced golden vectors for every
stage on \VEqCases{} windows: \VEqGoldClean{} clean validation windows from both datasets,
\VEqGoldNoisy{} noisy versions under \VEqGoldPerts{} perturbations (white noise at several SNRs and
synthetic sag and flicker power-quality disturbances), and \VEqGoldEdge{} edge cases (silence, clipping, synthetic
rotation outside the search range, and a constant offset with tiny noise). The HMM stages were
checked on \VEqGoldSeqs{} probability sequences under \VParNChangePen{} change penalties and
\VEqMinDurValues{} minimum durations (\VEqSeqCombos{} configurations). Tolerances for each stage
were derived from first-order rounding analysis and an independent numerical emulation of each
build, and frozen before the formal comparison (commit \texttt{\VEqTolCommit{}}, file hash
\texttt{\VEqTolSha{}}\ldots, \VEqTolCommitTime{}; comparison run \VEqRunTime{}); the hash is
recorded in every result row. The chronology matters: the C port had already been debugged
against the same golden vectors by its developer before the tolerances were frozen, so the
tolerances are frozen, not pre-registered. On \audHoldoutN{} windows outside the golden set,
unseen by the C developer, path and decision agree on \audHoldoutFloatPathSame{} and
\audHoldoutFloatDecSame{} in float32, no feature exceeds its tolerance (maximum ratio
\audHoldoutFloatRatioMax{}) and probabilities differ by at most \audHoldoutFloatDpMax{}. Tolerances for the edge cases
were calibrated on the emulation error for the same signals, a weaker standard, and some are wide
in float32 (e.g.\ the modulation prominence). Continuous stages are compared by maximum error relative to the stage's
scale; discrete outputs (rotation path, decision, HMM states) must agree exactly unless a tie is
demonstrated from the reference's own values. A failure at one stage marks later stages as
downstream. Both emulations passed all comparison rows (\VEqEmulDoublePass{} and \VEqEmulFloatPass{} of
\VEqEmulFloatRows{}) before the C code was compared.
\fi

\ifwip
\emph{Cost measurement.} We keep apart (1) the analytical arena estimate, (2) timings on
the PC (\VPcNIdmt{} timings cycling over \VPcDistinctIdmt{} distinct windows, each run once in a new
process, so warm-up does not warm the C code; workspace initialisation excluded), (3)
\texttt{idf.py size} of the firmware (ESP-IDF \VBuildIdf{}) and (4) measurements on an
\VChipModel{} without PSRAM (\VChipCpuMhz{}\,MHz; \VChipLatBldcN{} timed windows over UART after
\VChipLatWarmup{} warm-up runs, on-chip timer around the processing only). Energy is not measured.
\else
\emph{Cost measurement.} We keep four categories apart. (1) \emph{Analytical estimate}: the
arena size returned by the C layout function. (2) \emph{Measured on the PC}: \VPcNIdmt{} timings
per dataset cycling over \VPcDistinctIdmt{} distinct windows. The C executable is started anew and
timed once per window, so warm-up does not warm it, and its time excludes the workspace
initialisation, which the Python timing cannot exclude; the two effects partly offset.
(3) \emph{Firmware build}: \texttt{idf.py size} (ESP-IDF \VBuildIdf{}, GCC \VBuildGcc{}).
(4) \emph{Measured on the ESP32}: an \VChipModel{} (rev.\ \VChipRev{}, \VChipCpuMhz{}\,MHz,
\VChipFlashMb{}\,MB flash, no PSRAM, radios never started); windows arrive over UART (or, for the
stream, from flash) and are timed with the on-chip timer around each core call, with heap and
stack watermarks. Energy is not measured.
\fi

% ======================================================================
\section{Results}
\label{sec:results}

\subsection{Python--C equivalence}

\ifwip
Table~\ref{tab:equiv} summarises the verdicts. In double precision all \VEqStagesTotal{}
stage checks pass on all windows and sequence configurations; the rotation path, the decision and
every HMM output agree exactly. In float32 all \VEqRealCases{} windows of real audio pass every
stage. The failures are confined to the \VEqFloatBadCases{} synthetic constant-offset-plus-noise
windows, whose modulation spectrum lies at the float32 rounding floor in both implementations, so
the peak location is undetermined (the reference peak is \VEqDcRefPeakRatio{} times the reference
power at the bin chosen by the C code: not a tie). We count the modulation, feature and z-score
checks as failed there (\VEqFloatStagesPassStrict{} of \VEqStagesTotal{} checks pass on all
windows); the class decision agrees on \VEqFloatDecisionPass{} of \VEqCases{} windows. The named
variant \texttt{p11\_c\_sosd} passes every check. An independent audit reproduced the comparison
from a clean build; its corrections are applied here. These are PC results; the comparison on the
chip is in Section~\ref{sec:cost}.
\else
Table~\ref{tab:equiv} summarises the verdicts\ifwip\else{} and Fig.~\ref{fig:equiv} the worst
error-to-tolerance ratio per stage\fi. In double precision all \VEqStagesTotal{} stage checks pass
on all windows and sequence configurations (worst ratio \VEqDoubleMaxRatio{} of the tolerance).
The rotation path, the decision and every HMM output agree exactly. In float32, all
\VEqRealCases{} windows of real audio (clean and noisy) pass every stage. The
\VEqFloatBadCases{} remaining failures are the synthetic constant-offset-plus-noise edge case, once
per model. There the envelope is almost constant and its modulation spectrum lies at the float32
rounding floor in both implementations, so the peak location is undetermined at that precision:
the reference's peak is \VEqDcRefPeakRatio{} times the reference power at the bin chosen by the C
code, so this is not a tie. The modulation frequency moves by \VEqFloatFmodErr{}\,Hz, and we count
the envelope-modulation, feature and z-score checks as failed on these windows (the frozen
comparison rule had accepted the first two as a tie; we report the conservative reading), leaving
\VEqFloatStagesPassStrict{} of \VEqStagesTotal{} stage checks passing on all windows. Cosines,
probabilities and the class decision still pass on all windows
(\VEqFloatDecisionPass{}/\VEqCases{}). With the named variant \texttt{p11\_c\_sosd} every check
passes; it reduces the maximum error of the decimated envelope,
relative to its peak and taken as the median over windows, from \VEqEnvDecFloatMed{} to
\VEqEnvDecSosdMed{}, because the low-pass poles lie close to
the unit circle and float32 rounding accumulates in the filter state. These results were obtained
on the PC; the same comparison on the chip is
\ifchipmedido reported in Section~\ref{sec:cost}\else pending\fi. An independent audit
(report in the repository) reproduced the comparison from a clean build and its corrections are
applied here. We therefore
claim stage-wise agreement within tolerances frozen before the formal comparison, with the
exceptions above, not unconditional equivalence.
\fi

\begin{table}[t]
\centering
\scriptsize
\caption{Stage checks passed on the PC (\VEqCases{} windows; \VEqSeqCombos{} sequence
configurations). Conservative reading: a group fails if any stage fails on any window. All
float32 failures are the \VEqFloatBadCases{} synthetic constant-offset edge windows, where the
modulation peak is at the float32 noise floor; the decision is unchanged.}
\label{tab:equiv}
\input{../tables/latex/equivalencia_compacta.tex}
\end{table}

\ifwip\else
\begin{figure}[t]
\centering
\includegraphics[width=\columnwidth]{\FigEquiv}
\caption{Worst error-to-tolerance ratio per stage (PC). Ratios below one pass. F: failure on the
synthetic constant-offset edge windows (counted as failures even where the frozen rule accepted a
tie). Exact zeros are drawn at the bottom of the axis.}
\label{fig:equiv}
\end{figure}
\fi

\subsection{Cost}
\label{sec:cost}

\ifwip
\emph{Measured on the PC},
the C float32 build of the reference layout needs \VPcCIdmtMed{}\,ms per \VParSegIdmt{}-s window
and \VPcCBldcMed{}\,ms per \VParSegBldc{}-s window, against \VPcPyIdmtMed{}\,ms and
\VPcPyBldcMed{}\,ms for the Python reference (\VPcSpeedupIdmt{}$\times$, including interpreter
overhead); the STFT with harmonic sum takes \VPcCIdmtTrackPct{}\%, the rotation Viterbi
\VPcCIdmtVitPct{}\% and the envelope branch \VPcCIdmtEnvPct{}\% of the median time. These numbers
say nothing about the microcontroller.
\else
Table~\ref{tab:resources} lists the resources by category. The PC timings (column 2; the STFT
with harmonic sum takes \VPcCIdmtTrackPct{}\%, the rotation Viterbi \VPcCIdmtVitPct{}\% and the
envelope branch \VPcCIdmtEnvPct{}\% of the median) say nothing about the microcontroller.
\fi

The \emph{analytical estimate} of the working memory of the reference layout (one arena) is
\VEstWsFloatThreeKib{}\,KiB for \VParSegIdmt{}\,s and \VEstWsFloatOneKib{}\,KiB for
\VParSegBldc{}\,s, at least \VChipRefOverPoolThree{} and \VChipRefOverPoolOne{} times the internal
RAM that the firmware can hand to the working memory on our board (\VChipPoolKib{}\,KiB, measured)
and at least \VBuildEspsWsRatio{} and \VBuildEspsWsRatioOne{} times the whole DIRAM of the
ESP32-S3 \emph{build} (\VBuildEspsDramTotalKb{}\,KiB; its IRAM is \VBuildEspsIramPct{}\% used):
without PSRAM the reference layout cannot run.

\ifchipmedido
\emph{Measured on the ESP32} (\VChipPsramMb{}\,MB PSRAM), \texttt{\VChipVariant} fits a
\VParSegBldc{}-s window: the pool peaks at \VChipWsBldcB{}\,bytes (\VChipWsPoolPct{}\% of it,
Fig.~\ref{fig:esplat}b), the minimum free internal heap is \VChipHeapIntBldcB{}\,bytes and the
stack headroom \VChipStackBldc{}\,bytes. The \VParSegIdmt{}-s IDMT window does \emph{not} fit, so
no IDMT figure exists on the chip. All \VChipEqCases{} golden \VParSegBldc{}-s windows
(\VChipEqCasesReal{} real, \VChipEqCasesEdge{} edge cases) pass, on the chip, the \VChipEqStages{}
final outputs the firmware returns (Viterbi path and score through decision), with \VChipEqFail{} failures and a worst error-to-tolerance
ratio of \VChipEqMaxRatio{}; intermediate stages are not sent back. On the PC the variant passes
all \VEqLmStages{} compared stages. The median latency is \VChipLatBldcMed{}\,ms per
\VParSegBldc{}-s window (p95 \VChipLatBldcPninetyfive{}, maximum \VChipLatBldcMax{}\,ms), a
real-time factor of \VChipRtfBldc{}: on this board the pipeline does \emph{not} run in real time.
The STFT with harmonic sum takes \VChipBldcTrackPct{}\% of it (Fig.~\ref{fig:esplat}a; on the PC,
with the table, \VPcCBldcTrackPct{}\%). We attribute this to recomputing the harmonic positions in
double precision, which the ESP32 only emulates in software (the variant was not timed on the PC);
a float32 variant of that step or a board with PSRAM is the way towards real time. The variant's
image occupies \VBuildLmBinKb{}\,KiB of flash (build).
\else
\emph{Measured on the ESP32}: not available when this version was generated.
\fi

\ifwip\else\input{../tables/latex/recursos.tex}\fi

\iffigesp
\begin{figure}[t]
\centering
\includegraphics[width=\columnwidth]{\FigEspLat}
\caption{Measured on the ESP32 without PSRAM (\texttt{\VChipVariant}, A2212, \VParSegBldc{}-s
windows). (a) Median time per stage over \VChipLatBldcN{} windows. (b) Internal RAM handed to the
pool and its peak use (measured) against the estimated reference arena (hatched).}
\label{fig:esplat}
\end{figure}
\fi

\subsection{Classification under three protocols}

\input{../tables/latex/desempenho_compacto.tex}

\begin{figure}[t]
\centering
\includegraphics[width=\columnwidth]{\FigBalacc}
\caption{Balanced accuracy on the same windows with \VDataCiLevel{}\% group-bootstrap
intervals. Dotted line: chance. *Intentional leakage. On IDMT class = motor = recording; on A2212
Bearing exists only on motor M4 (confounded with motor identity), so motor-out Bearing recall is
zero by construction.}
\label{fig:balacc}
\end{figure}

Table~\ref{tab:desempenho-compacto} and
Fig.~\ref{fig:balacc} compare the methods. On IDMT every method is near the ceiling under both
protocols: the RMS-normalised level baseline, which only sees the crest factor, reaches
\VDiagBalRmsNormIdmtBlk{} balanced accuracy on held-out blocks, and the pipeline and the
classical features reach \VDiagBalDecIdmtBlk{} and \VDiagBalClsIdmtBlk{}. Because each class is one recording of one motor,
these numbers measure recording identity and cannot rank the methods.

On A2212 with a grouped split by file, the classical features reach \VDiagBalClsBldcFile{}
[\VDiagBalClsBldcFileLo{}, \VDiagBalClsBldcFileHi{}], the pipeline \VDiagBalDecBldcFile{}
[\VDiagBalDecBldcFileLo{}, \VDiagBalDecBldcFileHi{}] and the raw level baseline
\VDiagBalRmsRawBldcFile{}. The pipeline errs mostly on Healthy (recall
\VDiagRecDecBldcFileHealthy{}). The random split inflates balanced accuracy by
\VDiagInflDecBldc{} for the pipeline and \VDiagInflClsBldc{} for the classical features. Holding a
motor out lowers every method: the classical features fall to \VDiagBalClsBldcMotor{} and the
pipeline to \VDiagBalDecBldcMotor{}; Bearing recall is \VDiagRecDecBldcMotorBearing{} by
construction, since the only bearing motor is then absent from training. Without the M4 fold,
on Healthy versus Propeller of unseen motors, the classical features reach
\VDiagBalClsBldcNoMfour{}, the RMS-normalised level baseline \VDiagBalRmsNormBldcNoMfour{} and the
pipeline \VDiagBalDecBldcNoMfour{}. Compared like for like (Healthy and Propeller of motors
M1--M3 under the file split), holding the motor out costs \VDiagDropHPClsBldc{} for the classical
features and \VDiagDropHPDecBldc{} for the pipeline; the much larger drop of the full motor-out
protocol is dominated by the Bearing class losing its only training motor. The pipeline never
outperforms the classical features (it ties at the IDMT ceiling), and with a motor held out its
point estimate also falls below the level baseline, with overlapping intervals.

% matriz de confusão: figures/pdf (FigConf), fora do texto por espaço

\subsection{Interpretability: the tracked rotation}

\ifwip
On IDMT the median tracked rotation differs by class (\VRotIdmtGoodMed{},
\VRotIdmtBrokenMed{}, \VRotIdmtHeavyMed{}\,Hz), but class equals motor there. On A2212 its rank
correlation with the speed label of the files is \VRotBldcSpearman{} pooled over motors; within
motors it is \VRotBldcSpearmanMOne{} to \VRotBldcSpearmanMFour{} but significant only for M4
(\VRotBldcFilesMFour{} files), so it should not be read as shaft speed on that dataset.
\else
The tracker exposes a physical quantity that a technician can inspect. On IDMT, the median tracked rotation is \VRotIdmtGoodMed{},
\VRotIdmtBrokenMed{} and \VRotIdmtHeavyMed{}\,Hz for good, broken and heavy-load windows, with
inter-quartile ranges of at most \VRotIdmtHeavyIqr{}\,Hz; \VRotIdmtGoodOctPct{}\% of the good
windows lie above \VRotOctThr{} times the class median (all at \VRotOctRatioMin{}--\VRotOctRatioMax{}
times), an octave error. Because class equals motor on IDMT, these differences cannot be
attributed to the operating state. On A2212, pooled over the \VDataBldcRecordings{} files, the rank
correlation between the tracked rotation and the speed label in the file names is
\VRotBldcSpearman{}; within motors it is \VRotBldcSpearmanMOne{} to \VRotBldcSpearmanMFour{}
(\VRotBldcFilesMThree{}--\VRotBldcFilesMTwo{} files each), significant only for M4
($p \geq \audRotBldcSpearmanPMTwo{}$ for M1--M3). The absolute level
is not comparable across motors and \VRotBldcCeilPct{}\% of the windows lie within
\VRotCeilPctThr{}\% of the upper limit of the search range. The range and the grid step were chosen
on the IDMT training data, so on A2212 the tracker output should not be read as the shaft speed.
\fi

% espectro: figures/pdf (FigSpec), fora do texto por espaço

\ifwip\else
\subsection{Streaming decision}

On a stream of \VPcSeqWindows{} clean A2212 golden windows (\VPcSeqPerClass{} per class:
Healthy, Propeller, Bearing) the online HMM with change penalty \VPcSeqCp{} detected both changes
with a delay of \VPcSeqDelayC{} windows in C and \VPcSeqDelayPy{} in Python on the PC (maximum
probability difference \VPcSeqMaxProbDiff{}).
\ifchipmedido On the ESP32, the same stream read from flash by \texttt{\VChipVariant} gave a delay
of \VChipSeqDelay{} windows and online decisions equal to Python's on \VChipSeqSame{} of
\VChipSeqWindows{} windows; one HMM step takes \VChipSeqHmmUs{}\,$\mu$s.\fi{} Bearing is motor M4,
so this checks the implementation, not the detector.
\fi

% ======================================================================
\section{Discussion}

\emph{What the datasets allow us to conclude.} IDMT, used without its test set, supports no
statement about operating-state discrimination: one recording per class makes class, motor and
session the same variable, and a two-number level baseline is as good as any model. A2212 with a
file-level split supports statements about the state of motors that were \emph{already seen} in
training, and there a cheap classical feature set is better than our pipeline. Generalisation to a
new motor is the relevant question for a deployed monitor, and with \VDataBldcMotors{} motors, a
bearing fault present on one of them and speed ranges that differ by class~\cite{estacio2023}, the
data can only show that every method degrades; it cannot rank methods or estimate how a monitor
would behave on another motor. None of these experiments concerns failure prediction, and none
involves industrial motors. In our runs the random split inflated A2212 balanced accuracy only
moderately (\VDiagInflClsBldc{} to \VDiagInflDecBldc{}), and holding a motor out cost, like for
like on Healthy versus Propeller, a similar amount (\VDiagDropHPClsBldc{} for the classical features,
\VDiagDropHPDecBldc{} for the pipeline); the large drop of the full motor-out protocol is dominated
by the Bearing class, which exists only on M4. Published near-perfect accuracies on these datasets
should be read against that confound, which no split by segment or by file removes.

\ifwip\else
\emph{What the embedded work shows.} The C port reproduces the reference stage by stage within
tolerances that were frozen before the formal comparison, including the discrete Viterbi path and
the HMM decisions; the residual float32 failures occur only on a synthetic edge case whose
modulation peak lies at the float32 noise floor, they leave the decision unchanged, and they
disappear when the envelope filter state is kept in double.
\ifchipmedido
On a board without PSRAM, \texttt{\VChipVariant} runs \VParSegBldc{}-s windows with on-chip outputs
within those tolerances, but at a real-time factor of \VChipRtfBldc{}, and the \VParSegIdmt{}-s
window does not fit. We claim neither real-time operation nor low power consumption.
\else
Whether it runs in real time on the chip is still open: the PC timings do not transfer.
\fi
\fi

\emph{Limitations and future work.} Energy was not measured (no current sensor), and we do not
estimate it from datasheets. All audio came from the datasets through the serial link; a real
I2S microphone, a motor on a test bench with a tachometer to validate the rotation estimate, and
the official IDMT test set (unopened by design) are future work. An independent scientific audit
of datasets, equivalence and traceability (report in the repository) was carried out before the
on-chip measurements and does not cover them.

% ======================================================================
\section{Conclusion}

A music-transcription-inspired pipeline for acoustic motor-state classification was ported to
portable C and verified against its Python reference stage by stage, with tolerances frozen
before the formal comparison: all checks pass in double; in float32 every window of real audio
passes every stage, and only \VEqFloatBadCases{} synthetic edge windows fail the modulation,
feature and z-score checks, with an unchanged decision. On an ESP32 without PSRAM the reference
layout does not fit; the named variant \texttt{\VChipVariant} fits only \VParSegBldc{}-s windows,
matches the reference within the frozen tolerances on the chip, and is \VChipRtfBldc{} times slower
than real time, dominated by the table-free harmonic sum in emulated double precision; a float32
variant of that step or a board with PSRAM is the way forward. On the two public
datasets it does not beat simple baselines, and the datasets themselves cannot separate operating
state from motor identity. The value of the work lies in the verified, measurable implementation and
in making these limits explicit.

\bibliographystyle{IEEEtran}
\bibliography{references}

\end{document}
"
\ifdefined\usewip\wiptrue\else\wipfalse\fi

% Gerado por embarcados/src/python/gerar_artigo_numeros.py. NÃO EDITAR À MÃO.
% commit 9df3b2c280b69b02703d2c61d277e2bdf02cafad (árvore suja em embarcados/: True)
% Uma macro por número. Comentário: [categoria] fonte :: coluna.
% Prefixos: VPar código | VData/VDiag/VRot dados e validação | VEq equivalência (PC)
%           VEst 1·estimativa analítica | VPc 2·medição no PC | VBuild 3·build do firmware | VChip 4·medição no ESP32
\newif\ifchipmedido
\chipmedidofalse  % true só se results/manifests/bancada.json e esp32_lat_* existirem

\newcommand{\VParSr}{16{,}000}% [codigo] pipelines/_decifrae.py :: PreDecifrae.sr
\newcommand{\VParSrKhz}{16}% [codigo] pipelines/_decifrae.py :: PreDecifrae.sr/1000
\newcommand{\VParNfft}{8{,}192}% [codigo] pipelines/_decifrae.py :: PreDecifrae.n_fft
\newcommand{\VParHop}{1{,}024}% [codigo] pipelines/_decifrae.py :: PreDecifrae.hop
\newcommand{\VParFmin}{15}% [codigo] pipelines/_decifrae.py :: PreDecifrae.fmin
\newcommand{\VParFmax}{60}% [codigo] pipelines/_decifrae.py :: PreDecifrae.fmax
\newcommand{\VParStepCents}{2}% [codigo] pipelines/_decifrae.py :: PreDecifrae.passo_cents
\newcommand{\VParJ}{1{,}201}% [codigo] motorsound/rotacao.py :: len(rotacao.grade(fmin,fmax,passo))
\newcommand{\VParNbins}{4{,}097}% [codigo] motorsound/rotacao.py :: n_fft//2+1 (rfft)
\newcommand{\VParDf}{1.95}% [codigo] motorsound/rotacao.py :: sr/n_fft
\newcommand{\VParNharm}{48}% [codigo] pipelines/_decifrae.py :: N_HARM
\newcommand{\VParGroupOneLo}{1}% [codigo] pipelines/_decifrae.py :: GRUPOS[0][0]
\newcommand{\VParGroupOneHi}{4}% [codigo] pipelines/_decifrae.py :: GRUPOS[0][1]
\newcommand{\VParGroupTwoLo}{5}% [codigo] pipelines/_decifrae.py :: GRUPOS[1][0]
\newcommand{\VParGroupTwoHi}{12}% [codigo] pipelines/_decifrae.py :: GRUPOS[1][1]
\newcommand{\VParGroupThreeLo}{13}% [codigo] pipelines/_decifrae.py :: GRUPOS[2][0]
\newcommand{\VParGroupThreeHi}{24}% [codigo] pipelines/_decifrae.py :: GRUPOS[2][1]
\newcommand{\VParGroupFourLo}{25}% [codigo] pipelines/_decifrae.py :: GRUPOS[3][0]
\newcommand{\VParGroupFourHi}{48}% [codigo] pipelines/_decifrae.py :: GRUPOS[3][1]
\newcommand{\VParMaxJump}{50}% [codigo] pipelines/_decifrae.py :: PreDecifrae.max_salto
\newcommand{\VParMaxJumpCents}{100}% [codigo] pipelines/_decifrae.py :: max_salto*passo_cents
\newcommand{\VParOctStates}{600}% [codigo] motorsound/rotacao.py :: viterbi_rotacao: k_oct = round(1200/passo)
\newcommand{\VParWStep}{0.55}% [codigo] motorsound/rotacao.py :: VIT_W_STEP
\newcommand{\VParWOct}{7.0}% [codigo] motorsound/rotacao.py :: VIT_W_OCTAVE
\newcommand{\VParWSil}{3.2}% [codigo] motorsound/rotacao.py :: VIT_W_SIL
\newcommand{\VParWEnergy}{0.45}% [codigo] motorsound/rotacao.py :: emissao_baixo: peso de e_norm
\newcommand{\VParEmisEps}{\ensuremath{10^{-3}}}% [codigo] motorsound/rotacao.py :: emissao_baixo: eps do -log
\newcommand{\VParTempCos}{0.1}% [codigo] pipelines/_decifrae.py :: PostDecifrae.temp_cos
\newcommand{\VParButterOrder}{4}% [codigo] pipelines/_decifrae.py :: _envelope_mod: butter(ordem, fc)
\newcommand{\VParButterFc}{200}% [codigo] pipelines/_decifrae.py :: _envelope_mod: butter(ordem, fc)
\newcommand{\VParModLo}{2}% [codigo] pipelines/_decifrae.py :: _envelope_mod: banda
\newcommand{\VParModHi}{100}% [codigo] pipelines/_decifrae.py :: _envelope_mod: banda
\newcommand{\VParEnvRateHz}{1{,}000}% [codigo] pipelines/_decifrae.py :: _envelope_mod: [:: sr//1000]
\newcommand{\VParDecim}{16}% [codigo] pipelines/_decifrae.py :: sr // 1000
\newcommand{\VParBandLo}{10}% [codigo] pipelines/_decifrae.py :: extrair: banda das features
\newcommand{\VParBandHi}{4{,}000}% [codigo] pipelines/_decifrae.py :: extrair: banda das features
\newcommand{\VParSegIdmt}{3}% [codigo] motorsound/config.py :: AVAL_IDMT.seg_s
\newcommand{\VParSegBldc}{1}% [codigo] motorsound/config.py :: AVAL_A2212.seg_s
\newcommand{\VParRmsTarget}{0.1}% [codigo] motorsound/config.py :: AvaliacaoConfig.rms_alvo
\newcommand{\VParFramesIdmt}{39}% [codigo] motorsound/rotacao.py :: 1+(N-n_fft)//hop (stft_mag)
\newcommand{\VParFramesBldc}{8}% [codigo] motorsound/rotacao.py :: 1+(N-n_fft)//hop (stft_mag)
\newcommand{\VParNClassicFeat}{32}% [codigo] pipelines/_classicas.py :: N_FEATURES
\newcommand{\VParNChangePen}{6}% [codigo] embarcados/gold/seqs/idmt_seq00/params.json :: change_penalties
\newcommand{\VParCpMin}{0.5}% [codigo] embarcados/gold/seqs/idmt_seq00/params.json :: change_penalties
\newcommand{\VParCpMax}{16}% [codigo] embarcados/gold/seqs/idmt_seq00/params.json :: change_penalties
\newcommand{\VParNLogprob}{-50}% [codigo] embarcados/gold/seqs/idmt_seq00/params.json :: n_logprob
\newcommand{\VParHmmTemp}{1}% [codigo] embarcados/gold/seqs/idmt_seq00/params.json :: temp
\newcommand{\VDataIdmtWindows}{357}% [dados] embarcados/results/predictions/diagnostico_janelas_20261023_222831.csv :: len(unique(rec_id,idx)) da p11
\newcommand{\VDataIdmtRecordings}{3}% [dados] embarcados/results/predictions/diagnostico_janelas_20261023_222831.csv :: nunique(rec_id)
\newcommand{\VDataIdmtMotors}{3}% [dados] embarcados/results/predictions/diagnostico_janelas_20261023_222831.csv :: nunique(motor)
\newcommand{\VDataIdmtClasses}{3}% [dados] embarcados/results/predictions/diagnostico_janelas_20261023_222831.csv :: nunique(classe)
\newcommand{\VDataBldcWindows}{430}% [dados] embarcados/results/predictions/diagnostico_janelas_20261023_222831.csv :: len(unique(rec_id,idx)) da p11
\newcommand{\VDataBldcRecordings}{43}% [dados] embarcados/results/predictions/diagnostico_janelas_20261023_222831.csv :: nunique(rec_id)
\newcommand{\VDataBldcMotors}{4}% [dados] embarcados/results/predictions/diagnostico_janelas_20261023_222831.csv :: nunique(motor)
\newcommand{\VDataBldcClasses}{3}% [dados] embarcados/results/predictions/diagnostico_janelas_20261023_222831.csv :: nunique(classe)
\newcommand{\VDataBldcFilesHealthy}{22}% [dados] embarcados/results/predictions/diagnostico_janelas_20261023_222831.csv :: nunique(rec_id) por classe
\newcommand{\VDataBldcFilesPropeller}{14}% [dados] embarcados/results/predictions/diagnostico_janelas_20261023_222831.csv :: nunique(rec_id) por classe
\newcommand{\VDataBldcFilesBearing}{7}% [dados] embarcados/results/predictions/diagnostico_janelas_20261023_222831.csv :: nunique(rec_id) por classe
\newcommand{\VDataIdmtRandGroups}{18}% [dados] embarcados/results/csv/diagnostico_20261023_222831.csv :: n_grupos
\newcommand{\VDataIdmtBlkGroups}{15}% [dados] embarcados/results/csv/diagnostico_20261023_222831.csv :: n_grupos
\newcommand{\VDataBldcRandGroups}{43}% [dados] embarcados/results/csv/diagnostico_20261023_222831.csv :: n_grupos
\newcommand{\VDataBldcFileGroups}{43}% [dados] embarcados/results/csv/diagnostico_20261023_222831.csv :: n_grupos
\newcommand{\VDataBldcMotorGroups}{43}% [dados] embarcados/results/csv/diagnostico_20261023_222831.csv :: n_grupos
\newcommand{\VDataBldcNoMfourGroups}{36}% [dados] embarcados/results/csv/diagnostico_20261023_222831.csv :: n_grupos
\newcommand{\VDataNBoot}{2{,}000}% [dados] embarcados/results/csv/diagnostico_20261023_222831_meta.json :: n_boot
\newcommand{\VDataCiLevel}{95}% [dados] embarcados/results/csv/diagnostico_20261023_222831.csv :: nível do IC no nome da coluna ic<nível>_lo
\newcommand{\VDataSeedBoot}{20261023}% [dados] embarcados/results/csv/diagnostico_20261023_222831_meta.json :: seed_boot
\newcommand{\VDiagBalRmsNormIdmtRand}{0.992}% [diagnostico] embarcados/results/csv/diagnostico_20261023_222831.csv :: valor[b00_rmspico_norm,idmt,aleatorio,acuracia_balanceada,_macro]
\newcommand{\VDiagBalRmsNormIdmtRandLo}{0.983}% [diagnostico] embarcados/results/csv/diagnostico_20261023_222831.csv :: ic95_lo[b00_rmspico_norm,idmt,aleatorio,acuracia_balanceada,_macro]
\newcommand{\VDiagBalRmsNormIdmtRandHi}{1.000}% [diagnostico] embarcados/results/csv/diagnostico_20261023_222831.csv :: ic95_hi[b00_rmspico_norm,idmt,aleatorio,acuracia_balanceada,_macro]
\newcommand{\VDiagAucRmsNormIdmtRand}{0.998}% [diagnostico] embarcados/results/csv/diagnostico_20261023_222831.csv :: valor[b00_rmspico_norm,idmt,aleatorio,roc_auc_ovr_macro,_macro]
\newcommand{\VDiagFoneRmsNormIdmtRand}{0.992}% [diagnostico] embarcados/results/csv/diagnostico_20261023_222831.csv :: valor[b00_rmspico_norm,idmt,aleatorio,f1_macro,_macro]
\newcommand{\VDiagBalRmsRawIdmtRand}{0.997}% [diagnostico] embarcados/results/csv/diagnostico_20261023_222831.csv :: valor[b00_rmspico_bruto,idmt,aleatorio,acuracia_balanceada,_macro]
\newcommand{\VDiagBalRmsRawIdmtRandLo}{0.990}% [diagnostico] embarcados/results/csv/diagnostico_20261023_222831.csv :: ic95_lo[b00_rmspico_bruto,idmt,aleatorio,acuracia_balanceada,_macro]
\newcommand{\VDiagBalRmsRawIdmtRandHi}{1.000}% [diagnostico] embarcados/results/csv/diagnostico_20261023_222831.csv :: ic95_hi[b00_rmspico_bruto,idmt,aleatorio,acuracia_balanceada,_macro]
\newcommand{\VDiagAucRmsRawIdmtRand}{1.000}% [diagnostico] embarcados/results/csv/diagnostico_20261023_222831.csv :: valor[b00_rmspico_bruto,idmt,aleatorio,roc_auc_ovr_macro,_macro]
\newcommand{\VDiagFoneRmsRawIdmtRand}{0.997}% [diagnostico] embarcados/results/csv/diagnostico_20261023_222831.csv :: valor[b00_rmspico_bruto,idmt,aleatorio,f1_macro,_macro]
\newcommand{\VDiagBalClsIdmtRand}{1.000}% [diagnostico] embarcados/results/csv/diagnostico_20261023_222831.csv :: valor[p04_classicas,idmt,aleatorio,acuracia_balanceada,_macro]
\newcommand{\VDiagBalClsIdmtRandLo}{1.000}% [diagnostico] embarcados/results/csv/diagnostico_20261023_222831.csv :: ic95_lo[p04_classicas,idmt,aleatorio,acuracia_balanceada,_macro]
\newcommand{\VDiagBalClsIdmtRandHi}{1.000}% [diagnostico] embarcados/results/csv/diagnostico_20261023_222831.csv :: ic95_hi[p04_classicas,idmt,aleatorio,acuracia_balanceada,_macro]
\newcommand{\VDiagAucClsIdmtRand}{1.000}% [diagnostico] embarcados/results/csv/diagnostico_20261023_222831.csv :: valor[p04_classicas,idmt,aleatorio,roc_auc_ovr_macro,_macro]
\newcommand{\VDiagFoneClsIdmtRand}{1.000}% [diagnostico] embarcados/results/csv/diagnostico_20261023_222831.csv :: valor[p04_classicas,idmt,aleatorio,f1_macro,_macro]
\newcommand{\VDiagBalDecIdmtRand}{1.000}% [diagnostico] embarcados/results/csv/diagnostico_20261023_222831.csv :: valor[p11_decifrae_harmsum,idmt,aleatorio,acuracia_balanceada,_macro]
\newcommand{\VDiagBalDecIdmtRandLo}{1.000}% [diagnostico] embarcados/results/csv/diagnostico_20261023_222831.csv :: ic95_lo[p11_decifrae_harmsum,idmt,aleatorio,acuracia_balanceada,_macro]
\newcommand{\VDiagBalDecIdmtRandHi}{1.000}% [diagnostico] embarcados/results/csv/diagnostico_20261023_222831.csv :: ic95_hi[p11_decifrae_harmsum,idmt,aleatorio,acuracia_balanceada,_macro]
\newcommand{\VDiagAucDecIdmtRand}{1.000}% [diagnostico] embarcados/results/csv/diagnostico_20261023_222831.csv :: valor[p11_decifrae_harmsum,idmt,aleatorio,roc_auc_ovr_macro,_macro]
\newcommand{\VDiagFoneDecIdmtRand}{1.000}% [diagnostico] embarcados/results/csv/diagnostico_20261023_222831.csv :: valor[p11_decifrae_harmsum,idmt,aleatorio,f1_macro,_macro]
\newcommand{\VDiagBalRmsNormIdmtBlk}{0.991}% [diagnostico] embarcados/results/csv/diagnostico_20261023_222831.csv :: valor[b00_rmspico_norm,idmt,blocos,acuracia_balanceada,_macro]
\newcommand{\VDiagBalRmsNormIdmtBlkLo}{0.982}% [diagnostico] embarcados/results/csv/diagnostico_20261023_222831.csv :: ic95_lo[b00_rmspico_norm,idmt,blocos,acuracia_balanceada,_macro]
\newcommand{\VDiagBalRmsNormIdmtBlkHi}{0.997}% [diagnostico] embarcados/results/csv/diagnostico_20261023_222831.csv :: ic95_hi[b00_rmspico_norm,idmt,blocos,acuracia_balanceada,_macro]
\newcommand{\VDiagAucRmsNormIdmtBlk}{0.998}% [diagnostico] embarcados/results/csv/diagnostico_20261023_222831.csv :: valor[b00_rmspico_norm,idmt,blocos,roc_auc_ovr_macro,_macro]
\newcommand{\VDiagFoneRmsNormIdmtBlk}{0.991}% [diagnostico] embarcados/results/csv/diagnostico_20261023_222831.csv :: valor[b00_rmspico_norm,idmt,blocos,f1_macro,_macro]
\newcommand{\VDiagBalRmsRawIdmtBlk}{0.993}% [diagnostico] embarcados/results/csv/diagnostico_20261023_222831.csv :: valor[b00_rmspico_bruto,idmt,blocos,acuracia_balanceada,_macro]
\newcommand{\VDiagBalRmsRawIdmtBlkLo}{0.987}% [diagnostico] embarcados/results/csv/diagnostico_20261023_222831.csv :: ic95_lo[b00_rmspico_bruto,idmt,blocos,acuracia_balanceada,_macro]
\newcommand{\VDiagBalRmsRawIdmtBlkHi}{1.000}% [diagnostico] embarcados/results/csv/diagnostico_20261023_222831.csv :: ic95_hi[b00_rmspico_bruto,idmt,blocos,acuracia_balanceada,_macro]
\newcommand{\VDiagAucRmsRawIdmtBlk}{1.000}% [diagnostico] embarcados/results/csv/diagnostico_20261023_222831.csv :: valor[b00_rmspico_bruto,idmt,blocos,roc_auc_ovr_macro,_macro]
\newcommand{\VDiagFoneRmsRawIdmtBlk}{0.994}% [diagnostico] embarcados/results/csv/diagnostico_20261023_222831.csv :: valor[b00_rmspico_bruto,idmt,blocos,f1_macro,_macro]
\newcommand{\VDiagBalClsIdmtBlk}{1.000}% [diagnostico] embarcados/results/csv/diagnostico_20261023_222831.csv :: valor[p04_classicas,idmt,blocos,acuracia_balanceada,_macro]
\newcommand{\VDiagBalClsIdmtBlkLo}{1.000}% [diagnostico] embarcados/results/csv/diagnostico_20261023_222831.csv :: ic95_lo[p04_classicas,idmt,blocos,acuracia_balanceada,_macro]
\newcommand{\VDiagBalClsIdmtBlkHi}{1.000}% [diagnostico] embarcados/results/csv/diagnostico_20261023_222831.csv :: ic95_hi[p04_classicas,idmt,blocos,acuracia_balanceada,_macro]
\newcommand{\VDiagAucClsIdmtBlk}{1.000}% [diagnostico] embarcados/results/csv/diagnostico_20261023_222831.csv :: valor[p04_classicas,idmt,blocos,roc_auc_ovr_macro,_macro]
\newcommand{\VDiagFoneClsIdmtBlk}{1.000}% [diagnostico] embarcados/results/csv/diagnostico_20261023_222831.csv :: valor[p04_classicas,idmt,blocos,f1_macro,_macro]
\newcommand{\VDiagBalDecIdmtBlk}{1.000}% [diagnostico] embarcados/results/csv/diagnostico_20261023_222831.csv :: valor[p11_decifrae_harmsum,idmt,blocos,acuracia_balanceada,_macro]
\newcommand{\VDiagBalDecIdmtBlkLo}{1.000}% [diagnostico] embarcados/results/csv/diagnostico_20261023_222831.csv :: ic95_lo[p11_decifrae_harmsum,idmt,blocos,acuracia_balanceada,_macro]
\newcommand{\VDiagBalDecIdmtBlkHi}{1.000}% [diagnostico] embarcados/results/csv/diagnostico_20261023_222831.csv :: ic95_hi[p11_decifrae_harmsum,idmt,blocos,acuracia_balanceada,_macro]
\newcommand{\VDiagAucDecIdmtBlk}{1.000}% [diagnostico] embarcados/results/csv/diagnostico_20261023_222831.csv :: valor[p11_decifrae_harmsum,idmt,blocos,roc_auc_ovr_macro,_macro]
\newcommand{\VDiagFoneDecIdmtBlk}{1.000}% [diagnostico] embarcados/results/csv/diagnostico_20261023_222831.csv :: valor[p11_decifrae_harmsum,idmt,blocos,f1_macro,_macro]
\newcommand{\VDiagBalRmsNormBldcRand}{0.564}% [diagnostico] embarcados/results/csv/diagnostico_20261023_222831.csv :: valor[b00_rmspico_norm,a2212,aleatorio,acuracia_balanceada,_macro]
\newcommand{\VDiagBalRmsNormBldcRandLo}{0.515}% [diagnostico] embarcados/results/csv/diagnostico_20261023_222831.csv :: ic95_lo[b00_rmspico_norm,a2212,aleatorio,acuracia_balanceada,_macro]
\newcommand{\VDiagBalRmsNormBldcRandHi}{0.607}% [diagnostico] embarcados/results/csv/diagnostico_20261023_222831.csv :: ic95_hi[b00_rmspico_norm,a2212,aleatorio,acuracia_balanceada,_macro]
\newcommand{\VDiagAucRmsNormBldcRand}{0.784}% [diagnostico] embarcados/results/csv/diagnostico_20261023_222831.csv :: valor[b00_rmspico_norm,a2212,aleatorio,roc_auc_ovr_macro,_macro]
\newcommand{\VDiagFoneRmsNormBldcRand}{0.533}% [diagnostico] embarcados/results/csv/diagnostico_20261023_222831.csv :: valor[b00_rmspico_norm,a2212,aleatorio,f1_macro,_macro]
\newcommand{\VDiagBalRmsRawBldcRand}{0.625}% [diagnostico] embarcados/results/csv/diagnostico_20261023_222831.csv :: valor[b00_rmspico_bruto,a2212,aleatorio,acuracia_balanceada,_macro]
\newcommand{\VDiagBalRmsRawBldcRandLo}{0.571}% [diagnostico] embarcados/results/csv/diagnostico_20261023_222831.csv :: ic95_lo[b00_rmspico_bruto,a2212,aleatorio,acuracia_balanceada,_macro]
\newcommand{\VDiagBalRmsRawBldcRandHi}{0.660}% [diagnostico] embarcados/results/csv/diagnostico_20261023_222831.csv :: ic95_hi[b00_rmspico_bruto,a2212,aleatorio,acuracia_balanceada,_macro]
\newcommand{\VDiagAucRmsRawBldcRand}{0.786}% [diagnostico] embarcados/results/csv/diagnostico_20261023_222831.csv :: valor[b00_rmspico_bruto,a2212,aleatorio,roc_auc_ovr_macro,_macro]
\newcommand{\VDiagFoneRmsRawBldcRand}{0.588}% [diagnostico] embarcados/results/csv/diagnostico_20261023_222831.csv :: valor[b00_rmspico_bruto,a2212,aleatorio,f1_macro,_macro]
\newcommand{\VDiagBalClsBldcRand}{0.998}% [diagnostico] embarcados/results/csv/diagnostico_20261023_222831.csv :: valor[p04_classicas,a2212,aleatorio,acuracia_balanceada,_macro]
\newcommand{\VDiagBalClsBldcRandLo}{0.995}% [diagnostico] embarcados/results/csv/diagnostico_20261023_222831.csv :: ic95_lo[p04_classicas,a2212,aleatorio,acuracia_balanceada,_macro]
\newcommand{\VDiagBalClsBldcRandHi}{1.000}% [diagnostico] embarcados/results/csv/diagnostico_20261023_222831.csv :: ic95_hi[p04_classicas,a2212,aleatorio,acuracia_balanceada,_macro]
\newcommand{\VDiagAucClsBldcRand}{1.000}% [diagnostico] embarcados/results/csv/diagnostico_20261023_222831.csv :: valor[p04_classicas,a2212,aleatorio,roc_auc_ovr_macro,_macro]
\newcommand{\VDiagFoneClsBldcRand}{0.998}% [diagnostico] embarcados/results/csv/diagnostico_20261023_222831.csv :: valor[p04_classicas,a2212,aleatorio,f1_macro,_macro]
\newcommand{\VDiagBalDecBldcRand}{0.849}% [diagnostico] embarcados/results/csv/diagnostico_20261023_222831.csv :: valor[p11_decifrae_harmsum,a2212,aleatorio,acuracia_balanceada,_macro]
\newcommand{\VDiagBalDecBldcRandLo}{0.790}% [diagnostico] embarcados/results/csv/diagnostico_20261023_222831.csv :: ic95_lo[p11_decifrae_harmsum,a2212,aleatorio,acuracia_balanceada,_macro]
\newcommand{\VDiagBalDecBldcRandHi}{0.905}% [diagnostico] embarcados/results/csv/diagnostico_20261023_222831.csv :: ic95_hi[p11_decifrae_harmsum,a2212,aleatorio,acuracia_balanceada,_macro]
\newcommand{\VDiagAucDecBldcRand}{0.955}% [diagnostico] embarcados/results/csv/diagnostico_20261023_222831.csv :: valor[p11_decifrae_harmsum,a2212,aleatorio,roc_auc_ovr_macro,_macro]
\newcommand{\VDiagFoneDecBldcRand}{0.789}% [diagnostico] embarcados/results/csv/diagnostico_20261023_222831.csv :: valor[p11_decifrae_harmsum,a2212,aleatorio,f1_macro,_macro]
\newcommand{\VDiagBalRmsNormBldcFile}{0.562}% [diagnostico] embarcados/results/csv/diagnostico_20261023_222831.csv :: valor[b00_rmspico_norm,a2212,arquivo,acuracia_balanceada,_macro]
\newcommand{\VDiagBalRmsNormBldcFileLo}{0.512}% [diagnostico] embarcados/results/csv/diagnostico_20261023_222831.csv :: ic95_lo[b00_rmspico_norm,a2212,arquivo,acuracia_balanceada,_macro]
\newcommand{\VDiagBalRmsNormBldcFileHi}{0.606}% [diagnostico] embarcados/results/csv/diagnostico_20261023_222831.csv :: ic95_hi[b00_rmspico_norm,a2212,arquivo,acuracia_balanceada,_macro]
\newcommand{\VDiagAucRmsNormBldcFile}{0.759}% [diagnostico] embarcados/results/csv/diagnostico_20261023_222831.csv :: valor[b00_rmspico_norm,a2212,arquivo,roc_auc_ovr_macro,_macro]
\newcommand{\VDiagFoneRmsNormBldcFile}{0.532}% [diagnostico] embarcados/results/csv/diagnostico_20261023_222831.csv :: valor[b00_rmspico_norm,a2212,arquivo,f1_macro,_macro]
\newcommand{\VDiagBalRmsRawBldcFile}{0.608}% [diagnostico] embarcados/results/csv/diagnostico_20261023_222831.csv :: valor[b00_rmspico_bruto,a2212,arquivo,acuracia_balanceada,_macro]
\newcommand{\VDiagBalRmsRawBldcFileLo}{0.544}% [diagnostico] embarcados/results/csv/diagnostico_20261023_222831.csv :: ic95_lo[b00_rmspico_bruto,a2212,arquivo,acuracia_balanceada,_macro]
\newcommand{\VDiagBalRmsRawBldcFileHi}{0.656}% [diagnostico] embarcados/results/csv/diagnostico_20261023_222831.csv :: ic95_hi[b00_rmspico_bruto,a2212,arquivo,acuracia_balanceada,_macro]
\newcommand{\VDiagAucRmsRawBldcFile}{0.706}% [diagnostico] embarcados/results/csv/diagnostico_20261023_222831.csv :: valor[b00_rmspico_bruto,a2212,arquivo,roc_auc_ovr_macro,_macro]
\newcommand{\VDiagFoneRmsRawBldcFile}{0.576}% [diagnostico] embarcados/results/csv/diagnostico_20261023_222831.csv :: valor[b00_rmspico_bruto,a2212,arquivo,f1_macro,_macro]
\newcommand{\VDiagBalClsBldcFile}{0.981}% [diagnostico] embarcados/results/csv/diagnostico_20261023_222831.csv :: valor[p04_classicas,a2212,arquivo,acuracia_balanceada,_macro]
\newcommand{\VDiagBalClsBldcFileLo}{0.950}% [diagnostico] embarcados/results/csv/diagnostico_20261023_222831.csv :: ic95_lo[p04_classicas,a2212,arquivo,acuracia_balanceada,_macro]
\newcommand{\VDiagBalClsBldcFileHi}{0.998}% [diagnostico] embarcados/results/csv/diagnostico_20261023_222831.csv :: ic95_hi[p04_classicas,a2212,arquivo,acuracia_balanceada,_macro]
\newcommand{\VDiagAucClsBldcFile}{0.996}% [diagnostico] embarcados/results/csv/diagnostico_20261023_222831.csv :: valor[p04_classicas,a2212,arquivo,roc_auc_ovr_macro,_macro]
\newcommand{\VDiagFoneClsBldcFile}{0.982}% [diagnostico] embarcados/results/csv/diagnostico_20261023_222831.csv :: valor[p04_classicas,a2212,arquivo,f1_macro,_macro]
\newcommand{\VDiagBalDecBldcFile}{0.816}% [diagnostico] embarcados/results/csv/diagnostico_20261023_222831.csv :: valor[p11_decifrae_harmsum,a2212,arquivo,acuracia_balanceada,_macro]
\newcommand{\VDiagBalDecBldcFileLo}{0.747}% [diagnostico] embarcados/results/csv/diagnostico_20261023_222831.csv :: ic95_lo[p11_decifrae_harmsum,a2212,arquivo,acuracia_balanceada,_macro]
\newcommand{\VDiagBalDecBldcFileHi}{0.881}% [diagnostico] embarcados/results/csv/diagnostico_20261023_222831.csv :: ic95_hi[p11_decifrae_harmsum,a2212,arquivo,acuracia_balanceada,_macro]
\newcommand{\VDiagAucDecBldcFile}{0.935}% [diagnostico] embarcados/results/csv/diagnostico_20261023_222831.csv :: valor[p11_decifrae_harmsum,a2212,arquivo,roc_auc_ovr_macro,_macro]
\newcommand{\VDiagFoneDecBldcFile}{0.753}% [diagnostico] embarcados/results/csv/diagnostico_20261023_222831.csv :: valor[p11_decifrae_harmsum,a2212,arquivo,f1_macro,_macro]
\newcommand{\VDiagBalRmsNormBldcMotor}{0.542}% [diagnostico] embarcados/results/csv/diagnostico_20261023_222831.csv :: valor[b00_rmspico_norm,a2212,motor,acuracia_balanceada,_macro]
\newcommand{\VDiagBalRmsNormBldcMotorLo}{0.485}% [diagnostico] embarcados/results/csv/diagnostico_20261023_222831.csv :: ic95_lo[b00_rmspico_norm,a2212,motor,acuracia_balanceada,_macro]
\newcommand{\VDiagBalRmsNormBldcMotorHi}{0.596}% [diagnostico] embarcados/results/csv/diagnostico_20261023_222831.csv :: ic95_hi[b00_rmspico_norm,a2212,motor,acuracia_balanceada,_macro]
\newcommand{\VDiagAucRmsNormBldcMotor}{0.489}% [diagnostico] embarcados/results/csv/diagnostico_20261023_222831.csv :: valor[b00_rmspico_norm,a2212,motor,roc_auc_ovr_macro,_macro]
\newcommand{\VDiagFoneRmsNormBldcMotor}{0.509}% [diagnostico] embarcados/results/csv/diagnostico_20261023_222831.csv :: valor[b00_rmspico_norm,a2212,motor,f1_macro,_macro]
\newcommand{\VDiagBalRmsRawBldcMotor}{0.503}% [diagnostico] embarcados/results/csv/diagnostico_20261023_222831.csv :: valor[b00_rmspico_bruto,a2212,motor,acuracia_balanceada,_macro]
\newcommand{\VDiagBalRmsRawBldcMotorLo}{0.412}% [diagnostico] embarcados/results/csv/diagnostico_20261023_222831.csv :: ic95_lo[b00_rmspico_bruto,a2212,motor,acuracia_balanceada,_macro]
\newcommand{\VDiagBalRmsRawBldcMotorHi}{0.580}% [diagnostico] embarcados/results/csv/diagnostico_20261023_222831.csv :: ic95_hi[b00_rmspico_bruto,a2212,motor,acuracia_balanceada,_macro]
\newcommand{\VDiagAucRmsRawBldcMotor}{0.508}% [diagnostico] embarcados/results/csv/diagnostico_20261023_222831.csv :: valor[b00_rmspico_bruto,a2212,motor,roc_auc_ovr_macro,_macro]
\newcommand{\VDiagFoneRmsRawBldcMotor}{0.509}% [diagnostico] embarcados/results/csv/diagnostico_20261023_222831.csv :: valor[b00_rmspico_bruto,a2212,motor,f1_macro,_macro]
\newcommand{\VDiagBalClsBldcMotor}{0.613}% [diagnostico] embarcados/results/csv/diagnostico_20261023_222831.csv :: valor[p04_classicas,a2212,motor,acuracia_balanceada,_macro]
\newcommand{\VDiagBalClsBldcMotorLo}{0.579}% [diagnostico] embarcados/results/csv/diagnostico_20261023_222831.csv :: ic95_lo[p04_classicas,a2212,motor,acuracia_balanceada,_macro]
\newcommand{\VDiagBalClsBldcMotorHi}{0.641}% [diagnostico] embarcados/results/csv/diagnostico_20261023_222831.csv :: ic95_hi[p04_classicas,a2212,motor,acuracia_balanceada,_macro]
\newcommand{\VDiagAucClsBldcMotor}{0.618}% [diagnostico] embarcados/results/csv/diagnostico_20261023_222831.csv :: valor[p04_classicas,a2212,motor,roc_auc_ovr_macro,_macro]
\newcommand{\VDiagFoneClsBldcMotor}{0.563}% [diagnostico] embarcados/results/csv/diagnostico_20261023_222831.csv :: valor[p04_classicas,a2212,motor,f1_macro,_macro]
\newcommand{\VDiagBalDecBldcMotor}{0.456}% [diagnostico] embarcados/results/csv/diagnostico_20261023_222831.csv :: valor[p11_decifrae_harmsum,a2212,motor,acuracia_balanceada,_macro]
\newcommand{\VDiagBalDecBldcMotorLo}{0.397}% [diagnostico] embarcados/results/csv/diagnostico_20261023_222831.csv :: ic95_lo[p11_decifrae_harmsum,a2212,motor,acuracia_balanceada,_macro]
\newcommand{\VDiagBalDecBldcMotorHi}{0.519}% [diagnostico] embarcados/results/csv/diagnostico_20261023_222831.csv :: ic95_hi[p11_decifrae_harmsum,a2212,motor,acuracia_balanceada,_macro]
\newcommand{\VDiagAucDecBldcMotor}{0.484}% [diagnostico] embarcados/results/csv/diagnostico_20261023_222831.csv :: valor[p11_decifrae_harmsum,a2212,motor,roc_auc_ovr_macro,_macro]
\newcommand{\VDiagFoneDecBldcMotor}{0.418}% [diagnostico] embarcados/results/csv/diagnostico_20261023_222831.csv :: valor[p11_decifrae_harmsum,a2212,motor,f1_macro,_macro]
\newcommand{\VDiagBalRmsNormBldcNoMfour}{0.814}% [diagnostico] embarcados/results/csv/diagnostico_20261023_222831.csv :: valor[b00_rmspico_norm,a2212,motor_sem_M4,acuracia_balanceada,_macro]
\newcommand{\VDiagBalRmsNormBldcNoMfourLo}{0.720}% [diagnostico] embarcados/results/csv/diagnostico_20261023_222831.csv :: ic95_lo[b00_rmspico_norm,a2212,motor_sem_M4,acuracia_balanceada,_macro]
\newcommand{\VDiagBalRmsNormBldcNoMfourHi}{0.887}% [diagnostico] embarcados/results/csv/diagnostico_20261023_222831.csv :: ic95_hi[b00_rmspico_norm,a2212,motor_sem_M4,acuracia_balanceada,_macro]
\newcommand{\VDiagAucRmsNormBldcNoMfour}{0.871}% [diagnostico] embarcados/results/csv/diagnostico_20261023_222831.csv :: valor[b00_rmspico_norm,a2212,motor_sem_M4,roc_auc_ovr_macro,_macro]
\newcommand{\VDiagFoneRmsNormBldcNoMfour}{0.824}% [diagnostico] embarcados/results/csv/diagnostico_20261023_222831.csv :: valor[b00_rmspico_norm,a2212,motor_sem_M4,f1_macro,_macro]
\newcommand{\VDiagBalRmsRawBldcNoMfour}{0.754}% [diagnostico] embarcados/results/csv/diagnostico_20261023_222831.csv :: valor[b00_rmspico_bruto,a2212,motor_sem_M4,acuracia_balanceada,_macro]
\newcommand{\VDiagBalRmsRawBldcNoMfourLo}{0.623}% [diagnostico] embarcados/results/csv/diagnostico_20261023_222831.csv :: ic95_lo[b00_rmspico_bruto,a2212,motor_sem_M4,acuracia_balanceada,_macro]
\newcommand{\VDiagBalRmsRawBldcNoMfourHi}{0.873}% [diagnostico] embarcados/results/csv/diagnostico_20261023_222831.csv :: ic95_hi[b00_rmspico_bruto,a2212,motor_sem_M4,acuracia_balanceada,_macro]
\newcommand{\VDiagAucRmsRawBldcNoMfour}{0.912}% [diagnostico] embarcados/results/csv/diagnostico_20261023_222831.csv :: valor[b00_rmspico_bruto,a2212,motor_sem_M4,roc_auc_ovr_macro,_macro]
\newcommand{\VDiagFoneRmsRawBldcNoMfour}{0.823}% [diagnostico] embarcados/results/csv/diagnostico_20261023_222831.csv :: valor[b00_rmspico_bruto,a2212,motor_sem_M4,f1_macro,_macro]
\newcommand{\VDiagBalClsBldcNoMfour}{0.920}% [diagnostico] embarcados/results/csv/diagnostico_20261023_222831.csv :: valor[p04_classicas,a2212,motor_sem_M4,acuracia_balanceada,_macro]
\newcommand{\VDiagBalClsBldcNoMfourLo}{0.868}% [diagnostico] embarcados/results/csv/diagnostico_20261023_222831.csv :: ic95_lo[p04_classicas,a2212,motor_sem_M4,acuracia_balanceada,_macro]
\newcommand{\VDiagBalClsBldcNoMfourHi}{0.962}% [diagnostico] embarcados/results/csv/diagnostico_20261023_222831.csv :: ic95_hi[p04_classicas,a2212,motor_sem_M4,acuracia_balanceada,_macro]
\newcommand{\VDiagAucClsBldcNoMfour}{0.984}% [diagnostico] embarcados/results/csv/diagnostico_20261023_222831.csv :: valor[p04_classicas,a2212,motor_sem_M4,roc_auc_ovr_macro,_macro]
\newcommand{\VDiagFoneClsBldcNoMfour}{0.923}% [diagnostico] embarcados/results/csv/diagnostico_20261023_222831.csv :: valor[p04_classicas,a2212,motor_sem_M4,f1_macro,_macro]
\newcommand{\VDiagBalDecBldcNoMfour}{0.685}% [diagnostico] embarcados/results/csv/diagnostico_20261023_222831.csv :: valor[p11_decifrae_harmsum,a2212,motor_sem_M4,acuracia_balanceada,_macro]
\newcommand{\VDiagBalDecBldcNoMfourLo}{0.593}% [diagnostico] embarcados/results/csv/diagnostico_20261023_222831.csv :: ic95_lo[p11_decifrae_harmsum,a2212,motor_sem_M4,acuracia_balanceada,_macro]
\newcommand{\VDiagBalDecBldcNoMfourHi}{0.774}% [diagnostico] embarcados/results/csv/diagnostico_20261023_222831.csv :: ic95_hi[p11_decifrae_harmsum,a2212,motor_sem_M4,acuracia_balanceada,_macro]
\newcommand{\VDiagAucDecBldcNoMfour}{0.822}% [diagnostico] embarcados/results/csv/diagnostico_20261023_222831.csv :: valor[p11_decifrae_harmsum,a2212,motor_sem_M4,roc_auc_ovr_macro,_macro]
\newcommand{\VDiagFoneDecBldcNoMfour}{0.685}% [diagnostico] embarcados/results/csv/diagnostico_20261023_222831.csv :: valor[p11_decifrae_harmsum,a2212,motor_sem_M4,f1_macro,_macro]
\newcommand{\VDiagRecDecBldcFileHealthy}{0.577}% [diagnostico] embarcados/results/csv/diagnostico_20261023_222831.csv :: valor[p11,a2212,arquivo,sensibilidade,Healthy]
\newcommand{\VDiagRecDecBldcFilePropeller}{0.929}% [diagnostico] embarcados/results/csv/diagnostico_20261023_222831.csv :: valor[p11,a2212,arquivo,sensibilidade,Propeller]
\newcommand{\VDiagRecDecBldcFileBearing}{0.943}% [diagnostico] embarcados/results/csv/diagnostico_20261023_222831.csv :: valor[p11,a2212,arquivo,sensibilidade,Bearing]
\newcommand{\VDiagRecDecBldcMotorHealthy}{0.441}% [diagnostico] embarcados/results/csv/diagnostico_20261023_222831.csv :: valor[p11,a2212,motor,sensibilidade,Healthy]
\newcommand{\VDiagRecDecBldcMotorPropeller}{0.929}% [diagnostico] embarcados/results/csv/diagnostico_20261023_222831.csv :: valor[p11,a2212,motor,sensibilidade,Propeller]
\newcommand{\VDiagRecDecBldcMotorBearing}{0.000}% [diagnostico] embarcados/results/csv/diagnostico_20261023_222831.csv :: valor[p11,a2212,motor,sensibilidade,Bearing]
\newcommand{\VDiagInflRmsNormBldc}{0.002}% [diagnostico] embarcados/results/csv/diagnostico_20261023_222831.csv :: Bal[aleatorio]-Bal[arquivo] (a2212)
\newcommand{\VDiagDropRmsNormBldc}{0.019}% [diagnostico] embarcados/results/csv/diagnostico_20261023_222831.csv :: Bal[arquivo]-Bal[motor] (a2212)
\newcommand{\VDiagInflRmsRawBldc}{0.016}% [diagnostico] embarcados/results/csv/diagnostico_20261023_222831.csv :: Bal[aleatorio]-Bal[arquivo] (a2212)
\newcommand{\VDiagDropRmsRawBldc}{0.106}% [diagnostico] embarcados/results/csv/diagnostico_20261023_222831.csv :: Bal[arquivo]-Bal[motor] (a2212)
\newcommand{\VDiagInflClsBldc}{0.017}% [diagnostico] embarcados/results/csv/diagnostico_20261023_222831.csv :: Bal[aleatorio]-Bal[arquivo] (a2212)
\newcommand{\VDiagDropClsBldc}{0.368}% [diagnostico] embarcados/results/csv/diagnostico_20261023_222831.csv :: Bal[arquivo]-Bal[motor] (a2212)
\newcommand{\VDiagInflDecBldc}{0.033}% [diagnostico] embarcados/results/csv/diagnostico_20261023_222831.csv :: Bal[aleatorio]-Bal[arquivo] (a2212)
\newcommand{\VDiagDropDecBldc}{0.360}% [diagnostico] embarcados/results/csv/diagnostico_20261023_222831.csv :: Bal[arquivo]-Bal[motor] (a2212)
\newcommand{\VParNFeat}{13}% [codigo] embarcados/results/csv/p11_features_20261023_222831.csv :: nº de colunas f0..f12
\newcommand{\VRotIdmtGoodMed}{25.76}% [dados] embarcados/results/csv/p11_features_20261023_222831.csv :: mediana de f0 (idmt, good)
\newcommand{\VRotIdmtGoodIqr}{0.18}% [dados] embarcados/results/csv/p11_features_20261023_222831.csv :: IQR de f0 (idmt, good)
\newcommand{\VRotIdmtBrokenMed}{27.99}% [dados] embarcados/results/csv/p11_features_20261023_222831.csv :: mediana de f0 (idmt, broken)
\newcommand{\VRotIdmtBrokenIqr}{0.03}% [dados] embarcados/results/csv/p11_features_20261023_222831.csv :: IQR de f0 (idmt, broken)
\newcommand{\VRotIdmtHeavyMed}{24.65}% [dados] embarcados/results/csv/p11_features_20261023_222831.csv :: mediana de f0 (idmt, heavyload)
\newcommand{\VRotIdmtHeavyIqr}{0.20}% [dados] embarcados/results/csv/p11_features_20261023_222831.csv :: IQR de f0 (idmt, heavyload)
\newcommand{\VRotOctThr}{1.5}% [dados] embarcados/results/csv/p11_features_20261023_222831.csv :: limiar: f0 > LIM_OIT·mediana da classe
\newcommand{\VRotCeilPctThr}{1}% [dados] embarcados/results/csv/p11_features_20261023_222831.csv :: limiar: f0 >= LIM_TETO·fmax
\newcommand{\VRotIdmtGoodOctPct}{8.6}% [dados] embarcados/results/csv/p11_features_20261023_222831.csv :: % janelas good com f0 > LIM_OIT·mediana da classe (salto de oitava)
\newcommand{\VRotBldcCeilPct}{3.7}% [dados] embarcados/results/csv/p11_features_20261023_222831.csv :: % janelas A2212 com f0 >= 0,99·fmax
\newcommand{\VRotBldcSpearman}{0.02}% [dados] embarcados/results/csv/p11_features_20261023_222831.csv + embarcados/results/predictions/diagnostico_janelas_20261023_222831.csv :: Spearman(velocidade nominal do arquivo, mediana de f0 do arquivo), 43 arquivos
\newcommand{\VRotBldcSpearmanMOne}{0.30}% [dados] embarcados/results/csv/p11_features_20261023_222831.csv + embarcados/results/predictions/diagnostico_janelas_20261023_222831.csv :: Spearman por motor M1 (14 arquivos)
\newcommand{\VRotBldcFilesMOne}{14}% [dados] embarcados/results/csv/p11_features_20261023_222831.csv + embarcados/results/predictions/diagnostico_janelas_20261023_222831.csv :: nº de arquivos do M1
\newcommand{\VRotBldcSpearmanMTwo}{0.30}% [dados] embarcados/results/csv/p11_features_20261023_222831.csv + embarcados/results/predictions/diagnostico_janelas_20261023_222831.csv :: Spearman por motor M2 (16 arquivos)
\newcommand{\VRotBldcFilesMTwo}{16}% [dados] embarcados/results/csv/p11_features_20261023_222831.csv + embarcados/results/predictions/diagnostico_janelas_20261023_222831.csv :: nº de arquivos do M2
\newcommand{\VRotBldcSpearmanMThree}{0.54}% [dados] embarcados/results/csv/p11_features_20261023_222831.csv + embarcados/results/predictions/diagnostico_janelas_20261023_222831.csv :: Spearman por motor M3 (6 arquivos)
\newcommand{\VRotBldcFilesMThree}{6}% [dados] embarcados/results/csv/p11_features_20261023_222831.csv + embarcados/results/predictions/diagnostico_janelas_20261023_222831.csv :: nº de arquivos do M3
\newcommand{\VRotBldcSpearmanMFour}{1.00}% [dados] embarcados/results/csv/p11_features_20261023_222831.csv + embarcados/results/predictions/diagnostico_janelas_20261023_222831.csv :: Spearman por motor M4 (7 arquivos)
\newcommand{\VRotBldcFilesMFour}{7}% [dados] embarcados/results/csv/p11_features_20261023_222831.csv + embarcados/results/predictions/diagnostico_janelas_20261023_222831.csv :: nº de arquivos do M4
\newcommand{\VEqStagesTotal}{26}% [equivalencia] embarcados/results/csv/equivalencia_resumo_por_build_20261024-091729.csv :: nº de etapas por build
\newcommand{\VEqStagesCase}{20}% [equivalencia] embarcados/results/csv/equivalencia_resumo_por_build_20261024-091729.csv :: etapas de janela
\newcommand{\VEqStagesSeq}{6}% [equivalencia] embarcados/results/csv/equivalencia_resumo_por_build_20261024-091729.csv :: etapas de sequência
\newcommand{\VEqCases}{164}% [equivalencia] embarcados/results/csv/equivalencia_resumo_por_build_20261024-091729.csv :: n[y_prep]
\newcommand{\VEqSeqCombos}{240}% [equivalencia] embarcados/results/csv/equivalencia_resumo_por_build_20261024-091729.csv :: n[absorb]
\newcommand{\VEqDoubleStagesPass}{26}% [equivalencia] embarcados/results/csv/equivalencia_resumo_por_build_20261024-091729.csv :: nº de etapas com veredito PASS (double)
\newcommand{\VEqDoubleMaxRatio}{0.55}% [equivalencia] embarcados/results/csv/equivalencia_resumo_por_build_20261024-091729.csv :: máx razao_tol_max nas etapas sem empate/FAIL (double)
\newcommand{\VEqDoubleDecisionPass}{164}% [equivalencia] embarcados/results/csv/equivalencia_resumo_por_build_20261024-091729.csv :: pass+pass_empate[decisao] (double)
\newcommand{\VEqFloatStagesPass}{25}% [equivalencia] embarcados/results/csv/equivalencia_resumo_por_build_20261024-091729.csv :: nº de etapas com veredito PASS (float)
\newcommand{\VEqFloatMaxRatio}{0.34}% [equivalencia] embarcados/results/csv/equivalencia_resumo_por_build_20261024-091729.csv :: máx razao_tol_max nas etapas sem empate/FAIL (float)
\newcommand{\VEqFloatDecisionPass}{164}% [equivalencia] embarcados/results/csv/equivalencia_resumo_por_build_20261024-091729.csv :: pass+pass_empate[decisao] (float)
\newcommand{\VEqSosdStagesPass}{26}% [equivalencia] embarcados/results/csv/equivalencia_resumo_por_build_20261024-091729.csv :: nº de etapas com veredito PASS (float-sosd)
\newcommand{\VEqSosdMaxRatio}{0.23}% [equivalencia] embarcados/results/csv/equivalencia_resumo_por_build_20261024-091729.csv :: máx razao_tol_max nas etapas sem empate/FAIL (float-sosd)
\newcommand{\VEqSosdDecisionPass}{164}% [equivalencia] embarcados/results/csv/equivalencia_resumo_por_build_20261024-091729.csv :: pass+pass_empate[decisao] (float-sosd)
\newcommand{\VEqFloatZFail}{2}% [equivalencia] embarcados/results/csv/equivalencia_resumo_por_build_20261024-091729.csv :: fail[float,z]
\newcommand{\VEqFloatZPass}{162}% [equivalencia] embarcados/results/csv/equivalencia_resumo_por_build_20261024-091729.csv :: pass[float,z]
\newcommand{\VEqFloatFailStages}{z}% [equivalencia] embarcados/results/csv/equivalencia_resumo_por_build_20261024-091729.csv :: etapas sem PASS (float)
\newcommand{\VEqFloatTieStages}{2}% [equivalencia] embarcados/results/csv/equivalencia_resumo_por_build_20261024-091729.csv :: nº de etapas com PASS_EMPATE (float)
\newcommand{\VEqFloatTieCases}{2}% [equivalencia] embarcados/results/csv/equivalencia_float_20261024-091122.csv :: ids com PASS_EMPATE em env_mod
\newcommand{\VEqFloatTieCaseIds}{borda\_dc\_ruido3s\_ma2212, borda\_dc\_ruido3s\_midmt}% [equivalencia] embarcados/results/csv/equivalencia_float_20261024-091122.csv :: id
\newcommand{\VEqFloatTieGap}{\ensuremath{9.4\times10^{-16}}}% [equivalencia] embarcados/results/csv/equivalencia_float_20261024-091122.csv :: detalhe[env_mod]: P_ref[ref]-P_ref[C]
\newcommand{\VEqFloatTieTol}{\ensuremath{6.3\times10^{-11}}}% [equivalencia] embarcados/results/csv/equivalencia_float_20261024-091122.csv :: detalhe[env_mod]: 2·tol
\newcommand{\VEqFloatFmodErr}{37.1}% [equivalencia] embarcados/results/csv/equivalencia_float_20261024-091122.csv :: err_max[features[9]:f_mod] (empate)
\newcommand{\VEqFloatZErrMax}{1.61}% [equivalencia] embarcados/results/csv/equivalencia_float_20261024-091122.csv :: err_max[z] (FAIL)
\newcommand{\VEqFloatZFailIsEdge}{1}% [equivalencia] embarcados/results/csv/equivalencia_float_20261024-091122.csv :: todos os FAIL de z são casos de borda (1=sim)
\newcommand{\VEqEnvDecFloatMed}{\ensuremath{1.2\times10^{-5}}}% [equivalencia] embarcados/results/csv/equivalencia_float_20261024-091122.csv :: mediana err_rel_max[env_dec]
\newcommand{\VEqEnvDecSosdMed}{\ensuremath{1.2\times10^{-7}}}% [equivalencia] embarcados/results/csv/equivalencia_float-sosd_20261024-091222.csv :: mediana err_rel_max[env_dec]
\newcommand{\VEqEnvDecDoubleMed}{\ensuremath{1.1\times10^{-7}}}% [equivalencia] embarcados/results/csv/equivalencia_double_20261024-091106.csv :: mediana err_rel_max[env_dec]
\newcommand{\VEqTolSha}{d74dae800ceb}% [equivalencia] embarcados/configs/embedded/tolerancias.json :: sha256 (12 primeiros)
\newcommand{\VEqTolShaRec}{d74dae800ceb}% [equivalencia] embarcados/results/csv/equivalencia_resumo_por_build_20261024-091729.csv :: sha256_tolerancias gravado no _meta da comparação
\newcommand{\VEqTolCommit}{447a7c2}% [equivalencia] embarcados/configs/embedded/tolerancias.json :: git log -1 do tolerancias.json
\newcommand{\VEqTolCommitTime}{2026-10-24 09:10}% [equivalencia] embarcados/configs/embedded/tolerancias.json :: data do commit do tolerancias.json
\newcommand{\VEqRunTime}{2026-10-24 09:11}% [equivalencia] embarcados/results/csv/equivalencia_resumo_por_build_20261024-091729.csv :: data da comparação formal (float)
\newcommand{\VEqGoldClean}{60}% [equivalencia] embarcados/gold/manifest.json :: casos tipo limpo
\newcommand{\VEqGoldNoisy}{90}% [equivalencia] embarcados/gold/manifest.json :: casos tipo ruidoso
\newcommand{\VEqGoldEdge}{14}% [equivalencia] embarcados/gold/manifest.json :: casos tipo borda
\newcommand{\VEqGoldSeqs}{20}% [equivalencia] embarcados/gold/manifest.json :: n_seqs
\newcommand{\VEqGoldPerts}{5}% [equivalencia] embarcados/gold/manifest.json :: nº de perturbações
\newcommand{\VEqMinDurValues}{2}% [equivalencia] embarcados/gold/seqs/idmt_seq00/params.json :: len(min_durs_s)
\newcommand{\VEqEmulDoubleRows}{6{,}852}% [equivalencia] embarcados/results/csv/equivalencia_emul_double_20261024-090737_meta.json :: n_linhas
\newcommand{\VEqEmulDoublePass}{6{,}852}% [equivalencia] embarcados/results/csv/equivalencia_emul_double_20261024-090737_meta.json :: contagem_status.PASS
\newcommand{\VEqEmulFloatRows}{6{,}852}% [equivalencia] embarcados/results/csv/equivalencia_emul_float_20261024-090743_meta.json :: n_linhas
\newcommand{\VEqEmulFloatPass}{6{,}852}% [equivalencia] embarcados/results/csv/equivalencia_emul_float_20261024-090743_meta.json :: contagem_status.PASS
\newcommand{\VPcCIdmtMed}{13.29}% [pc] embarcados/results/metrics/bench_pc_20261024-091911.json :: resumo.c_float_pc/idmt.total_ms.mediana
\newcommand{\VPcPyIdmtMed}{95.3}% [pc] embarcados/results/metrics/bench_pc_20261024-091911.json :: resumo.python_ref/idmt.total_ms.mediana
\newcommand{\VPcCIdmtPninetyfive}{15.00}% [pc] embarcados/results/metrics/bench_pc_20261024-091911.json :: resumo.c_float_pc/idmt.total_ms.p95
\newcommand{\VPcPyIdmtPninetyfive}{100.2}% [pc] embarcados/results/metrics/bench_pc_20261024-091911.json :: resumo.python_ref/idmt.total_ms.p95
\newcommand{\VPcCIdmtMax}{15.80}% [pc] embarcados/results/metrics/bench_pc_20261024-091911.json :: resumo.c_float_pc/idmt.total_ms.max
\newcommand{\VPcPyIdmtMax}{110.2}% [pc] embarcados/results/metrics/bench_pc_20261024-091911.json :: resumo.python_ref/idmt.total_ms.max
\newcommand{\VPcCIdmtRtf}{0.0044}% [pc] embarcados/results/metrics/bench_pc_20261024-091911.json :: fator_tempo_real_mediana
\newcommand{\VPcPyIdmtRtf}{0.032}% [pc] embarcados/results/metrics/bench_pc_20261024-091911.json :: fator_tempo_real_mediana
\newcommand{\VPcSpeedupIdmt}{7.2}% [pc] embarcados/results/metrics/bench_pc_20261024-091911.json :: mediana Python / mediana C
\newcommand{\VPcCIdmtTrackMed}{5.45}% [pc] embarcados/results/metrics/bench_pc_20261024-091911.json :: rastrear_quadros_ms.mediana
\newcommand{\VPcCIdmtTrackPct}{41}% [pc] embarcados/results/metrics/bench_pc_20261024-091911.json :: 100·rastrear_quadros_ms.mediana/total_ms.mediana
\newcommand{\VPcCIdmtVitMed}{5.35}% [pc] embarcados/results/metrics/bench_pc_20261024-091911.json :: viterbi_ms.mediana
\newcommand{\VPcCIdmtVitPct}{40}% [pc] embarcados/results/metrics/bench_pc_20261024-091911.json :: 100·viterbi_ms.mediana/total_ms.mediana
\newcommand{\VPcCIdmtEnvMed}{2.27}% [pc] embarcados/results/metrics/bench_pc_20261024-091911.json :: envelope_mod_ms.mediana
\newcommand{\VPcCIdmtEnvPct}{17}% [pc] embarcados/results/metrics/bench_pc_20261024-091911.json :: 100·envelope_mod_ms.mediana/total_ms.mediana
\newcommand{\VPcCIdmtFeatMed}{0.013}% [pc] embarcados/results/metrics/bench_pc_20261024-091911.json :: features_ms.mediana
\newcommand{\VPcCIdmtFeatPct}{0}% [pc] embarcados/results/metrics/bench_pc_20261024-091911.json :: 100·features_ms.mediana/total_ms.mediana
\newcommand{\VPcCIdmtCosMed}{0.002}% [pc] embarcados/results/metrics/bench_pc_20261024-091911.json :: cosseno_ms.mediana
\newcommand{\VPcCIdmtCosPct}{0}% [pc] embarcados/results/metrics/bench_pc_20261024-091911.json :: 100·cosseno_ms.mediana/total_ms.mediana
\newcommand{\VPcNIdmt}{200}% [pc] embarcados/results/metrics/bench_pc_20261024-091911.json :: total_ms.n
\newcommand{\VPcCBldcMed}{3.18}% [pc] embarcados/results/metrics/bench_pc_20261024-091911.json :: resumo.c_float_pc/a2212.total_ms.mediana
\newcommand{\VPcPyBldcMed}{22.1}% [pc] embarcados/results/metrics/bench_pc_20261024-091911.json :: resumo.python_ref/a2212.total_ms.mediana
\newcommand{\VPcCBldcPninetyfive}{3.84}% [pc] embarcados/results/metrics/bench_pc_20261024-091911.json :: resumo.c_float_pc/a2212.total_ms.p95
\newcommand{\VPcPyBldcPninetyfive}{23.1}% [pc] embarcados/results/metrics/bench_pc_20261024-091911.json :: resumo.python_ref/a2212.total_ms.p95
\newcommand{\VPcCBldcMax}{4.45}% [pc] embarcados/results/metrics/bench_pc_20261024-091911.json :: resumo.c_float_pc/a2212.total_ms.max
\newcommand{\VPcPyBldcMax}{25.3}% [pc] embarcados/results/metrics/bench_pc_20261024-091911.json :: resumo.python_ref/a2212.total_ms.max
\newcommand{\VPcCBldcRtf}{0.0032}% [pc] embarcados/results/metrics/bench_pc_20261024-091911.json :: fator_tempo_real_mediana
\newcommand{\VPcPyBldcRtf}{0.022}% [pc] embarcados/results/metrics/bench_pc_20261024-091911.json :: fator_tempo_real_mediana
\newcommand{\VPcSpeedupBldc}{6.9}% [pc] embarcados/results/metrics/bench_pc_20261024-091911.json :: mediana Python / mediana C
\newcommand{\VPcCBldcTrackMed}{1.24}% [pc] embarcados/results/metrics/bench_pc_20261024-091911.json :: rastrear_quadros_ms.mediana
\newcommand{\VPcCBldcTrackPct}{39}% [pc] embarcados/results/metrics/bench_pc_20261024-091911.json :: 100·rastrear_quadros_ms.mediana/total_ms.mediana
\newcommand{\VPcCBldcVitMed}{1.14}% [pc] embarcados/results/metrics/bench_pc_20261024-091911.json :: viterbi_ms.mediana
\newcommand{\VPcCBldcVitPct}{36}% [pc] embarcados/results/metrics/bench_pc_20261024-091911.json :: 100·viterbi_ms.mediana/total_ms.mediana
\newcommand{\VPcCBldcEnvMed}{0.67}% [pc] embarcados/results/metrics/bench_pc_20261024-091911.json :: envelope_mod_ms.mediana
\newcommand{\VPcCBldcEnvPct}{21}% [pc] embarcados/results/metrics/bench_pc_20261024-091911.json :: 100·envelope_mod_ms.mediana/total_ms.mediana
\newcommand{\VPcCBldcFeatMed}{0.013}% [pc] embarcados/results/metrics/bench_pc_20261024-091911.json :: features_ms.mediana
\newcommand{\VPcCBldcFeatPct}{0}% [pc] embarcados/results/metrics/bench_pc_20261024-091911.json :: 100·features_ms.mediana/total_ms.mediana
\newcommand{\VPcCBldcCosMed}{0.002}% [pc] embarcados/results/metrics/bench_pc_20261024-091911.json :: cosseno_ms.mediana
\newcommand{\VPcCBldcCosPct}{0}% [pc] embarcados/results/metrics/bench_pc_20261024-091911.json :: 100·cosseno_ms.mediana/total_ms.mediana
\newcommand{\VPcNBldc}{200}% [pc] embarcados/results/metrics/bench_pc_20261024-091911.json :: total_ms.n
\newcommand{\VPcWarmup}{5}% [pc] embarcados/results/metrics/bench_pc_20261024-091911.json :: meta.argv.aquecimento
\newcommand{\VPcDistinctIdmt}{75}% [pc] embarcados/results/csv/bench_pc_20261024-091911.csv :: n_janelas_distintas
\newcommand{\VPcDistinctBldc}{75}% [pc] embarcados/results/csv/bench_pc_20261024-091911.csv :: n_janelas_distintas
\newcommand{\VPcCompiler}{MSVC 19.44 x64 \slash{}O2 \slash{}fp:precise}% [pc] embarcados/results/metrics/bench_pc_20261024-091911.json :: meta.compilador
\newcommand{\VPcNcpu}{16}% [pc] embarcados/results/metrics/bench_pc_20261024-091911.json :: meta.maquina.n_cpu
\newcommand{\VPcOs}{Windows 11}% [pc] embarcados/results/metrics/bench_pc_20261024-091911.json :: meta.maquina.plataforma
\newcommand{\VPcSeqWindows}{24}% [pc] embarcados/results/csv/seq_pc_20261024-092103.csv :: nº de linhas (janelas)
\newcommand{\VPcSeqSwitches}{2}% [pc] embarcados/results/csv/seq_pc_20261024-092103_meta.json :: len(deteccao_online_c)
\newcommand{\VPcSeqDelayC}{0}% [pc] embarcados/results/csv/seq_pc_20261024-092103_meta.json :: max atraso_janelas (C)
\newcommand{\VPcSeqDelayPy}{0}% [pc] embarcados/results/csv/seq_pc_20261024-092103_meta.json :: max atraso_janelas (Python)
\newcommand{\VPcSeqSameDecisions}{1}% [pc] embarcados/results/csv/seq_pc_20261024-092103_meta.json :: decisoes_iguais_python_c (1=sim)
\newcommand{\VPcSeqMaxProbDiff}{\ensuremath{3.2\times10^{-6}}}% [pc] embarcados/results/csv/seq_pc_20261024-092103_meta.json :: max_abs_dif_probs
\newcommand{\VPcSeqCp}{2}% [pc] embarcados/results/csv/seq_pc_20261024-092103_meta.json :: cp
\newcommand{\VPcSeqPerClass}{8}% [pc] embarcados/results/csv/seq_pc_20261024-092103.csv :: janelas por classe
\newcommand{\VEstWsFloatOneMb}{0.96}% [estimativa] embarcados/paper/fontes/ws_float_16000.json :: total_com_tabela/1e6
\newcommand{\VEstWsFloatOneNoTabMb}{0.61}% [estimativa] embarcados/paper/fontes/ws_float_16000.json :: total_sem_tabela/1e6
\newcommand{\VEstWsFloatOnePhaseBPct}{44}% [estimativa] embarcados/paper/fontes/ws_float_16000.json :: faseB/(persistente+max(faseA,faseB))
\newcommand{\VEstWsFloatThreeMb}{1.86}% [estimativa] embarcados/paper/fontes/ws_float_48000.json :: total_com_tabela/1e6
\newcommand{\VEstWsFloatThreeNoTabMb}{1.52}% [estimativa] embarcados/paper/fontes/ws_float_48000.json :: total_sem_tabela/1e6
\newcommand{\VEstWsFloatThreePhaseBPct}{64}% [estimativa] embarcados/paper/fontes/ws_float_48000.json :: faseB/(persistente+max(faseA,faseB))
\newcommand{\VEstWsDoubleOneMb}{1.79}% [estimativa] embarcados/paper/fontes/ws_double_16000.json :: total_com_tabela/1e6
\newcommand{\VEstWsDoubleOneNoTabMb}{1.21}% [estimativa] embarcados/paper/fontes/ws_double_16000.json :: total_sem_tabela/1e6
\newcommand{\VEstWsDoubleOnePhaseBPct}{47}% [estimativa] embarcados/paper/fontes/ws_double_16000.json :: faseB/(persistente+max(faseA,faseB))
\newcommand{\VEstWsDoubleThreeMb}{3.60}% [estimativa] embarcados/paper/fontes/ws_double_48000.json :: total_com_tabela/1e6
\newcommand{\VEstWsDoubleThreeNoTabMb}{3.02}% [estimativa] embarcados/paper/fontes/ws_double_48000.json :: total_sem_tabela/1e6
\newcommand{\VEstWsDoubleThreePhaseBPct}{67}% [estimativa] embarcados/paper/fontes/ws_double_48000.json :: faseB/(persistente+max(faseA,faseB))
\newcommand{\VEstHarmTableMb}{0.35}% [estimativa] embarcados/paper/fontes/ws_float_16000.json :: total_com_tabela - total_sem_tabela
\newcommand{\VBuildIdf}{v5.5.6}% [build] C:/Espressif/build_p11_esp32s3/project_description.json :: git_revision
\newcommand{\VBuildGcc}{14.2.0}% [build] C:/Espressif/build_p11_esp32s3/project_description.json :: c_compiler (versão do toolchain esp-)
\newcommand{\VBuildEspBinKb}{209}% [build] C:/Espressif/build_p11_esp32/size.txt :: binary size (hex)/1024
\newcommand{\VBuildEspBinPct}{20}% [build] C:/Espressif/build_p11_esp32/size.txt :: bin/partição app
\newcommand{\VBuildEspFlashCodeKb}{124}% [build] C:/Espressif/build_p11_esp32/size.txt :: Flash Code Used/1024
\newcommand{\VBuildEspFlashDataKb}{19}% [build] C:/Espressif/build_p11_esp32/size.txt :: Flash Data Used/1024
\newcommand{\VBuildEspDramKb}{15}% [build] C:/Espressif/build_p11_esp32/size.txt :: DRAM Used/1024
\newcommand{\VBuildEspDramTotalKb}{176}% [build] C:/Espressif/build_p11_esp32/size.txt :: DRAM Total/1024
\newcommand{\VBuildEspIramKb}{58}% [build] C:/Espressif/build_p11_esp32/size.txt :: IRAM Used/1024
\newcommand{\VBuildEspDramName}{DRAM}% [build] C:/Espressif/build_p11_esp32/size.txt :: nome da região de dados interna
\newcommand{\VBuildEspWsRatio}{10.3}% [estimativa] C:/Espressif/build_p11_esp32/size.txt :: estimativa ws float 3 s / DRAM total
\newcommand{\VBuildEspWsRatioOne}{5.3}% [estimativa] C:/Espressif/build_p11_esp32/size.txt :: estimativa ws float 1 s / DRAM total
\newcommand{\VBuildEspsBinKb}{233}% [build] C:/Espressif/build_p11_esp32s3/size.txt :: binary size (hex)/1024
\newcommand{\VBuildEspsBinPct}{23}% [build] C:/Espressif/build_p11_esp32s3/size.txt :: bin/partição app
\newcommand{\VBuildEspsFlashCodeKb}{139}% [build] C:/Espressif/build_p11_esp32s3/size.txt :: Flash Code Used/1024
\newcommand{\VBuildEspsFlashDataKb}{30}% [build] C:/Espressif/build_p11_esp32s3/size.txt :: Flash Data Used/1024
\newcommand{\VBuildEspsDramKb}{54}% [build] C:/Espressif/build_p11_esp32s3/size.txt :: DIRAM Used/1024
\newcommand{\VBuildEspsDramTotalKb}{334}% [build] C:/Espressif/build_p11_esp32s3/size.txt :: DIRAM Total/1024
\newcommand{\VBuildEspsIramKb}{16}% [build] C:/Espressif/build_p11_esp32s3/size.txt :: IRAM Used/1024
\newcommand{\VBuildEspsDramName}{DIRAM}% [build] C:/Espressif/build_p11_esp32s3/size.txt :: nome da região de dados interna
\newcommand{\VBuildEspsWsRatio}{5.5}% [estimativa] C:/Espressif/build_p11_esp32s3/size.txt :: estimativa ws float 3 s / DIRAM total
\newcommand{\VBuildEspsWsRatioOne}{2.8}% [estimativa] C:/Espressif/build_p11_esp32s3/size.txt :: estimativa ws float 1 s / DIRAM total
\newcommand{\VChipModel}{\textbf{[pending]}}% [esp32] embarcados/results/manifests/bancada.json :: chip_modelo
\newcommand{\VChipCpuMhz}{\textbf{[pending]}}% [esp32] embarcados/results/manifests/bancada.json :: cpu_mhz
\newcommand{\VChipPsramMb}{\textbf{[pending]}}% [esp32] embarcados/results/manifests/bancada.json :: psram_mb
\newcommand{\VChipLatIdmtMed}{\textbf{[pending]}}% [esp32] embarcados/results/manifests/bancada.json :: resumo.idmt.total_ms.mediana
\newcommand{\VChipLatIdmtPninetyfive}{\textbf{[pending]}}% [esp32] embarcados/results/manifests/bancada.json :: resumo.idmt.total_ms.p95
\newcommand{\VChipLatIdmtMax}{\textbf{[pending]}}% [esp32] embarcados/results/manifests/bancada.json :: resumo.idmt.total_ms.max
\newcommand{\VChipLatIdmtN}{\textbf{[pending]}}% [esp32] embarcados/results/manifests/bancada.json :: resumo.idmt.total_ms.n
\newcommand{\VChipRtfIdmt}{\textbf{[pending]}}% [esp32] embarcados/results/manifests/bancada.json :: fator_tempo_real_mediana
\newcommand{\VChipRtfMaxIdmt}{\textbf{[pending]}}% [esp32] embarcados/results/manifests/bancada.json :: fator_tempo_real_max
\newcommand{\VChipHeapIntIdmtKb}{\textbf{[pending]}}% [esp32] embarcados/results/manifests/bancada.json :: heap_min_livre_interna/1024
\newcommand{\VChipStackIdmt}{\textbf{[pending]}}% [esp32] embarcados/results/manifests/bancada.json :: pilha_folga_bytes_min
\newcommand{\VChipWsPsramIdmt}{\textbf{[pending]}}% [esp32] embarcados/results/manifests/bancada.json :: ws_na_psram
\newcommand{\VChipWsIdmtMb}{\textbf{[pending]}}% [esp32] embarcados/results/manifests/bancada.json :: ws_bytes/1e6
\newcommand{\VChipLatBldcMed}{\textbf{[pending]}}% [esp32] embarcados/results/manifests/bancada.json :: resumo.a2212.total_ms.mediana
\newcommand{\VChipLatBldcPninetyfive}{\textbf{[pending]}}% [esp32] embarcados/results/manifests/bancada.json :: resumo.a2212.total_ms.p95
\newcommand{\VChipLatBldcMax}{\textbf{[pending]}}% [esp32] embarcados/results/manifests/bancada.json :: resumo.a2212.total_ms.max
\newcommand{\VChipLatBldcN}{\textbf{[pending]}}% [esp32] embarcados/results/manifests/bancada.json :: resumo.a2212.total_ms.n
\newcommand{\VChipRtfBldc}{\textbf{[pending]}}% [esp32] embarcados/results/manifests/bancada.json :: fator_tempo_real_mediana
\newcommand{\VChipRtfMaxBldc}{\textbf{[pending]}}% [esp32] embarcados/results/manifests/bancada.json :: fator_tempo_real_max
\newcommand{\VChipHeapIntBldcKb}{\textbf{[pending]}}% [esp32] embarcados/results/manifests/bancada.json :: heap_min_livre_interna/1024
\newcommand{\VChipStackBldc}{\textbf{[pending]}}% [esp32] embarcados/results/manifests/bancada.json :: pilha_folga_bytes_min
\newcommand{\VChipWsPsramBldc}{\textbf{[pending]}}% [esp32] embarcados/results/manifests/bancada.json :: ws_na_psram
\newcommand{\VChipWsBldcMb}{\textbf{[pending]}}% [esp32] embarcados/results/manifests/bancada.json :: ws_bytes/1e6
\newcommand{\VChipEqStagesPass}{\textbf{[pending]}}% [esp32] embarcados/results/manifests/bancada.json :: equivalência no chip não rodou
\newcommand{\VChipEqStages}{\textbf{[pending]}}% [esp32] embarcados/results/manifests/bancada.json :: equivalência no chip não rodou
\newcommand{\VChipEqDecisionPass}{\textbf{[pending]}}% [esp32] embarcados/results/manifests/bancada.json :: equivalência no chip não rodou
\newcommand{\VChipSeqDelay}{\textbf{[pending]}}% [esp32] embarcados/results/manifests/bancada.json :: fluxo contínuo no chip não rodou

% Gerado por embarcados/src/python/gerar_artigo_numeros_aud.py. NÃO EDITAR À MÃO.
% commit 9fb82f016b4482c1a5b58f1b9bd08db149799b78 (árvore suja em embarcados/: True)
% Macros da auditoria independente (M2 holdout, M3 motor igual por igual, M4 rotação).
% Uma macro por número. Comentário: [categoria] fonte :: coluna.

\newcommand{\audHoldoutSeed}{7}% [equivalencia] embarcados/results/csv/aud_holdout_20261025-140558_meta.json :: seed
\newcommand{\audHoldoutN}{24}% [equivalencia] embarcados/results/csv/aud_holdout_20261025-140558_resumo.csv :: n_janelas[float,_todas]
\newcommand{\audHoldoutNIdmt}{12}% [equivalencia] embarcados/results/csv/aud_holdout_20261025-140558_resumo.csv :: n_janelas[float,idmt]
\newcommand{\audHoldoutNBldc}{12}% [equivalencia] embarcados/results/csv/aud_holdout_20261025-140558_resumo.csv :: n_janelas[float,a2212]
\newcommand{\audHoldoutFloatPathSame}{24}% [equivalencia] embarcados/results/csv/aud_holdout_20261025-140558_resumo.csv :: caminho_identico[float,_todas]
\newcommand{\audHoldoutFloatDecSame}{24}% [equivalencia] embarcados/results/csv/aud_holdout_20261025-140558_resumo.csv :: decisao_igual[float,_todas]
\newcommand{\audHoldoutFloatFeatOut}{0}% [equivalencia] embarcados/results/csv/aud_holdout_20261025-140558_resumo.csv :: n_features_fora_tol[float,_todas]
\newcommand{\audHoldoutFloatStageFail}{0}% [equivalencia] embarcados/results/csv/aud_holdout_20261025-140558_resumo.csv :: n_etapas_fail[float,_todas]
\newcommand{\audHoldoutFloatRatioMax}{0.20}% [equivalencia] embarcados/results/csv/aud_holdout_20261025-140558_resumo.csv :: razao_tol_max_features[float,_todas]
\newcommand{\audHoldoutFloatDpMax}{\ensuremath{8.7\times10^{-6}}}% [equivalencia] embarcados/results/csv/aud_holdout_20261025-140558_resumo.csv :: max_abs_dp[float,_todas]
\newcommand{\audHoldoutDoublePathSame}{24}% [equivalencia] embarcados/results/csv/aud_holdout_20261025-140558_resumo.csv :: caminho_identico[double,_todas]
\newcommand{\audHoldoutDoubleDecSame}{24}% [equivalencia] embarcados/results/csv/aud_holdout_20261025-140558_resumo.csv :: decisao_igual[double,_todas]
\newcommand{\audHoldoutDoubleFeatOut}{0}% [equivalencia] embarcados/results/csv/aud_holdout_20261025-140558_resumo.csv :: n_features_fora_tol[double,_todas]
\newcommand{\audHoldoutDoubleStageFail}{0}% [equivalencia] embarcados/results/csv/aud_holdout_20261025-140558_resumo.csv :: n_etapas_fail[double,_todas]
\newcommand{\audHoldoutDoubleRatioMax}{0.05}% [equivalencia] embarcados/results/csv/aud_holdout_20261025-140558_resumo.csv :: razao_tol_max_features[double,_todas]
\newcommand{\audHoldoutDoubleDpMax}{\ensuremath{1.9\times10^{-7}}}% [equivalencia] embarcados/results/csv/aud_holdout_20261025-140558_resumo.csv :: max_abs_dp[double,_todas]
\newcommand{\audRotBldcSpearman}{0.02}% [dados] embarcados/results/csv/aud_rotacao_motor_20261025-140355.csv :: spearman_rho[global]
\newcommand{\audRotBldcFiles}{43}% [dados] embarcados/results/csv/aud_rotacao_motor_20261025-140355.csv :: n_arquivos[global]
\newcommand{\audRotBldcWindows}{430}% [dados] embarcados/results/csv/aud_rotacao_motor_20261025-140355.csv :: n_janelas[global]
\newcommand{\audRotBldcCeilPct}{3.7}% [dados] embarcados/results/csv/aud_rotacao_motor_20261025-140355.csv :: 100·frac_janelas_teto[global]
\newcommand{\audRotBldcCeilThrPct}{99}% [dados] embarcados/results/csv/aud_rotacao_motor_20261025-140355.csv :: 100·lim_teto_rel
\newcommand{\audRotBldcSpearmanMOne}{0.30}% [dados] embarcados/results/csv/aud_rotacao_motor_20261025-140355.csv :: spearman_rho[M1]
\newcommand{\audRotBldcSpearmanPMOne}{0.29}% [dados] embarcados/results/csv/aud_rotacao_motor_20261025-140355.csv :: spearman_p[M1]
\newcommand{\audRotBldcFilesMOne}{14}% [dados] embarcados/results/csv/aud_rotacao_motor_20261025-140355.csv :: n_arquivos[M1]
\newcommand{\audRotBldcSpearmanMTwo}{0.30}% [dados] embarcados/results/csv/aud_rotacao_motor_20261025-140355.csv :: spearman_rho[M2]
\newcommand{\audRotBldcSpearmanPMTwo}{0.26}% [dados] embarcados/results/csv/aud_rotacao_motor_20261025-140355.csv :: spearman_p[M2]
\newcommand{\audRotBldcFilesMTwo}{16}% [dados] embarcados/results/csv/aud_rotacao_motor_20261025-140355.csv :: n_arquivos[M2]
\newcommand{\audRotBldcSpearmanMThree}{0.54}% [dados] embarcados/results/csv/aud_rotacao_motor_20261025-140355.csv :: spearman_rho[M3]
\newcommand{\audRotBldcSpearmanPMThree}{0.27}% [dados] embarcados/results/csv/aud_rotacao_motor_20261025-140355.csv :: spearman_p[M3]
\newcommand{\audRotBldcFilesMThree}{6}% [dados] embarcados/results/csv/aud_rotacao_motor_20261025-140355.csv :: n_arquivos[M3]
\newcommand{\audRotBldcSpearmanMFour}{1.00}% [dados] embarcados/results/csv/aud_rotacao_motor_20261025-140355.csv :: spearman_rho[M4]
\newcommand{\audRotBldcSpearmanPMFour}{0.00}% [dados] embarcados/results/csv/aud_rotacao_motor_20261025-140355.csv :: spearman_p[M4]
\newcommand{\audRotBldcFilesMFour}{7}% [dados] embarcados/results/csv/aud_rotacao_motor_20261025-140355.csv :: n_arquivos[M4]
\newcommand{\audRotBldcSpearmanMotorMin}{0.30}% [dados] embarcados/results/csv/aud_rotacao_motor_20261025-140355.csv :: min spearman_rho[M1..M4]
\newcommand{\audRotBldcSpearmanMotorMax}{1.00}% [dados] embarcados/results/csv/aud_rotacao_motor_20261025-140355.csv :: max spearman_rho[M1..M4]
\newcommand{\audRotBldcFilesMotorMin}{6}% [dados] embarcados/results/csv/aud_rotacao_motor_20261025-140355.csv :: min n_arquivos[M1..M4]
\newcommand{\audRotBldcFilesMotorMax}{16}% [dados] embarcados/results/csv/aud_rotacao_motor_20261025-140355.csv :: max n_arquivos[M1..M4]
\newcommand{\audMotorBalClsFileHP}{0.972}% [diagnostico] embarcados/results/csv/aud_motor_igual_20261025-140432.csv :: acuracia_balanceada[protocolo,p04_classicas,arquivo,HP_M1M3]
\newcommand{\audMotorBalClsFileHPLo}{0.924}% [diagnostico] embarcados/results/csv/aud_motor_igual_20261025-140432.csv :: ic95_lo[protocolo,p04_classicas,arquivo,HP_M1M3]
\newcommand{\audMotorBalClsFileHPHi}{0.998}% [diagnostico] embarcados/results/csv/aud_motor_igual_20261025-140432.csv :: ic95_hi[protocolo,p04_classicas,arquivo,HP_M1M3]
\newcommand{\audMotorBalClsMotorHP}{0.920}% [diagnostico] embarcados/results/csv/aud_motor_igual_20261025-140432.csv :: acuracia_balanceada[protocolo,p04_classicas,motor,HP_M1M3]
\newcommand{\audMotorBalClsMotorHPLo}{0.868}% [diagnostico] embarcados/results/csv/aud_motor_igual_20261025-140432.csv :: ic95_lo[protocolo,p04_classicas,motor,HP_M1M3]
\newcommand{\audMotorBalClsMotorHPHi}{0.962}% [diagnostico] embarcados/results/csv/aud_motor_igual_20261025-140432.csv :: ic95_hi[protocolo,p04_classicas,motor,HP_M1M3]
\newcommand{\audMotorDropClsHP}{0.052}% [diagnostico] embarcados/results/csv/aud_motor_igual_20261025-140432.csv :: acuracia_balanceada[queda,p04_classicas,HP_M1M3]
\newcommand{\audMotorDropClsHPLo}{0.023}% [diagnostico] embarcados/results/csv/aud_motor_igual_20261025-140432.csv :: ic95_lo[queda,p04_classicas,HP_M1M3]
\newcommand{\audMotorDropClsHPHi}{0.086}% [diagnostico] embarcados/results/csv/aud_motor_igual_20261025-140432.csv :: ic95_hi[queda,p04_classicas,HP_M1M3]
\newcommand{\audMotorDropClsFull}{0.368}% [diagnostico] embarcados/results/csv/aud_motor_igual_20261025-140432.csv :: queda_completa[queda,p04_classicas,HP_M1M3]
\newcommand{\audMotorBalDecFileHP}{0.753}% [diagnostico] embarcados/results/csv/aud_motor_igual_20261025-140432.csv :: acuracia_balanceada[protocolo,p11_decifrae_harmsum,arquivo,HP_M1M3]
\newcommand{\audMotorBalDecFileHPLo}{0.662}% [diagnostico] embarcados/results/csv/aud_motor_igual_20261025-140432.csv :: ic95_lo[protocolo,p11_decifrae_harmsum,arquivo,HP_M1M3]
\newcommand{\audMotorBalDecFileHPHi}{0.842}% [diagnostico] embarcados/results/csv/aud_motor_igual_20261025-140432.csv :: ic95_hi[protocolo,p11_decifrae_harmsum,arquivo,HP_M1M3]
\newcommand{\audMotorBalDecMotorHP}{0.685}% [diagnostico] embarcados/results/csv/aud_motor_igual_20261025-140432.csv :: acuracia_balanceada[protocolo,p11_decifrae_harmsum,motor,HP_M1M3]
\newcommand{\audMotorBalDecMotorHPLo}{0.593}% [diagnostico] embarcados/results/csv/aud_motor_igual_20261025-140432.csv :: ic95_lo[protocolo,p11_decifrae_harmsum,motor,HP_M1M3]
\newcommand{\audMotorBalDecMotorHPHi}{0.774}% [diagnostico] embarcados/results/csv/aud_motor_igual_20261025-140432.csv :: ic95_hi[protocolo,p11_decifrae_harmsum,motor,HP_M1M3]
\newcommand{\audMotorDropDecHP}{0.068}% [diagnostico] embarcados/results/csv/aud_motor_igual_20261025-140432.csv :: acuracia_balanceada[queda,p11_decifrae_harmsum,HP_M1M3]
\newcommand{\audMotorDropDecHPLo}{0.031}% [diagnostico] embarcados/results/csv/aud_motor_igual_20261025-140432.csv :: ic95_lo[queda,p11_decifrae_harmsum,HP_M1M3]
\newcommand{\audMotorDropDecHPHi}{0.116}% [diagnostico] embarcados/results/csv/aud_motor_igual_20261025-140432.csv :: ic95_hi[queda,p11_decifrae_harmsum,HP_M1M3]
\newcommand{\audMotorDropDecFull}{0.360}% [diagnostico] embarcados/results/csv/aud_motor_igual_20261025-140432.csv :: queda_completa[queda,p11_decifrae_harmsum,HP_M1M3]
\newcommand{\audMotorBalRmsNormFileHP}{0.843}% [diagnostico] embarcados/results/csv/aud_motor_igual_20261025-140432.csv :: acuracia_balanceada[protocolo,b00_rmspico_norm,arquivo,HP_M1M3]
\newcommand{\audMotorBalRmsNormFileHPLo}{0.760}% [diagnostico] embarcados/results/csv/aud_motor_igual_20261025-140432.csv :: ic95_lo[protocolo,b00_rmspico_norm,arquivo,HP_M1M3]
\newcommand{\audMotorBalRmsNormFileHPHi}{0.909}% [diagnostico] embarcados/results/csv/aud_motor_igual_20261025-140432.csv :: ic95_hi[protocolo,b00_rmspico_norm,arquivo,HP_M1M3]
\newcommand{\audMotorBalRmsNormMotorHP}{0.814}% [diagnostico] embarcados/results/csv/aud_motor_igual_20261025-140432.csv :: acuracia_balanceada[protocolo,b00_rmspico_norm,motor,HP_M1M3]
\newcommand{\audMotorBalRmsNormMotorHPLo}{0.720}% [diagnostico] embarcados/results/csv/aud_motor_igual_20261025-140432.csv :: ic95_lo[protocolo,b00_rmspico_norm,motor,HP_M1M3]
\newcommand{\audMotorBalRmsNormMotorHPHi}{0.887}% [diagnostico] embarcados/results/csv/aud_motor_igual_20261025-140432.csv :: ic95_hi[protocolo,b00_rmspico_norm,motor,HP_M1M3]
\newcommand{\audMotorDropRmsNormHP}{0.029}% [diagnostico] embarcados/results/csv/aud_motor_igual_20261025-140432.csv :: acuracia_balanceada[queda,b00_rmspico_norm,HP_M1M3]
\newcommand{\audMotorDropRmsNormHPLo}{-0.002}% [diagnostico] embarcados/results/csv/aud_motor_igual_20261025-140432.csv :: ic95_lo[queda,b00_rmspico_norm,HP_M1M3]
\newcommand{\audMotorDropRmsNormHPHi}{0.066}% [diagnostico] embarcados/results/csv/aud_motor_igual_20261025-140432.csv :: ic95_hi[queda,b00_rmspico_norm,HP_M1M3]
\newcommand{\audMotorDropRmsNormFull}{0.019}% [diagnostico] embarcados/results/csv/aud_motor_igual_20261025-140432.csv :: queda_completa[queda,b00_rmspico_norm,HP_M1M3]
\newcommand{\audMotorBalRmsRawFileHP}{0.913}% [diagnostico] embarcados/results/csv/aud_motor_igual_20261025-140432.csv :: acuracia_balanceada[protocolo,b00_rmspico_bruto,arquivo,HP_M1M3]
\newcommand{\audMotorBalRmsRawFileHPLo}{0.821}% [diagnostico] embarcados/results/csv/aud_motor_igual_20261025-140432.csv :: ic95_lo[protocolo,b00_rmspico_bruto,arquivo,HP_M1M3]
\newcommand{\audMotorBalRmsRawFileHPHi}{0.983}% [diagnostico] embarcados/results/csv/aud_motor_igual_20261025-140432.csv :: ic95_hi[protocolo,b00_rmspico_bruto,arquivo,HP_M1M3]
\newcommand{\audMotorBalRmsRawMotorHP}{0.754}% [diagnostico] embarcados/results/csv/aud_motor_igual_20261025-140432.csv :: acuracia_balanceada[protocolo,b00_rmspico_bruto,motor,HP_M1M3]
\newcommand{\audMotorBalRmsRawMotorHPLo}{0.623}% [diagnostico] embarcados/results/csv/aud_motor_igual_20261025-140432.csv :: ic95_lo[protocolo,b00_rmspico_bruto,motor,HP_M1M3]
\newcommand{\audMotorBalRmsRawMotorHPHi}{0.873}% [diagnostico] embarcados/results/csv/aud_motor_igual_20261025-140432.csv :: ic95_hi[protocolo,b00_rmspico_bruto,motor,HP_M1M3]
\newcommand{\audMotorDropRmsRawHP}{0.159}% [diagnostico] embarcados/results/csv/aud_motor_igual_20261025-140432.csv :: acuracia_balanceada[queda,b00_rmspico_bruto,HP_M1M3]
\newcommand{\audMotorDropRmsRawHPLo}{0.078}% [diagnostico] embarcados/results/csv/aud_motor_igual_20261025-140432.csv :: ic95_lo[queda,b00_rmspico_bruto,HP_M1M3]
\newcommand{\audMotorDropRmsRawHPHi}{0.251}% [diagnostico] embarcados/results/csv/aud_motor_igual_20261025-140432.csv :: ic95_hi[queda,b00_rmspico_bruto,HP_M1M3]
\newcommand{\audMotorDropRmsRawFull}{0.106}% [diagnostico] embarcados/results/csv/aud_motor_igual_20261025-140432.csv :: queda_completa[queda,b00_rmspico_bruto,HP_M1M3]
\newcommand{\audMotorHPFiles}{36}% [diagnostico] embarcados/results/csv/aud_motor_igual_20261025-140432.csv :: n_arquivos[queda]
\newcommand{\audMotorHPWindows}{360}% [diagnostico] embarcados/results/csv/aud_motor_igual_20261025-140432.csv :: n_janelas[queda]

% Gerado por embarcados/src/python/gerar_artigo_figuras.py. NÃO EDITAR À MÃO.
% Caminhos relativos a embarcados/paper/.
\newif\iffigesp
\figesptrue  % true só se a figura do ESP32 existir
\newcommand{\FigArq}{../figures/pdf/arquitetura.pdf}
\newcommand{\FigBalacc}{../figures/pdf/diag_balacc_en_20261023_222831.pdf}
\newcommand{\FigConf}{../figures/pdf/diag_confusao_p11_en_20261023_222831.pdf}
\newcommand{\FigEquiv}{../figures/pdf/equivalencia_erro_etapas_en.pdf}
\newcommand{\FigSpec}{../figures/pdf/diag_espectro_en_20261023_222831.pdf}
\newcommand{\FigEspLat}{../figures/pdf/esp32_lat_esp32_20261025-143026_en.pdf}


\begin{document}

\title{Porting a Harmonic-Sum Rotation Tracker for Acoustic Motor-State Monitoring to a
Microcontroller: Stage-Wise Verified C Port, Measured Cost, and What Two Public Datasets
Allow Us to Conclude\ifwip\ (Work in Progress)\fi}

% TODO(autores): nomes e afiliações a preencher pelos autores.
\author{\IEEEauthorblockN{[Authors to be completed]}
\IEEEauthorblockA{[Affiliation]}}

\maketitle

\begin{abstract}
We port to portable C a pipeline that classifies the operating state of small electric motors
from sound by borrowing ideas from bass-guitar transcription: a harmonic-sum rotation tracker
decoded by a Viterbi search, \VParNFeat{} physical features, cosine-to-prototype scoring and an
online hidden Markov model (HMM). The port is checked against the Python reference stage by
stage, on \VEqCases{} windows and \VEqSeqCombos{} sequence configurations, with tolerances
frozen before the formal comparison. In double precision all \VEqStagesTotal{} stage checks
pass. In float32, the build intended for the microcontroller, all \VEqRealCases{} windows of real
audio pass every stage; on \VEqFloatBadCases{} synthetic constant-offset-plus-noise windows the
envelope-modulation peak lies at the float32 noise floor and the modulation, feature and z-score
checks fail, while the class decision is identical on \VEqFloatDecisionPass{} of \VEqCases{}
windows. On an ESP32 without PSRAM the reference memory layout (estimated at
\VEstWsFloatThreeMb{}\,MB for a \VParSegIdmt{}-s window) does not fit. A named low-memory
variant fits \VParSegBldc{}-s windows, and on the chip its outputs agree with the reference within
the frozen tolerances on all \VChipEqCases{} windows compared; measured on the chip it is
\VChipRtfBldc{} times slower than real time, dominated by a table-free harmonic sum in
software-emulated double precision, and the \VParSegIdmt{}-s window does not fit. The pipeline does not outperform simple baselines:
on the IDMT dataset every method, including a two-number level baseline, saturates because
class, motor and recording coincide; on the A2212 dataset classical features reach
\VDiagBalClsBldcFile{} balanced accuracy against \VDiagBalDecBldcFile{} for the pipeline, and
with a motor held out all methods drop. The contribution is therefore the verified embedded
implementation, its cost, an interpretable rotation estimate, and an explicit account of what
these datasets can and cannot support.
\end{abstract}

\begin{IEEEkeywords}
acoustic condition monitoring, embedded signal processing, ESP32, rotation tracking,
Viterbi decoding, reproducibility, data leakage
\end{IEEEkeywords}

% ======================================================================
\section{Introduction}

Acoustic monitoring of electric motors is attractive because a microphone is cheap,
non-intrusive and needs no electrical connection to the machine. Several recent works report
near-perfect accuracy on two public datasets of small laboratory motors, the
IDMT-ISA-ELECTRIC-ENGINE set~\cite{grollmisch2019} and the A2212 brushless set~\cite{estacio2023},
often with deep networks~\cite{melcnn2025,ertarin2025} or large hand-crafted feature banks~\cite{etnext2026,tekin2025}.
Two aspects receive little attention. First, most of these results split recordings into
segments before the train/test split, so neighbouring segments of one recording fall on both
sides. Second, few works measure the cost on a microcontroller; an exception is the edge
classifier of Strantzalis et al.~\cite{strantzalis2022}, while other TinyML studies stop at a model
size and never run on hardware~\cite{almaini2026}.

This paper starts from a pipeline that we call Decifrae (p11). It reuses components of a
music-transcription system: the motor is treated as a ``one-string instrument'' whose
fundamental (the rotation frequency) is tracked through a sum over its harmonics and decoded with
a Viterbi search that has a silence state and an octave-jump cost, as a bass transcriber would do.
The analogy only motivates the \emph{representation}; we do not claim that motor sound is a
musical note, nor that music-transcription performance transfers to machines.

The question we answer is whether this pipeline fits and runs on an ESP32-class
microcontroller with an implementation that is demonstrably the same computation as the Python
reference, and what the two public datasets allow us to conclude about it. In short, it fits and
runs on an ESP32 without PSRAM only for \VParSegBldc{}-s windows, through a named low-memory
variant, and \VChipRtfBldc{} times slower than real time on that board.
\ifwip
We contribute a C port verified stage by stage against the reference with tolerances frozen before
the formal comparison, a cost account that keeps estimate, PC measurement, firmware build and
on-chip measurement apart, the tracked rotation as an interpretable output, and an evaluation
under three protocols showing that the pipeline does not beat simple baselines and that the
datasets cannot separate operating state from motor identity.
\else
Our contributions are:
\begin{itemize}
\item a portable C port without dynamic allocation, verified against the Python reference stage
by stage with tolerances frozen before the formal comparison, in double and in float32, with
every residual discrepancy listed;
\item a cost account that keeps four categories apart: analytical estimate, measurement on a PC,
firmware build (\texttt{idf.py size}), and measurement on an ESP32 without PSRAM, including the
low-memory variant needed there;
\item an interpretable intermediate output, the tracked rotation frequency, together with a check
of how far it can be trusted on each dataset;
\item an evaluation of the pipeline and of simple baselines on the same windows under three
protocols (random windows, grouped by recording or file, and motor held out), showing that the
pipeline does not outperform the baselines and that the datasets cannot separate operating
state from motor identity.
\end{itemize}
\fi
Three distinctions are kept throughout. We classify the \emph{current operating state} of a
motor; we do not predict future failures. Both datasets contain small laboratory motors, not
industrial machines. And an evaluation in which the test motor was also seen in training says
nothing about generalisation to a new motor.

% ======================================================================
\section{Related Work}

\emph{Datasets and splits.} Grollmisch et al.~\cite{grollmisch2019} released the IDMT set with
an official split (training on clean recordings, testing on recordings with added noise) and a
frame-level network. Each class was recorded on a different motor, so class and physical unit
coincide. Later works on IDMT and A2212 report near-perfect accuracy with mel-spectrogram
CNNs~\cite{melcnn2025}, image encodings~\cite{ertarin2025}, wavelet pattern features with
SVM/kNN~\cite{etnext2026} and wavelet statistics with bagged trees~\cite{tekin2025}; their splits
operate on segments rather than on recordings. The A2212 data paper~\cite{estacio2023} reports
no classifier; it documents that the bearing fault exists on a single motor.
\ifwip\else
The classical-feature baseline used here follows the feature list described in the abstract of
Vanaraj et al.~\cite{vanaraj2026}; we only read the abstract, so it is an inspired
re-implementation, not a reproduction. Glowacz~\cite{glowacz2018} classified a commutator motor
with a nearest-mean rule on selected spectral lines, the closest precedent of our
cosine-to-prototype decision.
\fi

\ifwip
\emph{Rotation tracking without a tachometer} has been reviewed on vibration
data~\cite{peeters2019}, and Viterbi decoding over a time--frequency representation predates our
work~\cite{xi2015}; we add its verified embedded implementation on microphone audio.
\else
\emph{Rotation tracking without a tachometer.} Peeters et al.~\cite{peeters2019} review and
compare tacholess speed estimators on vibration data, including multi-order probabilistic
approaches that combine several harmonics. Viterbi decoding over a time--frequency
representation for speed estimation predates our work~\cite{xi2015}, and harmonic weighting has
been proposed to suppress weak orders~\cite{yoo2023}. Wheat et al.~\cite{wheat2025} build
features on a harmonic grid given a known operating speed. Our tracker is therefore not new as an
idea; what we add is its embedded, verified implementation on microphone audio.
\fi

\ifwip\else
\emph{Embedded and anomaly-oriented work.} Strantzalis et al.~\cite{strantzalis2022} deploy a
small CNN on a Cortex-M4 and measure latency, RAM and response time to state changes on IDMT,
but with a random split of overlapping spectrograms. Almaini et al.~\cite{almaini2026} quantise an
MFCC network trained on urban sounds without deploying it. Unsupervised anomaly detection on
machine sounds (autoencoder baselines of the DCASE challenges) addresses novelty rather than state
classification and is not compared here.
\fi

% ======================================================================
\section{Method}
\label{sec:method}

Fig.~\ref{fig:arq} summarises the pipeline. Each analysis window (\VParSegIdmt{}\,s for IDMT,
\VParSegBldc{}\,s for A2212, both at \VParSrKhz{}\,kHz) is RMS-normalised and processed
independently; only the HMM carries state across windows.

\begin{figure}[t]
\centering
\includegraphics[width=\columnwidth]{\FigArq}
\caption{System overview. Everything inside the dashed box is the C port that runs on the
microcontroller; audio currently reaches the chip over UART (processor-in-the-loop) or from its
flash. Working memory: analytical estimate of the reference layout, and the pool peak of the
low-memory variant measured on the ESP32.}
\label{fig:arq}
\end{figure}

\emph{Rotation tracker.} A short-time Fourier transform with a symmetric Hann window of
\VParNfft{} samples and hop \VParHop{} gives magnitude frames $|X_t(k)|$. For each of
$J=\VParJ{}$ candidate rotation frequencies $f_j$ on a logarithmic grid from \VParFmin{} to
\VParFmax{}\,Hz (step \VParStepCents{}~cents), the saliency
$s_t(j)=\sum_{h=1}^{K}|X_t(h f_j/\Delta f)|$ sums $K=\VParNharm{}$ harmonics with linear
interpolation between bins, and is normalised per frame to $\tilde e_t(j)\in[0,1]$. The emission
cost is $-\log(\VParWEnergy{}\,\tilde e_t(j)+\VParEmisEps{})$; a silence state receives a cost
derived from the frame energy relative to the loudest frame of the window. A Viterbi search
allows jumps of up to \VParMaxJump{} states (\VParMaxJumpCents{} cents) with a cost proportional
to the jump, octave jumps of $\pm$\VParOctStates{} states with a fixed cost, and transitions into
and out of silence. The decoded path gives $f_{rot}(t)$.

\emph{Features and decision.} From the path and the mean power spectrum we compute \VParNFeat{}
features\ifwip\else{} (Table~\ref{tab:params})\fi: the median and inter-quartile range of $f_{rot}$, the silent-frame
fraction, the saliency along the path, a harmonic-energy ratio, the relative energy of four groups
of harmonic orders, the frequency and prominence of the main envelope-modulation peak (Hilbert
envelope, Butterworth low-pass, decimation to \VParEnvRateHz{}\,Hz, Welch spectrum), the spectral
centroid and the spectral flatness. Features are z-scored with training statistics and compared
by cosine with one prototype per class; a softmax with temperature \VParTempCos{} gives class
probabilities. For streams, an HMM with a fixed change penalty and an extra ``no decision''
state is run as a causal forward filter, one window at a time.

\ifwip\else
% Gerado por embarcados/src/python/gerar_artigo_tabelas.py a partir de embarcados/paper/numeros.json.
% NÃO EDITAR À MÃO.

\begin{table}[t]
\centering
\scriptsize
\setlength{\tabcolsep}{3pt}
\caption{Parameters of the p11 reference, read from the Python code (\texttt{PreDecifrae}, \texttt{rotacao}, \texttt{\_envelope\_mod}), and the 13 features in the order of the reference.}
\label{tab:params}
\begin{tabular}{@{}p{1.55cm}p{6.6cm}@{}}
\toprule
Sampling rate; window & 16\,kHz; 3\,s (IDMT), 1\,s (A2212); RMS $\rightarrow$ 0.1 \\
STFT & symmetric Hann, $N$=8{,}192, hop 1{,}024, no centring; 39 / 8 frames \\
Rotation grid & 15--60\,Hz, 2-cent step, $J$=1{,}201 states + silence \\
Harmonic sum & $K$=48 harmonics, unit weights, linear interpolation \\
Emission cost & $-\log(0.45\,\tilde e + \ensuremath{10^{-3}})$; silence from frame energy \\
Viterbi & step 0.55/semitone, jumps $\le$50 states, octave $\pm$600 states (cost 7.0), silence 3.2 \\
Envelope & $|\mathcal{H}(y)|$, Butterworth 4, 200\,Hz (causal \texttt{sosfilt}), $\downarrow$16, Welch; peak in 2--100\,Hz \\
Classifier & z-score, cosine to class means, softmax $T$=0.1 \\
Online HMM & change penalty swept (6 values, 0.5--16), $\log p_N$=-50 \\
\midrule
Features & (0) median $f_{rot}$; (1) IQR of $f_{rot}$; (2) silent-frame fraction; (3) mean normalised saliency on the path; (4) harmonic ratio ($k f_{rot}\pm1$ bin / 10--4{,}000\,Hz); (5) energy of orders 1--4 (dB); (6) energy of orders 5--12 (dB); (7) energy of orders 13--24 (dB); (8) energy of orders 25--48 (dB); (9) envelope-modulation peak frequency; (10) its prominence (dB); (11) spectral centroid; (12) spectral flatness \\
\bottomrule
\end{tabular}
\end{table}

\fi

\ifwip
\emph{C port.} The C code follows the Python line by line, compiles unchanged for the PC
and the ESP-IDF, processes the spectrogram frame by frame, includes a mixed-radix FFT for the
full-length Hilbert transform, and uses a caller-provided arena instead of dynamic allocation.
The real type is float32 on the microcontroller and double for verification. Any change of
arithmetic or memory layout is a named variant: \texttt{p11\_c\_sosd} keeps the envelope filter
state in double (more precise, not an approximation), and \texttt{\VChipVariant}, the one run on
our board without PSRAM, takes its memory from a pool of internal DRAM and IRAM blocks, drops the
harmonic table (positions recomputed in double) and computes the Hilbert FFT in four steps. An
ESP-DSP option (\texttt{p11\_c\_espdsp}) was neither compared nor measured and is left out.
\else
\emph{C port.} The C code follows the Python line by line and compiles unchanged for the PC and
for the ESP-IDF. It never stores the full spectrogram: each frame is windowed, transformed by a
real FFT, accumulated into the mean power spectrum and reduced to its saliency row. The Viterbi
search needs the whole window (the silence cost uses the maximum frame energy of the window), so
the normalised saliency and the back-pointers (short integers) are kept; the Hilbert envelope needs an FFT
of the full window length, which is not a power of two, so the port includes a mixed-radix FFT.
All buffers live in a caller-provided arena whose size is computed by the same layout function that
assigns the pointers. The real type is selected at compile time (float32 for the microcontroller,
double for verification). Index decisions that change integers (grid, medians, rounding of
harmonic positions) are computed in double in both builds, reproducing the reference's
round-half-to-even and first-index tie-breaking. Any change of arithmetic is a \emph{named
variant}, compared under the same tolerances. \texttt{p11\_c\_sosd} keeps the envelope filter
state in double, as SciPy does (more precise, not an approximation). \texttt{\VChipVariant}, run on
our board without PSRAM, calls the same core functions but carves its memory phase by phase from
a pool of internal DRAM and IRAM blocks (IRAM, word access only, holds only float vectors such as
the FFT twiddles), drops the harmonic-interpolation table (bin positions recomputed in double for
every frame) and computes the Hilbert FFT in four steps, the same DFT in another order. An ESP-DSP
option (\texttt{p11\_c\_espdsp}) was neither compared nor measured and is left out.
\fi

% ======================================================================
\section{Experimental Setup}
\label{sec:setup}

\emph{Data and protocols.} We use the \texttt{train/} part of IDMT (the official test set was not
opened) and all of A2212\ifwip\else{} (Table~\ref{tab:protocols})\fi. All methods are evaluated on exactly the
same windows and folds. On IDMT, the random-window protocol leaks by construction, and the
``blocks'' protocol holds out contiguous blocks of each recording; since there is one recording
per class, it still cannot separate class from motor. On A2212 we compare random windows, a
grouped split by file, and leave-one-motor-out; because the bearing fault exists only on motor
M4, we also report leave-one-motor-out without the M4 fold.

\ifwip\else% Gerado por embarcados/src/python/gerar_artigo_tabelas.py a partir de embarcados/paper/numeros.json.
% NÃO EDITAR À MÃO.

\begin{table}[t]
\centering
\scriptsize
\setlength{\tabcolsep}{3pt}
\caption{Datasets and evaluation protocols. IDMT: \texttt{train/} only, 3 recordings of 3 motors, 357 windows of 3\,s. A2212: 43 files of 4 motors (22 Healthy, 14 Propeller, 7 Bearing), 430 windows of 1\,s. $g$: groups used by the bootstrap. $^\dagger$Intentional leakage, reported only to measure inflation.}
\label{tab:protocols}
\begin{tabular}{@{}p{1.55cm}p{1.35cm}r p{3.75cm}@{}}
\toprule
Protocol & Split unit & $g$ & What it can / cannot support \\
\midrule
IDMT random$^\dagger$ & window & 18 & Upper bound only: neighbouring windows of one recording on both sides. \\
IDMT blocks & contiguous block & 15 & One recording per class: class = motor = recording. Measures recording identity, not operating state. \\
A2212 random$^\dagger$ & window & 43 & Upper bound only: windows of one file on both sides. \\
A2212 file & file & 43 & State of a \emph{known} motor (same motors in training and test). Bearing exists only on motor M4. \\
A2212 motor-out & motor & 43 & Unseen motor. With M4 held out, Bearing has no training data (recall 0.000 by construction). Only 4 motors: no motor-level CI. \\
A2212 motor-out w/o M4 & motor (M1--M3) & 36 & Healthy vs.\ Propeller on an unseen motor; 3 folds. \\
\bottomrule
\end{tabular}
\end{table}
\fi

\emph{Baselines.} A level baseline uses only the RMS and peak of the window with logistic
regression, once on RMS-normalised audio (only the crest factor remains) and once on raw audio. A
classical baseline (p04) uses \VParNClassicFeat{} frame statistics (zero-crossing rate, RMS,
spectral centroid and MFCC means and deviations) with a soft-voting ensemble, inspired
by~\cite{vanaraj2026}.

\emph{Metrics.} We report balanced accuracy, macro-F1, macro one-vs-rest ROC-AUC from continuous
scores, and per-class recall. Confidence intervals (\VDataCiLevel{}\%) come from a bootstrap that
resamples groups (recording blocks or files), stratified by class, with \VDataNBoot{} resamples;
windows of one group are never treated as independent. For the motor-out protocols the
independent unit is the motor, of which there are only \VDataBldcMotors{}, so no honest interval
exists at that level; the file-level interval shown is optimistic.

\ifwip
\emph{Equivalence protocol.} Golden vectors of every stage were produced by the
instrumented reference on \VEqCases{} windows (\VEqGoldClean{} clean, \VEqGoldNoisy{} noisy,
\VEqGoldEdge{} edge cases) and \VEqSeqCombos{} HMM sequence configurations. Per-stage tolerances,
derived from rounding analysis and an independent emulation, were frozen (\texttt{\VEqTolCommit{}})
before the formal comparison, although the C port had already been debugged against the same
golden vectors by its developer; discrete outputs must agree exactly unless a tie is demonstrated.
\else
\emph{Equivalence protocol.} The instrumented Python reference produced golden vectors for every
stage on \VEqCases{} windows: \VEqGoldClean{} clean validation windows from both datasets,
\VEqGoldNoisy{} noisy versions under \VEqGoldPerts{} perturbations (white noise at several SNRs and
synthetic sag and flicker power-quality disturbances), and \VEqGoldEdge{} edge cases (silence, clipping, synthetic
rotation outside the search range, and a constant offset with tiny noise). The HMM stages were
checked on \VEqGoldSeqs{} probability sequences under \VParNChangePen{} change penalties and
\VEqMinDurValues{} minimum durations (\VEqSeqCombos{} configurations). Tolerances for each stage
were derived from first-order rounding analysis and an independent numerical emulation of each
build, and frozen before the formal comparison (commit \texttt{\VEqTolCommit{}}, file hash
\texttt{\VEqTolSha{}}\ldots, \VEqTolCommitTime{}; comparison run \VEqRunTime{}); the hash is
recorded in every result row. The chronology matters: the C port had already been debugged
against the same golden vectors by its developer before the tolerances were frozen, so the
tolerances are frozen, not pre-registered. On \audHoldoutN{} windows outside the golden set,
unseen by the C developer, path and decision agree on \audHoldoutFloatPathSame{} and
\audHoldoutFloatDecSame{} in float32, no feature exceeds its tolerance (maximum ratio
\audHoldoutFloatRatioMax{}) and probabilities differ by at most \audHoldoutFloatDpMax{}. Tolerances for the edge cases
were calibrated on the emulation error for the same signals, a weaker standard, and some are wide
in float32 (e.g.\ the modulation prominence). Continuous stages are compared by maximum error relative to the stage's
scale; discrete outputs (rotation path, decision, HMM states) must agree exactly unless a tie is
demonstrated from the reference's own values. A failure at one stage marks later stages as
downstream. Both emulations passed all comparison rows (\VEqEmulDoublePass{} and \VEqEmulFloatPass{} of
\VEqEmulFloatRows{}) before the C code was compared.
\fi

\ifwip
\emph{Cost measurement.} We keep apart (1) the analytical arena estimate, (2) timings on
the PC (\VPcNIdmt{} timings cycling over \VPcDistinctIdmt{} distinct windows, each run once in a new
process, so warm-up does not warm the C code; workspace initialisation excluded), (3)
\texttt{idf.py size} of the firmware (ESP-IDF \VBuildIdf{}) and (4) measurements on an
\VChipModel{} without PSRAM (\VChipCpuMhz{}\,MHz; \VChipLatBldcN{} timed windows over UART after
\VChipLatWarmup{} warm-up runs, on-chip timer around the processing only). Energy is not measured.
\else
\emph{Cost measurement.} We keep four categories apart. (1) \emph{Analytical estimate}: the
arena size returned by the C layout function. (2) \emph{Measured on the PC}: \VPcNIdmt{} timings
per dataset cycling over \VPcDistinctIdmt{} distinct windows. The C executable is started anew and
timed once per window, so warm-up does not warm it, and its time excludes the workspace
initialisation, which the Python timing cannot exclude; the two effects partly offset.
(3) \emph{Firmware build}: \texttt{idf.py size} (ESP-IDF \VBuildIdf{}, GCC \VBuildGcc{}).
(4) \emph{Measured on the ESP32}: an \VChipModel{} (rev.\ \VChipRev{}, \VChipCpuMhz{}\,MHz,
\VChipFlashMb{}\,MB flash, no PSRAM, radios never started); windows arrive over UART (or, for the
stream, from flash) and are timed with the on-chip timer around each core call, with heap and
stack watermarks. Energy is not measured.
\fi

% ======================================================================
\section{Results}
\label{sec:results}

\subsection{Python--C equivalence}

\ifwip
Table~\ref{tab:equiv} summarises the verdicts. In double precision all \VEqStagesTotal{}
stage checks pass on all windows and sequence configurations; the rotation path, the decision and
every HMM output agree exactly. In float32 all \VEqRealCases{} windows of real audio pass every
stage. The failures are confined to the \VEqFloatBadCases{} synthetic constant-offset-plus-noise
windows, whose modulation spectrum lies at the float32 rounding floor in both implementations, so
the peak location is undetermined (the reference peak is \VEqDcRefPeakRatio{} times the reference
power at the bin chosen by the C code: not a tie). We count the modulation, feature and z-score
checks as failed there (\VEqFloatStagesPassStrict{} of \VEqStagesTotal{} checks pass on all
windows); the class decision agrees on \VEqFloatDecisionPass{} of \VEqCases{} windows. The named
variant \texttt{p11\_c\_sosd} passes every check. An independent audit reproduced the comparison
from a clean build; its corrections are applied here. These are PC results; the comparison on the
chip is in Section~\ref{sec:cost}.
\else
Table~\ref{tab:equiv} summarises the verdicts\ifwip\else{} and Fig.~\ref{fig:equiv} the worst
error-to-tolerance ratio per stage\fi. In double precision all \VEqStagesTotal{} stage checks pass
on all windows and sequence configurations (worst ratio \VEqDoubleMaxRatio{} of the tolerance).
The rotation path, the decision and every HMM output agree exactly. In float32, all
\VEqRealCases{} windows of real audio (clean and noisy) pass every stage. The
\VEqFloatBadCases{} remaining failures are the synthetic constant-offset-plus-noise edge case, once
per model. There the envelope is almost constant and its modulation spectrum lies at the float32
rounding floor in both implementations, so the peak location is undetermined at that precision:
the reference's peak is \VEqDcRefPeakRatio{} times the reference power at the bin chosen by the C
code, so this is not a tie. The modulation frequency moves by \VEqFloatFmodErr{}\,Hz, and we count
the envelope-modulation, feature and z-score checks as failed on these windows (the frozen
comparison rule had accepted the first two as a tie; we report the conservative reading), leaving
\VEqFloatStagesPassStrict{} of \VEqStagesTotal{} stage checks passing on all windows. Cosines,
probabilities and the class decision still pass on all windows
(\VEqFloatDecisionPass{}/\VEqCases{}). With the named variant \texttt{p11\_c\_sosd} every check
passes; it reduces the maximum error of the decimated envelope,
relative to its peak and taken as the median over windows, from \VEqEnvDecFloatMed{} to
\VEqEnvDecSosdMed{}, because the low-pass poles lie close to
the unit circle and float32 rounding accumulates in the filter state. These results were obtained
on the PC; the same comparison on the chip is
\ifchipmedido reported in Section~\ref{sec:cost}\else pending\fi. An independent audit
(report in the repository) reproduced the comparison from a clean build and its corrections are
applied here. We therefore
claim stage-wise agreement within tolerances frozen before the formal comparison, with the
exceptions above, not unconditional equivalence.
\fi

\begin{table}[t]
\centering
\scriptsize
\caption{Stage checks passed on the PC (\VEqCases{} windows; \VEqSeqCombos{} sequence
configurations). Conservative reading: a group fails if any stage fails on any window. All
float32 failures are the \VEqFloatBadCases{} synthetic constant-offset edge windows, where the
modulation peak is at the float32 noise floor; the decision is unchanged.}
\label{tab:equiv}
% Gerado por embarcados/src/python/gerar_artigo_tabelas.py a partir de embarcados/results/csv/equivalencia_resumo_por_build_20261024-101821.csv (veredito por etapa).
% NÃO EDITAR À MÃO.

\begin{tabular}{@{}lccc@{}}
\toprule
Stage group (\#stages) & C double & C float32 & float32+SOS64 \\
\midrule
Front end (6) & 6/6 & 6/6 & 6/6 \\
Rotation Viterbi (3) & 3/3 & 3/3 & 3/3 \\
Spectrum, envelope (6) & 6/6 & 5/6 (FAIL in 2/164) & 6/6 \\
13 features (1) & 1/1 & 0/1 (FAIL in 2/164) & 1/1 \\
z-score (1) & 1/1 & 0/1 (FAIL in 2/164) & 1/1 \\
Cosine, softmax, decision (3) & 3/3 & 3/3 & 3/3 \\
Sequences (HMM) (6) & 6/6 & 6/6 & 6/6 \\
\bottomrule
\end{tabular}

\end{table}

\ifwip\else
\begin{figure}[t]
\centering
\includegraphics[width=\columnwidth]{\FigEquiv}
\caption{Worst error-to-tolerance ratio per stage (PC). Ratios below one pass. F: failure on the
synthetic constant-offset edge windows (counted as failures even where the frozen rule accepted a
tie). Exact zeros are drawn at the bottom of the axis.}
\label{fig:equiv}
\end{figure}
\fi

\subsection{Cost}
\label{sec:cost}

\ifwip
\emph{Measured on the PC},
the C float32 build of the reference layout needs \VPcCIdmtMed{}\,ms per \VParSegIdmt{}-s window
and \VPcCBldcMed{}\,ms per \VParSegBldc{}-s window, against \VPcPyIdmtMed{}\,ms and
\VPcPyBldcMed{}\,ms for the Python reference (\VPcSpeedupIdmt{}$\times$, including interpreter
overhead); the STFT with harmonic sum takes \VPcCIdmtTrackPct{}\%, the rotation Viterbi
\VPcCIdmtVitPct{}\% and the envelope branch \VPcCIdmtEnvPct{}\% of the median time. These numbers
say nothing about the microcontroller.
\else
Table~\ref{tab:resources} lists the resources by category. The PC timings (column 2; the STFT
with harmonic sum takes \VPcCIdmtTrackPct{}\%, the rotation Viterbi \VPcCIdmtVitPct{}\% and the
envelope branch \VPcCIdmtEnvPct{}\% of the median) say nothing about the microcontroller.
\fi

The \emph{analytical estimate} of the working memory of the reference layout (one arena) is
\VEstWsFloatThreeKib{}\,KiB for \VParSegIdmt{}\,s and \VEstWsFloatOneKib{}\,KiB for
\VParSegBldc{}\,s, at least \VChipRefOverPoolThree{} and \VChipRefOverPoolOne{} times the internal
RAM that the firmware can hand to the working memory on our board (\VChipPoolKib{}\,KiB, measured)
and at least \VBuildEspsWsRatio{} and \VBuildEspsWsRatioOne{} times the whole DIRAM of the
ESP32-S3 \emph{build} (\VBuildEspsDramTotalKb{}\,KiB; its IRAM is \VBuildEspsIramPct{}\% used):
without PSRAM the reference layout cannot run.

\ifchipmedido
\emph{Measured on the ESP32} (\VChipPsramMb{}\,MB PSRAM), \texttt{\VChipVariant} fits a
\VParSegBldc{}-s window: the pool peaks at \VChipWsBldcB{}\,bytes (\VChipWsPoolPct{}\% of it,
Fig.~\ref{fig:esplat}b), the minimum free internal heap is \VChipHeapIntBldcB{}\,bytes and the
stack headroom \VChipStackBldc{}\,bytes. The \VParSegIdmt{}-s IDMT window does \emph{not} fit, so
no IDMT figure exists on the chip. All \VChipEqCases{} golden \VParSegBldc{}-s windows
(\VChipEqCasesReal{} real, \VChipEqCasesEdge{} edge cases) pass, on the chip, the \VChipEqStages{}
final outputs the firmware returns (Viterbi path and score through decision), with \VChipEqFail{} failures and a worst error-to-tolerance
ratio of \VChipEqMaxRatio{}; intermediate stages are not sent back. On the PC the variant passes
all \VEqLmStages{} compared stages. The median latency is \VChipLatBldcMed{}\,ms per
\VParSegBldc{}-s window (p95 \VChipLatBldcPninetyfive{}, maximum \VChipLatBldcMax{}\,ms), a
real-time factor of \VChipRtfBldc{}: on this board the pipeline does \emph{not} run in real time.
The STFT with harmonic sum takes \VChipBldcTrackPct{}\% of it (Fig.~\ref{fig:esplat}a; on the PC,
with the table, \VPcCBldcTrackPct{}\%). We attribute this to recomputing the harmonic positions in
double precision, which the ESP32 only emulates in software (the variant was not timed on the PC);
a float32 variant of that step or a board with PSRAM is the way towards real time. The variant's
image occupies \VBuildLmBinKb{}\,KiB of flash (build).
\else
\emph{Measured on the ESP32}: not available when this version was generated.
\fi

\ifwip\else% Gerado por embarcados/src/python/gerar_artigo_tabelas.py a partir de embarcados/paper/numeros.json.
% NÃO EDITAR À MÃO.

\begin{table*}[t]
\centering
\scriptsize
\setlength{\tabcolsep}{4pt}
\caption{Resources by evidence category: (1) analytical estimate (\texttt{p11\_host ws}); (2) measured on the PC (Windows 11, MSVC 19.44 x64 \slash{}O2 \slash{}fp:precise; 200 timings per dataset over 75 distinct windows, one new process per window); (3) \texttt{idf.py size}, static only; (4) measured on the ESP32 without PSRAM (200 timed windows after 5 warm-up runs). Columns (1)--(2) use the reference layout \texttt{p11\_c\_ref}; column (4), the only measurement on the microcontroller, uses \texttt{p11\_c\_lowmem}.}
\label{tab:resources}
\begin{tabular}{@{}lllll@{}}
\toprule
Quantity & (1) Estimate & (2) PC & (3) Build & (4) ESP32 \\
\midrule
Latency, C float32, median [p95] (ms), 3\,s / 1\,s & -- & 13.29 [15.00] / 3.18 [3.84] & -- & \textit{did not fit} / 2647.5 [2647.6] \\
Latency, Python reference, median (ms), 3\,s / 1\,s & -- & 95.3 / 22.1 & -- & -- \\
Real-time factor (C, median), 3\,s / 1\,s; STFT + harm.\ sum share at 1\,s & -- & 0.0044 / 0.0032; 39\,\% & -- & \textit{did not fit} / 2.65; 89\,\% \\
Working memory (KiB), 3\,s / 1\,s & 1820 / 936 & -- & -- & \textit{did not fit} / 296 (pool peak) \\
Internal RAM for the pool; min.\ free heap (KiB); stack headroom (B) & -- & -- & -- & 308; 7.4; 20{,}536 \\
App binary (KiB): ESP32, -S3 (\texttt{p11\_c\_ref}); ESP32 (\texttt{p11\_c\_lowmem}) & -- & -- & 209, 233; 215 & -- \\
Static DRAM, IRAM used / region (KiB), ESP32, \texttt{p11\_c\_lowmem} & -- & -- & 35/176, 58/128 & -- \\
\bottomrule
\end{tabular}
\end{table*}
\fi

\iffigesp
\begin{figure}[t]
\centering
\includegraphics[width=\columnwidth]{\FigEspLat}
\caption{Measured on the ESP32 without PSRAM (\texttt{\VChipVariant}, A2212, \VParSegBldc{}-s
windows). (a) Median time per stage over \VChipLatBldcN{} windows. (b) Internal RAM handed to the
pool and its peak use (measured) against the estimated reference arena (hatched).}
\label{fig:esplat}
\end{figure}
\fi

\subsection{Classification under three protocols}

% Gerado por embarcados/src/python/gerar_artigo_tabelas.py a partir de embarcados/paper/numeros.json.
% NÃO EDITAR À MÃO.

\begin{table}[t]
\centering
\scriptsize
\setlength{\tabcolsep}{3pt}
\caption{Balanced accuracy on the same windows (validation only); $g$ = groups. 95\% CIs from a group bootstrap (2{,}000 resamples) are shown in Fig.~\ref{fig:balacc}. $^\dagger$Intentional leakage.}
\label{tab:desempenho-compacto}
\begin{tabular}{@{}lrcccc@{}}
\toprule
Protocol & $g$ & RMS/pk (norm.) & RMS/pk (raw) & Classical & p11 \\
\midrule
IDMT random$^\dagger$ & 18 & 0.992 & 0.997 & 1.000 & 1.000 \\
IDMT blocks & 15 & 0.991 & 0.993 & 1.000 & 1.000 \\
A2212 random$^\dagger$ & 43 & 0.564 & 0.625 & 0.998 & 0.849 \\
A2212 file & 43 & 0.562 & 0.608 & 0.981 & 0.816 \\
A2212 motor-out & 43 & 0.542 & 0.503 & 0.613 & 0.456 \\
A2212 w/o M4 & 36 & 0.814 & 0.754 & 0.920 & 0.685 \\
\bottomrule
\end{tabular}
\end{table}


\begin{figure}[t]
\centering
\includegraphics[width=\columnwidth]{\FigBalacc}
\caption{Balanced accuracy on the same windows with \VDataCiLevel{}\% group-bootstrap
intervals. Dotted line: chance. *Intentional leakage. On IDMT class = motor = recording; on A2212
Bearing exists only on motor M4 (confounded with motor identity), so motor-out Bearing recall is
zero by construction.}
\label{fig:balacc}
\end{figure}

Table~\ref{tab:desempenho-compacto} and
Fig.~\ref{fig:balacc} compare the methods. On IDMT every method is near the ceiling under both
protocols: the RMS-normalised level baseline, which only sees the crest factor, reaches
\VDiagBalRmsNormIdmtBlk{} balanced accuracy on held-out blocks, and the pipeline and the
classical features reach \VDiagBalDecIdmtBlk{} and \VDiagBalClsIdmtBlk{}. Because each class is one recording of one motor,
these numbers measure recording identity and cannot rank the methods.

On A2212 with a grouped split by file, the classical features reach \VDiagBalClsBldcFile{}
[\VDiagBalClsBldcFileLo{}, \VDiagBalClsBldcFileHi{}], the pipeline \VDiagBalDecBldcFile{}
[\VDiagBalDecBldcFileLo{}, \VDiagBalDecBldcFileHi{}] and the raw level baseline
\VDiagBalRmsRawBldcFile{}. The pipeline errs mostly on Healthy (recall
\VDiagRecDecBldcFileHealthy{}). The random split inflates balanced accuracy by
\VDiagInflDecBldc{} for the pipeline and \VDiagInflClsBldc{} for the classical features. Holding a
motor out lowers every method: the classical features fall to \VDiagBalClsBldcMotor{} and the
pipeline to \VDiagBalDecBldcMotor{}; Bearing recall is \VDiagRecDecBldcMotorBearing{} by
construction, since the only bearing motor is then absent from training. Without the M4 fold,
on Healthy versus Propeller of unseen motors, the classical features reach
\VDiagBalClsBldcNoMfour{}, the RMS-normalised level baseline \VDiagBalRmsNormBldcNoMfour{} and the
pipeline \VDiagBalDecBldcNoMfour{}. Compared like for like (Healthy and Propeller of motors
M1--M3 under the file split), holding the motor out costs \VDiagDropHPClsBldc{} for the classical
features and \VDiagDropHPDecBldc{} for the pipeline; the much larger drop of the full motor-out
protocol is dominated by the Bearing class losing its only training motor. The pipeline never
outperforms the classical features (it ties at the IDMT ceiling), and with a motor held out its
point estimate also falls below the level baseline, with overlapping intervals.

% matriz de confusão: figures/pdf (FigConf), fora do texto por espaço

\subsection{Interpretability: the tracked rotation}

\ifwip
On IDMT the median tracked rotation differs by class (\VRotIdmtGoodMed{},
\VRotIdmtBrokenMed{}, \VRotIdmtHeavyMed{}\,Hz), but class equals motor there. On A2212 its rank
correlation with the speed label of the files is \VRotBldcSpearman{} pooled over motors; within
motors it is \VRotBldcSpearmanMOne{} to \VRotBldcSpearmanMFour{} but significant only for M4
(\VRotBldcFilesMFour{} files), so it should not be read as shaft speed on that dataset.
\else
The tracker exposes a physical quantity that a technician can inspect. On IDMT, the median tracked rotation is \VRotIdmtGoodMed{},
\VRotIdmtBrokenMed{} and \VRotIdmtHeavyMed{}\,Hz for good, broken and heavy-load windows, with
inter-quartile ranges of at most \VRotIdmtHeavyIqr{}\,Hz; \VRotIdmtGoodOctPct{}\% of the good
windows lie above \VRotOctThr{} times the class median (all at \VRotOctRatioMin{}--\VRotOctRatioMax{}
times), an octave error. Because class equals motor on IDMT, these differences cannot be
attributed to the operating state. On A2212, pooled over the \VDataBldcRecordings{} files, the rank
correlation between the tracked rotation and the speed label in the file names is
\VRotBldcSpearman{}; within motors it is \VRotBldcSpearmanMOne{} to \VRotBldcSpearmanMFour{}
(\VRotBldcFilesMThree{}--\VRotBldcFilesMTwo{} files each), significant only for M4
($p \geq \audRotBldcSpearmanPMTwo{}$ for M1--M3). The absolute level
is not comparable across motors and \VRotBldcCeilPct{}\% of the windows lie within
\VRotCeilPctThr{}\% of the upper limit of the search range. The range and the grid step were chosen
on the IDMT training data, so on A2212 the tracker output should not be read as the shaft speed.
\fi

% espectro: figures/pdf (FigSpec), fora do texto por espaço

\ifwip\else
\subsection{Streaming decision}

On a stream of \VPcSeqWindows{} clean A2212 golden windows (\VPcSeqPerClass{} per class:
Healthy, Propeller, Bearing) the online HMM with change penalty \VPcSeqCp{} detected both changes
with a delay of \VPcSeqDelayC{} windows in C and \VPcSeqDelayPy{} in Python on the PC (maximum
probability difference \VPcSeqMaxProbDiff{}).
\ifchipmedido On the ESP32, the same stream read from flash by \texttt{\VChipVariant} gave a delay
of \VChipSeqDelay{} windows and online decisions equal to Python's on \VChipSeqSame{} of
\VChipSeqWindows{} windows; one HMM step takes \VChipSeqHmmUs{}\,$\mu$s.\fi{} Bearing is motor M4,
so this checks the implementation, not the detector.
\fi

% ======================================================================
\section{Discussion}

\emph{What the datasets allow us to conclude.} IDMT, used without its test set, supports no
statement about operating-state discrimination: one recording per class makes class, motor and
session the same variable, and a two-number level baseline is as good as any model. A2212 with a
file-level split supports statements about the state of motors that were \emph{already seen} in
training, and there a cheap classical feature set is better than our pipeline. Generalisation to a
new motor is the relevant question for a deployed monitor, and with \VDataBldcMotors{} motors, a
bearing fault present on one of them and speed ranges that differ by class~\cite{estacio2023}, the
data can only show that every method degrades; it cannot rank methods or estimate how a monitor
would behave on another motor. None of these experiments concerns failure prediction, and none
involves industrial motors. In our runs the random split inflated A2212 balanced accuracy only
moderately (\VDiagInflClsBldc{} to \VDiagInflDecBldc{}), and holding a motor out cost, like for
like on Healthy versus Propeller, a similar amount (\VDiagDropHPClsBldc{} for the classical features,
\VDiagDropHPDecBldc{} for the pipeline); the large drop of the full motor-out protocol is dominated
by the Bearing class, which exists only on M4. Published near-perfect accuracies on these datasets
should be read against that confound, which no split by segment or by file removes.

\ifwip\else
\emph{What the embedded work shows.} The C port reproduces the reference stage by stage within
tolerances that were frozen before the formal comparison, including the discrete Viterbi path and
the HMM decisions; the residual float32 failures occur only on a synthetic edge case whose
modulation peak lies at the float32 noise floor, they leave the decision unchanged, and they
disappear when the envelope filter state is kept in double.
\ifchipmedido
On a board without PSRAM, \texttt{\VChipVariant} runs \VParSegBldc{}-s windows with on-chip outputs
within those tolerances, but at a real-time factor of \VChipRtfBldc{}, and the \VParSegIdmt{}-s
window does not fit. We claim neither real-time operation nor low power consumption.
\else
Whether it runs in real time on the chip is still open: the PC timings do not transfer.
\fi
\fi

\emph{Limitations and future work.} Energy was not measured (no current sensor), and we do not
estimate it from datasheets. All audio came from the datasets through the serial link; a real
I2S microphone, a motor on a test bench with a tachometer to validate the rotation estimate, and
the official IDMT test set (unopened by design) are future work. An independent scientific audit
of datasets, equivalence and traceability (report in the repository) was carried out before the
on-chip measurements and does not cover them.

% ======================================================================
\section{Conclusion}

A music-transcription-inspired pipeline for acoustic motor-state classification was ported to
portable C and verified against its Python reference stage by stage, with tolerances frozen
before the formal comparison: all checks pass in double; in float32 every window of real audio
passes every stage, and only \VEqFloatBadCases{} synthetic edge windows fail the modulation,
feature and z-score checks, with an unchanged decision. On an ESP32 without PSRAM the reference
layout does not fit; the named variant \texttt{\VChipVariant} fits only \VParSegBldc{}-s windows,
matches the reference within the frozen tolerances on the chip, and is \VChipRtfBldc{} times slower
than real time, dominated by the table-free harmonic sum in emulated double precision; a float32
variant of that step or a board with PSRAM is the way forward. On the two public
datasets it does not beat simple baselines, and the datasets themselves cannot separate operating
state from motor identity. The value of the work lies in the verified, measurable implementation and
in making these limits explicit.

\bibliographystyle{IEEEtran}
\bibliography{references}

\end{document}
"
\ifdefined\usewip\wiptrue\else\wipfalse\fi

% Gerado por embarcados/src/python/gerar_artigo_numeros.py. NÃO EDITAR À MÃO.
% commit 9df3b2c280b69b02703d2c61d277e2bdf02cafad (árvore suja em embarcados/: True)
% Uma macro por número. Comentário: [categoria] fonte :: coluna.
% Prefixos: VPar código | VData/VDiag/VRot dados e validação | VEq equivalência (PC)
%           VEst 1·estimativa analítica | VPc 2·medição no PC | VBuild 3·build do firmware | VChip 4·medição no ESP32
\newif\ifchipmedido
\chipmedidofalse  % true só se results/manifests/bancada.json e esp32_lat_* existirem

\newcommand{\VParSr}{16{,}000}% [codigo] pipelines/_decifrae.py :: PreDecifrae.sr
\newcommand{\VParSrKhz}{16}% [codigo] pipelines/_decifrae.py :: PreDecifrae.sr/1000
\newcommand{\VParNfft}{8{,}192}% [codigo] pipelines/_decifrae.py :: PreDecifrae.n_fft
\newcommand{\VParHop}{1{,}024}% [codigo] pipelines/_decifrae.py :: PreDecifrae.hop
\newcommand{\VParFmin}{15}% [codigo] pipelines/_decifrae.py :: PreDecifrae.fmin
\newcommand{\VParFmax}{60}% [codigo] pipelines/_decifrae.py :: PreDecifrae.fmax
\newcommand{\VParStepCents}{2}% [codigo] pipelines/_decifrae.py :: PreDecifrae.passo_cents
\newcommand{\VParJ}{1{,}201}% [codigo] motorsound/rotacao.py :: len(rotacao.grade(fmin,fmax,passo))
\newcommand{\VParNbins}{4{,}097}% [codigo] motorsound/rotacao.py :: n_fft//2+1 (rfft)
\newcommand{\VParDf}{1.95}% [codigo] motorsound/rotacao.py :: sr/n_fft
\newcommand{\VParNharm}{48}% [codigo] pipelines/_decifrae.py :: N_HARM
\newcommand{\VParGroupOneLo}{1}% [codigo] pipelines/_decifrae.py :: GRUPOS[0][0]
\newcommand{\VParGroupOneHi}{4}% [codigo] pipelines/_decifrae.py :: GRUPOS[0][1]
\newcommand{\VParGroupTwoLo}{5}% [codigo] pipelines/_decifrae.py :: GRUPOS[1][0]
\newcommand{\VParGroupTwoHi}{12}% [codigo] pipelines/_decifrae.py :: GRUPOS[1][1]
\newcommand{\VParGroupThreeLo}{13}% [codigo] pipelines/_decifrae.py :: GRUPOS[2][0]
\newcommand{\VParGroupThreeHi}{24}% [codigo] pipelines/_decifrae.py :: GRUPOS[2][1]
\newcommand{\VParGroupFourLo}{25}% [codigo] pipelines/_decifrae.py :: GRUPOS[3][0]
\newcommand{\VParGroupFourHi}{48}% [codigo] pipelines/_decifrae.py :: GRUPOS[3][1]
\newcommand{\VParMaxJump}{50}% [codigo] pipelines/_decifrae.py :: PreDecifrae.max_salto
\newcommand{\VParMaxJumpCents}{100}% [codigo] pipelines/_decifrae.py :: max_salto*passo_cents
\newcommand{\VParOctStates}{600}% [codigo] motorsound/rotacao.py :: viterbi_rotacao: k_oct = round(1200/passo)
\newcommand{\VParWStep}{0.55}% [codigo] motorsound/rotacao.py :: VIT_W_STEP
\newcommand{\VParWOct}{7.0}% [codigo] motorsound/rotacao.py :: VIT_W_OCTAVE
\newcommand{\VParWSil}{3.2}% [codigo] motorsound/rotacao.py :: VIT_W_SIL
\newcommand{\VParWEnergy}{0.45}% [codigo] motorsound/rotacao.py :: emissao_baixo: peso de e_norm
\newcommand{\VParEmisEps}{\ensuremath{10^{-3}}}% [codigo] motorsound/rotacao.py :: emissao_baixo: eps do -log
\newcommand{\VParTempCos}{0.1}% [codigo] pipelines/_decifrae.py :: PostDecifrae.temp_cos
\newcommand{\VParButterOrder}{4}% [codigo] pipelines/_decifrae.py :: _envelope_mod: butter(ordem, fc)
\newcommand{\VParButterFc}{200}% [codigo] pipelines/_decifrae.py :: _envelope_mod: butter(ordem, fc)
\newcommand{\VParModLo}{2}% [codigo] pipelines/_decifrae.py :: _envelope_mod: banda
\newcommand{\VParModHi}{100}% [codigo] pipelines/_decifrae.py :: _envelope_mod: banda
\newcommand{\VParEnvRateHz}{1{,}000}% [codigo] pipelines/_decifrae.py :: _envelope_mod: [:: sr//1000]
\newcommand{\VParDecim}{16}% [codigo] pipelines/_decifrae.py :: sr // 1000
\newcommand{\VParBandLo}{10}% [codigo] pipelines/_decifrae.py :: extrair: banda das features
\newcommand{\VParBandHi}{4{,}000}% [codigo] pipelines/_decifrae.py :: extrair: banda das features
\newcommand{\VParSegIdmt}{3}% [codigo] motorsound/config.py :: AVAL_IDMT.seg_s
\newcommand{\VParSegBldc}{1}% [codigo] motorsound/config.py :: AVAL_A2212.seg_s
\newcommand{\VParRmsTarget}{0.1}% [codigo] motorsound/config.py :: AvaliacaoConfig.rms_alvo
\newcommand{\VParFramesIdmt}{39}% [codigo] motorsound/rotacao.py :: 1+(N-n_fft)//hop (stft_mag)
\newcommand{\VParFramesBldc}{8}% [codigo] motorsound/rotacao.py :: 1+(N-n_fft)//hop (stft_mag)
\newcommand{\VParNClassicFeat}{32}% [codigo] pipelines/_classicas.py :: N_FEATURES
\newcommand{\VParNChangePen}{6}% [codigo] embarcados/gold/seqs/idmt_seq00/params.json :: change_penalties
\newcommand{\VParCpMin}{0.5}% [codigo] embarcados/gold/seqs/idmt_seq00/params.json :: change_penalties
\newcommand{\VParCpMax}{16}% [codigo] embarcados/gold/seqs/idmt_seq00/params.json :: change_penalties
\newcommand{\VParNLogprob}{-50}% [codigo] embarcados/gold/seqs/idmt_seq00/params.json :: n_logprob
\newcommand{\VParHmmTemp}{1}% [codigo] embarcados/gold/seqs/idmt_seq00/params.json :: temp
\newcommand{\VDataIdmtWindows}{357}% [dados] embarcados/results/predictions/diagnostico_janelas_20261023_222831.csv :: len(unique(rec_id,idx)) da p11
\newcommand{\VDataIdmtRecordings}{3}% [dados] embarcados/results/predictions/diagnostico_janelas_20261023_222831.csv :: nunique(rec_id)
\newcommand{\VDataIdmtMotors}{3}% [dados] embarcados/results/predictions/diagnostico_janelas_20261023_222831.csv :: nunique(motor)
\newcommand{\VDataIdmtClasses}{3}% [dados] embarcados/results/predictions/diagnostico_janelas_20261023_222831.csv :: nunique(classe)
\newcommand{\VDataBldcWindows}{430}% [dados] embarcados/results/predictions/diagnostico_janelas_20261023_222831.csv :: len(unique(rec_id,idx)) da p11
\newcommand{\VDataBldcRecordings}{43}% [dados] embarcados/results/predictions/diagnostico_janelas_20261023_222831.csv :: nunique(rec_id)
\newcommand{\VDataBldcMotors}{4}% [dados] embarcados/results/predictions/diagnostico_janelas_20261023_222831.csv :: nunique(motor)
\newcommand{\VDataBldcClasses}{3}% [dados] embarcados/results/predictions/diagnostico_janelas_20261023_222831.csv :: nunique(classe)
\newcommand{\VDataBldcFilesHealthy}{22}% [dados] embarcados/results/predictions/diagnostico_janelas_20261023_222831.csv :: nunique(rec_id) por classe
\newcommand{\VDataBldcFilesPropeller}{14}% [dados] embarcados/results/predictions/diagnostico_janelas_20261023_222831.csv :: nunique(rec_id) por classe
\newcommand{\VDataBldcFilesBearing}{7}% [dados] embarcados/results/predictions/diagnostico_janelas_20261023_222831.csv :: nunique(rec_id) por classe
\newcommand{\VDataIdmtRandGroups}{18}% [dados] embarcados/results/csv/diagnostico_20261023_222831.csv :: n_grupos
\newcommand{\VDataIdmtBlkGroups}{15}% [dados] embarcados/results/csv/diagnostico_20261023_222831.csv :: n_grupos
\newcommand{\VDataBldcRandGroups}{43}% [dados] embarcados/results/csv/diagnostico_20261023_222831.csv :: n_grupos
\newcommand{\VDataBldcFileGroups}{43}% [dados] embarcados/results/csv/diagnostico_20261023_222831.csv :: n_grupos
\newcommand{\VDataBldcMotorGroups}{43}% [dados] embarcados/results/csv/diagnostico_20261023_222831.csv :: n_grupos
\newcommand{\VDataBldcNoMfourGroups}{36}% [dados] embarcados/results/csv/diagnostico_20261023_222831.csv :: n_grupos
\newcommand{\VDataNBoot}{2{,}000}% [dados] embarcados/results/csv/diagnostico_20261023_222831_meta.json :: n_boot
\newcommand{\VDataCiLevel}{95}% [dados] embarcados/results/csv/diagnostico_20261023_222831.csv :: nível do IC no nome da coluna ic<nível>_lo
\newcommand{\VDataSeedBoot}{20261023}% [dados] embarcados/results/csv/diagnostico_20261023_222831_meta.json :: seed_boot
\newcommand{\VDiagBalRmsNormIdmtRand}{0.992}% [diagnostico] embarcados/results/csv/diagnostico_20261023_222831.csv :: valor[b00_rmspico_norm,idmt,aleatorio,acuracia_balanceada,_macro]
\newcommand{\VDiagBalRmsNormIdmtRandLo}{0.983}% [diagnostico] embarcados/results/csv/diagnostico_20261023_222831.csv :: ic95_lo[b00_rmspico_norm,idmt,aleatorio,acuracia_balanceada,_macro]
\newcommand{\VDiagBalRmsNormIdmtRandHi}{1.000}% [diagnostico] embarcados/results/csv/diagnostico_20261023_222831.csv :: ic95_hi[b00_rmspico_norm,idmt,aleatorio,acuracia_balanceada,_macro]
\newcommand{\VDiagAucRmsNormIdmtRand}{0.998}% [diagnostico] embarcados/results/csv/diagnostico_20261023_222831.csv :: valor[b00_rmspico_norm,idmt,aleatorio,roc_auc_ovr_macro,_macro]
\newcommand{\VDiagFoneRmsNormIdmtRand}{0.992}% [diagnostico] embarcados/results/csv/diagnostico_20261023_222831.csv :: valor[b00_rmspico_norm,idmt,aleatorio,f1_macro,_macro]
\newcommand{\VDiagBalRmsRawIdmtRand}{0.997}% [diagnostico] embarcados/results/csv/diagnostico_20261023_222831.csv :: valor[b00_rmspico_bruto,idmt,aleatorio,acuracia_balanceada,_macro]
\newcommand{\VDiagBalRmsRawIdmtRandLo}{0.990}% [diagnostico] embarcados/results/csv/diagnostico_20261023_222831.csv :: ic95_lo[b00_rmspico_bruto,idmt,aleatorio,acuracia_balanceada,_macro]
\newcommand{\VDiagBalRmsRawIdmtRandHi}{1.000}% [diagnostico] embarcados/results/csv/diagnostico_20261023_222831.csv :: ic95_hi[b00_rmspico_bruto,idmt,aleatorio,acuracia_balanceada,_macro]
\newcommand{\VDiagAucRmsRawIdmtRand}{1.000}% [diagnostico] embarcados/results/csv/diagnostico_20261023_222831.csv :: valor[b00_rmspico_bruto,idmt,aleatorio,roc_auc_ovr_macro,_macro]
\newcommand{\VDiagFoneRmsRawIdmtRand}{0.997}% [diagnostico] embarcados/results/csv/diagnostico_20261023_222831.csv :: valor[b00_rmspico_bruto,idmt,aleatorio,f1_macro,_macro]
\newcommand{\VDiagBalClsIdmtRand}{1.000}% [diagnostico] embarcados/results/csv/diagnostico_20261023_222831.csv :: valor[p04_classicas,idmt,aleatorio,acuracia_balanceada,_macro]
\newcommand{\VDiagBalClsIdmtRandLo}{1.000}% [diagnostico] embarcados/results/csv/diagnostico_20261023_222831.csv :: ic95_lo[p04_classicas,idmt,aleatorio,acuracia_balanceada,_macro]
\newcommand{\VDiagBalClsIdmtRandHi}{1.000}% [diagnostico] embarcados/results/csv/diagnostico_20261023_222831.csv :: ic95_hi[p04_classicas,idmt,aleatorio,acuracia_balanceada,_macro]
\newcommand{\VDiagAucClsIdmtRand}{1.000}% [diagnostico] embarcados/results/csv/diagnostico_20261023_222831.csv :: valor[p04_classicas,idmt,aleatorio,roc_auc_ovr_macro,_macro]
\newcommand{\VDiagFoneClsIdmtRand}{1.000}% [diagnostico] embarcados/results/csv/diagnostico_20261023_222831.csv :: valor[p04_classicas,idmt,aleatorio,f1_macro,_macro]
\newcommand{\VDiagBalDecIdmtRand}{1.000}% [diagnostico] embarcados/results/csv/diagnostico_20261023_222831.csv :: valor[p11_decifrae_harmsum,idmt,aleatorio,acuracia_balanceada,_macro]
\newcommand{\VDiagBalDecIdmtRandLo}{1.000}% [diagnostico] embarcados/results/csv/diagnostico_20261023_222831.csv :: ic95_lo[p11_decifrae_harmsum,idmt,aleatorio,acuracia_balanceada,_macro]
\newcommand{\VDiagBalDecIdmtRandHi}{1.000}% [diagnostico] embarcados/results/csv/diagnostico_20261023_222831.csv :: ic95_hi[p11_decifrae_harmsum,idmt,aleatorio,acuracia_balanceada,_macro]
\newcommand{\VDiagAucDecIdmtRand}{1.000}% [diagnostico] embarcados/results/csv/diagnostico_20261023_222831.csv :: valor[p11_decifrae_harmsum,idmt,aleatorio,roc_auc_ovr_macro,_macro]
\newcommand{\VDiagFoneDecIdmtRand}{1.000}% [diagnostico] embarcados/results/csv/diagnostico_20261023_222831.csv :: valor[p11_decifrae_harmsum,idmt,aleatorio,f1_macro,_macro]
\newcommand{\VDiagBalRmsNormIdmtBlk}{0.991}% [diagnostico] embarcados/results/csv/diagnostico_20261023_222831.csv :: valor[b00_rmspico_norm,idmt,blocos,acuracia_balanceada,_macro]
\newcommand{\VDiagBalRmsNormIdmtBlkLo}{0.982}% [diagnostico] embarcados/results/csv/diagnostico_20261023_222831.csv :: ic95_lo[b00_rmspico_norm,idmt,blocos,acuracia_balanceada,_macro]
\newcommand{\VDiagBalRmsNormIdmtBlkHi}{0.997}% [diagnostico] embarcados/results/csv/diagnostico_20261023_222831.csv :: ic95_hi[b00_rmspico_norm,idmt,blocos,acuracia_balanceada,_macro]
\newcommand{\VDiagAucRmsNormIdmtBlk}{0.998}% [diagnostico] embarcados/results/csv/diagnostico_20261023_222831.csv :: valor[b00_rmspico_norm,idmt,blocos,roc_auc_ovr_macro,_macro]
\newcommand{\VDiagFoneRmsNormIdmtBlk}{0.991}% [diagnostico] embarcados/results/csv/diagnostico_20261023_222831.csv :: valor[b00_rmspico_norm,idmt,blocos,f1_macro,_macro]
\newcommand{\VDiagBalRmsRawIdmtBlk}{0.993}% [diagnostico] embarcados/results/csv/diagnostico_20261023_222831.csv :: valor[b00_rmspico_bruto,idmt,blocos,acuracia_balanceada,_macro]
\newcommand{\VDiagBalRmsRawIdmtBlkLo}{0.987}% [diagnostico] embarcados/results/csv/diagnostico_20261023_222831.csv :: ic95_lo[b00_rmspico_bruto,idmt,blocos,acuracia_balanceada,_macro]
\newcommand{\VDiagBalRmsRawIdmtBlkHi}{1.000}% [diagnostico] embarcados/results/csv/diagnostico_20261023_222831.csv :: ic95_hi[b00_rmspico_bruto,idmt,blocos,acuracia_balanceada,_macro]
\newcommand{\VDiagAucRmsRawIdmtBlk}{1.000}% [diagnostico] embarcados/results/csv/diagnostico_20261023_222831.csv :: valor[b00_rmspico_bruto,idmt,blocos,roc_auc_ovr_macro,_macro]
\newcommand{\VDiagFoneRmsRawIdmtBlk}{0.994}% [diagnostico] embarcados/results/csv/diagnostico_20261023_222831.csv :: valor[b00_rmspico_bruto,idmt,blocos,f1_macro,_macro]
\newcommand{\VDiagBalClsIdmtBlk}{1.000}% [diagnostico] embarcados/results/csv/diagnostico_20261023_222831.csv :: valor[p04_classicas,idmt,blocos,acuracia_balanceada,_macro]
\newcommand{\VDiagBalClsIdmtBlkLo}{1.000}% [diagnostico] embarcados/results/csv/diagnostico_20261023_222831.csv :: ic95_lo[p04_classicas,idmt,blocos,acuracia_balanceada,_macro]
\newcommand{\VDiagBalClsIdmtBlkHi}{1.000}% [diagnostico] embarcados/results/csv/diagnostico_20261023_222831.csv :: ic95_hi[p04_classicas,idmt,blocos,acuracia_balanceada,_macro]
\newcommand{\VDiagAucClsIdmtBlk}{1.000}% [diagnostico] embarcados/results/csv/diagnostico_20261023_222831.csv :: valor[p04_classicas,idmt,blocos,roc_auc_ovr_macro,_macro]
\newcommand{\VDiagFoneClsIdmtBlk}{1.000}% [diagnostico] embarcados/results/csv/diagnostico_20261023_222831.csv :: valor[p04_classicas,idmt,blocos,f1_macro,_macro]
\newcommand{\VDiagBalDecIdmtBlk}{1.000}% [diagnostico] embarcados/results/csv/diagnostico_20261023_222831.csv :: valor[p11_decifrae_harmsum,idmt,blocos,acuracia_balanceada,_macro]
\newcommand{\VDiagBalDecIdmtBlkLo}{1.000}% [diagnostico] embarcados/results/csv/diagnostico_20261023_222831.csv :: ic95_lo[p11_decifrae_harmsum,idmt,blocos,acuracia_balanceada,_macro]
\newcommand{\VDiagBalDecIdmtBlkHi}{1.000}% [diagnostico] embarcados/results/csv/diagnostico_20261023_222831.csv :: ic95_hi[p11_decifrae_harmsum,idmt,blocos,acuracia_balanceada,_macro]
\newcommand{\VDiagAucDecIdmtBlk}{1.000}% [diagnostico] embarcados/results/csv/diagnostico_20261023_222831.csv :: valor[p11_decifrae_harmsum,idmt,blocos,roc_auc_ovr_macro,_macro]
\newcommand{\VDiagFoneDecIdmtBlk}{1.000}% [diagnostico] embarcados/results/csv/diagnostico_20261023_222831.csv :: valor[p11_decifrae_harmsum,idmt,blocos,f1_macro,_macro]
\newcommand{\VDiagBalRmsNormBldcRand}{0.564}% [diagnostico] embarcados/results/csv/diagnostico_20261023_222831.csv :: valor[b00_rmspico_norm,a2212,aleatorio,acuracia_balanceada,_macro]
\newcommand{\VDiagBalRmsNormBldcRandLo}{0.515}% [diagnostico] embarcados/results/csv/diagnostico_20261023_222831.csv :: ic95_lo[b00_rmspico_norm,a2212,aleatorio,acuracia_balanceada,_macro]
\newcommand{\VDiagBalRmsNormBldcRandHi}{0.607}% [diagnostico] embarcados/results/csv/diagnostico_20261023_222831.csv :: ic95_hi[b00_rmspico_norm,a2212,aleatorio,acuracia_balanceada,_macro]
\newcommand{\VDiagAucRmsNormBldcRand}{0.784}% [diagnostico] embarcados/results/csv/diagnostico_20261023_222831.csv :: valor[b00_rmspico_norm,a2212,aleatorio,roc_auc_ovr_macro,_macro]
\newcommand{\VDiagFoneRmsNormBldcRand}{0.533}% [diagnostico] embarcados/results/csv/diagnostico_20261023_222831.csv :: valor[b00_rmspico_norm,a2212,aleatorio,f1_macro,_macro]
\newcommand{\VDiagBalRmsRawBldcRand}{0.625}% [diagnostico] embarcados/results/csv/diagnostico_20261023_222831.csv :: valor[b00_rmspico_bruto,a2212,aleatorio,acuracia_balanceada,_macro]
\newcommand{\VDiagBalRmsRawBldcRandLo}{0.571}% [diagnostico] embarcados/results/csv/diagnostico_20261023_222831.csv :: ic95_lo[b00_rmspico_bruto,a2212,aleatorio,acuracia_balanceada,_macro]
\newcommand{\VDiagBalRmsRawBldcRandHi}{0.660}% [diagnostico] embarcados/results/csv/diagnostico_20261023_222831.csv :: ic95_hi[b00_rmspico_bruto,a2212,aleatorio,acuracia_balanceada,_macro]
\newcommand{\VDiagAucRmsRawBldcRand}{0.786}% [diagnostico] embarcados/results/csv/diagnostico_20261023_222831.csv :: valor[b00_rmspico_bruto,a2212,aleatorio,roc_auc_ovr_macro,_macro]
\newcommand{\VDiagFoneRmsRawBldcRand}{0.588}% [diagnostico] embarcados/results/csv/diagnostico_20261023_222831.csv :: valor[b00_rmspico_bruto,a2212,aleatorio,f1_macro,_macro]
\newcommand{\VDiagBalClsBldcRand}{0.998}% [diagnostico] embarcados/results/csv/diagnostico_20261023_222831.csv :: valor[p04_classicas,a2212,aleatorio,acuracia_balanceada,_macro]
\newcommand{\VDiagBalClsBldcRandLo}{0.995}% [diagnostico] embarcados/results/csv/diagnostico_20261023_222831.csv :: ic95_lo[p04_classicas,a2212,aleatorio,acuracia_balanceada,_macro]
\newcommand{\VDiagBalClsBldcRandHi}{1.000}% [diagnostico] embarcados/results/csv/diagnostico_20261023_222831.csv :: ic95_hi[p04_classicas,a2212,aleatorio,acuracia_balanceada,_macro]
\newcommand{\VDiagAucClsBldcRand}{1.000}% [diagnostico] embarcados/results/csv/diagnostico_20261023_222831.csv :: valor[p04_classicas,a2212,aleatorio,roc_auc_ovr_macro,_macro]
\newcommand{\VDiagFoneClsBldcRand}{0.998}% [diagnostico] embarcados/results/csv/diagnostico_20261023_222831.csv :: valor[p04_classicas,a2212,aleatorio,f1_macro,_macro]
\newcommand{\VDiagBalDecBldcRand}{0.849}% [diagnostico] embarcados/results/csv/diagnostico_20261023_222831.csv :: valor[p11_decifrae_harmsum,a2212,aleatorio,acuracia_balanceada,_macro]
\newcommand{\VDiagBalDecBldcRandLo}{0.790}% [diagnostico] embarcados/results/csv/diagnostico_20261023_222831.csv :: ic95_lo[p11_decifrae_harmsum,a2212,aleatorio,acuracia_balanceada,_macro]
\newcommand{\VDiagBalDecBldcRandHi}{0.905}% [diagnostico] embarcados/results/csv/diagnostico_20261023_222831.csv :: ic95_hi[p11_decifrae_harmsum,a2212,aleatorio,acuracia_balanceada,_macro]
\newcommand{\VDiagAucDecBldcRand}{0.955}% [diagnostico] embarcados/results/csv/diagnostico_20261023_222831.csv :: valor[p11_decifrae_harmsum,a2212,aleatorio,roc_auc_ovr_macro,_macro]
\newcommand{\VDiagFoneDecBldcRand}{0.789}% [diagnostico] embarcados/results/csv/diagnostico_20261023_222831.csv :: valor[p11_decifrae_harmsum,a2212,aleatorio,f1_macro,_macro]
\newcommand{\VDiagBalRmsNormBldcFile}{0.562}% [diagnostico] embarcados/results/csv/diagnostico_20261023_222831.csv :: valor[b00_rmspico_norm,a2212,arquivo,acuracia_balanceada,_macro]
\newcommand{\VDiagBalRmsNormBldcFileLo}{0.512}% [diagnostico] embarcados/results/csv/diagnostico_20261023_222831.csv :: ic95_lo[b00_rmspico_norm,a2212,arquivo,acuracia_balanceada,_macro]
\newcommand{\VDiagBalRmsNormBldcFileHi}{0.606}% [diagnostico] embarcados/results/csv/diagnostico_20261023_222831.csv :: ic95_hi[b00_rmspico_norm,a2212,arquivo,acuracia_balanceada,_macro]
\newcommand{\VDiagAucRmsNormBldcFile}{0.759}% [diagnostico] embarcados/results/csv/diagnostico_20261023_222831.csv :: valor[b00_rmspico_norm,a2212,arquivo,roc_auc_ovr_macro,_macro]
\newcommand{\VDiagFoneRmsNormBldcFile}{0.532}% [diagnostico] embarcados/results/csv/diagnostico_20261023_222831.csv :: valor[b00_rmspico_norm,a2212,arquivo,f1_macro,_macro]
\newcommand{\VDiagBalRmsRawBldcFile}{0.608}% [diagnostico] embarcados/results/csv/diagnostico_20261023_222831.csv :: valor[b00_rmspico_bruto,a2212,arquivo,acuracia_balanceada,_macro]
\newcommand{\VDiagBalRmsRawBldcFileLo}{0.544}% [diagnostico] embarcados/results/csv/diagnostico_20261023_222831.csv :: ic95_lo[b00_rmspico_bruto,a2212,arquivo,acuracia_balanceada,_macro]
\newcommand{\VDiagBalRmsRawBldcFileHi}{0.656}% [diagnostico] embarcados/results/csv/diagnostico_20261023_222831.csv :: ic95_hi[b00_rmspico_bruto,a2212,arquivo,acuracia_balanceada,_macro]
\newcommand{\VDiagAucRmsRawBldcFile}{0.706}% [diagnostico] embarcados/results/csv/diagnostico_20261023_222831.csv :: valor[b00_rmspico_bruto,a2212,arquivo,roc_auc_ovr_macro,_macro]
\newcommand{\VDiagFoneRmsRawBldcFile}{0.576}% [diagnostico] embarcados/results/csv/diagnostico_20261023_222831.csv :: valor[b00_rmspico_bruto,a2212,arquivo,f1_macro,_macro]
\newcommand{\VDiagBalClsBldcFile}{0.981}% [diagnostico] embarcados/results/csv/diagnostico_20261023_222831.csv :: valor[p04_classicas,a2212,arquivo,acuracia_balanceada,_macro]
\newcommand{\VDiagBalClsBldcFileLo}{0.950}% [diagnostico] embarcados/results/csv/diagnostico_20261023_222831.csv :: ic95_lo[p04_classicas,a2212,arquivo,acuracia_balanceada,_macro]
\newcommand{\VDiagBalClsBldcFileHi}{0.998}% [diagnostico] embarcados/results/csv/diagnostico_20261023_222831.csv :: ic95_hi[p04_classicas,a2212,arquivo,acuracia_balanceada,_macro]
\newcommand{\VDiagAucClsBldcFile}{0.996}% [diagnostico] embarcados/results/csv/diagnostico_20261023_222831.csv :: valor[p04_classicas,a2212,arquivo,roc_auc_ovr_macro,_macro]
\newcommand{\VDiagFoneClsBldcFile}{0.982}% [diagnostico] embarcados/results/csv/diagnostico_20261023_222831.csv :: valor[p04_classicas,a2212,arquivo,f1_macro,_macro]
\newcommand{\VDiagBalDecBldcFile}{0.816}% [diagnostico] embarcados/results/csv/diagnostico_20261023_222831.csv :: valor[p11_decifrae_harmsum,a2212,arquivo,acuracia_balanceada,_macro]
\newcommand{\VDiagBalDecBldcFileLo}{0.747}% [diagnostico] embarcados/results/csv/diagnostico_20261023_222831.csv :: ic95_lo[p11_decifrae_harmsum,a2212,arquivo,acuracia_balanceada,_macro]
\newcommand{\VDiagBalDecBldcFileHi}{0.881}% [diagnostico] embarcados/results/csv/diagnostico_20261023_222831.csv :: ic95_hi[p11_decifrae_harmsum,a2212,arquivo,acuracia_balanceada,_macro]
\newcommand{\VDiagAucDecBldcFile}{0.935}% [diagnostico] embarcados/results/csv/diagnostico_20261023_222831.csv :: valor[p11_decifrae_harmsum,a2212,arquivo,roc_auc_ovr_macro,_macro]
\newcommand{\VDiagFoneDecBldcFile}{0.753}% [diagnostico] embarcados/results/csv/diagnostico_20261023_222831.csv :: valor[p11_decifrae_harmsum,a2212,arquivo,f1_macro,_macro]
\newcommand{\VDiagBalRmsNormBldcMotor}{0.542}% [diagnostico] embarcados/results/csv/diagnostico_20261023_222831.csv :: valor[b00_rmspico_norm,a2212,motor,acuracia_balanceada,_macro]
\newcommand{\VDiagBalRmsNormBldcMotorLo}{0.485}% [diagnostico] embarcados/results/csv/diagnostico_20261023_222831.csv :: ic95_lo[b00_rmspico_norm,a2212,motor,acuracia_balanceada,_macro]
\newcommand{\VDiagBalRmsNormBldcMotorHi}{0.596}% [diagnostico] embarcados/results/csv/diagnostico_20261023_222831.csv :: ic95_hi[b00_rmspico_norm,a2212,motor,acuracia_balanceada,_macro]
\newcommand{\VDiagAucRmsNormBldcMotor}{0.489}% [diagnostico] embarcados/results/csv/diagnostico_20261023_222831.csv :: valor[b00_rmspico_norm,a2212,motor,roc_auc_ovr_macro,_macro]
\newcommand{\VDiagFoneRmsNormBldcMotor}{0.509}% [diagnostico] embarcados/results/csv/diagnostico_20261023_222831.csv :: valor[b00_rmspico_norm,a2212,motor,f1_macro,_macro]
\newcommand{\VDiagBalRmsRawBldcMotor}{0.503}% [diagnostico] embarcados/results/csv/diagnostico_20261023_222831.csv :: valor[b00_rmspico_bruto,a2212,motor,acuracia_balanceada,_macro]
\newcommand{\VDiagBalRmsRawBldcMotorLo}{0.412}% [diagnostico] embarcados/results/csv/diagnostico_20261023_222831.csv :: ic95_lo[b00_rmspico_bruto,a2212,motor,acuracia_balanceada,_macro]
\newcommand{\VDiagBalRmsRawBldcMotorHi}{0.580}% [diagnostico] embarcados/results/csv/diagnostico_20261023_222831.csv :: ic95_hi[b00_rmspico_bruto,a2212,motor,acuracia_balanceada,_macro]
\newcommand{\VDiagAucRmsRawBldcMotor}{0.508}% [diagnostico] embarcados/results/csv/diagnostico_20261023_222831.csv :: valor[b00_rmspico_bruto,a2212,motor,roc_auc_ovr_macro,_macro]
\newcommand{\VDiagFoneRmsRawBldcMotor}{0.509}% [diagnostico] embarcados/results/csv/diagnostico_20261023_222831.csv :: valor[b00_rmspico_bruto,a2212,motor,f1_macro,_macro]
\newcommand{\VDiagBalClsBldcMotor}{0.613}% [diagnostico] embarcados/results/csv/diagnostico_20261023_222831.csv :: valor[p04_classicas,a2212,motor,acuracia_balanceada,_macro]
\newcommand{\VDiagBalClsBldcMotorLo}{0.579}% [diagnostico] embarcados/results/csv/diagnostico_20261023_222831.csv :: ic95_lo[p04_classicas,a2212,motor,acuracia_balanceada,_macro]
\newcommand{\VDiagBalClsBldcMotorHi}{0.641}% [diagnostico] embarcados/results/csv/diagnostico_20261023_222831.csv :: ic95_hi[p04_classicas,a2212,motor,acuracia_balanceada,_macro]
\newcommand{\VDiagAucClsBldcMotor}{0.618}% [diagnostico] embarcados/results/csv/diagnostico_20261023_222831.csv :: valor[p04_classicas,a2212,motor,roc_auc_ovr_macro,_macro]
\newcommand{\VDiagFoneClsBldcMotor}{0.563}% [diagnostico] embarcados/results/csv/diagnostico_20261023_222831.csv :: valor[p04_classicas,a2212,motor,f1_macro,_macro]
\newcommand{\VDiagBalDecBldcMotor}{0.456}% [diagnostico] embarcados/results/csv/diagnostico_20261023_222831.csv :: valor[p11_decifrae_harmsum,a2212,motor,acuracia_balanceada,_macro]
\newcommand{\VDiagBalDecBldcMotorLo}{0.397}% [diagnostico] embarcados/results/csv/diagnostico_20261023_222831.csv :: ic95_lo[p11_decifrae_harmsum,a2212,motor,acuracia_balanceada,_macro]
\newcommand{\VDiagBalDecBldcMotorHi}{0.519}% [diagnostico] embarcados/results/csv/diagnostico_20261023_222831.csv :: ic95_hi[p11_decifrae_harmsum,a2212,motor,acuracia_balanceada,_macro]
\newcommand{\VDiagAucDecBldcMotor}{0.484}% [diagnostico] embarcados/results/csv/diagnostico_20261023_222831.csv :: valor[p11_decifrae_harmsum,a2212,motor,roc_auc_ovr_macro,_macro]
\newcommand{\VDiagFoneDecBldcMotor}{0.418}% [diagnostico] embarcados/results/csv/diagnostico_20261023_222831.csv :: valor[p11_decifrae_harmsum,a2212,motor,f1_macro,_macro]
\newcommand{\VDiagBalRmsNormBldcNoMfour}{0.814}% [diagnostico] embarcados/results/csv/diagnostico_20261023_222831.csv :: valor[b00_rmspico_norm,a2212,motor_sem_M4,acuracia_balanceada,_macro]
\newcommand{\VDiagBalRmsNormBldcNoMfourLo}{0.720}% [diagnostico] embarcados/results/csv/diagnostico_20261023_222831.csv :: ic95_lo[b00_rmspico_norm,a2212,motor_sem_M4,acuracia_balanceada,_macro]
\newcommand{\VDiagBalRmsNormBldcNoMfourHi}{0.887}% [diagnostico] embarcados/results/csv/diagnostico_20261023_222831.csv :: ic95_hi[b00_rmspico_norm,a2212,motor_sem_M4,acuracia_balanceada,_macro]
\newcommand{\VDiagAucRmsNormBldcNoMfour}{0.871}% [diagnostico] embarcados/results/csv/diagnostico_20261023_222831.csv :: valor[b00_rmspico_norm,a2212,motor_sem_M4,roc_auc_ovr_macro,_macro]
\newcommand{\VDiagFoneRmsNormBldcNoMfour}{0.824}% [diagnostico] embarcados/results/csv/diagnostico_20261023_222831.csv :: valor[b00_rmspico_norm,a2212,motor_sem_M4,f1_macro,_macro]
\newcommand{\VDiagBalRmsRawBldcNoMfour}{0.754}% [diagnostico] embarcados/results/csv/diagnostico_20261023_222831.csv :: valor[b00_rmspico_bruto,a2212,motor_sem_M4,acuracia_balanceada,_macro]
\newcommand{\VDiagBalRmsRawBldcNoMfourLo}{0.623}% [diagnostico] embarcados/results/csv/diagnostico_20261023_222831.csv :: ic95_lo[b00_rmspico_bruto,a2212,motor_sem_M4,acuracia_balanceada,_macro]
\newcommand{\VDiagBalRmsRawBldcNoMfourHi}{0.873}% [diagnostico] embarcados/results/csv/diagnostico_20261023_222831.csv :: ic95_hi[b00_rmspico_bruto,a2212,motor_sem_M4,acuracia_balanceada,_macro]
\newcommand{\VDiagAucRmsRawBldcNoMfour}{0.912}% [diagnostico] embarcados/results/csv/diagnostico_20261023_222831.csv :: valor[b00_rmspico_bruto,a2212,motor_sem_M4,roc_auc_ovr_macro,_macro]
\newcommand{\VDiagFoneRmsRawBldcNoMfour}{0.823}% [diagnostico] embarcados/results/csv/diagnostico_20261023_222831.csv :: valor[b00_rmspico_bruto,a2212,motor_sem_M4,f1_macro,_macro]
\newcommand{\VDiagBalClsBldcNoMfour}{0.920}% [diagnostico] embarcados/results/csv/diagnostico_20261023_222831.csv :: valor[p04_classicas,a2212,motor_sem_M4,acuracia_balanceada,_macro]
\newcommand{\VDiagBalClsBldcNoMfourLo}{0.868}% [diagnostico] embarcados/results/csv/diagnostico_20261023_222831.csv :: ic95_lo[p04_classicas,a2212,motor_sem_M4,acuracia_balanceada,_macro]
\newcommand{\VDiagBalClsBldcNoMfourHi}{0.962}% [diagnostico] embarcados/results/csv/diagnostico_20261023_222831.csv :: ic95_hi[p04_classicas,a2212,motor_sem_M4,acuracia_balanceada,_macro]
\newcommand{\VDiagAucClsBldcNoMfour}{0.984}% [diagnostico] embarcados/results/csv/diagnostico_20261023_222831.csv :: valor[p04_classicas,a2212,motor_sem_M4,roc_auc_ovr_macro,_macro]
\newcommand{\VDiagFoneClsBldcNoMfour}{0.923}% [diagnostico] embarcados/results/csv/diagnostico_20261023_222831.csv :: valor[p04_classicas,a2212,motor_sem_M4,f1_macro,_macro]
\newcommand{\VDiagBalDecBldcNoMfour}{0.685}% [diagnostico] embarcados/results/csv/diagnostico_20261023_222831.csv :: valor[p11_decifrae_harmsum,a2212,motor_sem_M4,acuracia_balanceada,_macro]
\newcommand{\VDiagBalDecBldcNoMfourLo}{0.593}% [diagnostico] embarcados/results/csv/diagnostico_20261023_222831.csv :: ic95_lo[p11_decifrae_harmsum,a2212,motor_sem_M4,acuracia_balanceada,_macro]
\newcommand{\VDiagBalDecBldcNoMfourHi}{0.774}% [diagnostico] embarcados/results/csv/diagnostico_20261023_222831.csv :: ic95_hi[p11_decifrae_harmsum,a2212,motor_sem_M4,acuracia_balanceada,_macro]
\newcommand{\VDiagAucDecBldcNoMfour}{0.822}% [diagnostico] embarcados/results/csv/diagnostico_20261023_222831.csv :: valor[p11_decifrae_harmsum,a2212,motor_sem_M4,roc_auc_ovr_macro,_macro]
\newcommand{\VDiagFoneDecBldcNoMfour}{0.685}% [diagnostico] embarcados/results/csv/diagnostico_20261023_222831.csv :: valor[p11_decifrae_harmsum,a2212,motor_sem_M4,f1_macro,_macro]
\newcommand{\VDiagRecDecBldcFileHealthy}{0.577}% [diagnostico] embarcados/results/csv/diagnostico_20261023_222831.csv :: valor[p11,a2212,arquivo,sensibilidade,Healthy]
\newcommand{\VDiagRecDecBldcFilePropeller}{0.929}% [diagnostico] embarcados/results/csv/diagnostico_20261023_222831.csv :: valor[p11,a2212,arquivo,sensibilidade,Propeller]
\newcommand{\VDiagRecDecBldcFileBearing}{0.943}% [diagnostico] embarcados/results/csv/diagnostico_20261023_222831.csv :: valor[p11,a2212,arquivo,sensibilidade,Bearing]
\newcommand{\VDiagRecDecBldcMotorHealthy}{0.441}% [diagnostico] embarcados/results/csv/diagnostico_20261023_222831.csv :: valor[p11,a2212,motor,sensibilidade,Healthy]
\newcommand{\VDiagRecDecBldcMotorPropeller}{0.929}% [diagnostico] embarcados/results/csv/diagnostico_20261023_222831.csv :: valor[p11,a2212,motor,sensibilidade,Propeller]
\newcommand{\VDiagRecDecBldcMotorBearing}{0.000}% [diagnostico] embarcados/results/csv/diagnostico_20261023_222831.csv :: valor[p11,a2212,motor,sensibilidade,Bearing]
\newcommand{\VDiagInflRmsNormBldc}{0.002}% [diagnostico] embarcados/results/csv/diagnostico_20261023_222831.csv :: Bal[aleatorio]-Bal[arquivo] (a2212)
\newcommand{\VDiagDropRmsNormBldc}{0.019}% [diagnostico] embarcados/results/csv/diagnostico_20261023_222831.csv :: Bal[arquivo]-Bal[motor] (a2212)
\newcommand{\VDiagInflRmsRawBldc}{0.016}% [diagnostico] embarcados/results/csv/diagnostico_20261023_222831.csv :: Bal[aleatorio]-Bal[arquivo] (a2212)
\newcommand{\VDiagDropRmsRawBldc}{0.106}% [diagnostico] embarcados/results/csv/diagnostico_20261023_222831.csv :: Bal[arquivo]-Bal[motor] (a2212)
\newcommand{\VDiagInflClsBldc}{0.017}% [diagnostico] embarcados/results/csv/diagnostico_20261023_222831.csv :: Bal[aleatorio]-Bal[arquivo] (a2212)
\newcommand{\VDiagDropClsBldc}{0.368}% [diagnostico] embarcados/results/csv/diagnostico_20261023_222831.csv :: Bal[arquivo]-Bal[motor] (a2212)
\newcommand{\VDiagInflDecBldc}{0.033}% [diagnostico] embarcados/results/csv/diagnostico_20261023_222831.csv :: Bal[aleatorio]-Bal[arquivo] (a2212)
\newcommand{\VDiagDropDecBldc}{0.360}% [diagnostico] embarcados/results/csv/diagnostico_20261023_222831.csv :: Bal[arquivo]-Bal[motor] (a2212)
\newcommand{\VParNFeat}{13}% [codigo] embarcados/results/csv/p11_features_20261023_222831.csv :: nº de colunas f0..f12
\newcommand{\VRotIdmtGoodMed}{25.76}% [dados] embarcados/results/csv/p11_features_20261023_222831.csv :: mediana de f0 (idmt, good)
\newcommand{\VRotIdmtGoodIqr}{0.18}% [dados] embarcados/results/csv/p11_features_20261023_222831.csv :: IQR de f0 (idmt, good)
\newcommand{\VRotIdmtBrokenMed}{27.99}% [dados] embarcados/results/csv/p11_features_20261023_222831.csv :: mediana de f0 (idmt, broken)
\newcommand{\VRotIdmtBrokenIqr}{0.03}% [dados] embarcados/results/csv/p11_features_20261023_222831.csv :: IQR de f0 (idmt, broken)
\newcommand{\VRotIdmtHeavyMed}{24.65}% [dados] embarcados/results/csv/p11_features_20261023_222831.csv :: mediana de f0 (idmt, heavyload)
\newcommand{\VRotIdmtHeavyIqr}{0.20}% [dados] embarcados/results/csv/p11_features_20261023_222831.csv :: IQR de f0 (idmt, heavyload)
\newcommand{\VRotOctThr}{1.5}% [dados] embarcados/results/csv/p11_features_20261023_222831.csv :: limiar: f0 > LIM_OIT·mediana da classe
\newcommand{\VRotCeilPctThr}{1}% [dados] embarcados/results/csv/p11_features_20261023_222831.csv :: limiar: f0 >= LIM_TETO·fmax
\newcommand{\VRotIdmtGoodOctPct}{8.6}% [dados] embarcados/results/csv/p11_features_20261023_222831.csv :: % janelas good com f0 > LIM_OIT·mediana da classe (salto de oitava)
\newcommand{\VRotBldcCeilPct}{3.7}% [dados] embarcados/results/csv/p11_features_20261023_222831.csv :: % janelas A2212 com f0 >= 0,99·fmax
\newcommand{\VRotBldcSpearman}{0.02}% [dados] embarcados/results/csv/p11_features_20261023_222831.csv + embarcados/results/predictions/diagnostico_janelas_20261023_222831.csv :: Spearman(velocidade nominal do arquivo, mediana de f0 do arquivo), 43 arquivos
\newcommand{\VRotBldcSpearmanMOne}{0.30}% [dados] embarcados/results/csv/p11_features_20261023_222831.csv + embarcados/results/predictions/diagnostico_janelas_20261023_222831.csv :: Spearman por motor M1 (14 arquivos)
\newcommand{\VRotBldcFilesMOne}{14}% [dados] embarcados/results/csv/p11_features_20261023_222831.csv + embarcados/results/predictions/diagnostico_janelas_20261023_222831.csv :: nº de arquivos do M1
\newcommand{\VRotBldcSpearmanMTwo}{0.30}% [dados] embarcados/results/csv/p11_features_20261023_222831.csv + embarcados/results/predictions/diagnostico_janelas_20261023_222831.csv :: Spearman por motor M2 (16 arquivos)
\newcommand{\VRotBldcFilesMTwo}{16}% [dados] embarcados/results/csv/p11_features_20261023_222831.csv + embarcados/results/predictions/diagnostico_janelas_20261023_222831.csv :: nº de arquivos do M2
\newcommand{\VRotBldcSpearmanMThree}{0.54}% [dados] embarcados/results/csv/p11_features_20261023_222831.csv + embarcados/results/predictions/diagnostico_janelas_20261023_222831.csv :: Spearman por motor M3 (6 arquivos)
\newcommand{\VRotBldcFilesMThree}{6}% [dados] embarcados/results/csv/p11_features_20261023_222831.csv + embarcados/results/predictions/diagnostico_janelas_20261023_222831.csv :: nº de arquivos do M3
\newcommand{\VRotBldcSpearmanMFour}{1.00}% [dados] embarcados/results/csv/p11_features_20261023_222831.csv + embarcados/results/predictions/diagnostico_janelas_20261023_222831.csv :: Spearman por motor M4 (7 arquivos)
\newcommand{\VRotBldcFilesMFour}{7}% [dados] embarcados/results/csv/p11_features_20261023_222831.csv + embarcados/results/predictions/diagnostico_janelas_20261023_222831.csv :: nº de arquivos do M4
\newcommand{\VEqStagesTotal}{26}% [equivalencia] embarcados/results/csv/equivalencia_resumo_por_build_20261024-091729.csv :: nº de etapas por build
\newcommand{\VEqStagesCase}{20}% [equivalencia] embarcados/results/csv/equivalencia_resumo_por_build_20261024-091729.csv :: etapas de janela
\newcommand{\VEqStagesSeq}{6}% [equivalencia] embarcados/results/csv/equivalencia_resumo_por_build_20261024-091729.csv :: etapas de sequência
\newcommand{\VEqCases}{164}% [equivalencia] embarcados/results/csv/equivalencia_resumo_por_build_20261024-091729.csv :: n[y_prep]
\newcommand{\VEqSeqCombos}{240}% [equivalencia] embarcados/results/csv/equivalencia_resumo_por_build_20261024-091729.csv :: n[absorb]
\newcommand{\VEqDoubleStagesPass}{26}% [equivalencia] embarcados/results/csv/equivalencia_resumo_por_build_20261024-091729.csv :: nº de etapas com veredito PASS (double)
\newcommand{\VEqDoubleMaxRatio}{0.55}% [equivalencia] embarcados/results/csv/equivalencia_resumo_por_build_20261024-091729.csv :: máx razao_tol_max nas etapas sem empate/FAIL (double)
\newcommand{\VEqDoubleDecisionPass}{164}% [equivalencia] embarcados/results/csv/equivalencia_resumo_por_build_20261024-091729.csv :: pass+pass_empate[decisao] (double)
\newcommand{\VEqFloatStagesPass}{25}% [equivalencia] embarcados/results/csv/equivalencia_resumo_por_build_20261024-091729.csv :: nº de etapas com veredito PASS (float)
\newcommand{\VEqFloatMaxRatio}{0.34}% [equivalencia] embarcados/results/csv/equivalencia_resumo_por_build_20261024-091729.csv :: máx razao_tol_max nas etapas sem empate/FAIL (float)
\newcommand{\VEqFloatDecisionPass}{164}% [equivalencia] embarcados/results/csv/equivalencia_resumo_por_build_20261024-091729.csv :: pass+pass_empate[decisao] (float)
\newcommand{\VEqSosdStagesPass}{26}% [equivalencia] embarcados/results/csv/equivalencia_resumo_por_build_20261024-091729.csv :: nº de etapas com veredito PASS (float-sosd)
\newcommand{\VEqSosdMaxRatio}{0.23}% [equivalencia] embarcados/results/csv/equivalencia_resumo_por_build_20261024-091729.csv :: máx razao_tol_max nas etapas sem empate/FAIL (float-sosd)
\newcommand{\VEqSosdDecisionPass}{164}% [equivalencia] embarcados/results/csv/equivalencia_resumo_por_build_20261024-091729.csv :: pass+pass_empate[decisao] (float-sosd)
\newcommand{\VEqFloatZFail}{2}% [equivalencia] embarcados/results/csv/equivalencia_resumo_por_build_20261024-091729.csv :: fail[float,z]
\newcommand{\VEqFloatZPass}{162}% [equivalencia] embarcados/results/csv/equivalencia_resumo_por_build_20261024-091729.csv :: pass[float,z]
\newcommand{\VEqFloatFailStages}{z}% [equivalencia] embarcados/results/csv/equivalencia_resumo_por_build_20261024-091729.csv :: etapas sem PASS (float)
\newcommand{\VEqFloatTieStages}{2}% [equivalencia] embarcados/results/csv/equivalencia_resumo_por_build_20261024-091729.csv :: nº de etapas com PASS_EMPATE (float)
\newcommand{\VEqFloatTieCases}{2}% [equivalencia] embarcados/results/csv/equivalencia_float_20261024-091122.csv :: ids com PASS_EMPATE em env_mod
\newcommand{\VEqFloatTieCaseIds}{borda\_dc\_ruido3s\_ma2212, borda\_dc\_ruido3s\_midmt}% [equivalencia] embarcados/results/csv/equivalencia_float_20261024-091122.csv :: id
\newcommand{\VEqFloatTieGap}{\ensuremath{9.4\times10^{-16}}}% [equivalencia] embarcados/results/csv/equivalencia_float_20261024-091122.csv :: detalhe[env_mod]: P_ref[ref]-P_ref[C]
\newcommand{\VEqFloatTieTol}{\ensuremath{6.3\times10^{-11}}}% [equivalencia] embarcados/results/csv/equivalencia_float_20261024-091122.csv :: detalhe[env_mod]: 2·tol
\newcommand{\VEqFloatFmodErr}{37.1}% [equivalencia] embarcados/results/csv/equivalencia_float_20261024-091122.csv :: err_max[features[9]:f_mod] (empate)
\newcommand{\VEqFloatZErrMax}{1.61}% [equivalencia] embarcados/results/csv/equivalencia_float_20261024-091122.csv :: err_max[z] (FAIL)
\newcommand{\VEqFloatZFailIsEdge}{1}% [equivalencia] embarcados/results/csv/equivalencia_float_20261024-091122.csv :: todos os FAIL de z são casos de borda (1=sim)
\newcommand{\VEqEnvDecFloatMed}{\ensuremath{1.2\times10^{-5}}}% [equivalencia] embarcados/results/csv/equivalencia_float_20261024-091122.csv :: mediana err_rel_max[env_dec]
\newcommand{\VEqEnvDecSosdMed}{\ensuremath{1.2\times10^{-7}}}% [equivalencia] embarcados/results/csv/equivalencia_float-sosd_20261024-091222.csv :: mediana err_rel_max[env_dec]
\newcommand{\VEqEnvDecDoubleMed}{\ensuremath{1.1\times10^{-7}}}% [equivalencia] embarcados/results/csv/equivalencia_double_20261024-091106.csv :: mediana err_rel_max[env_dec]
\newcommand{\VEqTolSha}{d74dae800ceb}% [equivalencia] embarcados/configs/embedded/tolerancias.json :: sha256 (12 primeiros)
\newcommand{\VEqTolShaRec}{d74dae800ceb}% [equivalencia] embarcados/results/csv/equivalencia_resumo_por_build_20261024-091729.csv :: sha256_tolerancias gravado no _meta da comparação
\newcommand{\VEqTolCommit}{447a7c2}% [equivalencia] embarcados/configs/embedded/tolerancias.json :: git log -1 do tolerancias.json
\newcommand{\VEqTolCommitTime}{2026-10-24 09:10}% [equivalencia] embarcados/configs/embedded/tolerancias.json :: data do commit do tolerancias.json
\newcommand{\VEqRunTime}{2026-10-24 09:11}% [equivalencia] embarcados/results/csv/equivalencia_resumo_por_build_20261024-091729.csv :: data da comparação formal (float)
\newcommand{\VEqGoldClean}{60}% [equivalencia] embarcados/gold/manifest.json :: casos tipo limpo
\newcommand{\VEqGoldNoisy}{90}% [equivalencia] embarcados/gold/manifest.json :: casos tipo ruidoso
\newcommand{\VEqGoldEdge}{14}% [equivalencia] embarcados/gold/manifest.json :: casos tipo borda
\newcommand{\VEqGoldSeqs}{20}% [equivalencia] embarcados/gold/manifest.json :: n_seqs
\newcommand{\VEqGoldPerts}{5}% [equivalencia] embarcados/gold/manifest.json :: nº de perturbações
\newcommand{\VEqMinDurValues}{2}% [equivalencia] embarcados/gold/seqs/idmt_seq00/params.json :: len(min_durs_s)
\newcommand{\VEqEmulDoubleRows}{6{,}852}% [equivalencia] embarcados/results/csv/equivalencia_emul_double_20261024-090737_meta.json :: n_linhas
\newcommand{\VEqEmulDoublePass}{6{,}852}% [equivalencia] embarcados/results/csv/equivalencia_emul_double_20261024-090737_meta.json :: contagem_status.PASS
\newcommand{\VEqEmulFloatRows}{6{,}852}% [equivalencia] embarcados/results/csv/equivalencia_emul_float_20261024-090743_meta.json :: n_linhas
\newcommand{\VEqEmulFloatPass}{6{,}852}% [equivalencia] embarcados/results/csv/equivalencia_emul_float_20261024-090743_meta.json :: contagem_status.PASS
\newcommand{\VPcCIdmtMed}{13.29}% [pc] embarcados/results/metrics/bench_pc_20261024-091911.json :: resumo.c_float_pc/idmt.total_ms.mediana
\newcommand{\VPcPyIdmtMed}{95.3}% [pc] embarcados/results/metrics/bench_pc_20261024-091911.json :: resumo.python_ref/idmt.total_ms.mediana
\newcommand{\VPcCIdmtPninetyfive}{15.00}% [pc] embarcados/results/metrics/bench_pc_20261024-091911.json :: resumo.c_float_pc/idmt.total_ms.p95
\newcommand{\VPcPyIdmtPninetyfive}{100.2}% [pc] embarcados/results/metrics/bench_pc_20261024-091911.json :: resumo.python_ref/idmt.total_ms.p95
\newcommand{\VPcCIdmtMax}{15.80}% [pc] embarcados/results/metrics/bench_pc_20261024-091911.json :: resumo.c_float_pc/idmt.total_ms.max
\newcommand{\VPcPyIdmtMax}{110.2}% [pc] embarcados/results/metrics/bench_pc_20261024-091911.json :: resumo.python_ref/idmt.total_ms.max
\newcommand{\VPcCIdmtRtf}{0.0044}% [pc] embarcados/results/metrics/bench_pc_20261024-091911.json :: fator_tempo_real_mediana
\newcommand{\VPcPyIdmtRtf}{0.032}% [pc] embarcados/results/metrics/bench_pc_20261024-091911.json :: fator_tempo_real_mediana
\newcommand{\VPcSpeedupIdmt}{7.2}% [pc] embarcados/results/metrics/bench_pc_20261024-091911.json :: mediana Python / mediana C
\newcommand{\VPcCIdmtTrackMed}{5.45}% [pc] embarcados/results/metrics/bench_pc_20261024-091911.json :: rastrear_quadros_ms.mediana
\newcommand{\VPcCIdmtTrackPct}{41}% [pc] embarcados/results/metrics/bench_pc_20261024-091911.json :: 100·rastrear_quadros_ms.mediana/total_ms.mediana
\newcommand{\VPcCIdmtVitMed}{5.35}% [pc] embarcados/results/metrics/bench_pc_20261024-091911.json :: viterbi_ms.mediana
\newcommand{\VPcCIdmtVitPct}{40}% [pc] embarcados/results/metrics/bench_pc_20261024-091911.json :: 100·viterbi_ms.mediana/total_ms.mediana
\newcommand{\VPcCIdmtEnvMed}{2.27}% [pc] embarcados/results/metrics/bench_pc_20261024-091911.json :: envelope_mod_ms.mediana
\newcommand{\VPcCIdmtEnvPct}{17}% [pc] embarcados/results/metrics/bench_pc_20261024-091911.json :: 100·envelope_mod_ms.mediana/total_ms.mediana
\newcommand{\VPcCIdmtFeatMed}{0.013}% [pc] embarcados/results/metrics/bench_pc_20261024-091911.json :: features_ms.mediana
\newcommand{\VPcCIdmtFeatPct}{0}% [pc] embarcados/results/metrics/bench_pc_20261024-091911.json :: 100·features_ms.mediana/total_ms.mediana
\newcommand{\VPcCIdmtCosMed}{0.002}% [pc] embarcados/results/metrics/bench_pc_20261024-091911.json :: cosseno_ms.mediana
\newcommand{\VPcCIdmtCosPct}{0}% [pc] embarcados/results/metrics/bench_pc_20261024-091911.json :: 100·cosseno_ms.mediana/total_ms.mediana
\newcommand{\VPcNIdmt}{200}% [pc] embarcados/results/metrics/bench_pc_20261024-091911.json :: total_ms.n
\newcommand{\VPcCBldcMed}{3.18}% [pc] embarcados/results/metrics/bench_pc_20261024-091911.json :: resumo.c_float_pc/a2212.total_ms.mediana
\newcommand{\VPcPyBldcMed}{22.1}% [pc] embarcados/results/metrics/bench_pc_20261024-091911.json :: resumo.python_ref/a2212.total_ms.mediana
\newcommand{\VPcCBldcPninetyfive}{3.84}% [pc] embarcados/results/metrics/bench_pc_20261024-091911.json :: resumo.c_float_pc/a2212.total_ms.p95
\newcommand{\VPcPyBldcPninetyfive}{23.1}% [pc] embarcados/results/metrics/bench_pc_20261024-091911.json :: resumo.python_ref/a2212.total_ms.p95
\newcommand{\VPcCBldcMax}{4.45}% [pc] embarcados/results/metrics/bench_pc_20261024-091911.json :: resumo.c_float_pc/a2212.total_ms.max
\newcommand{\VPcPyBldcMax}{25.3}% [pc] embarcados/results/metrics/bench_pc_20261024-091911.json :: resumo.python_ref/a2212.total_ms.max
\newcommand{\VPcCBldcRtf}{0.0032}% [pc] embarcados/results/metrics/bench_pc_20261024-091911.json :: fator_tempo_real_mediana
\newcommand{\VPcPyBldcRtf}{0.022}% [pc] embarcados/results/metrics/bench_pc_20261024-091911.json :: fator_tempo_real_mediana
\newcommand{\VPcSpeedupBldc}{6.9}% [pc] embarcados/results/metrics/bench_pc_20261024-091911.json :: mediana Python / mediana C
\newcommand{\VPcCBldcTrackMed}{1.24}% [pc] embarcados/results/metrics/bench_pc_20261024-091911.json :: rastrear_quadros_ms.mediana
\newcommand{\VPcCBldcTrackPct}{39}% [pc] embarcados/results/metrics/bench_pc_20261024-091911.json :: 100·rastrear_quadros_ms.mediana/total_ms.mediana
\newcommand{\VPcCBldcVitMed}{1.14}% [pc] embarcados/results/metrics/bench_pc_20261024-091911.json :: viterbi_ms.mediana
\newcommand{\VPcCBldcVitPct}{36}% [pc] embarcados/results/metrics/bench_pc_20261024-091911.json :: 100·viterbi_ms.mediana/total_ms.mediana
\newcommand{\VPcCBldcEnvMed}{0.67}% [pc] embarcados/results/metrics/bench_pc_20261024-091911.json :: envelope_mod_ms.mediana
\newcommand{\VPcCBldcEnvPct}{21}% [pc] embarcados/results/metrics/bench_pc_20261024-091911.json :: 100·envelope_mod_ms.mediana/total_ms.mediana
\newcommand{\VPcCBldcFeatMed}{0.013}% [pc] embarcados/results/metrics/bench_pc_20261024-091911.json :: features_ms.mediana
\newcommand{\VPcCBldcFeatPct}{0}% [pc] embarcados/results/metrics/bench_pc_20261024-091911.json :: 100·features_ms.mediana/total_ms.mediana
\newcommand{\VPcCBldcCosMed}{0.002}% [pc] embarcados/results/metrics/bench_pc_20261024-091911.json :: cosseno_ms.mediana
\newcommand{\VPcCBldcCosPct}{0}% [pc] embarcados/results/metrics/bench_pc_20261024-091911.json :: 100·cosseno_ms.mediana/total_ms.mediana
\newcommand{\VPcNBldc}{200}% [pc] embarcados/results/metrics/bench_pc_20261024-091911.json :: total_ms.n
\newcommand{\VPcWarmup}{5}% [pc] embarcados/results/metrics/bench_pc_20261024-091911.json :: meta.argv.aquecimento
\newcommand{\VPcDistinctIdmt}{75}% [pc] embarcados/results/csv/bench_pc_20261024-091911.csv :: n_janelas_distintas
\newcommand{\VPcDistinctBldc}{75}% [pc] embarcados/results/csv/bench_pc_20261024-091911.csv :: n_janelas_distintas
\newcommand{\VPcCompiler}{MSVC 19.44 x64 \slash{}O2 \slash{}fp:precise}% [pc] embarcados/results/metrics/bench_pc_20261024-091911.json :: meta.compilador
\newcommand{\VPcNcpu}{16}% [pc] embarcados/results/metrics/bench_pc_20261024-091911.json :: meta.maquina.n_cpu
\newcommand{\VPcOs}{Windows 11}% [pc] embarcados/results/metrics/bench_pc_20261024-091911.json :: meta.maquina.plataforma
\newcommand{\VPcSeqWindows}{24}% [pc] embarcados/results/csv/seq_pc_20261024-092103.csv :: nº de linhas (janelas)
\newcommand{\VPcSeqSwitches}{2}% [pc] embarcados/results/csv/seq_pc_20261024-092103_meta.json :: len(deteccao_online_c)
\newcommand{\VPcSeqDelayC}{0}% [pc] embarcados/results/csv/seq_pc_20261024-092103_meta.json :: max atraso_janelas (C)
\newcommand{\VPcSeqDelayPy}{0}% [pc] embarcados/results/csv/seq_pc_20261024-092103_meta.json :: max atraso_janelas (Python)
\newcommand{\VPcSeqSameDecisions}{1}% [pc] embarcados/results/csv/seq_pc_20261024-092103_meta.json :: decisoes_iguais_python_c (1=sim)
\newcommand{\VPcSeqMaxProbDiff}{\ensuremath{3.2\times10^{-6}}}% [pc] embarcados/results/csv/seq_pc_20261024-092103_meta.json :: max_abs_dif_probs
\newcommand{\VPcSeqCp}{2}% [pc] embarcados/results/csv/seq_pc_20261024-092103_meta.json :: cp
\newcommand{\VPcSeqPerClass}{8}% [pc] embarcados/results/csv/seq_pc_20261024-092103.csv :: janelas por classe
\newcommand{\VEstWsFloatOneMb}{0.96}% [estimativa] embarcados/paper/fontes/ws_float_16000.json :: total_com_tabela/1e6
\newcommand{\VEstWsFloatOneNoTabMb}{0.61}% [estimativa] embarcados/paper/fontes/ws_float_16000.json :: total_sem_tabela/1e6
\newcommand{\VEstWsFloatOnePhaseBPct}{44}% [estimativa] embarcados/paper/fontes/ws_float_16000.json :: faseB/(persistente+max(faseA,faseB))
\newcommand{\VEstWsFloatThreeMb}{1.86}% [estimativa] embarcados/paper/fontes/ws_float_48000.json :: total_com_tabela/1e6
\newcommand{\VEstWsFloatThreeNoTabMb}{1.52}% [estimativa] embarcados/paper/fontes/ws_float_48000.json :: total_sem_tabela/1e6
\newcommand{\VEstWsFloatThreePhaseBPct}{64}% [estimativa] embarcados/paper/fontes/ws_float_48000.json :: faseB/(persistente+max(faseA,faseB))
\newcommand{\VEstWsDoubleOneMb}{1.79}% [estimativa] embarcados/paper/fontes/ws_double_16000.json :: total_com_tabela/1e6
\newcommand{\VEstWsDoubleOneNoTabMb}{1.21}% [estimativa] embarcados/paper/fontes/ws_double_16000.json :: total_sem_tabela/1e6
\newcommand{\VEstWsDoubleOnePhaseBPct}{47}% [estimativa] embarcados/paper/fontes/ws_double_16000.json :: faseB/(persistente+max(faseA,faseB))
\newcommand{\VEstWsDoubleThreeMb}{3.60}% [estimativa] embarcados/paper/fontes/ws_double_48000.json :: total_com_tabela/1e6
\newcommand{\VEstWsDoubleThreeNoTabMb}{3.02}% [estimativa] embarcados/paper/fontes/ws_double_48000.json :: total_sem_tabela/1e6
\newcommand{\VEstWsDoubleThreePhaseBPct}{67}% [estimativa] embarcados/paper/fontes/ws_double_48000.json :: faseB/(persistente+max(faseA,faseB))
\newcommand{\VEstHarmTableMb}{0.35}% [estimativa] embarcados/paper/fontes/ws_float_16000.json :: total_com_tabela - total_sem_tabela
\newcommand{\VBuildIdf}{v5.5.6}% [build] C:/Espressif/build_p11_esp32s3/project_description.json :: git_revision
\newcommand{\VBuildGcc}{14.2.0}% [build] C:/Espressif/build_p11_esp32s3/project_description.json :: c_compiler (versão do toolchain esp-)
\newcommand{\VBuildEspBinKb}{209}% [build] C:/Espressif/build_p11_esp32/size.txt :: binary size (hex)/1024
\newcommand{\VBuildEspBinPct}{20}% [build] C:/Espressif/build_p11_esp32/size.txt :: bin/partição app
\newcommand{\VBuildEspFlashCodeKb}{124}% [build] C:/Espressif/build_p11_esp32/size.txt :: Flash Code Used/1024
\newcommand{\VBuildEspFlashDataKb}{19}% [build] C:/Espressif/build_p11_esp32/size.txt :: Flash Data Used/1024
\newcommand{\VBuildEspDramKb}{15}% [build] C:/Espressif/build_p11_esp32/size.txt :: DRAM Used/1024
\newcommand{\VBuildEspDramTotalKb}{176}% [build] C:/Espressif/build_p11_esp32/size.txt :: DRAM Total/1024
\newcommand{\VBuildEspIramKb}{58}% [build] C:/Espressif/build_p11_esp32/size.txt :: IRAM Used/1024
\newcommand{\VBuildEspDramName}{DRAM}% [build] C:/Espressif/build_p11_esp32/size.txt :: nome da região de dados interna
\newcommand{\VBuildEspWsRatio}{10.3}% [estimativa] C:/Espressif/build_p11_esp32/size.txt :: estimativa ws float 3 s / DRAM total
\newcommand{\VBuildEspWsRatioOne}{5.3}% [estimativa] C:/Espressif/build_p11_esp32/size.txt :: estimativa ws float 1 s / DRAM total
\newcommand{\VBuildEspsBinKb}{233}% [build] C:/Espressif/build_p11_esp32s3/size.txt :: binary size (hex)/1024
\newcommand{\VBuildEspsBinPct}{23}% [build] C:/Espressif/build_p11_esp32s3/size.txt :: bin/partição app
\newcommand{\VBuildEspsFlashCodeKb}{139}% [build] C:/Espressif/build_p11_esp32s3/size.txt :: Flash Code Used/1024
\newcommand{\VBuildEspsFlashDataKb}{30}% [build] C:/Espressif/build_p11_esp32s3/size.txt :: Flash Data Used/1024
\newcommand{\VBuildEspsDramKb}{54}% [build] C:/Espressif/build_p11_esp32s3/size.txt :: DIRAM Used/1024
\newcommand{\VBuildEspsDramTotalKb}{334}% [build] C:/Espressif/build_p11_esp32s3/size.txt :: DIRAM Total/1024
\newcommand{\VBuildEspsIramKb}{16}% [build] C:/Espressif/build_p11_esp32s3/size.txt :: IRAM Used/1024
\newcommand{\VBuildEspsDramName}{DIRAM}% [build] C:/Espressif/build_p11_esp32s3/size.txt :: nome da região de dados interna
\newcommand{\VBuildEspsWsRatio}{5.5}% [estimativa] C:/Espressif/build_p11_esp32s3/size.txt :: estimativa ws float 3 s / DIRAM total
\newcommand{\VBuildEspsWsRatioOne}{2.8}% [estimativa] C:/Espressif/build_p11_esp32s3/size.txt :: estimativa ws float 1 s / DIRAM total
\newcommand{\VChipModel}{\textbf{[pending]}}% [esp32] embarcados/results/manifests/bancada.json :: chip_modelo
\newcommand{\VChipCpuMhz}{\textbf{[pending]}}% [esp32] embarcados/results/manifests/bancada.json :: cpu_mhz
\newcommand{\VChipPsramMb}{\textbf{[pending]}}% [esp32] embarcados/results/manifests/bancada.json :: psram_mb
\newcommand{\VChipLatIdmtMed}{\textbf{[pending]}}% [esp32] embarcados/results/manifests/bancada.json :: resumo.idmt.total_ms.mediana
\newcommand{\VChipLatIdmtPninetyfive}{\textbf{[pending]}}% [esp32] embarcados/results/manifests/bancada.json :: resumo.idmt.total_ms.p95
\newcommand{\VChipLatIdmtMax}{\textbf{[pending]}}% [esp32] embarcados/results/manifests/bancada.json :: resumo.idmt.total_ms.max
\newcommand{\VChipLatIdmtN}{\textbf{[pending]}}% [esp32] embarcados/results/manifests/bancada.json :: resumo.idmt.total_ms.n
\newcommand{\VChipRtfIdmt}{\textbf{[pending]}}% [esp32] embarcados/results/manifests/bancada.json :: fator_tempo_real_mediana
\newcommand{\VChipRtfMaxIdmt}{\textbf{[pending]}}% [esp32] embarcados/results/manifests/bancada.json :: fator_tempo_real_max
\newcommand{\VChipHeapIntIdmtKb}{\textbf{[pending]}}% [esp32] embarcados/results/manifests/bancada.json :: heap_min_livre_interna/1024
\newcommand{\VChipStackIdmt}{\textbf{[pending]}}% [esp32] embarcados/results/manifests/bancada.json :: pilha_folga_bytes_min
\newcommand{\VChipWsPsramIdmt}{\textbf{[pending]}}% [esp32] embarcados/results/manifests/bancada.json :: ws_na_psram
\newcommand{\VChipWsIdmtMb}{\textbf{[pending]}}% [esp32] embarcados/results/manifests/bancada.json :: ws_bytes/1e6
\newcommand{\VChipLatBldcMed}{\textbf{[pending]}}% [esp32] embarcados/results/manifests/bancada.json :: resumo.a2212.total_ms.mediana
\newcommand{\VChipLatBldcPninetyfive}{\textbf{[pending]}}% [esp32] embarcados/results/manifests/bancada.json :: resumo.a2212.total_ms.p95
\newcommand{\VChipLatBldcMax}{\textbf{[pending]}}% [esp32] embarcados/results/manifests/bancada.json :: resumo.a2212.total_ms.max
\newcommand{\VChipLatBldcN}{\textbf{[pending]}}% [esp32] embarcados/results/manifests/bancada.json :: resumo.a2212.total_ms.n
\newcommand{\VChipRtfBldc}{\textbf{[pending]}}% [esp32] embarcados/results/manifests/bancada.json :: fator_tempo_real_mediana
\newcommand{\VChipRtfMaxBldc}{\textbf{[pending]}}% [esp32] embarcados/results/manifests/bancada.json :: fator_tempo_real_max
\newcommand{\VChipHeapIntBldcKb}{\textbf{[pending]}}% [esp32] embarcados/results/manifests/bancada.json :: heap_min_livre_interna/1024
\newcommand{\VChipStackBldc}{\textbf{[pending]}}% [esp32] embarcados/results/manifests/bancada.json :: pilha_folga_bytes_min
\newcommand{\VChipWsPsramBldc}{\textbf{[pending]}}% [esp32] embarcados/results/manifests/bancada.json :: ws_na_psram
\newcommand{\VChipWsBldcMb}{\textbf{[pending]}}% [esp32] embarcados/results/manifests/bancada.json :: ws_bytes/1e6
\newcommand{\VChipEqStagesPass}{\textbf{[pending]}}% [esp32] embarcados/results/manifests/bancada.json :: equivalência no chip não rodou
\newcommand{\VChipEqStages}{\textbf{[pending]}}% [esp32] embarcados/results/manifests/bancada.json :: equivalência no chip não rodou
\newcommand{\VChipEqDecisionPass}{\textbf{[pending]}}% [esp32] embarcados/results/manifests/bancada.json :: equivalência no chip não rodou
\newcommand{\VChipSeqDelay}{\textbf{[pending]}}% [esp32] embarcados/results/manifests/bancada.json :: fluxo contínuo no chip não rodou

% Gerado por embarcados/src/python/gerar_artigo_numeros_aud.py. NÃO EDITAR À MÃO.
% commit 9fb82f016b4482c1a5b58f1b9bd08db149799b78 (árvore suja em embarcados/: True)
% Macros da auditoria independente (M2 holdout, M3 motor igual por igual, M4 rotação).
% Uma macro por número. Comentário: [categoria] fonte :: coluna.

\newcommand{\audHoldoutSeed}{7}% [equivalencia] embarcados/results/csv/aud_holdout_20261025-140558_meta.json :: seed
\newcommand{\audHoldoutN}{24}% [equivalencia] embarcados/results/csv/aud_holdout_20261025-140558_resumo.csv :: n_janelas[float,_todas]
\newcommand{\audHoldoutNIdmt}{12}% [equivalencia] embarcados/results/csv/aud_holdout_20261025-140558_resumo.csv :: n_janelas[float,idmt]
\newcommand{\audHoldoutNBldc}{12}% [equivalencia] embarcados/results/csv/aud_holdout_20261025-140558_resumo.csv :: n_janelas[float,a2212]
\newcommand{\audHoldoutFloatPathSame}{24}% [equivalencia] embarcados/results/csv/aud_holdout_20261025-140558_resumo.csv :: caminho_identico[float,_todas]
\newcommand{\audHoldoutFloatDecSame}{24}% [equivalencia] embarcados/results/csv/aud_holdout_20261025-140558_resumo.csv :: decisao_igual[float,_todas]
\newcommand{\audHoldoutFloatFeatOut}{0}% [equivalencia] embarcados/results/csv/aud_holdout_20261025-140558_resumo.csv :: n_features_fora_tol[float,_todas]
\newcommand{\audHoldoutFloatStageFail}{0}% [equivalencia] embarcados/results/csv/aud_holdout_20261025-140558_resumo.csv :: n_etapas_fail[float,_todas]
\newcommand{\audHoldoutFloatRatioMax}{0.20}% [equivalencia] embarcados/results/csv/aud_holdout_20261025-140558_resumo.csv :: razao_tol_max_features[float,_todas]
\newcommand{\audHoldoutFloatDpMax}{\ensuremath{8.7\times10^{-6}}}% [equivalencia] embarcados/results/csv/aud_holdout_20261025-140558_resumo.csv :: max_abs_dp[float,_todas]
\newcommand{\audHoldoutDoublePathSame}{24}% [equivalencia] embarcados/results/csv/aud_holdout_20261025-140558_resumo.csv :: caminho_identico[double,_todas]
\newcommand{\audHoldoutDoubleDecSame}{24}% [equivalencia] embarcados/results/csv/aud_holdout_20261025-140558_resumo.csv :: decisao_igual[double,_todas]
\newcommand{\audHoldoutDoubleFeatOut}{0}% [equivalencia] embarcados/results/csv/aud_holdout_20261025-140558_resumo.csv :: n_features_fora_tol[double,_todas]
\newcommand{\audHoldoutDoubleStageFail}{0}% [equivalencia] embarcados/results/csv/aud_holdout_20261025-140558_resumo.csv :: n_etapas_fail[double,_todas]
\newcommand{\audHoldoutDoubleRatioMax}{0.05}% [equivalencia] embarcados/results/csv/aud_holdout_20261025-140558_resumo.csv :: razao_tol_max_features[double,_todas]
\newcommand{\audHoldoutDoubleDpMax}{\ensuremath{1.9\times10^{-7}}}% [equivalencia] embarcados/results/csv/aud_holdout_20261025-140558_resumo.csv :: max_abs_dp[double,_todas]
\newcommand{\audRotBldcSpearman}{0.02}% [dados] embarcados/results/csv/aud_rotacao_motor_20261025-140355.csv :: spearman_rho[global]
\newcommand{\audRotBldcFiles}{43}% [dados] embarcados/results/csv/aud_rotacao_motor_20261025-140355.csv :: n_arquivos[global]
\newcommand{\audRotBldcWindows}{430}% [dados] embarcados/results/csv/aud_rotacao_motor_20261025-140355.csv :: n_janelas[global]
\newcommand{\audRotBldcCeilPct}{3.7}% [dados] embarcados/results/csv/aud_rotacao_motor_20261025-140355.csv :: 100·frac_janelas_teto[global]
\newcommand{\audRotBldcCeilThrPct}{99}% [dados] embarcados/results/csv/aud_rotacao_motor_20261025-140355.csv :: 100·lim_teto_rel
\newcommand{\audRotBldcSpearmanMOne}{0.30}% [dados] embarcados/results/csv/aud_rotacao_motor_20261025-140355.csv :: spearman_rho[M1]
\newcommand{\audRotBldcSpearmanPMOne}{0.29}% [dados] embarcados/results/csv/aud_rotacao_motor_20261025-140355.csv :: spearman_p[M1]
\newcommand{\audRotBldcFilesMOne}{14}% [dados] embarcados/results/csv/aud_rotacao_motor_20261025-140355.csv :: n_arquivos[M1]
\newcommand{\audRotBldcSpearmanMTwo}{0.30}% [dados] embarcados/results/csv/aud_rotacao_motor_20261025-140355.csv :: spearman_rho[M2]
\newcommand{\audRotBldcSpearmanPMTwo}{0.26}% [dados] embarcados/results/csv/aud_rotacao_motor_20261025-140355.csv :: spearman_p[M2]
\newcommand{\audRotBldcFilesMTwo}{16}% [dados] embarcados/results/csv/aud_rotacao_motor_20261025-140355.csv :: n_arquivos[M2]
\newcommand{\audRotBldcSpearmanMThree}{0.54}% [dados] embarcados/results/csv/aud_rotacao_motor_20261025-140355.csv :: spearman_rho[M3]
\newcommand{\audRotBldcSpearmanPMThree}{0.27}% [dados] embarcados/results/csv/aud_rotacao_motor_20261025-140355.csv :: spearman_p[M3]
\newcommand{\audRotBldcFilesMThree}{6}% [dados] embarcados/results/csv/aud_rotacao_motor_20261025-140355.csv :: n_arquivos[M3]
\newcommand{\audRotBldcSpearmanMFour}{1.00}% [dados] embarcados/results/csv/aud_rotacao_motor_20261025-140355.csv :: spearman_rho[M4]
\newcommand{\audRotBldcSpearmanPMFour}{0.00}% [dados] embarcados/results/csv/aud_rotacao_motor_20261025-140355.csv :: spearman_p[M4]
\newcommand{\audRotBldcFilesMFour}{7}% [dados] embarcados/results/csv/aud_rotacao_motor_20261025-140355.csv :: n_arquivos[M4]
\newcommand{\audRotBldcSpearmanMotorMin}{0.30}% [dados] embarcados/results/csv/aud_rotacao_motor_20261025-140355.csv :: min spearman_rho[M1..M4]
\newcommand{\audRotBldcSpearmanMotorMax}{1.00}% [dados] embarcados/results/csv/aud_rotacao_motor_20261025-140355.csv :: max spearman_rho[M1..M4]
\newcommand{\audRotBldcFilesMotorMin}{6}% [dados] embarcados/results/csv/aud_rotacao_motor_20261025-140355.csv :: min n_arquivos[M1..M4]
\newcommand{\audRotBldcFilesMotorMax}{16}% [dados] embarcados/results/csv/aud_rotacao_motor_20261025-140355.csv :: max n_arquivos[M1..M4]
\newcommand{\audMotorBalClsFileHP}{0.972}% [diagnostico] embarcados/results/csv/aud_motor_igual_20261025-140432.csv :: acuracia_balanceada[protocolo,p04_classicas,arquivo,HP_M1M3]
\newcommand{\audMotorBalClsFileHPLo}{0.924}% [diagnostico] embarcados/results/csv/aud_motor_igual_20261025-140432.csv :: ic95_lo[protocolo,p04_classicas,arquivo,HP_M1M3]
\newcommand{\audMotorBalClsFileHPHi}{0.998}% [diagnostico] embarcados/results/csv/aud_motor_igual_20261025-140432.csv :: ic95_hi[protocolo,p04_classicas,arquivo,HP_M1M3]
\newcommand{\audMotorBalClsMotorHP}{0.920}% [diagnostico] embarcados/results/csv/aud_motor_igual_20261025-140432.csv :: acuracia_balanceada[protocolo,p04_classicas,motor,HP_M1M3]
\newcommand{\audMotorBalClsMotorHPLo}{0.868}% [diagnostico] embarcados/results/csv/aud_motor_igual_20261025-140432.csv :: ic95_lo[protocolo,p04_classicas,motor,HP_M1M3]
\newcommand{\audMotorBalClsMotorHPHi}{0.962}% [diagnostico] embarcados/results/csv/aud_motor_igual_20261025-140432.csv :: ic95_hi[protocolo,p04_classicas,motor,HP_M1M3]
\newcommand{\audMotorDropClsHP}{0.052}% [diagnostico] embarcados/results/csv/aud_motor_igual_20261025-140432.csv :: acuracia_balanceada[queda,p04_classicas,HP_M1M3]
\newcommand{\audMotorDropClsHPLo}{0.023}% [diagnostico] embarcados/results/csv/aud_motor_igual_20261025-140432.csv :: ic95_lo[queda,p04_classicas,HP_M1M3]
\newcommand{\audMotorDropClsHPHi}{0.086}% [diagnostico] embarcados/results/csv/aud_motor_igual_20261025-140432.csv :: ic95_hi[queda,p04_classicas,HP_M1M3]
\newcommand{\audMotorDropClsFull}{0.368}% [diagnostico] embarcados/results/csv/aud_motor_igual_20261025-140432.csv :: queda_completa[queda,p04_classicas,HP_M1M3]
\newcommand{\audMotorBalDecFileHP}{0.753}% [diagnostico] embarcados/results/csv/aud_motor_igual_20261025-140432.csv :: acuracia_balanceada[protocolo,p11_decifrae_harmsum,arquivo,HP_M1M3]
\newcommand{\audMotorBalDecFileHPLo}{0.662}% [diagnostico] embarcados/results/csv/aud_motor_igual_20261025-140432.csv :: ic95_lo[protocolo,p11_decifrae_harmsum,arquivo,HP_M1M3]
\newcommand{\audMotorBalDecFileHPHi}{0.842}% [diagnostico] embarcados/results/csv/aud_motor_igual_20261025-140432.csv :: ic95_hi[protocolo,p11_decifrae_harmsum,arquivo,HP_M1M3]
\newcommand{\audMotorBalDecMotorHP}{0.685}% [diagnostico] embarcados/results/csv/aud_motor_igual_20261025-140432.csv :: acuracia_balanceada[protocolo,p11_decifrae_harmsum,motor,HP_M1M3]
\newcommand{\audMotorBalDecMotorHPLo}{0.593}% [diagnostico] embarcados/results/csv/aud_motor_igual_20261025-140432.csv :: ic95_lo[protocolo,p11_decifrae_harmsum,motor,HP_M1M3]
\newcommand{\audMotorBalDecMotorHPHi}{0.774}% [diagnostico] embarcados/results/csv/aud_motor_igual_20261025-140432.csv :: ic95_hi[protocolo,p11_decifrae_harmsum,motor,HP_M1M3]
\newcommand{\audMotorDropDecHP}{0.068}% [diagnostico] embarcados/results/csv/aud_motor_igual_20261025-140432.csv :: acuracia_balanceada[queda,p11_decifrae_harmsum,HP_M1M3]
\newcommand{\audMotorDropDecHPLo}{0.031}% [diagnostico] embarcados/results/csv/aud_motor_igual_20261025-140432.csv :: ic95_lo[queda,p11_decifrae_harmsum,HP_M1M3]
\newcommand{\audMotorDropDecHPHi}{0.116}% [diagnostico] embarcados/results/csv/aud_motor_igual_20261025-140432.csv :: ic95_hi[queda,p11_decifrae_harmsum,HP_M1M3]
\newcommand{\audMotorDropDecFull}{0.360}% [diagnostico] embarcados/results/csv/aud_motor_igual_20261025-140432.csv :: queda_completa[queda,p11_decifrae_harmsum,HP_M1M3]
\newcommand{\audMotorBalRmsNormFileHP}{0.843}% [diagnostico] embarcados/results/csv/aud_motor_igual_20261025-140432.csv :: acuracia_balanceada[protocolo,b00_rmspico_norm,arquivo,HP_M1M3]
\newcommand{\audMotorBalRmsNormFileHPLo}{0.760}% [diagnostico] embarcados/results/csv/aud_motor_igual_20261025-140432.csv :: ic95_lo[protocolo,b00_rmspico_norm,arquivo,HP_M1M3]
\newcommand{\audMotorBalRmsNormFileHPHi}{0.909}% [diagnostico] embarcados/results/csv/aud_motor_igual_20261025-140432.csv :: ic95_hi[protocolo,b00_rmspico_norm,arquivo,HP_M1M3]
\newcommand{\audMotorBalRmsNormMotorHP}{0.814}% [diagnostico] embarcados/results/csv/aud_motor_igual_20261025-140432.csv :: acuracia_balanceada[protocolo,b00_rmspico_norm,motor,HP_M1M3]
\newcommand{\audMotorBalRmsNormMotorHPLo}{0.720}% [diagnostico] embarcados/results/csv/aud_motor_igual_20261025-140432.csv :: ic95_lo[protocolo,b00_rmspico_norm,motor,HP_M1M3]
\newcommand{\audMotorBalRmsNormMotorHPHi}{0.887}% [diagnostico] embarcados/results/csv/aud_motor_igual_20261025-140432.csv :: ic95_hi[protocolo,b00_rmspico_norm,motor,HP_M1M3]
\newcommand{\audMotorDropRmsNormHP}{0.029}% [diagnostico] embarcados/results/csv/aud_motor_igual_20261025-140432.csv :: acuracia_balanceada[queda,b00_rmspico_norm,HP_M1M3]
\newcommand{\audMotorDropRmsNormHPLo}{-0.002}% [diagnostico] embarcados/results/csv/aud_motor_igual_20261025-140432.csv :: ic95_lo[queda,b00_rmspico_norm,HP_M1M3]
\newcommand{\audMotorDropRmsNormHPHi}{0.066}% [diagnostico] embarcados/results/csv/aud_motor_igual_20261025-140432.csv :: ic95_hi[queda,b00_rmspico_norm,HP_M1M3]
\newcommand{\audMotorDropRmsNormFull}{0.019}% [diagnostico] embarcados/results/csv/aud_motor_igual_20261025-140432.csv :: queda_completa[queda,b00_rmspico_norm,HP_M1M3]
\newcommand{\audMotorBalRmsRawFileHP}{0.913}% [diagnostico] embarcados/results/csv/aud_motor_igual_20261025-140432.csv :: acuracia_balanceada[protocolo,b00_rmspico_bruto,arquivo,HP_M1M3]
\newcommand{\audMotorBalRmsRawFileHPLo}{0.821}% [diagnostico] embarcados/results/csv/aud_motor_igual_20261025-140432.csv :: ic95_lo[protocolo,b00_rmspico_bruto,arquivo,HP_M1M3]
\newcommand{\audMotorBalRmsRawFileHPHi}{0.983}% [diagnostico] embarcados/results/csv/aud_motor_igual_20261025-140432.csv :: ic95_hi[protocolo,b00_rmspico_bruto,arquivo,HP_M1M3]
\newcommand{\audMotorBalRmsRawMotorHP}{0.754}% [diagnostico] embarcados/results/csv/aud_motor_igual_20261025-140432.csv :: acuracia_balanceada[protocolo,b00_rmspico_bruto,motor,HP_M1M3]
\newcommand{\audMotorBalRmsRawMotorHPLo}{0.623}% [diagnostico] embarcados/results/csv/aud_motor_igual_20261025-140432.csv :: ic95_lo[protocolo,b00_rmspico_bruto,motor,HP_M1M3]
\newcommand{\audMotorBalRmsRawMotorHPHi}{0.873}% [diagnostico] embarcados/results/csv/aud_motor_igual_20261025-140432.csv :: ic95_hi[protocolo,b00_rmspico_bruto,motor,HP_M1M3]
\newcommand{\audMotorDropRmsRawHP}{0.159}% [diagnostico] embarcados/results/csv/aud_motor_igual_20261025-140432.csv :: acuracia_balanceada[queda,b00_rmspico_bruto,HP_M1M3]
\newcommand{\audMotorDropRmsRawHPLo}{0.078}% [diagnostico] embarcados/results/csv/aud_motor_igual_20261025-140432.csv :: ic95_lo[queda,b00_rmspico_bruto,HP_M1M3]
\newcommand{\audMotorDropRmsRawHPHi}{0.251}% [diagnostico] embarcados/results/csv/aud_motor_igual_20261025-140432.csv :: ic95_hi[queda,b00_rmspico_bruto,HP_M1M3]
\newcommand{\audMotorDropRmsRawFull}{0.106}% [diagnostico] embarcados/results/csv/aud_motor_igual_20261025-140432.csv :: queda_completa[queda,b00_rmspico_bruto,HP_M1M3]
\newcommand{\audMotorHPFiles}{36}% [diagnostico] embarcados/results/csv/aud_motor_igual_20261025-140432.csv :: n_arquivos[queda]
\newcommand{\audMotorHPWindows}{360}% [diagnostico] embarcados/results/csv/aud_motor_igual_20261025-140432.csv :: n_janelas[queda]

% Gerado por embarcados/src/python/gerar_artigo_figuras.py. NÃO EDITAR À MÃO.
% Caminhos relativos a embarcados/paper/.
\newif\iffigesp
\figesptrue  % true só se a figura do ESP32 existir
\newcommand{\FigArq}{../figures/pdf/arquitetura.pdf}
\newcommand{\FigBalacc}{../figures/pdf/diag_balacc_en_20261023_222831.pdf}
\newcommand{\FigConf}{../figures/pdf/diag_confusao_p11_en_20261023_222831.pdf}
\newcommand{\FigEquiv}{../figures/pdf/equivalencia_erro_etapas_en.pdf}
\newcommand{\FigSpec}{../figures/pdf/diag_espectro_en_20261023_222831.pdf}
\newcommand{\FigEspLat}{../figures/pdf/esp32_lat_esp32_20261025-143026_en.pdf}


\begin{document}

\title{Porting a Harmonic-Sum Rotation Tracker for Acoustic Motor-State Monitoring to a
Microcontroller: Stage-Wise Verified C Port, Measured Cost, and What Two Public Datasets
Allow Us to Conclude\ifwip\ (Work in Progress)\fi}

% TODO(autores): nomes e afiliações a preencher pelos autores.
\author{\IEEEauthorblockN{[Authors to be completed]}
\IEEEauthorblockA{[Affiliation]}}

\maketitle

\begin{abstract}
We port to portable C a pipeline that classifies the operating state of small electric motors
from sound by borrowing ideas from bass-guitar transcription: a harmonic-sum rotation tracker
decoded by a Viterbi search, \VParNFeat{} physical features, cosine-to-prototype scoring and an
online hidden Markov model (HMM). The port is checked against the Python reference stage by
stage, on \VEqCases{} windows and \VEqSeqCombos{} sequence configurations, with tolerances
frozen before the formal comparison. In double precision all \VEqStagesTotal{} stage checks
pass. In float32, the build intended for the microcontroller, all \VEqRealCases{} windows of real
audio pass every stage; on \VEqFloatBadCases{} synthetic constant-offset-plus-noise windows the
envelope-modulation peak lies at the float32 noise floor and the modulation, feature and z-score
checks fail, while the class decision is identical on \VEqFloatDecisionPass{} of \VEqCases{}
windows. On an ESP32 without PSRAM the reference memory layout (estimated at
\VEstWsFloatThreeMb{}\,MB for a \VParSegIdmt{}-s window) does not fit. A named low-memory
variant fits \VParSegBldc{}-s windows, and on the chip its outputs agree with the reference within
the frozen tolerances on all \VChipEqCases{} windows compared; measured on the chip it is
\VChipRtfBldc{} times slower than real time, dominated by a table-free harmonic sum in
software-emulated double precision, and the \VParSegIdmt{}-s window does not fit. The pipeline does not outperform simple baselines:
on the IDMT dataset every method, including a two-number level baseline, saturates because
class, motor and recording coincide; on the A2212 dataset classical features reach
\VDiagBalClsBldcFile{} balanced accuracy against \VDiagBalDecBldcFile{} for the pipeline, and
with a motor held out all methods drop. The contribution is therefore the verified embedded
implementation, its cost, an interpretable rotation estimate, and an explicit account of what
these datasets can and cannot support.
\end{abstract}

\begin{IEEEkeywords}
acoustic condition monitoring, embedded signal processing, ESP32, rotation tracking,
Viterbi decoding, reproducibility, data leakage
\end{IEEEkeywords}

% ======================================================================
\section{Introduction}

Acoustic monitoring of electric motors is attractive because a microphone is cheap,
non-intrusive and needs no electrical connection to the machine. Several recent works report
near-perfect accuracy on two public datasets of small laboratory motors, the
IDMT-ISA-ELECTRIC-ENGINE set~\cite{grollmisch2019} and the A2212 brushless set~\cite{estacio2023},
often with deep networks~\cite{melcnn2025,ertarin2025} or large hand-crafted feature banks~\cite{etnext2026,tekin2025}.
Two aspects receive little attention. First, most of these results split recordings into
segments before the train/test split, so neighbouring segments of one recording fall on both
sides. Second, few works measure the cost on a microcontroller; an exception is the edge
classifier of Strantzalis et al.~\cite{strantzalis2022}, while other TinyML studies stop at a model
size and never run on hardware~\cite{almaini2026}.

This paper starts from a pipeline that we call Decifrae (p11). It reuses components of a
music-transcription system: the motor is treated as a ``one-string instrument'' whose
fundamental (the rotation frequency) is tracked through a sum over its harmonics and decoded with
a Viterbi search that has a silence state and an octave-jump cost, as a bass transcriber would do.
The analogy only motivates the \emph{representation}; we do not claim that motor sound is a
musical note, nor that music-transcription performance transfers to machines.

The question we answer is whether this pipeline fits and runs on an ESP32-class
microcontroller with an implementation that is demonstrably the same computation as the Python
reference, and what the two public datasets allow us to conclude about it. In short, it fits and
runs on an ESP32 without PSRAM only for \VParSegBldc{}-s windows, through a named low-memory
variant, and \VChipRtfBldc{} times slower than real time on that board.
\ifwip
We contribute a C port verified stage by stage against the reference with tolerances frozen before
the formal comparison, a cost account that keeps estimate, PC measurement, firmware build and
on-chip measurement apart, the tracked rotation as an interpretable output, and an evaluation
under three protocols showing that the pipeline does not beat simple baselines and that the
datasets cannot separate operating state from motor identity.
\else
Our contributions are:
\begin{itemize}
\item a portable C port without dynamic allocation, verified against the Python reference stage
by stage with tolerances frozen before the formal comparison, in double and in float32, with
every residual discrepancy listed;
\item a cost account that keeps four categories apart: analytical estimate, measurement on a PC,
firmware build (\texttt{idf.py size}), and measurement on an ESP32 without PSRAM, including the
low-memory variant needed there;
\item an interpretable intermediate output, the tracked rotation frequency, together with a check
of how far it can be trusted on each dataset;
\item an evaluation of the pipeline and of simple baselines on the same windows under three
protocols (random windows, grouped by recording or file, and motor held out), showing that the
pipeline does not outperform the baselines and that the datasets cannot separate operating
state from motor identity.
\end{itemize}
\fi
Three distinctions are kept throughout. We classify the \emph{current operating state} of a
motor; we do not predict future failures. Both datasets contain small laboratory motors, not
industrial machines. And an evaluation in which the test motor was also seen in training says
nothing about generalisation to a new motor.

% ======================================================================
\section{Related Work}

\emph{Datasets and splits.} Grollmisch et al.~\cite{grollmisch2019} released the IDMT set with
an official split (training on clean recordings, testing on recordings with added noise) and a
frame-level network. Each class was recorded on a different motor, so class and physical unit
coincide. Later works on IDMT and A2212 report near-perfect accuracy with mel-spectrogram
CNNs~\cite{melcnn2025}, image encodings~\cite{ertarin2025}, wavelet pattern features with
SVM/kNN~\cite{etnext2026} and wavelet statistics with bagged trees~\cite{tekin2025}; their splits
operate on segments rather than on recordings. The A2212 data paper~\cite{estacio2023} reports
no classifier; it documents that the bearing fault exists on a single motor.
\ifwip\else
The classical-feature baseline used here follows the feature list described in the abstract of
Vanaraj et al.~\cite{vanaraj2026}; we only read the abstract, so it is an inspired
re-implementation, not a reproduction. Glowacz~\cite{glowacz2018} classified a commutator motor
with a nearest-mean rule on selected spectral lines, the closest precedent of our
cosine-to-prototype decision.
\fi

\ifwip
\emph{Rotation tracking without a tachometer} has been reviewed on vibration
data~\cite{peeters2019}, and Viterbi decoding over a time--frequency representation predates our
work~\cite{xi2015}; we add its verified embedded implementation on microphone audio.
\else
\emph{Rotation tracking without a tachometer.} Peeters et al.~\cite{peeters2019} review and
compare tacholess speed estimators on vibration data, including multi-order probabilistic
approaches that combine several harmonics. Viterbi decoding over a time--frequency
representation for speed estimation predates our work~\cite{xi2015}, and harmonic weighting has
been proposed to suppress weak orders~\cite{yoo2023}. Wheat et al.~\cite{wheat2025} build
features on a harmonic grid given a known operating speed. Our tracker is therefore not new as an
idea; what we add is its embedded, verified implementation on microphone audio.
\fi

\ifwip\else
\emph{Embedded and anomaly-oriented work.} Strantzalis et al.~\cite{strantzalis2022} deploy a
small CNN on a Cortex-M4 and measure latency, RAM and response time to state changes on IDMT,
but with a random split of overlapping spectrograms. Almaini et al.~\cite{almaini2026} quantise an
MFCC network trained on urban sounds without deploying it. Unsupervised anomaly detection on
machine sounds (autoencoder baselines of the DCASE challenges) addresses novelty rather than state
classification and is not compared here.
\fi

% ======================================================================
\section{Method}
\label{sec:method}

Fig.~\ref{fig:arq} summarises the pipeline. Each analysis window (\VParSegIdmt{}\,s for IDMT,
\VParSegBldc{}\,s for A2212, both at \VParSrKhz{}\,kHz) is RMS-normalised and processed
independently; only the HMM carries state across windows.

\begin{figure}[t]
\centering
\includegraphics[width=\columnwidth]{\FigArq}
\caption{System overview. Everything inside the dashed box is the C port that runs on the
microcontroller; audio currently reaches the chip over UART (processor-in-the-loop) or from its
flash. Working memory: analytical estimate of the reference layout, and the pool peak of the
low-memory variant measured on the ESP32.}
\label{fig:arq}
\end{figure}

\emph{Rotation tracker.} A short-time Fourier transform with a symmetric Hann window of
\VParNfft{} samples and hop \VParHop{} gives magnitude frames $|X_t(k)|$. For each of
$J=\VParJ{}$ candidate rotation frequencies $f_j$ on a logarithmic grid from \VParFmin{} to
\VParFmax{}\,Hz (step \VParStepCents{}~cents), the saliency
$s_t(j)=\sum_{h=1}^{K}|X_t(h f_j/\Delta f)|$ sums $K=\VParNharm{}$ harmonics with linear
interpolation between bins, and is normalised per frame to $\tilde e_t(j)\in[0,1]$. The emission
cost is $-\log(\VParWEnergy{}\,\tilde e_t(j)+\VParEmisEps{})$; a silence state receives a cost
derived from the frame energy relative to the loudest frame of the window. A Viterbi search
allows jumps of up to \VParMaxJump{} states (\VParMaxJumpCents{} cents) with a cost proportional
to the jump, octave jumps of $\pm$\VParOctStates{} states with a fixed cost, and transitions into
and out of silence. The decoded path gives $f_{rot}(t)$.

\emph{Features and decision.} From the path and the mean power spectrum we compute \VParNFeat{}
features\ifwip\else{} (Table~\ref{tab:params})\fi: the median and inter-quartile range of $f_{rot}$, the silent-frame
fraction, the saliency along the path, a harmonic-energy ratio, the relative energy of four groups
of harmonic orders, the frequency and prominence of the main envelope-modulation peak (Hilbert
envelope, Butterworth low-pass, decimation to \VParEnvRateHz{}\,Hz, Welch spectrum), the spectral
centroid and the spectral flatness. Features are z-scored with training statistics and compared
by cosine with one prototype per class; a softmax with temperature \VParTempCos{} gives class
probabilities. For streams, an HMM with a fixed change penalty and an extra ``no decision''
state is run as a causal forward filter, one window at a time.

\ifwip\else
% Gerado por embarcados/src/python/gerar_artigo_tabelas.py a partir de embarcados/paper/numeros.json.
% NÃO EDITAR À MÃO.

\begin{table}[t]
\centering
\scriptsize
\setlength{\tabcolsep}{3pt}
\caption{Parameters of the p11 reference, read from the Python code (\texttt{PreDecifrae}, \texttt{rotacao}, \texttt{\_envelope\_mod}), and the 13 features in the order of the reference.}
\label{tab:params}
\begin{tabular}{@{}p{1.55cm}p{6.6cm}@{}}
\toprule
Sampling rate; window & 16\,kHz; 3\,s (IDMT), 1\,s (A2212); RMS $\rightarrow$ 0.1 \\
STFT & symmetric Hann, $N$=8{,}192, hop 1{,}024, no centring; 39 / 8 frames \\
Rotation grid & 15--60\,Hz, 2-cent step, $J$=1{,}201 states + silence \\
Harmonic sum & $K$=48 harmonics, unit weights, linear interpolation \\
Emission cost & $-\log(0.45\,\tilde e + \ensuremath{10^{-3}})$; silence from frame energy \\
Viterbi & step 0.55/semitone, jumps $\le$50 states, octave $\pm$600 states (cost 7.0), silence 3.2 \\
Envelope & $|\mathcal{H}(y)|$, Butterworth 4, 200\,Hz (causal \texttt{sosfilt}), $\downarrow$16, Welch; peak in 2--100\,Hz \\
Classifier & z-score, cosine to class means, softmax $T$=0.1 \\
Online HMM & change penalty swept (6 values, 0.5--16), $\log p_N$=-50 \\
\midrule
Features & (0) median $f_{rot}$; (1) IQR of $f_{rot}$; (2) silent-frame fraction; (3) mean normalised saliency on the path; (4) harmonic ratio ($k f_{rot}\pm1$ bin / 10--4{,}000\,Hz); (5) energy of orders 1--4 (dB); (6) energy of orders 5--12 (dB); (7) energy of orders 13--24 (dB); (8) energy of orders 25--48 (dB); (9) envelope-modulation peak frequency; (10) its prominence (dB); (11) spectral centroid; (12) spectral flatness \\
\bottomrule
\end{tabular}
\end{table}

\fi

\ifwip
\emph{C port.} The C code follows the Python line by line, compiles unchanged for the PC
and the ESP-IDF, processes the spectrogram frame by frame, includes a mixed-radix FFT for the
full-length Hilbert transform, and uses a caller-provided arena instead of dynamic allocation.
The real type is float32 on the microcontroller and double for verification. Any change of
arithmetic or memory layout is a named variant: \texttt{p11\_c\_sosd} keeps the envelope filter
state in double (more precise, not an approximation), and \texttt{\VChipVariant}, the one run on
our board without PSRAM, takes its memory from a pool of internal DRAM and IRAM blocks, drops the
harmonic table (positions recomputed in double) and computes the Hilbert FFT in four steps. An
ESP-DSP option (\texttt{p11\_c\_espdsp}) was neither compared nor measured and is left out.
\else
\emph{C port.} The C code follows the Python line by line and compiles unchanged for the PC and
for the ESP-IDF. It never stores the full spectrogram: each frame is windowed, transformed by a
real FFT, accumulated into the mean power spectrum and reduced to its saliency row. The Viterbi
search needs the whole window (the silence cost uses the maximum frame energy of the window), so
the normalised saliency and the back-pointers (short integers) are kept; the Hilbert envelope needs an FFT
of the full window length, which is not a power of two, so the port includes a mixed-radix FFT.
All buffers live in a caller-provided arena whose size is computed by the same layout function that
assigns the pointers. The real type is selected at compile time (float32 for the microcontroller,
double for verification). Index decisions that change integers (grid, medians, rounding of
harmonic positions) are computed in double in both builds, reproducing the reference's
round-half-to-even and first-index tie-breaking. Any change of arithmetic is a \emph{named
variant}, compared under the same tolerances. \texttt{p11\_c\_sosd} keeps the envelope filter
state in double, as SciPy does (more precise, not an approximation). \texttt{\VChipVariant}, run on
our board without PSRAM, calls the same core functions but carves its memory phase by phase from
a pool of internal DRAM and IRAM blocks (IRAM, word access only, holds only float vectors such as
the FFT twiddles), drops the harmonic-interpolation table (bin positions recomputed in double for
every frame) and computes the Hilbert FFT in four steps, the same DFT in another order. An ESP-DSP
option (\texttt{p11\_c\_espdsp}) was neither compared nor measured and is left out.
\fi

% ======================================================================
\section{Experimental Setup}
\label{sec:setup}

\emph{Data and protocols.} We use the \texttt{train/} part of IDMT (the official test set was not
opened) and all of A2212\ifwip\else{} (Table~\ref{tab:protocols})\fi. All methods are evaluated on exactly the
same windows and folds. On IDMT, the random-window protocol leaks by construction, and the
``blocks'' protocol holds out contiguous blocks of each recording; since there is one recording
per class, it still cannot separate class from motor. On A2212 we compare random windows, a
grouped split by file, and leave-one-motor-out; because the bearing fault exists only on motor
M4, we also report leave-one-motor-out without the M4 fold.

\ifwip\else% Gerado por embarcados/src/python/gerar_artigo_tabelas.py a partir de embarcados/paper/numeros.json.
% NÃO EDITAR À MÃO.

\begin{table}[t]
\centering
\scriptsize
\setlength{\tabcolsep}{3pt}
\caption{Datasets and evaluation protocols. IDMT: \texttt{train/} only, 3 recordings of 3 motors, 357 windows of 3\,s. A2212: 43 files of 4 motors (22 Healthy, 14 Propeller, 7 Bearing), 430 windows of 1\,s. $g$: groups used by the bootstrap. $^\dagger$Intentional leakage, reported only to measure inflation.}
\label{tab:protocols}
\begin{tabular}{@{}p{1.55cm}p{1.35cm}r p{3.75cm}@{}}
\toprule
Protocol & Split unit & $g$ & What it can / cannot support \\
\midrule
IDMT random$^\dagger$ & window & 18 & Upper bound only: neighbouring windows of one recording on both sides. \\
IDMT blocks & contiguous block & 15 & One recording per class: class = motor = recording. Measures recording identity, not operating state. \\
A2212 random$^\dagger$ & window & 43 & Upper bound only: windows of one file on both sides. \\
A2212 file & file & 43 & State of a \emph{known} motor (same motors in training and test). Bearing exists only on motor M4. \\
A2212 motor-out & motor & 43 & Unseen motor. With M4 held out, Bearing has no training data (recall 0.000 by construction). Only 4 motors: no motor-level CI. \\
A2212 motor-out w/o M4 & motor (M1--M3) & 36 & Healthy vs.\ Propeller on an unseen motor; 3 folds. \\
\bottomrule
\end{tabular}
\end{table}
\fi

\emph{Baselines.} A level baseline uses only the RMS and peak of the window with logistic
regression, once on RMS-normalised audio (only the crest factor remains) and once on raw audio. A
classical baseline (p04) uses \VParNClassicFeat{} frame statistics (zero-crossing rate, RMS,
spectral centroid and MFCC means and deviations) with a soft-voting ensemble, inspired
by~\cite{vanaraj2026}.

\emph{Metrics.} We report balanced accuracy, macro-F1, macro one-vs-rest ROC-AUC from continuous
scores, and per-class recall. Confidence intervals (\VDataCiLevel{}\%) come from a bootstrap that
resamples groups (recording blocks or files), stratified by class, with \VDataNBoot{} resamples;
windows of one group are never treated as independent. For the motor-out protocols the
independent unit is the motor, of which there are only \VDataBldcMotors{}, so no honest interval
exists at that level; the file-level interval shown is optimistic.

\ifwip
\emph{Equivalence protocol.} Golden vectors of every stage were produced by the
instrumented reference on \VEqCases{} windows (\VEqGoldClean{} clean, \VEqGoldNoisy{} noisy,
\VEqGoldEdge{} edge cases) and \VEqSeqCombos{} HMM sequence configurations. Per-stage tolerances,
derived from rounding analysis and an independent emulation, were frozen (\texttt{\VEqTolCommit{}})
before the formal comparison, although the C port had already been debugged against the same
golden vectors by its developer; discrete outputs must agree exactly unless a tie is demonstrated.
\else
\emph{Equivalence protocol.} The instrumented Python reference produced golden vectors for every
stage on \VEqCases{} windows: \VEqGoldClean{} clean validation windows from both datasets,
\VEqGoldNoisy{} noisy versions under \VEqGoldPerts{} perturbations (white noise at several SNRs and
synthetic sag and flicker power-quality disturbances), and \VEqGoldEdge{} edge cases (silence, clipping, synthetic
rotation outside the search range, and a constant offset with tiny noise). The HMM stages were
checked on \VEqGoldSeqs{} probability sequences under \VParNChangePen{} change penalties and
\VEqMinDurValues{} minimum durations (\VEqSeqCombos{} configurations). Tolerances for each stage
were derived from first-order rounding analysis and an independent numerical emulation of each
build, and frozen before the formal comparison (commit \texttt{\VEqTolCommit{}}, file hash
\texttt{\VEqTolSha{}}\ldots, \VEqTolCommitTime{}; comparison run \VEqRunTime{}); the hash is
recorded in every result row. The chronology matters: the C port had already been debugged
against the same golden vectors by its developer before the tolerances were frozen, so the
tolerances are frozen, not pre-registered. On \audHoldoutN{} windows outside the golden set,
unseen by the C developer, path and decision agree on \audHoldoutFloatPathSame{} and
\audHoldoutFloatDecSame{} in float32, no feature exceeds its tolerance (maximum ratio
\audHoldoutFloatRatioMax{}) and probabilities differ by at most \audHoldoutFloatDpMax{}. Tolerances for the edge cases
were calibrated on the emulation error for the same signals, a weaker standard, and some are wide
in float32 (e.g.\ the modulation prominence). Continuous stages are compared by maximum error relative to the stage's
scale; discrete outputs (rotation path, decision, HMM states) must agree exactly unless a tie is
demonstrated from the reference's own values. A failure at one stage marks later stages as
downstream. Both emulations passed all comparison rows (\VEqEmulDoublePass{} and \VEqEmulFloatPass{} of
\VEqEmulFloatRows{}) before the C code was compared.
\fi

\ifwip
\emph{Cost measurement.} We keep apart (1) the analytical arena estimate, (2) timings on
the PC (\VPcNIdmt{} timings cycling over \VPcDistinctIdmt{} distinct windows, each run once in a new
process, so warm-up does not warm the C code; workspace initialisation excluded), (3)
\texttt{idf.py size} of the firmware (ESP-IDF \VBuildIdf{}) and (4) measurements on an
\VChipModel{} without PSRAM (\VChipCpuMhz{}\,MHz; \VChipLatBldcN{} timed windows over UART after
\VChipLatWarmup{} warm-up runs, on-chip timer around the processing only). Energy is not measured.
\else
\emph{Cost measurement.} We keep four categories apart. (1) \emph{Analytical estimate}: the
arena size returned by the C layout function. (2) \emph{Measured on the PC}: \VPcNIdmt{} timings
per dataset cycling over \VPcDistinctIdmt{} distinct windows. The C executable is started anew and
timed once per window, so warm-up does not warm it, and its time excludes the workspace
initialisation, which the Python timing cannot exclude; the two effects partly offset.
(3) \emph{Firmware build}: \texttt{idf.py size} (ESP-IDF \VBuildIdf{}, GCC \VBuildGcc{}).
(4) \emph{Measured on the ESP32}: an \VChipModel{} (rev.\ \VChipRev{}, \VChipCpuMhz{}\,MHz,
\VChipFlashMb{}\,MB flash, no PSRAM, radios never started); windows arrive over UART (or, for the
stream, from flash) and are timed with the on-chip timer around each core call, with heap and
stack watermarks. Energy is not measured.
\fi

% ======================================================================
\section{Results}
\label{sec:results}

\subsection{Python--C equivalence}

\ifwip
Table~\ref{tab:equiv} summarises the verdicts. In double precision all \VEqStagesTotal{}
stage checks pass on all windows and sequence configurations; the rotation path, the decision and
every HMM output agree exactly. In float32 all \VEqRealCases{} windows of real audio pass every
stage. The failures are confined to the \VEqFloatBadCases{} synthetic constant-offset-plus-noise
windows, whose modulation spectrum lies at the float32 rounding floor in both implementations, so
the peak location is undetermined (the reference peak is \VEqDcRefPeakRatio{} times the reference
power at the bin chosen by the C code: not a tie). We count the modulation, feature and z-score
checks as failed there (\VEqFloatStagesPassStrict{} of \VEqStagesTotal{} checks pass on all
windows); the class decision agrees on \VEqFloatDecisionPass{} of \VEqCases{} windows. The named
variant \texttt{p11\_c\_sosd} passes every check. An independent audit reproduced the comparison
from a clean build; its corrections are applied here. These are PC results; the comparison on the
chip is in Section~\ref{sec:cost}.
\else
Table~\ref{tab:equiv} summarises the verdicts\ifwip\else{} and Fig.~\ref{fig:equiv} the worst
error-to-tolerance ratio per stage\fi. In double precision all \VEqStagesTotal{} stage checks pass
on all windows and sequence configurations (worst ratio \VEqDoubleMaxRatio{} of the tolerance).
The rotation path, the decision and every HMM output agree exactly. In float32, all
\VEqRealCases{} windows of real audio (clean and noisy) pass every stage. The
\VEqFloatBadCases{} remaining failures are the synthetic constant-offset-plus-noise edge case, once
per model. There the envelope is almost constant and its modulation spectrum lies at the float32
rounding floor in both implementations, so the peak location is undetermined at that precision:
the reference's peak is \VEqDcRefPeakRatio{} times the reference power at the bin chosen by the C
code, so this is not a tie. The modulation frequency moves by \VEqFloatFmodErr{}\,Hz, and we count
the envelope-modulation, feature and z-score checks as failed on these windows (the frozen
comparison rule had accepted the first two as a tie; we report the conservative reading), leaving
\VEqFloatStagesPassStrict{} of \VEqStagesTotal{} stage checks passing on all windows. Cosines,
probabilities and the class decision still pass on all windows
(\VEqFloatDecisionPass{}/\VEqCases{}). With the named variant \texttt{p11\_c\_sosd} every check
passes; it reduces the maximum error of the decimated envelope,
relative to its peak and taken as the median over windows, from \VEqEnvDecFloatMed{} to
\VEqEnvDecSosdMed{}, because the low-pass poles lie close to
the unit circle and float32 rounding accumulates in the filter state. These results were obtained
on the PC; the same comparison on the chip is
\ifchipmedido reported in Section~\ref{sec:cost}\else pending\fi. An independent audit
(report in the repository) reproduced the comparison from a clean build and its corrections are
applied here. We therefore
claim stage-wise agreement within tolerances frozen before the formal comparison, with the
exceptions above, not unconditional equivalence.
\fi

\begin{table}[t]
\centering
\scriptsize
\caption{Stage checks passed on the PC (\VEqCases{} windows; \VEqSeqCombos{} sequence
configurations). Conservative reading: a group fails if any stage fails on any window. All
float32 failures are the \VEqFloatBadCases{} synthetic constant-offset edge windows, where the
modulation peak is at the float32 noise floor; the decision is unchanged.}
\label{tab:equiv}
% Gerado por embarcados/src/python/gerar_artigo_tabelas.py a partir de embarcados/results/csv/equivalencia_resumo_por_build_20261024-101821.csv (veredito por etapa).
% NÃO EDITAR À MÃO.

\begin{tabular}{@{}lccc@{}}
\toprule
Stage group (\#stages) & C double & C float32 & float32+SOS64 \\
\midrule
Front end (6) & 6/6 & 6/6 & 6/6 \\
Rotation Viterbi (3) & 3/3 & 3/3 & 3/3 \\
Spectrum, envelope (6) & 6/6 & 5/6 (FAIL in 2/164) & 6/6 \\
13 features (1) & 1/1 & 0/1 (FAIL in 2/164) & 1/1 \\
z-score (1) & 1/1 & 0/1 (FAIL in 2/164) & 1/1 \\
Cosine, softmax, decision (3) & 3/3 & 3/3 & 3/3 \\
Sequences (HMM) (6) & 6/6 & 6/6 & 6/6 \\
\bottomrule
\end{tabular}

\end{table}

\ifwip\else
\begin{figure}[t]
\centering
\includegraphics[width=\columnwidth]{\FigEquiv}
\caption{Worst error-to-tolerance ratio per stage (PC). Ratios below one pass. F: failure on the
synthetic constant-offset edge windows (counted as failures even where the frozen rule accepted a
tie). Exact zeros are drawn at the bottom of the axis.}
\label{fig:equiv}
\end{figure}
\fi

\subsection{Cost}
\label{sec:cost}

\ifwip
\emph{Measured on the PC},
the C float32 build of the reference layout needs \VPcCIdmtMed{}\,ms per \VParSegIdmt{}-s window
and \VPcCBldcMed{}\,ms per \VParSegBldc{}-s window, against \VPcPyIdmtMed{}\,ms and
\VPcPyBldcMed{}\,ms for the Python reference (\VPcSpeedupIdmt{}$\times$, including interpreter
overhead); the STFT with harmonic sum takes \VPcCIdmtTrackPct{}\%, the rotation Viterbi
\VPcCIdmtVitPct{}\% and the envelope branch \VPcCIdmtEnvPct{}\% of the median time. These numbers
say nothing about the microcontroller.
\else
Table~\ref{tab:resources} lists the resources by category. The PC timings (column 2; the STFT
with harmonic sum takes \VPcCIdmtTrackPct{}\%, the rotation Viterbi \VPcCIdmtVitPct{}\% and the
envelope branch \VPcCIdmtEnvPct{}\% of the median) say nothing about the microcontroller.
\fi

The \emph{analytical estimate} of the working memory of the reference layout (one arena) is
\VEstWsFloatThreeKib{}\,KiB for \VParSegIdmt{}\,s and \VEstWsFloatOneKib{}\,KiB for
\VParSegBldc{}\,s, at least \VChipRefOverPoolThree{} and \VChipRefOverPoolOne{} times the internal
RAM that the firmware can hand to the working memory on our board (\VChipPoolKib{}\,KiB, measured)
and at least \VBuildEspsWsRatio{} and \VBuildEspsWsRatioOne{} times the whole DIRAM of the
ESP32-S3 \emph{build} (\VBuildEspsDramTotalKb{}\,KiB; its IRAM is \VBuildEspsIramPct{}\% used):
without PSRAM the reference layout cannot run.

\ifchipmedido
\emph{Measured on the ESP32} (\VChipPsramMb{}\,MB PSRAM), \texttt{\VChipVariant} fits a
\VParSegBldc{}-s window: the pool peaks at \VChipWsBldcB{}\,bytes (\VChipWsPoolPct{}\% of it,
Fig.~\ref{fig:esplat}b), the minimum free internal heap is \VChipHeapIntBldcB{}\,bytes and the
stack headroom \VChipStackBldc{}\,bytes. The \VParSegIdmt{}-s IDMT window does \emph{not} fit, so
no IDMT figure exists on the chip. All \VChipEqCases{} golden \VParSegBldc{}-s windows
(\VChipEqCasesReal{} real, \VChipEqCasesEdge{} edge cases) pass, on the chip, the \VChipEqStages{}
final outputs the firmware returns (Viterbi path and score through decision), with \VChipEqFail{} failures and a worst error-to-tolerance
ratio of \VChipEqMaxRatio{}; intermediate stages are not sent back. On the PC the variant passes
all \VEqLmStages{} compared stages. The median latency is \VChipLatBldcMed{}\,ms per
\VParSegBldc{}-s window (p95 \VChipLatBldcPninetyfive{}, maximum \VChipLatBldcMax{}\,ms), a
real-time factor of \VChipRtfBldc{}: on this board the pipeline does \emph{not} run in real time.
The STFT with harmonic sum takes \VChipBldcTrackPct{}\% of it (Fig.~\ref{fig:esplat}a; on the PC,
with the table, \VPcCBldcTrackPct{}\%). We attribute this to recomputing the harmonic positions in
double precision, which the ESP32 only emulates in software (the variant was not timed on the PC);
a float32 variant of that step or a board with PSRAM is the way towards real time. The variant's
image occupies \VBuildLmBinKb{}\,KiB of flash (build).
\else
\emph{Measured on the ESP32}: not available when this version was generated.
\fi

\ifwip\else% Gerado por embarcados/src/python/gerar_artigo_tabelas.py a partir de embarcados/paper/numeros.json.
% NÃO EDITAR À MÃO.

\begin{table*}[t]
\centering
\scriptsize
\setlength{\tabcolsep}{4pt}
\caption{Resources by evidence category: (1) analytical estimate (\texttt{p11\_host ws}); (2) measured on the PC (Windows 11, MSVC 19.44 x64 \slash{}O2 \slash{}fp:precise; 200 timings per dataset over 75 distinct windows, one new process per window); (3) \texttt{idf.py size}, static only; (4) measured on the ESP32 without PSRAM (200 timed windows after 5 warm-up runs). Columns (1)--(2) use the reference layout \texttt{p11\_c\_ref}; column (4), the only measurement on the microcontroller, uses \texttt{p11\_c\_lowmem}.}
\label{tab:resources}
\begin{tabular}{@{}lllll@{}}
\toprule
Quantity & (1) Estimate & (2) PC & (3) Build & (4) ESP32 \\
\midrule
Latency, C float32, median [p95] (ms), 3\,s / 1\,s & -- & 13.29 [15.00] / 3.18 [3.84] & -- & \textit{did not fit} / 2647.5 [2647.6] \\
Latency, Python reference, median (ms), 3\,s / 1\,s & -- & 95.3 / 22.1 & -- & -- \\
Real-time factor (C, median), 3\,s / 1\,s; STFT + harm.\ sum share at 1\,s & -- & 0.0044 / 0.0032; 39\,\% & -- & \textit{did not fit} / 2.65; 89\,\% \\
Working memory (KiB), 3\,s / 1\,s & 1820 / 936 & -- & -- & \textit{did not fit} / 296 (pool peak) \\
Internal RAM for the pool; min.\ free heap (KiB); stack headroom (B) & -- & -- & -- & 308; 7.4; 20{,}536 \\
App binary (KiB): ESP32, -S3 (\texttt{p11\_c\_ref}); ESP32 (\texttt{p11\_c\_lowmem}) & -- & -- & 209, 233; 215 & -- \\
Static DRAM, IRAM used / region (KiB), ESP32, \texttt{p11\_c\_lowmem} & -- & -- & 35/176, 58/128 & -- \\
\bottomrule
\end{tabular}
\end{table*}
\fi

\iffigesp
\begin{figure}[t]
\centering
\includegraphics[width=\columnwidth]{\FigEspLat}
\caption{Measured on the ESP32 without PSRAM (\texttt{\VChipVariant}, A2212, \VParSegBldc{}-s
windows). (a) Median time per stage over \VChipLatBldcN{} windows. (b) Internal RAM handed to the
pool and its peak use (measured) against the estimated reference arena (hatched).}
\label{fig:esplat}
\end{figure}
\fi

\subsection{Classification under three protocols}

% Gerado por embarcados/src/python/gerar_artigo_tabelas.py a partir de embarcados/paper/numeros.json.
% NÃO EDITAR À MÃO.

\begin{table}[t]
\centering
\scriptsize
\setlength{\tabcolsep}{3pt}
\caption{Balanced accuracy on the same windows (validation only); $g$ = groups. 95\% CIs from a group bootstrap (2{,}000 resamples) are shown in Fig.~\ref{fig:balacc}. $^\dagger$Intentional leakage.}
\label{tab:desempenho-compacto}
\begin{tabular}{@{}lrcccc@{}}
\toprule
Protocol & $g$ & RMS/pk (norm.) & RMS/pk (raw) & Classical & p11 \\
\midrule
IDMT random$^\dagger$ & 18 & 0.992 & 0.997 & 1.000 & 1.000 \\
IDMT blocks & 15 & 0.991 & 0.993 & 1.000 & 1.000 \\
A2212 random$^\dagger$ & 43 & 0.564 & 0.625 & 0.998 & 0.849 \\
A2212 file & 43 & 0.562 & 0.608 & 0.981 & 0.816 \\
A2212 motor-out & 43 & 0.542 & 0.503 & 0.613 & 0.456 \\
A2212 w/o M4 & 36 & 0.814 & 0.754 & 0.920 & 0.685 \\
\bottomrule
\end{tabular}
\end{table}


\begin{figure}[t]
\centering
\includegraphics[width=\columnwidth]{\FigBalacc}
\caption{Balanced accuracy on the same windows with \VDataCiLevel{}\% group-bootstrap
intervals. Dotted line: chance. *Intentional leakage. On IDMT class = motor = recording; on A2212
Bearing exists only on motor M4 (confounded with motor identity), so motor-out Bearing recall is
zero by construction.}
\label{fig:balacc}
\end{figure}

Table~\ref{tab:desempenho-compacto} and
Fig.~\ref{fig:balacc} compare the methods. On IDMT every method is near the ceiling under both
protocols: the RMS-normalised level baseline, which only sees the crest factor, reaches
\VDiagBalRmsNormIdmtBlk{} balanced accuracy on held-out blocks, and the pipeline and the
classical features reach \VDiagBalDecIdmtBlk{} and \VDiagBalClsIdmtBlk{}. Because each class is one recording of one motor,
these numbers measure recording identity and cannot rank the methods.

On A2212 with a grouped split by file, the classical features reach \VDiagBalClsBldcFile{}
[\VDiagBalClsBldcFileLo{}, \VDiagBalClsBldcFileHi{}], the pipeline \VDiagBalDecBldcFile{}
[\VDiagBalDecBldcFileLo{}, \VDiagBalDecBldcFileHi{}] and the raw level baseline
\VDiagBalRmsRawBldcFile{}. The pipeline errs mostly on Healthy (recall
\VDiagRecDecBldcFileHealthy{}). The random split inflates balanced accuracy by
\VDiagInflDecBldc{} for the pipeline and \VDiagInflClsBldc{} for the classical features. Holding a
motor out lowers every method: the classical features fall to \VDiagBalClsBldcMotor{} and the
pipeline to \VDiagBalDecBldcMotor{}; Bearing recall is \VDiagRecDecBldcMotorBearing{} by
construction, since the only bearing motor is then absent from training. Without the M4 fold,
on Healthy versus Propeller of unseen motors, the classical features reach
\VDiagBalClsBldcNoMfour{}, the RMS-normalised level baseline \VDiagBalRmsNormBldcNoMfour{} and the
pipeline \VDiagBalDecBldcNoMfour{}. Compared like for like (Healthy and Propeller of motors
M1--M3 under the file split), holding the motor out costs \VDiagDropHPClsBldc{} for the classical
features and \VDiagDropHPDecBldc{} for the pipeline; the much larger drop of the full motor-out
protocol is dominated by the Bearing class losing its only training motor. The pipeline never
outperforms the classical features (it ties at the IDMT ceiling), and with a motor held out its
point estimate also falls below the level baseline, with overlapping intervals.

% matriz de confusão: figures/pdf (FigConf), fora do texto por espaço

\subsection{Interpretability: the tracked rotation}

\ifwip
On IDMT the median tracked rotation differs by class (\VRotIdmtGoodMed{},
\VRotIdmtBrokenMed{}, \VRotIdmtHeavyMed{}\,Hz), but class equals motor there. On A2212 its rank
correlation with the speed label of the files is \VRotBldcSpearman{} pooled over motors; within
motors it is \VRotBldcSpearmanMOne{} to \VRotBldcSpearmanMFour{} but significant only for M4
(\VRotBldcFilesMFour{} files), so it should not be read as shaft speed on that dataset.
\else
The tracker exposes a physical quantity that a technician can inspect. On IDMT, the median tracked rotation is \VRotIdmtGoodMed{},
\VRotIdmtBrokenMed{} and \VRotIdmtHeavyMed{}\,Hz for good, broken and heavy-load windows, with
inter-quartile ranges of at most \VRotIdmtHeavyIqr{}\,Hz; \VRotIdmtGoodOctPct{}\% of the good
windows lie above \VRotOctThr{} times the class median (all at \VRotOctRatioMin{}--\VRotOctRatioMax{}
times), an octave error. Because class equals motor on IDMT, these differences cannot be
attributed to the operating state. On A2212, pooled over the \VDataBldcRecordings{} files, the rank
correlation between the tracked rotation and the speed label in the file names is
\VRotBldcSpearman{}; within motors it is \VRotBldcSpearmanMOne{} to \VRotBldcSpearmanMFour{}
(\VRotBldcFilesMThree{}--\VRotBldcFilesMTwo{} files each), significant only for M4
($p \geq \audRotBldcSpearmanPMTwo{}$ for M1--M3). The absolute level
is not comparable across motors and \VRotBldcCeilPct{}\% of the windows lie within
\VRotCeilPctThr{}\% of the upper limit of the search range. The range and the grid step were chosen
on the IDMT training data, so on A2212 the tracker output should not be read as the shaft speed.
\fi

% espectro: figures/pdf (FigSpec), fora do texto por espaço

\ifwip\else
\subsection{Streaming decision}

On a stream of \VPcSeqWindows{} clean A2212 golden windows (\VPcSeqPerClass{} per class:
Healthy, Propeller, Bearing) the online HMM with change penalty \VPcSeqCp{} detected both changes
with a delay of \VPcSeqDelayC{} windows in C and \VPcSeqDelayPy{} in Python on the PC (maximum
probability difference \VPcSeqMaxProbDiff{}).
\ifchipmedido On the ESP32, the same stream read from flash by \texttt{\VChipVariant} gave a delay
of \VChipSeqDelay{} windows and online decisions equal to Python's on \VChipSeqSame{} of
\VChipSeqWindows{} windows; one HMM step takes \VChipSeqHmmUs{}\,$\mu$s.\fi{} Bearing is motor M4,
so this checks the implementation, not the detector.
\fi

% ======================================================================
\section{Discussion}

\emph{What the datasets allow us to conclude.} IDMT, used without its test set, supports no
statement about operating-state discrimination: one recording per class makes class, motor and
session the same variable, and a two-number level baseline is as good as any model. A2212 with a
file-level split supports statements about the state of motors that were \emph{already seen} in
training, and there a cheap classical feature set is better than our pipeline. Generalisation to a
new motor is the relevant question for a deployed monitor, and with \VDataBldcMotors{} motors, a
bearing fault present on one of them and speed ranges that differ by class~\cite{estacio2023}, the
data can only show that every method degrades; it cannot rank methods or estimate how a monitor
would behave on another motor. None of these experiments concerns failure prediction, and none
involves industrial motors. In our runs the random split inflated A2212 balanced accuracy only
moderately (\VDiagInflClsBldc{} to \VDiagInflDecBldc{}), and holding a motor out cost, like for
like on Healthy versus Propeller, a similar amount (\VDiagDropHPClsBldc{} for the classical features,
\VDiagDropHPDecBldc{} for the pipeline); the large drop of the full motor-out protocol is dominated
by the Bearing class, which exists only on M4. Published near-perfect accuracies on these datasets
should be read against that confound, which no split by segment or by file removes.

\ifwip\else
\emph{What the embedded work shows.} The C port reproduces the reference stage by stage within
tolerances that were frozen before the formal comparison, including the discrete Viterbi path and
the HMM decisions; the residual float32 failures occur only on a synthetic edge case whose
modulation peak lies at the float32 noise floor, they leave the decision unchanged, and they
disappear when the envelope filter state is kept in double.
\ifchipmedido
On a board without PSRAM, \texttt{\VChipVariant} runs \VParSegBldc{}-s windows with on-chip outputs
within those tolerances, but at a real-time factor of \VChipRtfBldc{}, and the \VParSegIdmt{}-s
window does not fit. We claim neither real-time operation nor low power consumption.
\else
Whether it runs in real time on the chip is still open: the PC timings do not transfer.
\fi
\fi

\emph{Limitations and future work.} Energy was not measured (no current sensor), and we do not
estimate it from datasheets. All audio came from the datasets through the serial link; a real
I2S microphone, a motor on a test bench with a tachometer to validate the rotation estimate, and
the official IDMT test set (unopened by design) are future work. An independent scientific audit
of datasets, equivalence and traceability (report in the repository) was carried out before the
on-chip measurements and does not cover them.

% ======================================================================
\section{Conclusion}

A music-transcription-inspired pipeline for acoustic motor-state classification was ported to
portable C and verified against its Python reference stage by stage, with tolerances frozen
before the formal comparison: all checks pass in double; in float32 every window of real audio
passes every stage, and only \VEqFloatBadCases{} synthetic edge windows fail the modulation,
feature and z-score checks, with an unchanged decision. On an ESP32 without PSRAM the reference
layout does not fit; the named variant \texttt{\VChipVariant} fits only \VParSegBldc{}-s windows,
matches the reference within the frozen tolerances on the chip, and is \VChipRtfBldc{} times slower
than real time, dominated by the table-free harmonic sum in emulated double precision; a float32
variant of that step or a board with PSRAM is the way forward. On the two public
datasets it does not beat simple baselines, and the datasets themselves cannot separate operating
state from motor identity. The value of the work lies in the verified, measurable implementation and
in making these limits explicit.

\bibliographystyle{IEEEtran}
\bibliography{references}

\end{document}
"
\ifdefined\usewip\wiptrue\else\wipfalse\fi

% Gerado por embarcados/src/python/gerar_artigo_numeros.py. NÃO EDITAR À MÃO.
% commit 9df3b2c280b69b02703d2c61d277e2bdf02cafad (árvore suja em embarcados/: True)
% Uma macro por número. Comentário: [categoria] fonte :: coluna.
% Prefixos: VPar código | VData/VDiag/VRot dados e validação | VEq equivalência (PC)
%           VEst 1·estimativa analítica | VPc 2·medição no PC | VBuild 3·build do firmware | VChip 4·medição no ESP32
\newif\ifchipmedido
\chipmedidofalse  % true só se results/manifests/bancada.json e esp32_lat_* existirem

\newcommand{\VParSr}{16{,}000}% [codigo] pipelines/_decifrae.py :: PreDecifrae.sr
\newcommand{\VParSrKhz}{16}% [codigo] pipelines/_decifrae.py :: PreDecifrae.sr/1000
\newcommand{\VParNfft}{8{,}192}% [codigo] pipelines/_decifrae.py :: PreDecifrae.n_fft
\newcommand{\VParHop}{1{,}024}% [codigo] pipelines/_decifrae.py :: PreDecifrae.hop
\newcommand{\VParFmin}{15}% [codigo] pipelines/_decifrae.py :: PreDecifrae.fmin
\newcommand{\VParFmax}{60}% [codigo] pipelines/_decifrae.py :: PreDecifrae.fmax
\newcommand{\VParStepCents}{2}% [codigo] pipelines/_decifrae.py :: PreDecifrae.passo_cents
\newcommand{\VParJ}{1{,}201}% [codigo] motorsound/rotacao.py :: len(rotacao.grade(fmin,fmax,passo))
\newcommand{\VParNbins}{4{,}097}% [codigo] motorsound/rotacao.py :: n_fft//2+1 (rfft)
\newcommand{\VParDf}{1.95}% [codigo] motorsound/rotacao.py :: sr/n_fft
\newcommand{\VParNharm}{48}% [codigo] pipelines/_decifrae.py :: N_HARM
\newcommand{\VParGroupOneLo}{1}% [codigo] pipelines/_decifrae.py :: GRUPOS[0][0]
\newcommand{\VParGroupOneHi}{4}% [codigo] pipelines/_decifrae.py :: GRUPOS[0][1]
\newcommand{\VParGroupTwoLo}{5}% [codigo] pipelines/_decifrae.py :: GRUPOS[1][0]
\newcommand{\VParGroupTwoHi}{12}% [codigo] pipelines/_decifrae.py :: GRUPOS[1][1]
\newcommand{\VParGroupThreeLo}{13}% [codigo] pipelines/_decifrae.py :: GRUPOS[2][0]
\newcommand{\VParGroupThreeHi}{24}% [codigo] pipelines/_decifrae.py :: GRUPOS[2][1]
\newcommand{\VParGroupFourLo}{25}% [codigo] pipelines/_decifrae.py :: GRUPOS[3][0]
\newcommand{\VParGroupFourHi}{48}% [codigo] pipelines/_decifrae.py :: GRUPOS[3][1]
\newcommand{\VParMaxJump}{50}% [codigo] pipelines/_decifrae.py :: PreDecifrae.max_salto
\newcommand{\VParMaxJumpCents}{100}% [codigo] pipelines/_decifrae.py :: max_salto*passo_cents
\newcommand{\VParOctStates}{600}% [codigo] motorsound/rotacao.py :: viterbi_rotacao: k_oct = round(1200/passo)
\newcommand{\VParWStep}{0.55}% [codigo] motorsound/rotacao.py :: VIT_W_STEP
\newcommand{\VParWOct}{7.0}% [codigo] motorsound/rotacao.py :: VIT_W_OCTAVE
\newcommand{\VParWSil}{3.2}% [codigo] motorsound/rotacao.py :: VIT_W_SIL
\newcommand{\VParWEnergy}{0.45}% [codigo] motorsound/rotacao.py :: emissao_baixo: peso de e_norm
\newcommand{\VParEmisEps}{\ensuremath{10^{-3}}}% [codigo] motorsound/rotacao.py :: emissao_baixo: eps do -log
\newcommand{\VParTempCos}{0.1}% [codigo] pipelines/_decifrae.py :: PostDecifrae.temp_cos
\newcommand{\VParButterOrder}{4}% [codigo] pipelines/_decifrae.py :: _envelope_mod: butter(ordem, fc)
\newcommand{\VParButterFc}{200}% [codigo] pipelines/_decifrae.py :: _envelope_mod: butter(ordem, fc)
\newcommand{\VParModLo}{2}% [codigo] pipelines/_decifrae.py :: _envelope_mod: banda
\newcommand{\VParModHi}{100}% [codigo] pipelines/_decifrae.py :: _envelope_mod: banda
\newcommand{\VParEnvRateHz}{1{,}000}% [codigo] pipelines/_decifrae.py :: _envelope_mod: [:: sr//1000]
\newcommand{\VParDecim}{16}% [codigo] pipelines/_decifrae.py :: sr // 1000
\newcommand{\VParBandLo}{10}% [codigo] pipelines/_decifrae.py :: extrair: banda das features
\newcommand{\VParBandHi}{4{,}000}% [codigo] pipelines/_decifrae.py :: extrair: banda das features
\newcommand{\VParSegIdmt}{3}% [codigo] motorsound/config.py :: AVAL_IDMT.seg_s
\newcommand{\VParSegBldc}{1}% [codigo] motorsound/config.py :: AVAL_A2212.seg_s
\newcommand{\VParRmsTarget}{0.1}% [codigo] motorsound/config.py :: AvaliacaoConfig.rms_alvo
\newcommand{\VParFramesIdmt}{39}% [codigo] motorsound/rotacao.py :: 1+(N-n_fft)//hop (stft_mag)
\newcommand{\VParFramesBldc}{8}% [codigo] motorsound/rotacao.py :: 1+(N-n_fft)//hop (stft_mag)
\newcommand{\VParNClassicFeat}{32}% [codigo] pipelines/_classicas.py :: N_FEATURES
\newcommand{\VParNChangePen}{6}% [codigo] embarcados/gold/seqs/idmt_seq00/params.json :: change_penalties
\newcommand{\VParCpMin}{0.5}% [codigo] embarcados/gold/seqs/idmt_seq00/params.json :: change_penalties
\newcommand{\VParCpMax}{16}% [codigo] embarcados/gold/seqs/idmt_seq00/params.json :: change_penalties
\newcommand{\VParNLogprob}{-50}% [codigo] embarcados/gold/seqs/idmt_seq00/params.json :: n_logprob
\newcommand{\VParHmmTemp}{1}% [codigo] embarcados/gold/seqs/idmt_seq00/params.json :: temp
\newcommand{\VDataIdmtWindows}{357}% [dados] embarcados/results/predictions/diagnostico_janelas_20261023_222831.csv :: len(unique(rec_id,idx)) da p11
\newcommand{\VDataIdmtRecordings}{3}% [dados] embarcados/results/predictions/diagnostico_janelas_20261023_222831.csv :: nunique(rec_id)
\newcommand{\VDataIdmtMotors}{3}% [dados] embarcados/results/predictions/diagnostico_janelas_20261023_222831.csv :: nunique(motor)
\newcommand{\VDataIdmtClasses}{3}% [dados] embarcados/results/predictions/diagnostico_janelas_20261023_222831.csv :: nunique(classe)
\newcommand{\VDataBldcWindows}{430}% [dados] embarcados/results/predictions/diagnostico_janelas_20261023_222831.csv :: len(unique(rec_id,idx)) da p11
\newcommand{\VDataBldcRecordings}{43}% [dados] embarcados/results/predictions/diagnostico_janelas_20261023_222831.csv :: nunique(rec_id)
\newcommand{\VDataBldcMotors}{4}% [dados] embarcados/results/predictions/diagnostico_janelas_20261023_222831.csv :: nunique(motor)
\newcommand{\VDataBldcClasses}{3}% [dados] embarcados/results/predictions/diagnostico_janelas_20261023_222831.csv :: nunique(classe)
\newcommand{\VDataBldcFilesHealthy}{22}% [dados] embarcados/results/predictions/diagnostico_janelas_20261023_222831.csv :: nunique(rec_id) por classe
\newcommand{\VDataBldcFilesPropeller}{14}% [dados] embarcados/results/predictions/diagnostico_janelas_20261023_222831.csv :: nunique(rec_id) por classe
\newcommand{\VDataBldcFilesBearing}{7}% [dados] embarcados/results/predictions/diagnostico_janelas_20261023_222831.csv :: nunique(rec_id) por classe
\newcommand{\VDataIdmtRandGroups}{18}% [dados] embarcados/results/csv/diagnostico_20261023_222831.csv :: n_grupos
\newcommand{\VDataIdmtBlkGroups}{15}% [dados] embarcados/results/csv/diagnostico_20261023_222831.csv :: n_grupos
\newcommand{\VDataBldcRandGroups}{43}% [dados] embarcados/results/csv/diagnostico_20261023_222831.csv :: n_grupos
\newcommand{\VDataBldcFileGroups}{43}% [dados] embarcados/results/csv/diagnostico_20261023_222831.csv :: n_grupos
\newcommand{\VDataBldcMotorGroups}{43}% [dados] embarcados/results/csv/diagnostico_20261023_222831.csv :: n_grupos
\newcommand{\VDataBldcNoMfourGroups}{36}% [dados] embarcados/results/csv/diagnostico_20261023_222831.csv :: n_grupos
\newcommand{\VDataNBoot}{2{,}000}% [dados] embarcados/results/csv/diagnostico_20261023_222831_meta.json :: n_boot
\newcommand{\VDataCiLevel}{95}% [dados] embarcados/results/csv/diagnostico_20261023_222831.csv :: nível do IC no nome da coluna ic<nível>_lo
\newcommand{\VDataSeedBoot}{20261023}% [dados] embarcados/results/csv/diagnostico_20261023_222831_meta.json :: seed_boot
\newcommand{\VDiagBalRmsNormIdmtRand}{0.992}% [diagnostico] embarcados/results/csv/diagnostico_20261023_222831.csv :: valor[b00_rmspico_norm,idmt,aleatorio,acuracia_balanceada,_macro]
\newcommand{\VDiagBalRmsNormIdmtRandLo}{0.983}% [diagnostico] embarcados/results/csv/diagnostico_20261023_222831.csv :: ic95_lo[b00_rmspico_norm,idmt,aleatorio,acuracia_balanceada,_macro]
\newcommand{\VDiagBalRmsNormIdmtRandHi}{1.000}% [diagnostico] embarcados/results/csv/diagnostico_20261023_222831.csv :: ic95_hi[b00_rmspico_norm,idmt,aleatorio,acuracia_balanceada,_macro]
\newcommand{\VDiagAucRmsNormIdmtRand}{0.998}% [diagnostico] embarcados/results/csv/diagnostico_20261023_222831.csv :: valor[b00_rmspico_norm,idmt,aleatorio,roc_auc_ovr_macro,_macro]
\newcommand{\VDiagFoneRmsNormIdmtRand}{0.992}% [diagnostico] embarcados/results/csv/diagnostico_20261023_222831.csv :: valor[b00_rmspico_norm,idmt,aleatorio,f1_macro,_macro]
\newcommand{\VDiagBalRmsRawIdmtRand}{0.997}% [diagnostico] embarcados/results/csv/diagnostico_20261023_222831.csv :: valor[b00_rmspico_bruto,idmt,aleatorio,acuracia_balanceada,_macro]
\newcommand{\VDiagBalRmsRawIdmtRandLo}{0.990}% [diagnostico] embarcados/results/csv/diagnostico_20261023_222831.csv :: ic95_lo[b00_rmspico_bruto,idmt,aleatorio,acuracia_balanceada,_macro]
\newcommand{\VDiagBalRmsRawIdmtRandHi}{1.000}% [diagnostico] embarcados/results/csv/diagnostico_20261023_222831.csv :: ic95_hi[b00_rmspico_bruto,idmt,aleatorio,acuracia_balanceada,_macro]
\newcommand{\VDiagAucRmsRawIdmtRand}{1.000}% [diagnostico] embarcados/results/csv/diagnostico_20261023_222831.csv :: valor[b00_rmspico_bruto,idmt,aleatorio,roc_auc_ovr_macro,_macro]
\newcommand{\VDiagFoneRmsRawIdmtRand}{0.997}% [diagnostico] embarcados/results/csv/diagnostico_20261023_222831.csv :: valor[b00_rmspico_bruto,idmt,aleatorio,f1_macro,_macro]
\newcommand{\VDiagBalClsIdmtRand}{1.000}% [diagnostico] embarcados/results/csv/diagnostico_20261023_222831.csv :: valor[p04_classicas,idmt,aleatorio,acuracia_balanceada,_macro]
\newcommand{\VDiagBalClsIdmtRandLo}{1.000}% [diagnostico] embarcados/results/csv/diagnostico_20261023_222831.csv :: ic95_lo[p04_classicas,idmt,aleatorio,acuracia_balanceada,_macro]
\newcommand{\VDiagBalClsIdmtRandHi}{1.000}% [diagnostico] embarcados/results/csv/diagnostico_20261023_222831.csv :: ic95_hi[p04_classicas,idmt,aleatorio,acuracia_balanceada,_macro]
\newcommand{\VDiagAucClsIdmtRand}{1.000}% [diagnostico] embarcados/results/csv/diagnostico_20261023_222831.csv :: valor[p04_classicas,idmt,aleatorio,roc_auc_ovr_macro,_macro]
\newcommand{\VDiagFoneClsIdmtRand}{1.000}% [diagnostico] embarcados/results/csv/diagnostico_20261023_222831.csv :: valor[p04_classicas,idmt,aleatorio,f1_macro,_macro]
\newcommand{\VDiagBalDecIdmtRand}{1.000}% [diagnostico] embarcados/results/csv/diagnostico_20261023_222831.csv :: valor[p11_decifrae_harmsum,idmt,aleatorio,acuracia_balanceada,_macro]
\newcommand{\VDiagBalDecIdmtRandLo}{1.000}% [diagnostico] embarcados/results/csv/diagnostico_20261023_222831.csv :: ic95_lo[p11_decifrae_harmsum,idmt,aleatorio,acuracia_balanceada,_macro]
\newcommand{\VDiagBalDecIdmtRandHi}{1.000}% [diagnostico] embarcados/results/csv/diagnostico_20261023_222831.csv :: ic95_hi[p11_decifrae_harmsum,idmt,aleatorio,acuracia_balanceada,_macro]
\newcommand{\VDiagAucDecIdmtRand}{1.000}% [diagnostico] embarcados/results/csv/diagnostico_20261023_222831.csv :: valor[p11_decifrae_harmsum,idmt,aleatorio,roc_auc_ovr_macro,_macro]
\newcommand{\VDiagFoneDecIdmtRand}{1.000}% [diagnostico] embarcados/results/csv/diagnostico_20261023_222831.csv :: valor[p11_decifrae_harmsum,idmt,aleatorio,f1_macro,_macro]
\newcommand{\VDiagBalRmsNormIdmtBlk}{0.991}% [diagnostico] embarcados/results/csv/diagnostico_20261023_222831.csv :: valor[b00_rmspico_norm,idmt,blocos,acuracia_balanceada,_macro]
\newcommand{\VDiagBalRmsNormIdmtBlkLo}{0.982}% [diagnostico] embarcados/results/csv/diagnostico_20261023_222831.csv :: ic95_lo[b00_rmspico_norm,idmt,blocos,acuracia_balanceada,_macro]
\newcommand{\VDiagBalRmsNormIdmtBlkHi}{0.997}% [diagnostico] embarcados/results/csv/diagnostico_20261023_222831.csv :: ic95_hi[b00_rmspico_norm,idmt,blocos,acuracia_balanceada,_macro]
\newcommand{\VDiagAucRmsNormIdmtBlk}{0.998}% [diagnostico] embarcados/results/csv/diagnostico_20261023_222831.csv :: valor[b00_rmspico_norm,idmt,blocos,roc_auc_ovr_macro,_macro]
\newcommand{\VDiagFoneRmsNormIdmtBlk}{0.991}% [diagnostico] embarcados/results/csv/diagnostico_20261023_222831.csv :: valor[b00_rmspico_norm,idmt,blocos,f1_macro,_macro]
\newcommand{\VDiagBalRmsRawIdmtBlk}{0.993}% [diagnostico] embarcados/results/csv/diagnostico_20261023_222831.csv :: valor[b00_rmspico_bruto,idmt,blocos,acuracia_balanceada,_macro]
\newcommand{\VDiagBalRmsRawIdmtBlkLo}{0.987}% [diagnostico] embarcados/results/csv/diagnostico_20261023_222831.csv :: ic95_lo[b00_rmspico_bruto,idmt,blocos,acuracia_balanceada,_macro]
\newcommand{\VDiagBalRmsRawIdmtBlkHi}{1.000}% [diagnostico] embarcados/results/csv/diagnostico_20261023_222831.csv :: ic95_hi[b00_rmspico_bruto,idmt,blocos,acuracia_balanceada,_macro]
\newcommand{\VDiagAucRmsRawIdmtBlk}{1.000}% [diagnostico] embarcados/results/csv/diagnostico_20261023_222831.csv :: valor[b00_rmspico_bruto,idmt,blocos,roc_auc_ovr_macro,_macro]
\newcommand{\VDiagFoneRmsRawIdmtBlk}{0.994}% [diagnostico] embarcados/results/csv/diagnostico_20261023_222831.csv :: valor[b00_rmspico_bruto,idmt,blocos,f1_macro,_macro]
\newcommand{\VDiagBalClsIdmtBlk}{1.000}% [diagnostico] embarcados/results/csv/diagnostico_20261023_222831.csv :: valor[p04_classicas,idmt,blocos,acuracia_balanceada,_macro]
\newcommand{\VDiagBalClsIdmtBlkLo}{1.000}% [diagnostico] embarcados/results/csv/diagnostico_20261023_222831.csv :: ic95_lo[p04_classicas,idmt,blocos,acuracia_balanceada,_macro]
\newcommand{\VDiagBalClsIdmtBlkHi}{1.000}% [diagnostico] embarcados/results/csv/diagnostico_20261023_222831.csv :: ic95_hi[p04_classicas,idmt,blocos,acuracia_balanceada,_macro]
\newcommand{\VDiagAucClsIdmtBlk}{1.000}% [diagnostico] embarcados/results/csv/diagnostico_20261023_222831.csv :: valor[p04_classicas,idmt,blocos,roc_auc_ovr_macro,_macro]
\newcommand{\VDiagFoneClsIdmtBlk}{1.000}% [diagnostico] embarcados/results/csv/diagnostico_20261023_222831.csv :: valor[p04_classicas,idmt,blocos,f1_macro,_macro]
\newcommand{\VDiagBalDecIdmtBlk}{1.000}% [diagnostico] embarcados/results/csv/diagnostico_20261023_222831.csv :: valor[p11_decifrae_harmsum,idmt,blocos,acuracia_balanceada,_macro]
\newcommand{\VDiagBalDecIdmtBlkLo}{1.000}% [diagnostico] embarcados/results/csv/diagnostico_20261023_222831.csv :: ic95_lo[p11_decifrae_harmsum,idmt,blocos,acuracia_balanceada,_macro]
\newcommand{\VDiagBalDecIdmtBlkHi}{1.000}% [diagnostico] embarcados/results/csv/diagnostico_20261023_222831.csv :: ic95_hi[p11_decifrae_harmsum,idmt,blocos,acuracia_balanceada,_macro]
\newcommand{\VDiagAucDecIdmtBlk}{1.000}% [diagnostico] embarcados/results/csv/diagnostico_20261023_222831.csv :: valor[p11_decifrae_harmsum,idmt,blocos,roc_auc_ovr_macro,_macro]
\newcommand{\VDiagFoneDecIdmtBlk}{1.000}% [diagnostico] embarcados/results/csv/diagnostico_20261023_222831.csv :: valor[p11_decifrae_harmsum,idmt,blocos,f1_macro,_macro]
\newcommand{\VDiagBalRmsNormBldcRand}{0.564}% [diagnostico] embarcados/results/csv/diagnostico_20261023_222831.csv :: valor[b00_rmspico_norm,a2212,aleatorio,acuracia_balanceada,_macro]
\newcommand{\VDiagBalRmsNormBldcRandLo}{0.515}% [diagnostico] embarcados/results/csv/diagnostico_20261023_222831.csv :: ic95_lo[b00_rmspico_norm,a2212,aleatorio,acuracia_balanceada,_macro]
\newcommand{\VDiagBalRmsNormBldcRandHi}{0.607}% [diagnostico] embarcados/results/csv/diagnostico_20261023_222831.csv :: ic95_hi[b00_rmspico_norm,a2212,aleatorio,acuracia_balanceada,_macro]
\newcommand{\VDiagAucRmsNormBldcRand}{0.784}% [diagnostico] embarcados/results/csv/diagnostico_20261023_222831.csv :: valor[b00_rmspico_norm,a2212,aleatorio,roc_auc_ovr_macro,_macro]
\newcommand{\VDiagFoneRmsNormBldcRand}{0.533}% [diagnostico] embarcados/results/csv/diagnostico_20261023_222831.csv :: valor[b00_rmspico_norm,a2212,aleatorio,f1_macro,_macro]
\newcommand{\VDiagBalRmsRawBldcRand}{0.625}% [diagnostico] embarcados/results/csv/diagnostico_20261023_222831.csv :: valor[b00_rmspico_bruto,a2212,aleatorio,acuracia_balanceada,_macro]
\newcommand{\VDiagBalRmsRawBldcRandLo}{0.571}% [diagnostico] embarcados/results/csv/diagnostico_20261023_222831.csv :: ic95_lo[b00_rmspico_bruto,a2212,aleatorio,acuracia_balanceada,_macro]
\newcommand{\VDiagBalRmsRawBldcRandHi}{0.660}% [diagnostico] embarcados/results/csv/diagnostico_20261023_222831.csv :: ic95_hi[b00_rmspico_bruto,a2212,aleatorio,acuracia_balanceada,_macro]
\newcommand{\VDiagAucRmsRawBldcRand}{0.786}% [diagnostico] embarcados/results/csv/diagnostico_20261023_222831.csv :: valor[b00_rmspico_bruto,a2212,aleatorio,roc_auc_ovr_macro,_macro]
\newcommand{\VDiagFoneRmsRawBldcRand}{0.588}% [diagnostico] embarcados/results/csv/diagnostico_20261023_222831.csv :: valor[b00_rmspico_bruto,a2212,aleatorio,f1_macro,_macro]
\newcommand{\VDiagBalClsBldcRand}{0.998}% [diagnostico] embarcados/results/csv/diagnostico_20261023_222831.csv :: valor[p04_classicas,a2212,aleatorio,acuracia_balanceada,_macro]
\newcommand{\VDiagBalClsBldcRandLo}{0.995}% [diagnostico] embarcados/results/csv/diagnostico_20261023_222831.csv :: ic95_lo[p04_classicas,a2212,aleatorio,acuracia_balanceada,_macro]
\newcommand{\VDiagBalClsBldcRandHi}{1.000}% [diagnostico] embarcados/results/csv/diagnostico_20261023_222831.csv :: ic95_hi[p04_classicas,a2212,aleatorio,acuracia_balanceada,_macro]
\newcommand{\VDiagAucClsBldcRand}{1.000}% [diagnostico] embarcados/results/csv/diagnostico_20261023_222831.csv :: valor[p04_classicas,a2212,aleatorio,roc_auc_ovr_macro,_macro]
\newcommand{\VDiagFoneClsBldcRand}{0.998}% [diagnostico] embarcados/results/csv/diagnostico_20261023_222831.csv :: valor[p04_classicas,a2212,aleatorio,f1_macro,_macro]
\newcommand{\VDiagBalDecBldcRand}{0.849}% [diagnostico] embarcados/results/csv/diagnostico_20261023_222831.csv :: valor[p11_decifrae_harmsum,a2212,aleatorio,acuracia_balanceada,_macro]
\newcommand{\VDiagBalDecBldcRandLo}{0.790}% [diagnostico] embarcados/results/csv/diagnostico_20261023_222831.csv :: ic95_lo[p11_decifrae_harmsum,a2212,aleatorio,acuracia_balanceada,_macro]
\newcommand{\VDiagBalDecBldcRandHi}{0.905}% [diagnostico] embarcados/results/csv/diagnostico_20261023_222831.csv :: ic95_hi[p11_decifrae_harmsum,a2212,aleatorio,acuracia_balanceada,_macro]
\newcommand{\VDiagAucDecBldcRand}{0.955}% [diagnostico] embarcados/results/csv/diagnostico_20261023_222831.csv :: valor[p11_decifrae_harmsum,a2212,aleatorio,roc_auc_ovr_macro,_macro]
\newcommand{\VDiagFoneDecBldcRand}{0.789}% [diagnostico] embarcados/results/csv/diagnostico_20261023_222831.csv :: valor[p11_decifrae_harmsum,a2212,aleatorio,f1_macro,_macro]
\newcommand{\VDiagBalRmsNormBldcFile}{0.562}% [diagnostico] embarcados/results/csv/diagnostico_20261023_222831.csv :: valor[b00_rmspico_norm,a2212,arquivo,acuracia_balanceada,_macro]
\newcommand{\VDiagBalRmsNormBldcFileLo}{0.512}% [diagnostico] embarcados/results/csv/diagnostico_20261023_222831.csv :: ic95_lo[b00_rmspico_norm,a2212,arquivo,acuracia_balanceada,_macro]
\newcommand{\VDiagBalRmsNormBldcFileHi}{0.606}% [diagnostico] embarcados/results/csv/diagnostico_20261023_222831.csv :: ic95_hi[b00_rmspico_norm,a2212,arquivo,acuracia_balanceada,_macro]
\newcommand{\VDiagAucRmsNormBldcFile}{0.759}% [diagnostico] embarcados/results/csv/diagnostico_20261023_222831.csv :: valor[b00_rmspico_norm,a2212,arquivo,roc_auc_ovr_macro,_macro]
\newcommand{\VDiagFoneRmsNormBldcFile}{0.532}% [diagnostico] embarcados/results/csv/diagnostico_20261023_222831.csv :: valor[b00_rmspico_norm,a2212,arquivo,f1_macro,_macro]
\newcommand{\VDiagBalRmsRawBldcFile}{0.608}% [diagnostico] embarcados/results/csv/diagnostico_20261023_222831.csv :: valor[b00_rmspico_bruto,a2212,arquivo,acuracia_balanceada,_macro]
\newcommand{\VDiagBalRmsRawBldcFileLo}{0.544}% [diagnostico] embarcados/results/csv/diagnostico_20261023_222831.csv :: ic95_lo[b00_rmspico_bruto,a2212,arquivo,acuracia_balanceada,_macro]
\newcommand{\VDiagBalRmsRawBldcFileHi}{0.656}% [diagnostico] embarcados/results/csv/diagnostico_20261023_222831.csv :: ic95_hi[b00_rmspico_bruto,a2212,arquivo,acuracia_balanceada,_macro]
\newcommand{\VDiagAucRmsRawBldcFile}{0.706}% [diagnostico] embarcados/results/csv/diagnostico_20261023_222831.csv :: valor[b00_rmspico_bruto,a2212,arquivo,roc_auc_ovr_macro,_macro]
\newcommand{\VDiagFoneRmsRawBldcFile}{0.576}% [diagnostico] embarcados/results/csv/diagnostico_20261023_222831.csv :: valor[b00_rmspico_bruto,a2212,arquivo,f1_macro,_macro]
\newcommand{\VDiagBalClsBldcFile}{0.981}% [diagnostico] embarcados/results/csv/diagnostico_20261023_222831.csv :: valor[p04_classicas,a2212,arquivo,acuracia_balanceada,_macro]
\newcommand{\VDiagBalClsBldcFileLo}{0.950}% [diagnostico] embarcados/results/csv/diagnostico_20261023_222831.csv :: ic95_lo[p04_classicas,a2212,arquivo,acuracia_balanceada,_macro]
\newcommand{\VDiagBalClsBldcFileHi}{0.998}% [diagnostico] embarcados/results/csv/diagnostico_20261023_222831.csv :: ic95_hi[p04_classicas,a2212,arquivo,acuracia_balanceada,_macro]
\newcommand{\VDiagAucClsBldcFile}{0.996}% [diagnostico] embarcados/results/csv/diagnostico_20261023_222831.csv :: valor[p04_classicas,a2212,arquivo,roc_auc_ovr_macro,_macro]
\newcommand{\VDiagFoneClsBldcFile}{0.982}% [diagnostico] embarcados/results/csv/diagnostico_20261023_222831.csv :: valor[p04_classicas,a2212,arquivo,f1_macro,_macro]
\newcommand{\VDiagBalDecBldcFile}{0.816}% [diagnostico] embarcados/results/csv/diagnostico_20261023_222831.csv :: valor[p11_decifrae_harmsum,a2212,arquivo,acuracia_balanceada,_macro]
\newcommand{\VDiagBalDecBldcFileLo}{0.747}% [diagnostico] embarcados/results/csv/diagnostico_20261023_222831.csv :: ic95_lo[p11_decifrae_harmsum,a2212,arquivo,acuracia_balanceada,_macro]
\newcommand{\VDiagBalDecBldcFileHi}{0.881}% [diagnostico] embarcados/results/csv/diagnostico_20261023_222831.csv :: ic95_hi[p11_decifrae_harmsum,a2212,arquivo,acuracia_balanceada,_macro]
\newcommand{\VDiagAucDecBldcFile}{0.935}% [diagnostico] embarcados/results/csv/diagnostico_20261023_222831.csv :: valor[p11_decifrae_harmsum,a2212,arquivo,roc_auc_ovr_macro,_macro]
\newcommand{\VDiagFoneDecBldcFile}{0.753}% [diagnostico] embarcados/results/csv/diagnostico_20261023_222831.csv :: valor[p11_decifrae_harmsum,a2212,arquivo,f1_macro,_macro]
\newcommand{\VDiagBalRmsNormBldcMotor}{0.542}% [diagnostico] embarcados/results/csv/diagnostico_20261023_222831.csv :: valor[b00_rmspico_norm,a2212,motor,acuracia_balanceada,_macro]
\newcommand{\VDiagBalRmsNormBldcMotorLo}{0.485}% [diagnostico] embarcados/results/csv/diagnostico_20261023_222831.csv :: ic95_lo[b00_rmspico_norm,a2212,motor,acuracia_balanceada,_macro]
\newcommand{\VDiagBalRmsNormBldcMotorHi}{0.596}% [diagnostico] embarcados/results/csv/diagnostico_20261023_222831.csv :: ic95_hi[b00_rmspico_norm,a2212,motor,acuracia_balanceada,_macro]
\newcommand{\VDiagAucRmsNormBldcMotor}{0.489}% [diagnostico] embarcados/results/csv/diagnostico_20261023_222831.csv :: valor[b00_rmspico_norm,a2212,motor,roc_auc_ovr_macro,_macro]
\newcommand{\VDiagFoneRmsNormBldcMotor}{0.509}% [diagnostico] embarcados/results/csv/diagnostico_20261023_222831.csv :: valor[b00_rmspico_norm,a2212,motor,f1_macro,_macro]
\newcommand{\VDiagBalRmsRawBldcMotor}{0.503}% [diagnostico] embarcados/results/csv/diagnostico_20261023_222831.csv :: valor[b00_rmspico_bruto,a2212,motor,acuracia_balanceada,_macro]
\newcommand{\VDiagBalRmsRawBldcMotorLo}{0.412}% [diagnostico] embarcados/results/csv/diagnostico_20261023_222831.csv :: ic95_lo[b00_rmspico_bruto,a2212,motor,acuracia_balanceada,_macro]
\newcommand{\VDiagBalRmsRawBldcMotorHi}{0.580}% [diagnostico] embarcados/results/csv/diagnostico_20261023_222831.csv :: ic95_hi[b00_rmspico_bruto,a2212,motor,acuracia_balanceada,_macro]
\newcommand{\VDiagAucRmsRawBldcMotor}{0.508}% [diagnostico] embarcados/results/csv/diagnostico_20261023_222831.csv :: valor[b00_rmspico_bruto,a2212,motor,roc_auc_ovr_macro,_macro]
\newcommand{\VDiagFoneRmsRawBldcMotor}{0.509}% [diagnostico] embarcados/results/csv/diagnostico_20261023_222831.csv :: valor[b00_rmspico_bruto,a2212,motor,f1_macro,_macro]
\newcommand{\VDiagBalClsBldcMotor}{0.613}% [diagnostico] embarcados/results/csv/diagnostico_20261023_222831.csv :: valor[p04_classicas,a2212,motor,acuracia_balanceada,_macro]
\newcommand{\VDiagBalClsBldcMotorLo}{0.579}% [diagnostico] embarcados/results/csv/diagnostico_20261023_222831.csv :: ic95_lo[p04_classicas,a2212,motor,acuracia_balanceada,_macro]
\newcommand{\VDiagBalClsBldcMotorHi}{0.641}% [diagnostico] embarcados/results/csv/diagnostico_20261023_222831.csv :: ic95_hi[p04_classicas,a2212,motor,acuracia_balanceada,_macro]
\newcommand{\VDiagAucClsBldcMotor}{0.618}% [diagnostico] embarcados/results/csv/diagnostico_20261023_222831.csv :: valor[p04_classicas,a2212,motor,roc_auc_ovr_macro,_macro]
\newcommand{\VDiagFoneClsBldcMotor}{0.563}% [diagnostico] embarcados/results/csv/diagnostico_20261023_222831.csv :: valor[p04_classicas,a2212,motor,f1_macro,_macro]
\newcommand{\VDiagBalDecBldcMotor}{0.456}% [diagnostico] embarcados/results/csv/diagnostico_20261023_222831.csv :: valor[p11_decifrae_harmsum,a2212,motor,acuracia_balanceada,_macro]
\newcommand{\VDiagBalDecBldcMotorLo}{0.397}% [diagnostico] embarcados/results/csv/diagnostico_20261023_222831.csv :: ic95_lo[p11_decifrae_harmsum,a2212,motor,acuracia_balanceada,_macro]
\newcommand{\VDiagBalDecBldcMotorHi}{0.519}% [diagnostico] embarcados/results/csv/diagnostico_20261023_222831.csv :: ic95_hi[p11_decifrae_harmsum,a2212,motor,acuracia_balanceada,_macro]
\newcommand{\VDiagAucDecBldcMotor}{0.484}% [diagnostico] embarcados/results/csv/diagnostico_20261023_222831.csv :: valor[p11_decifrae_harmsum,a2212,motor,roc_auc_ovr_macro,_macro]
\newcommand{\VDiagFoneDecBldcMotor}{0.418}% [diagnostico] embarcados/results/csv/diagnostico_20261023_222831.csv :: valor[p11_decifrae_harmsum,a2212,motor,f1_macro,_macro]
\newcommand{\VDiagBalRmsNormBldcNoMfour}{0.814}% [diagnostico] embarcados/results/csv/diagnostico_20261023_222831.csv :: valor[b00_rmspico_norm,a2212,motor_sem_M4,acuracia_balanceada,_macro]
\newcommand{\VDiagBalRmsNormBldcNoMfourLo}{0.720}% [diagnostico] embarcados/results/csv/diagnostico_20261023_222831.csv :: ic95_lo[b00_rmspico_norm,a2212,motor_sem_M4,acuracia_balanceada,_macro]
\newcommand{\VDiagBalRmsNormBldcNoMfourHi}{0.887}% [diagnostico] embarcados/results/csv/diagnostico_20261023_222831.csv :: ic95_hi[b00_rmspico_norm,a2212,motor_sem_M4,acuracia_balanceada,_macro]
\newcommand{\VDiagAucRmsNormBldcNoMfour}{0.871}% [diagnostico] embarcados/results/csv/diagnostico_20261023_222831.csv :: valor[b00_rmspico_norm,a2212,motor_sem_M4,roc_auc_ovr_macro,_macro]
\newcommand{\VDiagFoneRmsNormBldcNoMfour}{0.824}% [diagnostico] embarcados/results/csv/diagnostico_20261023_222831.csv :: valor[b00_rmspico_norm,a2212,motor_sem_M4,f1_macro,_macro]
\newcommand{\VDiagBalRmsRawBldcNoMfour}{0.754}% [diagnostico] embarcados/results/csv/diagnostico_20261023_222831.csv :: valor[b00_rmspico_bruto,a2212,motor_sem_M4,acuracia_balanceada,_macro]
\newcommand{\VDiagBalRmsRawBldcNoMfourLo}{0.623}% [diagnostico] embarcados/results/csv/diagnostico_20261023_222831.csv :: ic95_lo[b00_rmspico_bruto,a2212,motor_sem_M4,acuracia_balanceada,_macro]
\newcommand{\VDiagBalRmsRawBldcNoMfourHi}{0.873}% [diagnostico] embarcados/results/csv/diagnostico_20261023_222831.csv :: ic95_hi[b00_rmspico_bruto,a2212,motor_sem_M4,acuracia_balanceada,_macro]
\newcommand{\VDiagAucRmsRawBldcNoMfour}{0.912}% [diagnostico] embarcados/results/csv/diagnostico_20261023_222831.csv :: valor[b00_rmspico_bruto,a2212,motor_sem_M4,roc_auc_ovr_macro,_macro]
\newcommand{\VDiagFoneRmsRawBldcNoMfour}{0.823}% [diagnostico] embarcados/results/csv/diagnostico_20261023_222831.csv :: valor[b00_rmspico_bruto,a2212,motor_sem_M4,f1_macro,_macro]
\newcommand{\VDiagBalClsBldcNoMfour}{0.920}% [diagnostico] embarcados/results/csv/diagnostico_20261023_222831.csv :: valor[p04_classicas,a2212,motor_sem_M4,acuracia_balanceada,_macro]
\newcommand{\VDiagBalClsBldcNoMfourLo}{0.868}% [diagnostico] embarcados/results/csv/diagnostico_20261023_222831.csv :: ic95_lo[p04_classicas,a2212,motor_sem_M4,acuracia_balanceada,_macro]
\newcommand{\VDiagBalClsBldcNoMfourHi}{0.962}% [diagnostico] embarcados/results/csv/diagnostico_20261023_222831.csv :: ic95_hi[p04_classicas,a2212,motor_sem_M4,acuracia_balanceada,_macro]
\newcommand{\VDiagAucClsBldcNoMfour}{0.984}% [diagnostico] embarcados/results/csv/diagnostico_20261023_222831.csv :: valor[p04_classicas,a2212,motor_sem_M4,roc_auc_ovr_macro,_macro]
\newcommand{\VDiagFoneClsBldcNoMfour}{0.923}% [diagnostico] embarcados/results/csv/diagnostico_20261023_222831.csv :: valor[p04_classicas,a2212,motor_sem_M4,f1_macro,_macro]
\newcommand{\VDiagBalDecBldcNoMfour}{0.685}% [diagnostico] embarcados/results/csv/diagnostico_20261023_222831.csv :: valor[p11_decifrae_harmsum,a2212,motor_sem_M4,acuracia_balanceada,_macro]
\newcommand{\VDiagBalDecBldcNoMfourLo}{0.593}% [diagnostico] embarcados/results/csv/diagnostico_20261023_222831.csv :: ic95_lo[p11_decifrae_harmsum,a2212,motor_sem_M4,acuracia_balanceada,_macro]
\newcommand{\VDiagBalDecBldcNoMfourHi}{0.774}% [diagnostico] embarcados/results/csv/diagnostico_20261023_222831.csv :: ic95_hi[p11_decifrae_harmsum,a2212,motor_sem_M4,acuracia_balanceada,_macro]
\newcommand{\VDiagAucDecBldcNoMfour}{0.822}% [diagnostico] embarcados/results/csv/diagnostico_20261023_222831.csv :: valor[p11_decifrae_harmsum,a2212,motor_sem_M4,roc_auc_ovr_macro,_macro]
\newcommand{\VDiagFoneDecBldcNoMfour}{0.685}% [diagnostico] embarcados/results/csv/diagnostico_20261023_222831.csv :: valor[p11_decifrae_harmsum,a2212,motor_sem_M4,f1_macro,_macro]
\newcommand{\VDiagRecDecBldcFileHealthy}{0.577}% [diagnostico] embarcados/results/csv/diagnostico_20261023_222831.csv :: valor[p11,a2212,arquivo,sensibilidade,Healthy]
\newcommand{\VDiagRecDecBldcFilePropeller}{0.929}% [diagnostico] embarcados/results/csv/diagnostico_20261023_222831.csv :: valor[p11,a2212,arquivo,sensibilidade,Propeller]
\newcommand{\VDiagRecDecBldcFileBearing}{0.943}% [diagnostico] embarcados/results/csv/diagnostico_20261023_222831.csv :: valor[p11,a2212,arquivo,sensibilidade,Bearing]
\newcommand{\VDiagRecDecBldcMotorHealthy}{0.441}% [diagnostico] embarcados/results/csv/diagnostico_20261023_222831.csv :: valor[p11,a2212,motor,sensibilidade,Healthy]
\newcommand{\VDiagRecDecBldcMotorPropeller}{0.929}% [diagnostico] embarcados/results/csv/diagnostico_20261023_222831.csv :: valor[p11,a2212,motor,sensibilidade,Propeller]
\newcommand{\VDiagRecDecBldcMotorBearing}{0.000}% [diagnostico] embarcados/results/csv/diagnostico_20261023_222831.csv :: valor[p11,a2212,motor,sensibilidade,Bearing]
\newcommand{\VDiagInflRmsNormBldc}{0.002}% [diagnostico] embarcados/results/csv/diagnostico_20261023_222831.csv :: Bal[aleatorio]-Bal[arquivo] (a2212)
\newcommand{\VDiagDropRmsNormBldc}{0.019}% [diagnostico] embarcados/results/csv/diagnostico_20261023_222831.csv :: Bal[arquivo]-Bal[motor] (a2212)
\newcommand{\VDiagInflRmsRawBldc}{0.016}% [diagnostico] embarcados/results/csv/diagnostico_20261023_222831.csv :: Bal[aleatorio]-Bal[arquivo] (a2212)
\newcommand{\VDiagDropRmsRawBldc}{0.106}% [diagnostico] embarcados/results/csv/diagnostico_20261023_222831.csv :: Bal[arquivo]-Bal[motor] (a2212)
\newcommand{\VDiagInflClsBldc}{0.017}% [diagnostico] embarcados/results/csv/diagnostico_20261023_222831.csv :: Bal[aleatorio]-Bal[arquivo] (a2212)
\newcommand{\VDiagDropClsBldc}{0.368}% [diagnostico] embarcados/results/csv/diagnostico_20261023_222831.csv :: Bal[arquivo]-Bal[motor] (a2212)
\newcommand{\VDiagInflDecBldc}{0.033}% [diagnostico] embarcados/results/csv/diagnostico_20261023_222831.csv :: Bal[aleatorio]-Bal[arquivo] (a2212)
\newcommand{\VDiagDropDecBldc}{0.360}% [diagnostico] embarcados/results/csv/diagnostico_20261023_222831.csv :: Bal[arquivo]-Bal[motor] (a2212)
\newcommand{\VParNFeat}{13}% [codigo] embarcados/results/csv/p11_features_20261023_222831.csv :: nº de colunas f0..f12
\newcommand{\VRotIdmtGoodMed}{25.76}% [dados] embarcados/results/csv/p11_features_20261023_222831.csv :: mediana de f0 (idmt, good)
\newcommand{\VRotIdmtGoodIqr}{0.18}% [dados] embarcados/results/csv/p11_features_20261023_222831.csv :: IQR de f0 (idmt, good)
\newcommand{\VRotIdmtBrokenMed}{27.99}% [dados] embarcados/results/csv/p11_features_20261023_222831.csv :: mediana de f0 (idmt, broken)
\newcommand{\VRotIdmtBrokenIqr}{0.03}% [dados] embarcados/results/csv/p11_features_20261023_222831.csv :: IQR de f0 (idmt, broken)
\newcommand{\VRotIdmtHeavyMed}{24.65}% [dados] embarcados/results/csv/p11_features_20261023_222831.csv :: mediana de f0 (idmt, heavyload)
\newcommand{\VRotIdmtHeavyIqr}{0.20}% [dados] embarcados/results/csv/p11_features_20261023_222831.csv :: IQR de f0 (idmt, heavyload)
\newcommand{\VRotOctThr}{1.5}% [dados] embarcados/results/csv/p11_features_20261023_222831.csv :: limiar: f0 > LIM_OIT·mediana da classe
\newcommand{\VRotCeilPctThr}{1}% [dados] embarcados/results/csv/p11_features_20261023_222831.csv :: limiar: f0 >= LIM_TETO·fmax
\newcommand{\VRotIdmtGoodOctPct}{8.6}% [dados] embarcados/results/csv/p11_features_20261023_222831.csv :: % janelas good com f0 > LIM_OIT·mediana da classe (salto de oitava)
\newcommand{\VRotBldcCeilPct}{3.7}% [dados] embarcados/results/csv/p11_features_20261023_222831.csv :: % janelas A2212 com f0 >= 0,99·fmax
\newcommand{\VRotBldcSpearman}{0.02}% [dados] embarcados/results/csv/p11_features_20261023_222831.csv + embarcados/results/predictions/diagnostico_janelas_20261023_222831.csv :: Spearman(velocidade nominal do arquivo, mediana de f0 do arquivo), 43 arquivos
\newcommand{\VRotBldcSpearmanMOne}{0.30}% [dados] embarcados/results/csv/p11_features_20261023_222831.csv + embarcados/results/predictions/diagnostico_janelas_20261023_222831.csv :: Spearman por motor M1 (14 arquivos)
\newcommand{\VRotBldcFilesMOne}{14}% [dados] embarcados/results/csv/p11_features_20261023_222831.csv + embarcados/results/predictions/diagnostico_janelas_20261023_222831.csv :: nº de arquivos do M1
\newcommand{\VRotBldcSpearmanMTwo}{0.30}% [dados] embarcados/results/csv/p11_features_20261023_222831.csv + embarcados/results/predictions/diagnostico_janelas_20261023_222831.csv :: Spearman por motor M2 (16 arquivos)
\newcommand{\VRotBldcFilesMTwo}{16}% [dados] embarcados/results/csv/p11_features_20261023_222831.csv + embarcados/results/predictions/diagnostico_janelas_20261023_222831.csv :: nº de arquivos do M2
\newcommand{\VRotBldcSpearmanMThree}{0.54}% [dados] embarcados/results/csv/p11_features_20261023_222831.csv + embarcados/results/predictions/diagnostico_janelas_20261023_222831.csv :: Spearman por motor M3 (6 arquivos)
\newcommand{\VRotBldcFilesMThree}{6}% [dados] embarcados/results/csv/p11_features_20261023_222831.csv + embarcados/results/predictions/diagnostico_janelas_20261023_222831.csv :: nº de arquivos do M3
\newcommand{\VRotBldcSpearmanMFour}{1.00}% [dados] embarcados/results/csv/p11_features_20261023_222831.csv + embarcados/results/predictions/diagnostico_janelas_20261023_222831.csv :: Spearman por motor M4 (7 arquivos)
\newcommand{\VRotBldcFilesMFour}{7}% [dados] embarcados/results/csv/p11_features_20261023_222831.csv + embarcados/results/predictions/diagnostico_janelas_20261023_222831.csv :: nº de arquivos do M4
\newcommand{\VEqStagesTotal}{26}% [equivalencia] embarcados/results/csv/equivalencia_resumo_por_build_20261024-091729.csv :: nº de etapas por build
\newcommand{\VEqStagesCase}{20}% [equivalencia] embarcados/results/csv/equivalencia_resumo_por_build_20261024-091729.csv :: etapas de janela
\newcommand{\VEqStagesSeq}{6}% [equivalencia] embarcados/results/csv/equivalencia_resumo_por_build_20261024-091729.csv :: etapas de sequência
\newcommand{\VEqCases}{164}% [equivalencia] embarcados/results/csv/equivalencia_resumo_por_build_20261024-091729.csv :: n[y_prep]
\newcommand{\VEqSeqCombos}{240}% [equivalencia] embarcados/results/csv/equivalencia_resumo_por_build_20261024-091729.csv :: n[absorb]
\newcommand{\VEqDoubleStagesPass}{26}% [equivalencia] embarcados/results/csv/equivalencia_resumo_por_build_20261024-091729.csv :: nº de etapas com veredito PASS (double)
\newcommand{\VEqDoubleMaxRatio}{0.55}% [equivalencia] embarcados/results/csv/equivalencia_resumo_por_build_20261024-091729.csv :: máx razao_tol_max nas etapas sem empate/FAIL (double)
\newcommand{\VEqDoubleDecisionPass}{164}% [equivalencia] embarcados/results/csv/equivalencia_resumo_por_build_20261024-091729.csv :: pass+pass_empate[decisao] (double)
\newcommand{\VEqFloatStagesPass}{25}% [equivalencia] embarcados/results/csv/equivalencia_resumo_por_build_20261024-091729.csv :: nº de etapas com veredito PASS (float)
\newcommand{\VEqFloatMaxRatio}{0.34}% [equivalencia] embarcados/results/csv/equivalencia_resumo_por_build_20261024-091729.csv :: máx razao_tol_max nas etapas sem empate/FAIL (float)
\newcommand{\VEqFloatDecisionPass}{164}% [equivalencia] embarcados/results/csv/equivalencia_resumo_por_build_20261024-091729.csv :: pass+pass_empate[decisao] (float)
\newcommand{\VEqSosdStagesPass}{26}% [equivalencia] embarcados/results/csv/equivalencia_resumo_por_build_20261024-091729.csv :: nº de etapas com veredito PASS (float-sosd)
\newcommand{\VEqSosdMaxRatio}{0.23}% [equivalencia] embarcados/results/csv/equivalencia_resumo_por_build_20261024-091729.csv :: máx razao_tol_max nas etapas sem empate/FAIL (float-sosd)
\newcommand{\VEqSosdDecisionPass}{164}% [equivalencia] embarcados/results/csv/equivalencia_resumo_por_build_20261024-091729.csv :: pass+pass_empate[decisao] (float-sosd)
\newcommand{\VEqFloatZFail}{2}% [equivalencia] embarcados/results/csv/equivalencia_resumo_por_build_20261024-091729.csv :: fail[float,z]
\newcommand{\VEqFloatZPass}{162}% [equivalencia] embarcados/results/csv/equivalencia_resumo_por_build_20261024-091729.csv :: pass[float,z]
\newcommand{\VEqFloatFailStages}{z}% [equivalencia] embarcados/results/csv/equivalencia_resumo_por_build_20261024-091729.csv :: etapas sem PASS (float)
\newcommand{\VEqFloatTieStages}{2}% [equivalencia] embarcados/results/csv/equivalencia_resumo_por_build_20261024-091729.csv :: nº de etapas com PASS_EMPATE (float)
\newcommand{\VEqFloatTieCases}{2}% [equivalencia] embarcados/results/csv/equivalencia_float_20261024-091122.csv :: ids com PASS_EMPATE em env_mod
\newcommand{\VEqFloatTieCaseIds}{borda\_dc\_ruido3s\_ma2212, borda\_dc\_ruido3s\_midmt}% [equivalencia] embarcados/results/csv/equivalencia_float_20261024-091122.csv :: id
\newcommand{\VEqFloatTieGap}{\ensuremath{9.4\times10^{-16}}}% [equivalencia] embarcados/results/csv/equivalencia_float_20261024-091122.csv :: detalhe[env_mod]: P_ref[ref]-P_ref[C]
\newcommand{\VEqFloatTieTol}{\ensuremath{6.3\times10^{-11}}}% [equivalencia] embarcados/results/csv/equivalencia_float_20261024-091122.csv :: detalhe[env_mod]: 2·tol
\newcommand{\VEqFloatFmodErr}{37.1}% [equivalencia] embarcados/results/csv/equivalencia_float_20261024-091122.csv :: err_max[features[9]:f_mod] (empate)
\newcommand{\VEqFloatZErrMax}{1.61}% [equivalencia] embarcados/results/csv/equivalencia_float_20261024-091122.csv :: err_max[z] (FAIL)
\newcommand{\VEqFloatZFailIsEdge}{1}% [equivalencia] embarcados/results/csv/equivalencia_float_20261024-091122.csv :: todos os FAIL de z são casos de borda (1=sim)
\newcommand{\VEqEnvDecFloatMed}{\ensuremath{1.2\times10^{-5}}}% [equivalencia] embarcados/results/csv/equivalencia_float_20261024-091122.csv :: mediana err_rel_max[env_dec]
\newcommand{\VEqEnvDecSosdMed}{\ensuremath{1.2\times10^{-7}}}% [equivalencia] embarcados/results/csv/equivalencia_float-sosd_20261024-091222.csv :: mediana err_rel_max[env_dec]
\newcommand{\VEqEnvDecDoubleMed}{\ensuremath{1.1\times10^{-7}}}% [equivalencia] embarcados/results/csv/equivalencia_double_20261024-091106.csv :: mediana err_rel_max[env_dec]
\newcommand{\VEqTolSha}{d74dae800ceb}% [equivalencia] embarcados/configs/embedded/tolerancias.json :: sha256 (12 primeiros)
\newcommand{\VEqTolShaRec}{d74dae800ceb}% [equivalencia] embarcados/results/csv/equivalencia_resumo_por_build_20261024-091729.csv :: sha256_tolerancias gravado no _meta da comparação
\newcommand{\VEqTolCommit}{447a7c2}% [equivalencia] embarcados/configs/embedded/tolerancias.json :: git log -1 do tolerancias.json
\newcommand{\VEqTolCommitTime}{2026-10-24 09:10}% [equivalencia] embarcados/configs/embedded/tolerancias.json :: data do commit do tolerancias.json
\newcommand{\VEqRunTime}{2026-10-24 09:11}% [equivalencia] embarcados/results/csv/equivalencia_resumo_por_build_20261024-091729.csv :: data da comparação formal (float)
\newcommand{\VEqGoldClean}{60}% [equivalencia] embarcados/gold/manifest.json :: casos tipo limpo
\newcommand{\VEqGoldNoisy}{90}% [equivalencia] embarcados/gold/manifest.json :: casos tipo ruidoso
\newcommand{\VEqGoldEdge}{14}% [equivalencia] embarcados/gold/manifest.json :: casos tipo borda
\newcommand{\VEqGoldSeqs}{20}% [equivalencia] embarcados/gold/manifest.json :: n_seqs
\newcommand{\VEqGoldPerts}{5}% [equivalencia] embarcados/gold/manifest.json :: nº de perturbações
\newcommand{\VEqMinDurValues}{2}% [equivalencia] embarcados/gold/seqs/idmt_seq00/params.json :: len(min_durs_s)
\newcommand{\VEqEmulDoubleRows}{6{,}852}% [equivalencia] embarcados/results/csv/equivalencia_emul_double_20261024-090737_meta.json :: n_linhas
\newcommand{\VEqEmulDoublePass}{6{,}852}% [equivalencia] embarcados/results/csv/equivalencia_emul_double_20261024-090737_meta.json :: contagem_status.PASS
\newcommand{\VEqEmulFloatRows}{6{,}852}% [equivalencia] embarcados/results/csv/equivalencia_emul_float_20261024-090743_meta.json :: n_linhas
\newcommand{\VEqEmulFloatPass}{6{,}852}% [equivalencia] embarcados/results/csv/equivalencia_emul_float_20261024-090743_meta.json :: contagem_status.PASS
\newcommand{\VPcCIdmtMed}{13.29}% [pc] embarcados/results/metrics/bench_pc_20261024-091911.json :: resumo.c_float_pc/idmt.total_ms.mediana
\newcommand{\VPcPyIdmtMed}{95.3}% [pc] embarcados/results/metrics/bench_pc_20261024-091911.json :: resumo.python_ref/idmt.total_ms.mediana
\newcommand{\VPcCIdmtPninetyfive}{15.00}% [pc] embarcados/results/metrics/bench_pc_20261024-091911.json :: resumo.c_float_pc/idmt.total_ms.p95
\newcommand{\VPcPyIdmtPninetyfive}{100.2}% [pc] embarcados/results/metrics/bench_pc_20261024-091911.json :: resumo.python_ref/idmt.total_ms.p95
\newcommand{\VPcCIdmtMax}{15.80}% [pc] embarcados/results/metrics/bench_pc_20261024-091911.json :: resumo.c_float_pc/idmt.total_ms.max
\newcommand{\VPcPyIdmtMax}{110.2}% [pc] embarcados/results/metrics/bench_pc_20261024-091911.json :: resumo.python_ref/idmt.total_ms.max
\newcommand{\VPcCIdmtRtf}{0.0044}% [pc] embarcados/results/metrics/bench_pc_20261024-091911.json :: fator_tempo_real_mediana
\newcommand{\VPcPyIdmtRtf}{0.032}% [pc] embarcados/results/metrics/bench_pc_20261024-091911.json :: fator_tempo_real_mediana
\newcommand{\VPcSpeedupIdmt}{7.2}% [pc] embarcados/results/metrics/bench_pc_20261024-091911.json :: mediana Python / mediana C
\newcommand{\VPcCIdmtTrackMed}{5.45}% [pc] embarcados/results/metrics/bench_pc_20261024-091911.json :: rastrear_quadros_ms.mediana
\newcommand{\VPcCIdmtTrackPct}{41}% [pc] embarcados/results/metrics/bench_pc_20261024-091911.json :: 100·rastrear_quadros_ms.mediana/total_ms.mediana
\newcommand{\VPcCIdmtVitMed}{5.35}% [pc] embarcados/results/metrics/bench_pc_20261024-091911.json :: viterbi_ms.mediana
\newcommand{\VPcCIdmtVitPct}{40}% [pc] embarcados/results/metrics/bench_pc_20261024-091911.json :: 100·viterbi_ms.mediana/total_ms.mediana
\newcommand{\VPcCIdmtEnvMed}{2.27}% [pc] embarcados/results/metrics/bench_pc_20261024-091911.json :: envelope_mod_ms.mediana
\newcommand{\VPcCIdmtEnvPct}{17}% [pc] embarcados/results/metrics/bench_pc_20261024-091911.json :: 100·envelope_mod_ms.mediana/total_ms.mediana
\newcommand{\VPcCIdmtFeatMed}{0.013}% [pc] embarcados/results/metrics/bench_pc_20261024-091911.json :: features_ms.mediana
\newcommand{\VPcCIdmtFeatPct}{0}% [pc] embarcados/results/metrics/bench_pc_20261024-091911.json :: 100·features_ms.mediana/total_ms.mediana
\newcommand{\VPcCIdmtCosMed}{0.002}% [pc] embarcados/results/metrics/bench_pc_20261024-091911.json :: cosseno_ms.mediana
\newcommand{\VPcCIdmtCosPct}{0}% [pc] embarcados/results/metrics/bench_pc_20261024-091911.json :: 100·cosseno_ms.mediana/total_ms.mediana
\newcommand{\VPcNIdmt}{200}% [pc] embarcados/results/metrics/bench_pc_20261024-091911.json :: total_ms.n
\newcommand{\VPcCBldcMed}{3.18}% [pc] embarcados/results/metrics/bench_pc_20261024-091911.json :: resumo.c_float_pc/a2212.total_ms.mediana
\newcommand{\VPcPyBldcMed}{22.1}% [pc] embarcados/results/metrics/bench_pc_20261024-091911.json :: resumo.python_ref/a2212.total_ms.mediana
\newcommand{\VPcCBldcPninetyfive}{3.84}% [pc] embarcados/results/metrics/bench_pc_20261024-091911.json :: resumo.c_float_pc/a2212.total_ms.p95
\newcommand{\VPcPyBldcPninetyfive}{23.1}% [pc] embarcados/results/metrics/bench_pc_20261024-091911.json :: resumo.python_ref/a2212.total_ms.p95
\newcommand{\VPcCBldcMax}{4.45}% [pc] embarcados/results/metrics/bench_pc_20261024-091911.json :: resumo.c_float_pc/a2212.total_ms.max
\newcommand{\VPcPyBldcMax}{25.3}% [pc] embarcados/results/metrics/bench_pc_20261024-091911.json :: resumo.python_ref/a2212.total_ms.max
\newcommand{\VPcCBldcRtf}{0.0032}% [pc] embarcados/results/metrics/bench_pc_20261024-091911.json :: fator_tempo_real_mediana
\newcommand{\VPcPyBldcRtf}{0.022}% [pc] embarcados/results/metrics/bench_pc_20261024-091911.json :: fator_tempo_real_mediana
\newcommand{\VPcSpeedupBldc}{6.9}% [pc] embarcados/results/metrics/bench_pc_20261024-091911.json :: mediana Python / mediana C
\newcommand{\VPcCBldcTrackMed}{1.24}% [pc] embarcados/results/metrics/bench_pc_20261024-091911.json :: rastrear_quadros_ms.mediana
\newcommand{\VPcCBldcTrackPct}{39}% [pc] embarcados/results/metrics/bench_pc_20261024-091911.json :: 100·rastrear_quadros_ms.mediana/total_ms.mediana
\newcommand{\VPcCBldcVitMed}{1.14}% [pc] embarcados/results/metrics/bench_pc_20261024-091911.json :: viterbi_ms.mediana
\newcommand{\VPcCBldcVitPct}{36}% [pc] embarcados/results/metrics/bench_pc_20261024-091911.json :: 100·viterbi_ms.mediana/total_ms.mediana
\newcommand{\VPcCBldcEnvMed}{0.67}% [pc] embarcados/results/metrics/bench_pc_20261024-091911.json :: envelope_mod_ms.mediana
\newcommand{\VPcCBldcEnvPct}{21}% [pc] embarcados/results/metrics/bench_pc_20261024-091911.json :: 100·envelope_mod_ms.mediana/total_ms.mediana
\newcommand{\VPcCBldcFeatMed}{0.013}% [pc] embarcados/results/metrics/bench_pc_20261024-091911.json :: features_ms.mediana
\newcommand{\VPcCBldcFeatPct}{0}% [pc] embarcados/results/metrics/bench_pc_20261024-091911.json :: 100·features_ms.mediana/total_ms.mediana
\newcommand{\VPcCBldcCosMed}{0.002}% [pc] embarcados/results/metrics/bench_pc_20261024-091911.json :: cosseno_ms.mediana
\newcommand{\VPcCBldcCosPct}{0}% [pc] embarcados/results/metrics/bench_pc_20261024-091911.json :: 100·cosseno_ms.mediana/total_ms.mediana
\newcommand{\VPcNBldc}{200}% [pc] embarcados/results/metrics/bench_pc_20261024-091911.json :: total_ms.n
\newcommand{\VPcWarmup}{5}% [pc] embarcados/results/metrics/bench_pc_20261024-091911.json :: meta.argv.aquecimento
\newcommand{\VPcDistinctIdmt}{75}% [pc] embarcados/results/csv/bench_pc_20261024-091911.csv :: n_janelas_distintas
\newcommand{\VPcDistinctBldc}{75}% [pc] embarcados/results/csv/bench_pc_20261024-091911.csv :: n_janelas_distintas
\newcommand{\VPcCompiler}{MSVC 19.44 x64 \slash{}O2 \slash{}fp:precise}% [pc] embarcados/results/metrics/bench_pc_20261024-091911.json :: meta.compilador
\newcommand{\VPcNcpu}{16}% [pc] embarcados/results/metrics/bench_pc_20261024-091911.json :: meta.maquina.n_cpu
\newcommand{\VPcOs}{Windows 11}% [pc] embarcados/results/metrics/bench_pc_20261024-091911.json :: meta.maquina.plataforma
\newcommand{\VPcSeqWindows}{24}% [pc] embarcados/results/csv/seq_pc_20261024-092103.csv :: nº de linhas (janelas)
\newcommand{\VPcSeqSwitches}{2}% [pc] embarcados/results/csv/seq_pc_20261024-092103_meta.json :: len(deteccao_online_c)
\newcommand{\VPcSeqDelayC}{0}% [pc] embarcados/results/csv/seq_pc_20261024-092103_meta.json :: max atraso_janelas (C)
\newcommand{\VPcSeqDelayPy}{0}% [pc] embarcados/results/csv/seq_pc_20261024-092103_meta.json :: max atraso_janelas (Python)
\newcommand{\VPcSeqSameDecisions}{1}% [pc] embarcados/results/csv/seq_pc_20261024-092103_meta.json :: decisoes_iguais_python_c (1=sim)
\newcommand{\VPcSeqMaxProbDiff}{\ensuremath{3.2\times10^{-6}}}% [pc] embarcados/results/csv/seq_pc_20261024-092103_meta.json :: max_abs_dif_probs
\newcommand{\VPcSeqCp}{2}% [pc] embarcados/results/csv/seq_pc_20261024-092103_meta.json :: cp
\newcommand{\VPcSeqPerClass}{8}% [pc] embarcados/results/csv/seq_pc_20261024-092103.csv :: janelas por classe
\newcommand{\VEstWsFloatOneMb}{0.96}% [estimativa] embarcados/paper/fontes/ws_float_16000.json :: total_com_tabela/1e6
\newcommand{\VEstWsFloatOneNoTabMb}{0.61}% [estimativa] embarcados/paper/fontes/ws_float_16000.json :: total_sem_tabela/1e6
\newcommand{\VEstWsFloatOnePhaseBPct}{44}% [estimativa] embarcados/paper/fontes/ws_float_16000.json :: faseB/(persistente+max(faseA,faseB))
\newcommand{\VEstWsFloatThreeMb}{1.86}% [estimativa] embarcados/paper/fontes/ws_float_48000.json :: total_com_tabela/1e6
\newcommand{\VEstWsFloatThreeNoTabMb}{1.52}% [estimativa] embarcados/paper/fontes/ws_float_48000.json :: total_sem_tabela/1e6
\newcommand{\VEstWsFloatThreePhaseBPct}{64}% [estimativa] embarcados/paper/fontes/ws_float_48000.json :: faseB/(persistente+max(faseA,faseB))
\newcommand{\VEstWsDoubleOneMb}{1.79}% [estimativa] embarcados/paper/fontes/ws_double_16000.json :: total_com_tabela/1e6
\newcommand{\VEstWsDoubleOneNoTabMb}{1.21}% [estimativa] embarcados/paper/fontes/ws_double_16000.json :: total_sem_tabela/1e6
\newcommand{\VEstWsDoubleOnePhaseBPct}{47}% [estimativa] embarcados/paper/fontes/ws_double_16000.json :: faseB/(persistente+max(faseA,faseB))
\newcommand{\VEstWsDoubleThreeMb}{3.60}% [estimativa] embarcados/paper/fontes/ws_double_48000.json :: total_com_tabela/1e6
\newcommand{\VEstWsDoubleThreeNoTabMb}{3.02}% [estimativa] embarcados/paper/fontes/ws_double_48000.json :: total_sem_tabela/1e6
\newcommand{\VEstWsDoubleThreePhaseBPct}{67}% [estimativa] embarcados/paper/fontes/ws_double_48000.json :: faseB/(persistente+max(faseA,faseB))
\newcommand{\VEstHarmTableMb}{0.35}% [estimativa] embarcados/paper/fontes/ws_float_16000.json :: total_com_tabela - total_sem_tabela
\newcommand{\VBuildIdf}{v5.5.6}% [build] C:/Espressif/build_p11_esp32s3/project_description.json :: git_revision
\newcommand{\VBuildGcc}{14.2.0}% [build] C:/Espressif/build_p11_esp32s3/project_description.json :: c_compiler (versão do toolchain esp-)
\newcommand{\VBuildEspBinKb}{209}% [build] C:/Espressif/build_p11_esp32/size.txt :: binary size (hex)/1024
\newcommand{\VBuildEspBinPct}{20}% [build] C:/Espressif/build_p11_esp32/size.txt :: bin/partição app
\newcommand{\VBuildEspFlashCodeKb}{124}% [build] C:/Espressif/build_p11_esp32/size.txt :: Flash Code Used/1024
\newcommand{\VBuildEspFlashDataKb}{19}% [build] C:/Espressif/build_p11_esp32/size.txt :: Flash Data Used/1024
\newcommand{\VBuildEspDramKb}{15}% [build] C:/Espressif/build_p11_esp32/size.txt :: DRAM Used/1024
\newcommand{\VBuildEspDramTotalKb}{176}% [build] C:/Espressif/build_p11_esp32/size.txt :: DRAM Total/1024
\newcommand{\VBuildEspIramKb}{58}% [build] C:/Espressif/build_p11_esp32/size.txt :: IRAM Used/1024
\newcommand{\VBuildEspDramName}{DRAM}% [build] C:/Espressif/build_p11_esp32/size.txt :: nome da região de dados interna
\newcommand{\VBuildEspWsRatio}{10.3}% [estimativa] C:/Espressif/build_p11_esp32/size.txt :: estimativa ws float 3 s / DRAM total
\newcommand{\VBuildEspWsRatioOne}{5.3}% [estimativa] C:/Espressif/build_p11_esp32/size.txt :: estimativa ws float 1 s / DRAM total
\newcommand{\VBuildEspsBinKb}{233}% [build] C:/Espressif/build_p11_esp32s3/size.txt :: binary size (hex)/1024
\newcommand{\VBuildEspsBinPct}{23}% [build] C:/Espressif/build_p11_esp32s3/size.txt :: bin/partição app
\newcommand{\VBuildEspsFlashCodeKb}{139}% [build] C:/Espressif/build_p11_esp32s3/size.txt :: Flash Code Used/1024
\newcommand{\VBuildEspsFlashDataKb}{30}% [build] C:/Espressif/build_p11_esp32s3/size.txt :: Flash Data Used/1024
\newcommand{\VBuildEspsDramKb}{54}% [build] C:/Espressif/build_p11_esp32s3/size.txt :: DIRAM Used/1024
\newcommand{\VBuildEspsDramTotalKb}{334}% [build] C:/Espressif/build_p11_esp32s3/size.txt :: DIRAM Total/1024
\newcommand{\VBuildEspsIramKb}{16}% [build] C:/Espressif/build_p11_esp32s3/size.txt :: IRAM Used/1024
\newcommand{\VBuildEspsDramName}{DIRAM}% [build] C:/Espressif/build_p11_esp32s3/size.txt :: nome da região de dados interna
\newcommand{\VBuildEspsWsRatio}{5.5}% [estimativa] C:/Espressif/build_p11_esp32s3/size.txt :: estimativa ws float 3 s / DIRAM total
\newcommand{\VBuildEspsWsRatioOne}{2.8}% [estimativa] C:/Espressif/build_p11_esp32s3/size.txt :: estimativa ws float 1 s / DIRAM total
\newcommand{\VChipModel}{\textbf{[pending]}}% [esp32] embarcados/results/manifests/bancada.json :: chip_modelo
\newcommand{\VChipCpuMhz}{\textbf{[pending]}}% [esp32] embarcados/results/manifests/bancada.json :: cpu_mhz
\newcommand{\VChipPsramMb}{\textbf{[pending]}}% [esp32] embarcados/results/manifests/bancada.json :: psram_mb
\newcommand{\VChipLatIdmtMed}{\textbf{[pending]}}% [esp32] embarcados/results/manifests/bancada.json :: resumo.idmt.total_ms.mediana
\newcommand{\VChipLatIdmtPninetyfive}{\textbf{[pending]}}% [esp32] embarcados/results/manifests/bancada.json :: resumo.idmt.total_ms.p95
\newcommand{\VChipLatIdmtMax}{\textbf{[pending]}}% [esp32] embarcados/results/manifests/bancada.json :: resumo.idmt.total_ms.max
\newcommand{\VChipLatIdmtN}{\textbf{[pending]}}% [esp32] embarcados/results/manifests/bancada.json :: resumo.idmt.total_ms.n
\newcommand{\VChipRtfIdmt}{\textbf{[pending]}}% [esp32] embarcados/results/manifests/bancada.json :: fator_tempo_real_mediana
\newcommand{\VChipRtfMaxIdmt}{\textbf{[pending]}}% [esp32] embarcados/results/manifests/bancada.json :: fator_tempo_real_max
\newcommand{\VChipHeapIntIdmtKb}{\textbf{[pending]}}% [esp32] embarcados/results/manifests/bancada.json :: heap_min_livre_interna/1024
\newcommand{\VChipStackIdmt}{\textbf{[pending]}}% [esp32] embarcados/results/manifests/bancada.json :: pilha_folga_bytes_min
\newcommand{\VChipWsPsramIdmt}{\textbf{[pending]}}% [esp32] embarcados/results/manifests/bancada.json :: ws_na_psram
\newcommand{\VChipWsIdmtMb}{\textbf{[pending]}}% [esp32] embarcados/results/manifests/bancada.json :: ws_bytes/1e6
\newcommand{\VChipLatBldcMed}{\textbf{[pending]}}% [esp32] embarcados/results/manifests/bancada.json :: resumo.a2212.total_ms.mediana
\newcommand{\VChipLatBldcPninetyfive}{\textbf{[pending]}}% [esp32] embarcados/results/manifests/bancada.json :: resumo.a2212.total_ms.p95
\newcommand{\VChipLatBldcMax}{\textbf{[pending]}}% [esp32] embarcados/results/manifests/bancada.json :: resumo.a2212.total_ms.max
\newcommand{\VChipLatBldcN}{\textbf{[pending]}}% [esp32] embarcados/results/manifests/bancada.json :: resumo.a2212.total_ms.n
\newcommand{\VChipRtfBldc}{\textbf{[pending]}}% [esp32] embarcados/results/manifests/bancada.json :: fator_tempo_real_mediana
\newcommand{\VChipRtfMaxBldc}{\textbf{[pending]}}% [esp32] embarcados/results/manifests/bancada.json :: fator_tempo_real_max
\newcommand{\VChipHeapIntBldcKb}{\textbf{[pending]}}% [esp32] embarcados/results/manifests/bancada.json :: heap_min_livre_interna/1024
\newcommand{\VChipStackBldc}{\textbf{[pending]}}% [esp32] embarcados/results/manifests/bancada.json :: pilha_folga_bytes_min
\newcommand{\VChipWsPsramBldc}{\textbf{[pending]}}% [esp32] embarcados/results/manifests/bancada.json :: ws_na_psram
\newcommand{\VChipWsBldcMb}{\textbf{[pending]}}% [esp32] embarcados/results/manifests/bancada.json :: ws_bytes/1e6
\newcommand{\VChipEqStagesPass}{\textbf{[pending]}}% [esp32] embarcados/results/manifests/bancada.json :: equivalência no chip não rodou
\newcommand{\VChipEqStages}{\textbf{[pending]}}% [esp32] embarcados/results/manifests/bancada.json :: equivalência no chip não rodou
\newcommand{\VChipEqDecisionPass}{\textbf{[pending]}}% [esp32] embarcados/results/manifests/bancada.json :: equivalência no chip não rodou
\newcommand{\VChipSeqDelay}{\textbf{[pending]}}% [esp32] embarcados/results/manifests/bancada.json :: fluxo contínuo no chip não rodou

% Gerado por embarcados/src/python/gerar_artigo_numeros_aud.py. NÃO EDITAR À MÃO.
% commit 9fb82f016b4482c1a5b58f1b9bd08db149799b78 (árvore suja em embarcados/: True)
% Macros da auditoria independente (M2 holdout, M3 motor igual por igual, M4 rotação).
% Uma macro por número. Comentário: [categoria] fonte :: coluna.

\newcommand{\audHoldoutSeed}{7}% [equivalencia] embarcados/results/csv/aud_holdout_20261025-140558_meta.json :: seed
\newcommand{\audHoldoutN}{24}% [equivalencia] embarcados/results/csv/aud_holdout_20261025-140558_resumo.csv :: n_janelas[float,_todas]
\newcommand{\audHoldoutNIdmt}{12}% [equivalencia] embarcados/results/csv/aud_holdout_20261025-140558_resumo.csv :: n_janelas[float,idmt]
\newcommand{\audHoldoutNBldc}{12}% [equivalencia] embarcados/results/csv/aud_holdout_20261025-140558_resumo.csv :: n_janelas[float,a2212]
\newcommand{\audHoldoutFloatPathSame}{24}% [equivalencia] embarcados/results/csv/aud_holdout_20261025-140558_resumo.csv :: caminho_identico[float,_todas]
\newcommand{\audHoldoutFloatDecSame}{24}% [equivalencia] embarcados/results/csv/aud_holdout_20261025-140558_resumo.csv :: decisao_igual[float,_todas]
\newcommand{\audHoldoutFloatFeatOut}{0}% [equivalencia] embarcados/results/csv/aud_holdout_20261025-140558_resumo.csv :: n_features_fora_tol[float,_todas]
\newcommand{\audHoldoutFloatStageFail}{0}% [equivalencia] embarcados/results/csv/aud_holdout_20261025-140558_resumo.csv :: n_etapas_fail[float,_todas]
\newcommand{\audHoldoutFloatRatioMax}{0.20}% [equivalencia] embarcados/results/csv/aud_holdout_20261025-140558_resumo.csv :: razao_tol_max_features[float,_todas]
\newcommand{\audHoldoutFloatDpMax}{\ensuremath{8.7\times10^{-6}}}% [equivalencia] embarcados/results/csv/aud_holdout_20261025-140558_resumo.csv :: max_abs_dp[float,_todas]
\newcommand{\audHoldoutDoublePathSame}{24}% [equivalencia] embarcados/results/csv/aud_holdout_20261025-140558_resumo.csv :: caminho_identico[double,_todas]
\newcommand{\audHoldoutDoubleDecSame}{24}% [equivalencia] embarcados/results/csv/aud_holdout_20261025-140558_resumo.csv :: decisao_igual[double,_todas]
\newcommand{\audHoldoutDoubleFeatOut}{0}% [equivalencia] embarcados/results/csv/aud_holdout_20261025-140558_resumo.csv :: n_features_fora_tol[double,_todas]
\newcommand{\audHoldoutDoubleStageFail}{0}% [equivalencia] embarcados/results/csv/aud_holdout_20261025-140558_resumo.csv :: n_etapas_fail[double,_todas]
\newcommand{\audHoldoutDoubleRatioMax}{0.05}% [equivalencia] embarcados/results/csv/aud_holdout_20261025-140558_resumo.csv :: razao_tol_max_features[double,_todas]
\newcommand{\audHoldoutDoubleDpMax}{\ensuremath{1.9\times10^{-7}}}% [equivalencia] embarcados/results/csv/aud_holdout_20261025-140558_resumo.csv :: max_abs_dp[double,_todas]
\newcommand{\audRotBldcSpearman}{0.02}% [dados] embarcados/results/csv/aud_rotacao_motor_20261025-140355.csv :: spearman_rho[global]
\newcommand{\audRotBldcFiles}{43}% [dados] embarcados/results/csv/aud_rotacao_motor_20261025-140355.csv :: n_arquivos[global]
\newcommand{\audRotBldcWindows}{430}% [dados] embarcados/results/csv/aud_rotacao_motor_20261025-140355.csv :: n_janelas[global]
\newcommand{\audRotBldcCeilPct}{3.7}% [dados] embarcados/results/csv/aud_rotacao_motor_20261025-140355.csv :: 100·frac_janelas_teto[global]
\newcommand{\audRotBldcCeilThrPct}{99}% [dados] embarcados/results/csv/aud_rotacao_motor_20261025-140355.csv :: 100·lim_teto_rel
\newcommand{\audRotBldcSpearmanMOne}{0.30}% [dados] embarcados/results/csv/aud_rotacao_motor_20261025-140355.csv :: spearman_rho[M1]
\newcommand{\audRotBldcSpearmanPMOne}{0.29}% [dados] embarcados/results/csv/aud_rotacao_motor_20261025-140355.csv :: spearman_p[M1]
\newcommand{\audRotBldcFilesMOne}{14}% [dados] embarcados/results/csv/aud_rotacao_motor_20261025-140355.csv :: n_arquivos[M1]
\newcommand{\audRotBldcSpearmanMTwo}{0.30}% [dados] embarcados/results/csv/aud_rotacao_motor_20261025-140355.csv :: spearman_rho[M2]
\newcommand{\audRotBldcSpearmanPMTwo}{0.26}% [dados] embarcados/results/csv/aud_rotacao_motor_20261025-140355.csv :: spearman_p[M2]
\newcommand{\audRotBldcFilesMTwo}{16}% [dados] embarcados/results/csv/aud_rotacao_motor_20261025-140355.csv :: n_arquivos[M2]
\newcommand{\audRotBldcSpearmanMThree}{0.54}% [dados] embarcados/results/csv/aud_rotacao_motor_20261025-140355.csv :: spearman_rho[M3]
\newcommand{\audRotBldcSpearmanPMThree}{0.27}% [dados] embarcados/results/csv/aud_rotacao_motor_20261025-140355.csv :: spearman_p[M3]
\newcommand{\audRotBldcFilesMThree}{6}% [dados] embarcados/results/csv/aud_rotacao_motor_20261025-140355.csv :: n_arquivos[M3]
\newcommand{\audRotBldcSpearmanMFour}{1.00}% [dados] embarcados/results/csv/aud_rotacao_motor_20261025-140355.csv :: spearman_rho[M4]
\newcommand{\audRotBldcSpearmanPMFour}{0.00}% [dados] embarcados/results/csv/aud_rotacao_motor_20261025-140355.csv :: spearman_p[M4]
\newcommand{\audRotBldcFilesMFour}{7}% [dados] embarcados/results/csv/aud_rotacao_motor_20261025-140355.csv :: n_arquivos[M4]
\newcommand{\audRotBldcSpearmanMotorMin}{0.30}% [dados] embarcados/results/csv/aud_rotacao_motor_20261025-140355.csv :: min spearman_rho[M1..M4]
\newcommand{\audRotBldcSpearmanMotorMax}{1.00}% [dados] embarcados/results/csv/aud_rotacao_motor_20261025-140355.csv :: max spearman_rho[M1..M4]
\newcommand{\audRotBldcFilesMotorMin}{6}% [dados] embarcados/results/csv/aud_rotacao_motor_20261025-140355.csv :: min n_arquivos[M1..M4]
\newcommand{\audRotBldcFilesMotorMax}{16}% [dados] embarcados/results/csv/aud_rotacao_motor_20261025-140355.csv :: max n_arquivos[M1..M4]
\newcommand{\audMotorBalClsFileHP}{0.972}% [diagnostico] embarcados/results/csv/aud_motor_igual_20261025-140432.csv :: acuracia_balanceada[protocolo,p04_classicas,arquivo,HP_M1M3]
\newcommand{\audMotorBalClsFileHPLo}{0.924}% [diagnostico] embarcados/results/csv/aud_motor_igual_20261025-140432.csv :: ic95_lo[protocolo,p04_classicas,arquivo,HP_M1M3]
\newcommand{\audMotorBalClsFileHPHi}{0.998}% [diagnostico] embarcados/results/csv/aud_motor_igual_20261025-140432.csv :: ic95_hi[protocolo,p04_classicas,arquivo,HP_M1M3]
\newcommand{\audMotorBalClsMotorHP}{0.920}% [diagnostico] embarcados/results/csv/aud_motor_igual_20261025-140432.csv :: acuracia_balanceada[protocolo,p04_classicas,motor,HP_M1M3]
\newcommand{\audMotorBalClsMotorHPLo}{0.868}% [diagnostico] embarcados/results/csv/aud_motor_igual_20261025-140432.csv :: ic95_lo[protocolo,p04_classicas,motor,HP_M1M3]
\newcommand{\audMotorBalClsMotorHPHi}{0.962}% [diagnostico] embarcados/results/csv/aud_motor_igual_20261025-140432.csv :: ic95_hi[protocolo,p04_classicas,motor,HP_M1M3]
\newcommand{\audMotorDropClsHP}{0.052}% [diagnostico] embarcados/results/csv/aud_motor_igual_20261025-140432.csv :: acuracia_balanceada[queda,p04_classicas,HP_M1M3]
\newcommand{\audMotorDropClsHPLo}{0.023}% [diagnostico] embarcados/results/csv/aud_motor_igual_20261025-140432.csv :: ic95_lo[queda,p04_classicas,HP_M1M3]
\newcommand{\audMotorDropClsHPHi}{0.086}% [diagnostico] embarcados/results/csv/aud_motor_igual_20261025-140432.csv :: ic95_hi[queda,p04_classicas,HP_M1M3]
\newcommand{\audMotorDropClsFull}{0.368}% [diagnostico] embarcados/results/csv/aud_motor_igual_20261025-140432.csv :: queda_completa[queda,p04_classicas,HP_M1M3]
\newcommand{\audMotorBalDecFileHP}{0.753}% [diagnostico] embarcados/results/csv/aud_motor_igual_20261025-140432.csv :: acuracia_balanceada[protocolo,p11_decifrae_harmsum,arquivo,HP_M1M3]
\newcommand{\audMotorBalDecFileHPLo}{0.662}% [diagnostico] embarcados/results/csv/aud_motor_igual_20261025-140432.csv :: ic95_lo[protocolo,p11_decifrae_harmsum,arquivo,HP_M1M3]
\newcommand{\audMotorBalDecFileHPHi}{0.842}% [diagnostico] embarcados/results/csv/aud_motor_igual_20261025-140432.csv :: ic95_hi[protocolo,p11_decifrae_harmsum,arquivo,HP_M1M3]
\newcommand{\audMotorBalDecMotorHP}{0.685}% [diagnostico] embarcados/results/csv/aud_motor_igual_20261025-140432.csv :: acuracia_balanceada[protocolo,p11_decifrae_harmsum,motor,HP_M1M3]
\newcommand{\audMotorBalDecMotorHPLo}{0.593}% [diagnostico] embarcados/results/csv/aud_motor_igual_20261025-140432.csv :: ic95_lo[protocolo,p11_decifrae_harmsum,motor,HP_M1M3]
\newcommand{\audMotorBalDecMotorHPHi}{0.774}% [diagnostico] embarcados/results/csv/aud_motor_igual_20261025-140432.csv :: ic95_hi[protocolo,p11_decifrae_harmsum,motor,HP_M1M3]
\newcommand{\audMotorDropDecHP}{0.068}% [diagnostico] embarcados/results/csv/aud_motor_igual_20261025-140432.csv :: acuracia_balanceada[queda,p11_decifrae_harmsum,HP_M1M3]
\newcommand{\audMotorDropDecHPLo}{0.031}% [diagnostico] embarcados/results/csv/aud_motor_igual_20261025-140432.csv :: ic95_lo[queda,p11_decifrae_harmsum,HP_M1M3]
\newcommand{\audMotorDropDecHPHi}{0.116}% [diagnostico] embarcados/results/csv/aud_motor_igual_20261025-140432.csv :: ic95_hi[queda,p11_decifrae_harmsum,HP_M1M3]
\newcommand{\audMotorDropDecFull}{0.360}% [diagnostico] embarcados/results/csv/aud_motor_igual_20261025-140432.csv :: queda_completa[queda,p11_decifrae_harmsum,HP_M1M3]
\newcommand{\audMotorBalRmsNormFileHP}{0.843}% [diagnostico] embarcados/results/csv/aud_motor_igual_20261025-140432.csv :: acuracia_balanceada[protocolo,b00_rmspico_norm,arquivo,HP_M1M3]
\newcommand{\audMotorBalRmsNormFileHPLo}{0.760}% [diagnostico] embarcados/results/csv/aud_motor_igual_20261025-140432.csv :: ic95_lo[protocolo,b00_rmspico_norm,arquivo,HP_M1M3]
\newcommand{\audMotorBalRmsNormFileHPHi}{0.909}% [diagnostico] embarcados/results/csv/aud_motor_igual_20261025-140432.csv :: ic95_hi[protocolo,b00_rmspico_norm,arquivo,HP_M1M3]
\newcommand{\audMotorBalRmsNormMotorHP}{0.814}% [diagnostico] embarcados/results/csv/aud_motor_igual_20261025-140432.csv :: acuracia_balanceada[protocolo,b00_rmspico_norm,motor,HP_M1M3]
\newcommand{\audMotorBalRmsNormMotorHPLo}{0.720}% [diagnostico] embarcados/results/csv/aud_motor_igual_20261025-140432.csv :: ic95_lo[protocolo,b00_rmspico_norm,motor,HP_M1M3]
\newcommand{\audMotorBalRmsNormMotorHPHi}{0.887}% [diagnostico] embarcados/results/csv/aud_motor_igual_20261025-140432.csv :: ic95_hi[protocolo,b00_rmspico_norm,motor,HP_M1M3]
\newcommand{\audMotorDropRmsNormHP}{0.029}% [diagnostico] embarcados/results/csv/aud_motor_igual_20261025-140432.csv :: acuracia_balanceada[queda,b00_rmspico_norm,HP_M1M3]
\newcommand{\audMotorDropRmsNormHPLo}{-0.002}% [diagnostico] embarcados/results/csv/aud_motor_igual_20261025-140432.csv :: ic95_lo[queda,b00_rmspico_norm,HP_M1M3]
\newcommand{\audMotorDropRmsNormHPHi}{0.066}% [diagnostico] embarcados/results/csv/aud_motor_igual_20261025-140432.csv :: ic95_hi[queda,b00_rmspico_norm,HP_M1M3]
\newcommand{\audMotorDropRmsNormFull}{0.019}% [diagnostico] embarcados/results/csv/aud_motor_igual_20261025-140432.csv :: queda_completa[queda,b00_rmspico_norm,HP_M1M3]
\newcommand{\audMotorBalRmsRawFileHP}{0.913}% [diagnostico] embarcados/results/csv/aud_motor_igual_20261025-140432.csv :: acuracia_balanceada[protocolo,b00_rmspico_bruto,arquivo,HP_M1M3]
\newcommand{\audMotorBalRmsRawFileHPLo}{0.821}% [diagnostico] embarcados/results/csv/aud_motor_igual_20261025-140432.csv :: ic95_lo[protocolo,b00_rmspico_bruto,arquivo,HP_M1M3]
\newcommand{\audMotorBalRmsRawFileHPHi}{0.983}% [diagnostico] embarcados/results/csv/aud_motor_igual_20261025-140432.csv :: ic95_hi[protocolo,b00_rmspico_bruto,arquivo,HP_M1M3]
\newcommand{\audMotorBalRmsRawMotorHP}{0.754}% [diagnostico] embarcados/results/csv/aud_motor_igual_20261025-140432.csv :: acuracia_balanceada[protocolo,b00_rmspico_bruto,motor,HP_M1M3]
\newcommand{\audMotorBalRmsRawMotorHPLo}{0.623}% [diagnostico] embarcados/results/csv/aud_motor_igual_20261025-140432.csv :: ic95_lo[protocolo,b00_rmspico_bruto,motor,HP_M1M3]
\newcommand{\audMotorBalRmsRawMotorHPHi}{0.873}% [diagnostico] embarcados/results/csv/aud_motor_igual_20261025-140432.csv :: ic95_hi[protocolo,b00_rmspico_bruto,motor,HP_M1M3]
\newcommand{\audMotorDropRmsRawHP}{0.159}% [diagnostico] embarcados/results/csv/aud_motor_igual_20261025-140432.csv :: acuracia_balanceada[queda,b00_rmspico_bruto,HP_M1M3]
\newcommand{\audMotorDropRmsRawHPLo}{0.078}% [diagnostico] embarcados/results/csv/aud_motor_igual_20261025-140432.csv :: ic95_lo[queda,b00_rmspico_bruto,HP_M1M3]
\newcommand{\audMotorDropRmsRawHPHi}{0.251}% [diagnostico] embarcados/results/csv/aud_motor_igual_20261025-140432.csv :: ic95_hi[queda,b00_rmspico_bruto,HP_M1M3]
\newcommand{\audMotorDropRmsRawFull}{0.106}% [diagnostico] embarcados/results/csv/aud_motor_igual_20261025-140432.csv :: queda_completa[queda,b00_rmspico_bruto,HP_M1M3]
\newcommand{\audMotorHPFiles}{36}% [diagnostico] embarcados/results/csv/aud_motor_igual_20261025-140432.csv :: n_arquivos[queda]
\newcommand{\audMotorHPWindows}{360}% [diagnostico] embarcados/results/csv/aud_motor_igual_20261025-140432.csv :: n_janelas[queda]

% Gerado por embarcados/src/python/gerar_artigo_figuras.py. NÃO EDITAR À MÃO.
% Caminhos relativos a embarcados/paper/.
\newif\iffigesp
\figesptrue  % true só se a figura do ESP32 existir
\newcommand{\FigArq}{../figures/pdf/arquitetura.pdf}
\newcommand{\FigBalacc}{../figures/pdf/diag_balacc_en_20261023_222831.pdf}
\newcommand{\FigConf}{../figures/pdf/diag_confusao_p11_en_20261023_222831.pdf}
\newcommand{\FigEquiv}{../figures/pdf/equivalencia_erro_etapas_en.pdf}
\newcommand{\FigSpec}{../figures/pdf/diag_espectro_en_20261023_222831.pdf}
\newcommand{\FigEspLat}{../figures/pdf/esp32_lat_esp32_20261025-143026_en.pdf}


\begin{document}

\title{Porting a Harmonic-Sum Rotation Tracker for Acoustic Motor-State Monitoring to a
Microcontroller: Stage-Wise Verified C Port, Measured Cost, and What Two Public Datasets
Allow Us to Conclude\ifwip\ (Work in Progress)\fi}

% TODO(autores): nomes e afiliações a preencher pelos autores.
\author{\IEEEauthorblockN{[Authors to be completed]}
\IEEEauthorblockA{[Affiliation]}}

\maketitle

\begin{abstract}
We port to portable C a pipeline that classifies the operating state of small electric motors
from sound by borrowing ideas from bass-guitar transcription: a harmonic-sum rotation tracker
decoded by a Viterbi search, \VParNFeat{} physical features, cosine-to-prototype scoring and an
online hidden Markov model (HMM). The port is checked against the Python reference stage by
stage, on \VEqCases{} windows and \VEqSeqCombos{} sequence configurations, with tolerances
frozen before the formal comparison. In double precision all \VEqStagesTotal{} stage checks
pass. In float32, the build intended for the microcontroller, all \VEqRealCases{} windows of real
audio pass every stage; on \VEqFloatBadCases{} synthetic constant-offset-plus-noise windows the
envelope-modulation peak lies at the float32 noise floor and the modulation, feature and z-score
checks fail, while the class decision is identical on \VEqFloatDecisionPass{} of \VEqCases{}
windows. On an ESP32 without PSRAM the reference memory layout (estimated at
\VEstWsFloatThreeMb{}\,MB for a \VParSegIdmt{}-s window) does not fit. A named low-memory
variant fits \VParSegBldc{}-s windows, and on the chip its outputs agree with the reference within
the frozen tolerances on all \VChipEqCases{} windows compared; measured on the chip it is
\VChipRtfBldc{} times slower than real time, dominated by a table-free harmonic sum in
software-emulated double precision, and the \VParSegIdmt{}-s window does not fit. The pipeline does not outperform simple baselines:
on the IDMT dataset every method, including a two-number level baseline, saturates because
class, motor and recording coincide; on the A2212 dataset classical features reach
\VDiagBalClsBldcFile{} balanced accuracy against \VDiagBalDecBldcFile{} for the pipeline, and
with a motor held out all methods drop. The contribution is therefore the verified embedded
implementation, its cost, an interpretable rotation estimate, and an explicit account of what
these datasets can and cannot support.
\end{abstract}

\begin{IEEEkeywords}
acoustic condition monitoring, embedded signal processing, ESP32, rotation tracking,
Viterbi decoding, reproducibility, data leakage
\end{IEEEkeywords}

% ======================================================================
\section{Introduction}

Acoustic monitoring of electric motors is attractive because a microphone is cheap,
non-intrusive and needs no electrical connection to the machine. Several recent works report
near-perfect accuracy on two public datasets of small laboratory motors, the
IDMT-ISA-ELECTRIC-ENGINE set~\cite{grollmisch2019} and the A2212 brushless set~\cite{estacio2023},
often with deep networks~\cite{melcnn2025,ertarin2025} or large hand-crafted feature banks~\cite{etnext2026,tekin2025}.
Two aspects receive little attention. First, most of these results split recordings into
segments before the train/test split, so neighbouring segments of one recording fall on both
sides. Second, few works measure the cost on a microcontroller; an exception is the edge
classifier of Strantzalis et al.~\cite{strantzalis2022}, while other TinyML studies stop at a model
size and never run on hardware~\cite{almaini2026}.

This paper starts from a pipeline that we call Decifrae (p11). It reuses components of a
music-transcription system: the motor is treated as a ``one-string instrument'' whose
fundamental (the rotation frequency) is tracked through a sum over its harmonics and decoded with
a Viterbi search that has a silence state and an octave-jump cost, as a bass transcriber would do.
The analogy only motivates the \emph{representation}; we do not claim that motor sound is a
musical note, nor that music-transcription performance transfers to machines.

The question we answer is whether this pipeline fits and runs on an ESP32-class
microcontroller with an implementation that is demonstrably the same computation as the Python
reference, and what the two public datasets allow us to conclude about it. In short, it fits and
runs on an ESP32 without PSRAM only for \VParSegBldc{}-s windows, through a named low-memory
variant, and \VChipRtfBldc{} times slower than real time on that board.
\ifwip
We contribute a C port verified stage by stage against the reference with tolerances frozen before
the formal comparison, a cost account that keeps estimate, PC measurement, firmware build and
on-chip measurement apart, the tracked rotation as an interpretable output, and an evaluation
under three protocols showing that the pipeline does not beat simple baselines and that the
datasets cannot separate operating state from motor identity.
\else
Our contributions are:
\begin{itemize}
\item a portable C port without dynamic allocation, verified against the Python reference stage
by stage with tolerances frozen before the formal comparison, in double and in float32, with
every residual discrepancy listed;
\item a cost account that keeps four categories apart: analytical estimate, measurement on a PC,
firmware build (\texttt{idf.py size}), and measurement on an ESP32 without PSRAM, including the
low-memory variant needed there;
\item an interpretable intermediate output, the tracked rotation frequency, together with a check
of how far it can be trusted on each dataset;
\item an evaluation of the pipeline and of simple baselines on the same windows under three
protocols (random windows, grouped by recording or file, and motor held out), showing that the
pipeline does not outperform the baselines and that the datasets cannot separate operating
state from motor identity.
\end{itemize}
\fi
Three distinctions are kept throughout. We classify the \emph{current operating state} of a
motor; we do not predict future failures. Both datasets contain small laboratory motors, not
industrial machines. And an evaluation in which the test motor was also seen in training says
nothing about generalisation to a new motor.

% ======================================================================
\section{Related Work}

\emph{Datasets and splits.} Grollmisch et al.~\cite{grollmisch2019} released the IDMT set with
an official split (training on clean recordings, testing on recordings with added noise) and a
frame-level network. Each class was recorded on a different motor, so class and physical unit
coincide. Later works on IDMT and A2212 report near-perfect accuracy with mel-spectrogram
CNNs~\cite{melcnn2025}, image encodings~\cite{ertarin2025}, wavelet pattern features with
SVM/kNN~\cite{etnext2026} and wavelet statistics with bagged trees~\cite{tekin2025}; their splits
operate on segments rather than on recordings. The A2212 data paper~\cite{estacio2023} reports
no classifier; it documents that the bearing fault exists on a single motor.
\ifwip\else
The classical-feature baseline used here follows the feature list described in the abstract of
Vanaraj et al.~\cite{vanaraj2026}; we only read the abstract, so it is an inspired
re-implementation, not a reproduction. Glowacz~\cite{glowacz2018} classified a commutator motor
with a nearest-mean rule on selected spectral lines, the closest precedent of our
cosine-to-prototype decision.
\fi

\ifwip
\emph{Rotation tracking without a tachometer} has been reviewed on vibration
data~\cite{peeters2019}, and Viterbi decoding over a time--frequency representation predates our
work~\cite{xi2015}; we add its verified embedded implementation on microphone audio.
\else
\emph{Rotation tracking without a tachometer.} Peeters et al.~\cite{peeters2019} review and
compare tacholess speed estimators on vibration data, including multi-order probabilistic
approaches that combine several harmonics. Viterbi decoding over a time--frequency
representation for speed estimation predates our work~\cite{xi2015}, and harmonic weighting has
been proposed to suppress weak orders~\cite{yoo2023}. Wheat et al.~\cite{wheat2025} build
features on a harmonic grid given a known operating speed. Our tracker is therefore not new as an
idea; what we add is its embedded, verified implementation on microphone audio.
\fi

\ifwip\else
\emph{Embedded and anomaly-oriented work.} Strantzalis et al.~\cite{strantzalis2022} deploy a
small CNN on a Cortex-M4 and measure latency, RAM and response time to state changes on IDMT,
but with a random split of overlapping spectrograms. Almaini et al.~\cite{almaini2026} quantise an
MFCC network trained on urban sounds without deploying it. Unsupervised anomaly detection on
machine sounds (autoencoder baselines of the DCASE challenges) addresses novelty rather than state
classification and is not compared here.
\fi

% ======================================================================
\section{Method}
\label{sec:method}

Fig.~\ref{fig:arq} summarises the pipeline. Each analysis window (\VParSegIdmt{}\,s for IDMT,
\VParSegBldc{}\,s for A2212, both at \VParSrKhz{}\,kHz) is RMS-normalised and processed
independently; only the HMM carries state across windows.

\begin{figure}[t]
\centering
\includegraphics[width=\columnwidth]{\FigArq}
\caption{System overview. Everything inside the dashed box is the C port that runs on the
microcontroller; audio currently reaches the chip over UART (processor-in-the-loop) or from its
flash. Working memory: analytical estimate of the reference layout, and the pool peak of the
low-memory variant measured on the ESP32.}
\label{fig:arq}
\end{figure}

\emph{Rotation tracker.} A short-time Fourier transform with a symmetric Hann window of
\VParNfft{} samples and hop \VParHop{} gives magnitude frames $|X_t(k)|$. For each of
$J=\VParJ{}$ candidate rotation frequencies $f_j$ on a logarithmic grid from \VParFmin{} to
\VParFmax{}\,Hz (step \VParStepCents{}~cents), the saliency
$s_t(j)=\sum_{h=1}^{K}|X_t(h f_j/\Delta f)|$ sums $K=\VParNharm{}$ harmonics with linear
interpolation between bins, and is normalised per frame to $\tilde e_t(j)\in[0,1]$. The emission
cost is $-\log(\VParWEnergy{}\,\tilde e_t(j)+\VParEmisEps{})$; a silence state receives a cost
derived from the frame energy relative to the loudest frame of the window. A Viterbi search
allows jumps of up to \VParMaxJump{} states (\VParMaxJumpCents{} cents) with a cost proportional
to the jump, octave jumps of $\pm$\VParOctStates{} states with a fixed cost, and transitions into
and out of silence. The decoded path gives $f_{rot}(t)$.

\emph{Features and decision.} From the path and the mean power spectrum we compute \VParNFeat{}
features\ifwip\else{} (Table~\ref{tab:params})\fi: the median and inter-quartile range of $f_{rot}$, the silent-frame
fraction, the saliency along the path, a harmonic-energy ratio, the relative energy of four groups
of harmonic orders, the frequency and prominence of the main envelope-modulation peak (Hilbert
envelope, Butterworth low-pass, decimation to \VParEnvRateHz{}\,Hz, Welch spectrum), the spectral
centroid and the spectral flatness. Features are z-scored with training statistics and compared
by cosine with one prototype per class; a softmax with temperature \VParTempCos{} gives class
probabilities. For streams, an HMM with a fixed change penalty and an extra ``no decision''
state is run as a causal forward filter, one window at a time.

\ifwip\else
% Gerado por embarcados/src/python/gerar_artigo_tabelas.py a partir de embarcados/paper/numeros.json.
% NÃO EDITAR À MÃO.

\begin{table}[t]
\centering
\scriptsize
\setlength{\tabcolsep}{3pt}
\caption{Parameters of the p11 reference, read from the Python code (\texttt{PreDecifrae}, \texttt{rotacao}, \texttt{\_envelope\_mod}), and the 13 features in the order of the reference.}
\label{tab:params}
\begin{tabular}{@{}p{1.55cm}p{6.6cm}@{}}
\toprule
Sampling rate; window & 16\,kHz; 3\,s (IDMT), 1\,s (A2212); RMS $\rightarrow$ 0.1 \\
STFT & symmetric Hann, $N$=8{,}192, hop 1{,}024, no centring; 39 / 8 frames \\
Rotation grid & 15--60\,Hz, 2-cent step, $J$=1{,}201 states + silence \\
Harmonic sum & $K$=48 harmonics, unit weights, linear interpolation \\
Emission cost & $-\log(0.45\,\tilde e + \ensuremath{10^{-3}})$; silence from frame energy \\
Viterbi & step 0.55/semitone, jumps $\le$50 states, octave $\pm$600 states (cost 7.0), silence 3.2 \\
Envelope & $|\mathcal{H}(y)|$, Butterworth 4, 200\,Hz (causal \texttt{sosfilt}), $\downarrow$16, Welch; peak in 2--100\,Hz \\
Classifier & z-score, cosine to class means, softmax $T$=0.1 \\
Online HMM & change penalty swept (6 values, 0.5--16), $\log p_N$=-50 \\
\midrule
Features & (0) median $f_{rot}$; (1) IQR of $f_{rot}$; (2) silent-frame fraction; (3) mean normalised saliency on the path; (4) harmonic ratio ($k f_{rot}\pm1$ bin / 10--4{,}000\,Hz); (5) energy of orders 1--4 (dB); (6) energy of orders 5--12 (dB); (7) energy of orders 13--24 (dB); (8) energy of orders 25--48 (dB); (9) envelope-modulation peak frequency; (10) its prominence (dB); (11) spectral centroid; (12) spectral flatness \\
\bottomrule
\end{tabular}
\end{table}

\fi

\ifwip
\emph{C port.} The C code follows the Python line by line, compiles unchanged for the PC
and the ESP-IDF, processes the spectrogram frame by frame, includes a mixed-radix FFT for the
full-length Hilbert transform, and uses a caller-provided arena instead of dynamic allocation.
The real type is float32 on the microcontroller and double for verification. Any change of
arithmetic or memory layout is a named variant: \texttt{p11\_c\_sosd} keeps the envelope filter
state in double (more precise, not an approximation), and \texttt{\VChipVariant}, the one run on
our board without PSRAM, takes its memory from a pool of internal DRAM and IRAM blocks, drops the
harmonic table (positions recomputed in double) and computes the Hilbert FFT in four steps. An
ESP-DSP option (\texttt{p11\_c\_espdsp}) was neither compared nor measured and is left out.
\else
\emph{C port.} The C code follows the Python line by line and compiles unchanged for the PC and
for the ESP-IDF. It never stores the full spectrogram: each frame is windowed, transformed by a
real FFT, accumulated into the mean power spectrum and reduced to its saliency row. The Viterbi
search needs the whole window (the silence cost uses the maximum frame energy of the window), so
the normalised saliency and the back-pointers (short integers) are kept; the Hilbert envelope needs an FFT
of the full window length, which is not a power of two, so the port includes a mixed-radix FFT.
All buffers live in a caller-provided arena whose size is computed by the same layout function that
assigns the pointers. The real type is selected at compile time (float32 for the microcontroller,
double for verification). Index decisions that change integers (grid, medians, rounding of
harmonic positions) are computed in double in both builds, reproducing the reference's
round-half-to-even and first-index tie-breaking. Any change of arithmetic is a \emph{named
variant}, compared under the same tolerances. \texttt{p11\_c\_sosd} keeps the envelope filter
state in double, as SciPy does (more precise, not an approximation). \texttt{\VChipVariant}, run on
our board without PSRAM, calls the same core functions but carves its memory phase by phase from
a pool of internal DRAM and IRAM blocks (IRAM, word access only, holds only float vectors such as
the FFT twiddles), drops the harmonic-interpolation table (bin positions recomputed in double for
every frame) and computes the Hilbert FFT in four steps, the same DFT in another order. An ESP-DSP
option (\texttt{p11\_c\_espdsp}) was neither compared nor measured and is left out.
\fi

% ======================================================================
\section{Experimental Setup}
\label{sec:setup}

\emph{Data and protocols.} We use the \texttt{train/} part of IDMT (the official test set was not
opened) and all of A2212\ifwip\else{} (Table~\ref{tab:protocols})\fi. All methods are evaluated on exactly the
same windows and folds. On IDMT, the random-window protocol leaks by construction, and the
``blocks'' protocol holds out contiguous blocks of each recording; since there is one recording
per class, it still cannot separate class from motor. On A2212 we compare random windows, a
grouped split by file, and leave-one-motor-out; because the bearing fault exists only on motor
M4, we also report leave-one-motor-out without the M4 fold.

\ifwip\else% Gerado por embarcados/src/python/gerar_artigo_tabelas.py a partir de embarcados/paper/numeros.json.
% NÃO EDITAR À MÃO.

\begin{table}[t]
\centering
\scriptsize
\setlength{\tabcolsep}{3pt}
\caption{Datasets and evaluation protocols. IDMT: \texttt{train/} only, 3 recordings of 3 motors, 357 windows of 3\,s. A2212: 43 files of 4 motors (22 Healthy, 14 Propeller, 7 Bearing), 430 windows of 1\,s. $g$: groups used by the bootstrap. $^\dagger$Intentional leakage, reported only to measure inflation.}
\label{tab:protocols}
\begin{tabular}{@{}p{1.55cm}p{1.35cm}r p{3.75cm}@{}}
\toprule
Protocol & Split unit & $g$ & What it can / cannot support \\
\midrule
IDMT random$^\dagger$ & window & 18 & Upper bound only: neighbouring windows of one recording on both sides. \\
IDMT blocks & contiguous block & 15 & One recording per class: class = motor = recording. Measures recording identity, not operating state. \\
A2212 random$^\dagger$ & window & 43 & Upper bound only: windows of one file on both sides. \\
A2212 file & file & 43 & State of a \emph{known} motor (same motors in training and test). Bearing exists only on motor M4. \\
A2212 motor-out & motor & 43 & Unseen motor. With M4 held out, Bearing has no training data (recall 0.000 by construction). Only 4 motors: no motor-level CI. \\
A2212 motor-out w/o M4 & motor (M1--M3) & 36 & Healthy vs.\ Propeller on an unseen motor; 3 folds. \\
\bottomrule
\end{tabular}
\end{table}
\fi

\emph{Baselines.} A level baseline uses only the RMS and peak of the window with logistic
regression, once on RMS-normalised audio (only the crest factor remains) and once on raw audio. A
classical baseline (p04) uses \VParNClassicFeat{} frame statistics (zero-crossing rate, RMS,
spectral centroid and MFCC means and deviations) with a soft-voting ensemble, inspired
by~\cite{vanaraj2026}.

\emph{Metrics.} We report balanced accuracy, macro-F1, macro one-vs-rest ROC-AUC from continuous
scores, and per-class recall. Confidence intervals (\VDataCiLevel{}\%) come from a bootstrap that
resamples groups (recording blocks or files), stratified by class, with \VDataNBoot{} resamples;
windows of one group are never treated as independent. For the motor-out protocols the
independent unit is the motor, of which there are only \VDataBldcMotors{}, so no honest interval
exists at that level; the file-level interval shown is optimistic.

\ifwip
\emph{Equivalence protocol.} Golden vectors of every stage were produced by the
instrumented reference on \VEqCases{} windows (\VEqGoldClean{} clean, \VEqGoldNoisy{} noisy,
\VEqGoldEdge{} edge cases) and \VEqSeqCombos{} HMM sequence configurations. Per-stage tolerances,
derived from rounding analysis and an independent emulation, were frozen (\texttt{\VEqTolCommit{}})
before the formal comparison, although the C port had already been debugged against the same
golden vectors by its developer; discrete outputs must agree exactly unless a tie is demonstrated.
\else
\emph{Equivalence protocol.} The instrumented Python reference produced golden vectors for every
stage on \VEqCases{} windows: \VEqGoldClean{} clean validation windows from both datasets,
\VEqGoldNoisy{} noisy versions under \VEqGoldPerts{} perturbations (white noise at several SNRs and
synthetic sag and flicker power-quality disturbances), and \VEqGoldEdge{} edge cases (silence, clipping, synthetic
rotation outside the search range, and a constant offset with tiny noise). The HMM stages were
checked on \VEqGoldSeqs{} probability sequences under \VParNChangePen{} change penalties and
\VEqMinDurValues{} minimum durations (\VEqSeqCombos{} configurations). Tolerances for each stage
were derived from first-order rounding analysis and an independent numerical emulation of each
build, and frozen before the formal comparison (commit \texttt{\VEqTolCommit{}}, file hash
\texttt{\VEqTolSha{}}\ldots, \VEqTolCommitTime{}; comparison run \VEqRunTime{}); the hash is
recorded in every result row. The chronology matters: the C port had already been debugged
against the same golden vectors by its developer before the tolerances were frozen, so the
tolerances are frozen, not pre-registered. On \audHoldoutN{} windows outside the golden set,
unseen by the C developer, path and decision agree on \audHoldoutFloatPathSame{} and
\audHoldoutFloatDecSame{} in float32, no feature exceeds its tolerance (maximum ratio
\audHoldoutFloatRatioMax{}) and probabilities differ by at most \audHoldoutFloatDpMax{}. Tolerances for the edge cases
were calibrated on the emulation error for the same signals, a weaker standard, and some are wide
in float32 (e.g.\ the modulation prominence). Continuous stages are compared by maximum error relative to the stage's
scale; discrete outputs (rotation path, decision, HMM states) must agree exactly unless a tie is
demonstrated from the reference's own values. A failure at one stage marks later stages as
downstream. Both emulations passed all comparison rows (\VEqEmulDoublePass{} and \VEqEmulFloatPass{} of
\VEqEmulFloatRows{}) before the C code was compared.
\fi

\ifwip
\emph{Cost measurement.} We keep apart (1) the analytical arena estimate, (2) timings on
the PC (\VPcNIdmt{} timings cycling over \VPcDistinctIdmt{} distinct windows, each run once in a new
process, so warm-up does not warm the C code; workspace initialisation excluded), (3)
\texttt{idf.py size} of the firmware (ESP-IDF \VBuildIdf{}) and (4) measurements on an
\VChipModel{} without PSRAM (\VChipCpuMhz{}\,MHz; \VChipLatBldcN{} timed windows over UART after
\VChipLatWarmup{} warm-up runs, on-chip timer around the processing only). Energy is not measured.
\else
\emph{Cost measurement.} We keep four categories apart. (1) \emph{Analytical estimate}: the
arena size returned by the C layout function. (2) \emph{Measured on the PC}: \VPcNIdmt{} timings
per dataset cycling over \VPcDistinctIdmt{} distinct windows. The C executable is started anew and
timed once per window, so warm-up does not warm it, and its time excludes the workspace
initialisation, which the Python timing cannot exclude; the two effects partly offset.
(3) \emph{Firmware build}: \texttt{idf.py size} (ESP-IDF \VBuildIdf{}, GCC \VBuildGcc{}).
(4) \emph{Measured on the ESP32}: an \VChipModel{} (rev.\ \VChipRev{}, \VChipCpuMhz{}\,MHz,
\VChipFlashMb{}\,MB flash, no PSRAM, radios never started); windows arrive over UART (or, for the
stream, from flash) and are timed with the on-chip timer around each core call, with heap and
stack watermarks. Energy is not measured.
\fi

% ======================================================================
\section{Results}
\label{sec:results}

\subsection{Python--C equivalence}

\ifwip
Table~\ref{tab:equiv} summarises the verdicts. In double precision all \VEqStagesTotal{}
stage checks pass on all windows and sequence configurations; the rotation path, the decision and
every HMM output agree exactly. In float32 all \VEqRealCases{} windows of real audio pass every
stage. The failures are confined to the \VEqFloatBadCases{} synthetic constant-offset-plus-noise
windows, whose modulation spectrum lies at the float32 rounding floor in both implementations, so
the peak location is undetermined (the reference peak is \VEqDcRefPeakRatio{} times the reference
power at the bin chosen by the C code: not a tie). We count the modulation, feature and z-score
checks as failed there (\VEqFloatStagesPassStrict{} of \VEqStagesTotal{} checks pass on all
windows); the class decision agrees on \VEqFloatDecisionPass{} of \VEqCases{} windows. The named
variant \texttt{p11\_c\_sosd} passes every check. An independent audit reproduced the comparison
from a clean build; its corrections are applied here. These are PC results; the comparison on the
chip is in Section~\ref{sec:cost}.
\else
Table~\ref{tab:equiv} summarises the verdicts\ifwip\else{} and Fig.~\ref{fig:equiv} the worst
error-to-tolerance ratio per stage\fi. In double precision all \VEqStagesTotal{} stage checks pass
on all windows and sequence configurations (worst ratio \VEqDoubleMaxRatio{} of the tolerance).
The rotation path, the decision and every HMM output agree exactly. In float32, all
\VEqRealCases{} windows of real audio (clean and noisy) pass every stage. The
\VEqFloatBadCases{} remaining failures are the synthetic constant-offset-plus-noise edge case, once
per model. There the envelope is almost constant and its modulation spectrum lies at the float32
rounding floor in both implementations, so the peak location is undetermined at that precision:
the reference's peak is \VEqDcRefPeakRatio{} times the reference power at the bin chosen by the C
code, so this is not a tie. The modulation frequency moves by \VEqFloatFmodErr{}\,Hz, and we count
the envelope-modulation, feature and z-score checks as failed on these windows (the frozen
comparison rule had accepted the first two as a tie; we report the conservative reading), leaving
\VEqFloatStagesPassStrict{} of \VEqStagesTotal{} stage checks passing on all windows. Cosines,
probabilities and the class decision still pass on all windows
(\VEqFloatDecisionPass{}/\VEqCases{}). With the named variant \texttt{p11\_c\_sosd} every check
passes; it reduces the maximum error of the decimated envelope,
relative to its peak and taken as the median over windows, from \VEqEnvDecFloatMed{} to
\VEqEnvDecSosdMed{}, because the low-pass poles lie close to
the unit circle and float32 rounding accumulates in the filter state. These results were obtained
on the PC; the same comparison on the chip is
\ifchipmedido reported in Section~\ref{sec:cost}\else pending\fi. An independent audit
(report in the repository) reproduced the comparison from a clean build and its corrections are
applied here. We therefore
claim stage-wise agreement within tolerances frozen before the formal comparison, with the
exceptions above, not unconditional equivalence.
\fi

\begin{table}[t]
\centering
\scriptsize
\caption{Stage checks passed on the PC (\VEqCases{} windows; \VEqSeqCombos{} sequence
configurations). Conservative reading: a group fails if any stage fails on any window. All
float32 failures are the \VEqFloatBadCases{} synthetic constant-offset edge windows, where the
modulation peak is at the float32 noise floor; the decision is unchanged.}
\label{tab:equiv}
% Gerado por embarcados/src/python/gerar_artigo_tabelas.py a partir de embarcados/results/csv/equivalencia_resumo_por_build_20261024-101821.csv (veredito por etapa).
% NÃO EDITAR À MÃO.

\begin{tabular}{@{}lccc@{}}
\toprule
Stage group (\#stages) & C double & C float32 & float32+SOS64 \\
\midrule
Front end (6) & 6/6 & 6/6 & 6/6 \\
Rotation Viterbi (3) & 3/3 & 3/3 & 3/3 \\
Spectrum, envelope (6) & 6/6 & 5/6 (FAIL in 2/164) & 6/6 \\
13 features (1) & 1/1 & 0/1 (FAIL in 2/164) & 1/1 \\
z-score (1) & 1/1 & 0/1 (FAIL in 2/164) & 1/1 \\
Cosine, softmax, decision (3) & 3/3 & 3/3 & 3/3 \\
Sequences (HMM) (6) & 6/6 & 6/6 & 6/6 \\
\bottomrule
\end{tabular}

\end{table}

\ifwip\else
\begin{figure}[t]
\centering
\includegraphics[width=\columnwidth]{\FigEquiv}
\caption{Worst error-to-tolerance ratio per stage (PC). Ratios below one pass. F: failure on the
synthetic constant-offset edge windows (counted as failures even where the frozen rule accepted a
tie). Exact zeros are drawn at the bottom of the axis.}
\label{fig:equiv}
\end{figure}
\fi

\subsection{Cost}
\label{sec:cost}

\ifwip
\emph{Measured on the PC},
the C float32 build of the reference layout needs \VPcCIdmtMed{}\,ms per \VParSegIdmt{}-s window
and \VPcCBldcMed{}\,ms per \VParSegBldc{}-s window, against \VPcPyIdmtMed{}\,ms and
\VPcPyBldcMed{}\,ms for the Python reference (\VPcSpeedupIdmt{}$\times$, including interpreter
overhead); the STFT with harmonic sum takes \VPcCIdmtTrackPct{}\%, the rotation Viterbi
\VPcCIdmtVitPct{}\% and the envelope branch \VPcCIdmtEnvPct{}\% of the median time. These numbers
say nothing about the microcontroller.
\else
Table~\ref{tab:resources} lists the resources by category. The PC timings (column 2; the STFT
with harmonic sum takes \VPcCIdmtTrackPct{}\%, the rotation Viterbi \VPcCIdmtVitPct{}\% and the
envelope branch \VPcCIdmtEnvPct{}\% of the median) say nothing about the microcontroller.
\fi

The \emph{analytical estimate} of the working memory of the reference layout (one arena) is
\VEstWsFloatThreeKib{}\,KiB for \VParSegIdmt{}\,s and \VEstWsFloatOneKib{}\,KiB for
\VParSegBldc{}\,s, at least \VChipRefOverPoolThree{} and \VChipRefOverPoolOne{} times the internal
RAM that the firmware can hand to the working memory on our board (\VChipPoolKib{}\,KiB, measured)
and at least \VBuildEspsWsRatio{} and \VBuildEspsWsRatioOne{} times the whole DIRAM of the
ESP32-S3 \emph{build} (\VBuildEspsDramTotalKb{}\,KiB; its IRAM is \VBuildEspsIramPct{}\% used):
without PSRAM the reference layout cannot run.

\ifchipmedido
\emph{Measured on the ESP32} (\VChipPsramMb{}\,MB PSRAM), \texttt{\VChipVariant} fits a
\VParSegBldc{}-s window: the pool peaks at \VChipWsBldcB{}\,bytes (\VChipWsPoolPct{}\% of it,
Fig.~\ref{fig:esplat}b), the minimum free internal heap is \VChipHeapIntBldcB{}\,bytes and the
stack headroom \VChipStackBldc{}\,bytes. The \VParSegIdmt{}-s IDMT window does \emph{not} fit, so
no IDMT figure exists on the chip. All \VChipEqCases{} golden \VParSegBldc{}-s windows
(\VChipEqCasesReal{} real, \VChipEqCasesEdge{} edge cases) pass, on the chip, the \VChipEqStages{}
final outputs the firmware returns (Viterbi path and score through decision), with \VChipEqFail{} failures and a worst error-to-tolerance
ratio of \VChipEqMaxRatio{}; intermediate stages are not sent back. On the PC the variant passes
all \VEqLmStages{} compared stages. The median latency is \VChipLatBldcMed{}\,ms per
\VParSegBldc{}-s window (p95 \VChipLatBldcPninetyfive{}, maximum \VChipLatBldcMax{}\,ms), a
real-time factor of \VChipRtfBldc{}: on this board the pipeline does \emph{not} run in real time.
The STFT with harmonic sum takes \VChipBldcTrackPct{}\% of it (Fig.~\ref{fig:esplat}a; on the PC,
with the table, \VPcCBldcTrackPct{}\%). We attribute this to recomputing the harmonic positions in
double precision, which the ESP32 only emulates in software (the variant was not timed on the PC);
a float32 variant of that step or a board with PSRAM is the way towards real time. The variant's
image occupies \VBuildLmBinKb{}\,KiB of flash (build).
\else
\emph{Measured on the ESP32}: not available when this version was generated.
\fi

\ifwip\else% Gerado por embarcados/src/python/gerar_artigo_tabelas.py a partir de embarcados/paper/numeros.json.
% NÃO EDITAR À MÃO.

\begin{table*}[t]
\centering
\scriptsize
\setlength{\tabcolsep}{4pt}
\caption{Resources by evidence category: (1) analytical estimate (\texttt{p11\_host ws}); (2) measured on the PC (Windows 11, MSVC 19.44 x64 \slash{}O2 \slash{}fp:precise; 200 timings per dataset over 75 distinct windows, one new process per window); (3) \texttt{idf.py size}, static only; (4) measured on the ESP32 without PSRAM (200 timed windows after 5 warm-up runs). Columns (1)--(2) use the reference layout \texttt{p11\_c\_ref}; column (4), the only measurement on the microcontroller, uses \texttt{p11\_c\_lowmem}.}
\label{tab:resources}
\begin{tabular}{@{}lllll@{}}
\toprule
Quantity & (1) Estimate & (2) PC & (3) Build & (4) ESP32 \\
\midrule
Latency, C float32, median [p95] (ms), 3\,s / 1\,s & -- & 13.29 [15.00] / 3.18 [3.84] & -- & \textit{did not fit} / 2647.5 [2647.6] \\
Latency, Python reference, median (ms), 3\,s / 1\,s & -- & 95.3 / 22.1 & -- & -- \\
Real-time factor (C, median), 3\,s / 1\,s; STFT + harm.\ sum share at 1\,s & -- & 0.0044 / 0.0032; 39\,\% & -- & \textit{did not fit} / 2.65; 89\,\% \\
Working memory (KiB), 3\,s / 1\,s & 1820 / 936 & -- & -- & \textit{did not fit} / 296 (pool peak) \\
Internal RAM for the pool; min.\ free heap (KiB); stack headroom (B) & -- & -- & -- & 308; 7.4; 20{,}536 \\
App binary (KiB): ESP32, -S3 (\texttt{p11\_c\_ref}); ESP32 (\texttt{p11\_c\_lowmem}) & -- & -- & 209, 233; 215 & -- \\
Static DRAM, IRAM used / region (KiB), ESP32, \texttt{p11\_c\_lowmem} & -- & -- & 35/176, 58/128 & -- \\
\bottomrule
\end{tabular}
\end{table*}
\fi

\iffigesp
\begin{figure}[t]
\centering
\includegraphics[width=\columnwidth]{\FigEspLat}
\caption{Measured on the ESP32 without PSRAM (\texttt{\VChipVariant}, A2212, \VParSegBldc{}-s
windows). (a) Median time per stage over \VChipLatBldcN{} windows. (b) Internal RAM handed to the
pool and its peak use (measured) against the estimated reference arena (hatched).}
\label{fig:esplat}
\end{figure}
\fi

\subsection{Classification under three protocols}

% Gerado por embarcados/src/python/gerar_artigo_tabelas.py a partir de embarcados/paper/numeros.json.
% NÃO EDITAR À MÃO.

\begin{table}[t]
\centering
\scriptsize
\setlength{\tabcolsep}{3pt}
\caption{Balanced accuracy on the same windows (validation only); $g$ = groups. 95\% CIs from a group bootstrap (2{,}000 resamples) are shown in Fig.~\ref{fig:balacc}. $^\dagger$Intentional leakage.}
\label{tab:desempenho-compacto}
\begin{tabular}{@{}lrcccc@{}}
\toprule
Protocol & $g$ & RMS/pk (norm.) & RMS/pk (raw) & Classical & p11 \\
\midrule
IDMT random$^\dagger$ & 18 & 0.992 & 0.997 & 1.000 & 1.000 \\
IDMT blocks & 15 & 0.991 & 0.993 & 1.000 & 1.000 \\
A2212 random$^\dagger$ & 43 & 0.564 & 0.625 & 0.998 & 0.849 \\
A2212 file & 43 & 0.562 & 0.608 & 0.981 & 0.816 \\
A2212 motor-out & 43 & 0.542 & 0.503 & 0.613 & 0.456 \\
A2212 w/o M4 & 36 & 0.814 & 0.754 & 0.920 & 0.685 \\
\bottomrule
\end{tabular}
\end{table}


\begin{figure}[t]
\centering
\includegraphics[width=\columnwidth]{\FigBalacc}
\caption{Balanced accuracy on the same windows with \VDataCiLevel{}\% group-bootstrap
intervals. Dotted line: chance. *Intentional leakage. On IDMT class = motor = recording; on A2212
Bearing exists only on motor M4 (confounded with motor identity), so motor-out Bearing recall is
zero by construction.}
\label{fig:balacc}
\end{figure}

Table~\ref{tab:desempenho-compacto} and
Fig.~\ref{fig:balacc} compare the methods. On IDMT every method is near the ceiling under both
protocols: the RMS-normalised level baseline, which only sees the crest factor, reaches
\VDiagBalRmsNormIdmtBlk{} balanced accuracy on held-out blocks, and the pipeline and the
classical features reach \VDiagBalDecIdmtBlk{} and \VDiagBalClsIdmtBlk{}. Because each class is one recording of one motor,
these numbers measure recording identity and cannot rank the methods.

On A2212 with a grouped split by file, the classical features reach \VDiagBalClsBldcFile{}
[\VDiagBalClsBldcFileLo{}, \VDiagBalClsBldcFileHi{}], the pipeline \VDiagBalDecBldcFile{}
[\VDiagBalDecBldcFileLo{}, \VDiagBalDecBldcFileHi{}] and the raw level baseline
\VDiagBalRmsRawBldcFile{}. The pipeline errs mostly on Healthy (recall
\VDiagRecDecBldcFileHealthy{}). The random split inflates balanced accuracy by
\VDiagInflDecBldc{} for the pipeline and \VDiagInflClsBldc{} for the classical features. Holding a
motor out lowers every method: the classical features fall to \VDiagBalClsBldcMotor{} and the
pipeline to \VDiagBalDecBldcMotor{}; Bearing recall is \VDiagRecDecBldcMotorBearing{} by
construction, since the only bearing motor is then absent from training. Without the M4 fold,
on Healthy versus Propeller of unseen motors, the classical features reach
\VDiagBalClsBldcNoMfour{}, the RMS-normalised level baseline \VDiagBalRmsNormBldcNoMfour{} and the
pipeline \VDiagBalDecBldcNoMfour{}. Compared like for like (Healthy and Propeller of motors
M1--M3 under the file split), holding the motor out costs \VDiagDropHPClsBldc{} for the classical
features and \VDiagDropHPDecBldc{} for the pipeline; the much larger drop of the full motor-out
protocol is dominated by the Bearing class losing its only training motor. The pipeline never
outperforms the classical features (it ties at the IDMT ceiling), and with a motor held out its
point estimate also falls below the level baseline, with overlapping intervals.

% matriz de confusão: figures/pdf (FigConf), fora do texto por espaço

\subsection{Interpretability: the tracked rotation}

\ifwip
On IDMT the median tracked rotation differs by class (\VRotIdmtGoodMed{},
\VRotIdmtBrokenMed{}, \VRotIdmtHeavyMed{}\,Hz), but class equals motor there. On A2212 its rank
correlation with the speed label of the files is \VRotBldcSpearman{} pooled over motors; within
motors it is \VRotBldcSpearmanMOne{} to \VRotBldcSpearmanMFour{} but significant only for M4
(\VRotBldcFilesMFour{} files), so it should not be read as shaft speed on that dataset.
\else
The tracker exposes a physical quantity that a technician can inspect. On IDMT, the median tracked rotation is \VRotIdmtGoodMed{},
\VRotIdmtBrokenMed{} and \VRotIdmtHeavyMed{}\,Hz for good, broken and heavy-load windows, with
inter-quartile ranges of at most \VRotIdmtHeavyIqr{}\,Hz; \VRotIdmtGoodOctPct{}\% of the good
windows lie above \VRotOctThr{} times the class median (all at \VRotOctRatioMin{}--\VRotOctRatioMax{}
times), an octave error. Because class equals motor on IDMT, these differences cannot be
attributed to the operating state. On A2212, pooled over the \VDataBldcRecordings{} files, the rank
correlation between the tracked rotation and the speed label in the file names is
\VRotBldcSpearman{}; within motors it is \VRotBldcSpearmanMOne{} to \VRotBldcSpearmanMFour{}
(\VRotBldcFilesMThree{}--\VRotBldcFilesMTwo{} files each), significant only for M4
($p \geq \audRotBldcSpearmanPMTwo{}$ for M1--M3). The absolute level
is not comparable across motors and \VRotBldcCeilPct{}\% of the windows lie within
\VRotCeilPctThr{}\% of the upper limit of the search range. The range and the grid step were chosen
on the IDMT training data, so on A2212 the tracker output should not be read as the shaft speed.
\fi

% espectro: figures/pdf (FigSpec), fora do texto por espaço

\ifwip\else
\subsection{Streaming decision}

On a stream of \VPcSeqWindows{} clean A2212 golden windows (\VPcSeqPerClass{} per class:
Healthy, Propeller, Bearing) the online HMM with change penalty \VPcSeqCp{} detected both changes
with a delay of \VPcSeqDelayC{} windows in C and \VPcSeqDelayPy{} in Python on the PC (maximum
probability difference \VPcSeqMaxProbDiff{}).
\ifchipmedido On the ESP32, the same stream read from flash by \texttt{\VChipVariant} gave a delay
of \VChipSeqDelay{} windows and online decisions equal to Python's on \VChipSeqSame{} of
\VChipSeqWindows{} windows; one HMM step takes \VChipSeqHmmUs{}\,$\mu$s.\fi{} Bearing is motor M4,
so this checks the implementation, not the detector.
\fi

% ======================================================================
\section{Discussion}

\emph{What the datasets allow us to conclude.} IDMT, used without its test set, supports no
statement about operating-state discrimination: one recording per class makes class, motor and
session the same variable, and a two-number level baseline is as good as any model. A2212 with a
file-level split supports statements about the state of motors that were \emph{already seen} in
training, and there a cheap classical feature set is better than our pipeline. Generalisation to a
new motor is the relevant question for a deployed monitor, and with \VDataBldcMotors{} motors, a
bearing fault present on one of them and speed ranges that differ by class~\cite{estacio2023}, the
data can only show that every method degrades; it cannot rank methods or estimate how a monitor
would behave on another motor. None of these experiments concerns failure prediction, and none
involves industrial motors. In our runs the random split inflated A2212 balanced accuracy only
moderately (\VDiagInflClsBldc{} to \VDiagInflDecBldc{}), and holding a motor out cost, like for
like on Healthy versus Propeller, a similar amount (\VDiagDropHPClsBldc{} for the classical features,
\VDiagDropHPDecBldc{} for the pipeline); the large drop of the full motor-out protocol is dominated
by the Bearing class, which exists only on M4. Published near-perfect accuracies on these datasets
should be read against that confound, which no split by segment or by file removes.

\ifwip\else
\emph{What the embedded work shows.} The C port reproduces the reference stage by stage within
tolerances that were frozen before the formal comparison, including the discrete Viterbi path and
the HMM decisions; the residual float32 failures occur only on a synthetic edge case whose
modulation peak lies at the float32 noise floor, they leave the decision unchanged, and they
disappear when the envelope filter state is kept in double.
\ifchipmedido
On a board without PSRAM, \texttt{\VChipVariant} runs \VParSegBldc{}-s windows with on-chip outputs
within those tolerances, but at a real-time factor of \VChipRtfBldc{}, and the \VParSegIdmt{}-s
window does not fit. We claim neither real-time operation nor low power consumption.
\else
Whether it runs in real time on the chip is still open: the PC timings do not transfer.
\fi
\fi

\emph{Limitations and future work.} Energy was not measured (no current sensor), and we do not
estimate it from datasheets. All audio came from the datasets through the serial link; a real
I2S microphone, a motor on a test bench with a tachometer to validate the rotation estimate, and
the official IDMT test set (unopened by design) are future work. An independent scientific audit
of datasets, equivalence and traceability (report in the repository) was carried out before the
on-chip measurements and does not cover them.

% ======================================================================
\section{Conclusion}

A music-transcription-inspired pipeline for acoustic motor-state classification was ported to
portable C and verified against its Python reference stage by stage, with tolerances frozen
before the formal comparison: all checks pass in double; in float32 every window of real audio
passes every stage, and only \VEqFloatBadCases{} synthetic edge windows fail the modulation,
feature and z-score checks, with an unchanged decision. On an ESP32 without PSRAM the reference
layout does not fit; the named variant \texttt{\VChipVariant} fits only \VParSegBldc{}-s windows,
matches the reference within the frozen tolerances on the chip, and is \VChipRtfBldc{} times slower
than real time, dominated by the table-free harmonic sum in emulated double precision; a float32
variant of that step or a board with PSRAM is the way forward. On the two public
datasets it does not beat simple baselines, and the datasets themselves cannot separate operating
state from motor identity. The value of the work lies in the verified, measurable implementation and
in making these limits explicit.

\bibliographystyle{IEEEtran}
\bibliography{references}

\end{document}
